% !TeX program = xelatex
% Standalone English volume. Use XeLaTeX until the references stabilize.
\documentclass[a4paper,11pt]{article}
\usepackage[T1]{fontenc}
\usepackage{lmodern}
\usepackage[english]{babel}
\usepackage[left=26mm,right=26mm,top=25mm,bottom=25mm,headheight=15pt,headsep=7mm,footskip=11mm]{geometry}
\usepackage{amsmath,amssymb,mathtools,amsthm}
\newtheorem{proposition}{Proposition}[section]
\newtheorem{lemma}[proposition]{Lemma}
\newtheorem{theorem}[proposition]{Theorem}
\newtheorem{corollary}[proposition]{Corollary}
\newtheorem{maintheorem}[proposition]{Main Theorem}
\theoremstyle{definition}
\newtheorem{definition}[proposition]{Definition}
\theoremstyle{remark}
\newtheorem{remark}[proposition]{Remark}
\makeatletter
\renewcommand*\l@section{\@dottedtocline{1}{0em}{2.4em}}
\renewcommand*\l@subsection{\@dottedtocline{2}{1.45em}{3.45em}}
\renewcommand*\l@subsubsection{\@dottedtocline{3}{3.8em}{4.55em}}
\makeatother
\usepackage{xcolor,fancyhdr,titlesec,needspace,xurl}
\usepackage{longtable,multicol}
\usepackage{hyperref}
\definecolor{ink}{HTML}{1A3542}
\definecolor{muted}{HTML}{4F606A}
\hypersetup{unicode=true,colorlinks=true,linkcolor=ink,filecolor=ink,urlcolor=ink,citecolor=ink,bookmarksnumbered=true,
 pdfdisplaydoctitle=true,pdflang={en},
 pdftitle={The generalized Riemann hypothesis for Dirichlet L-functions. Edition 930, Volume XVIII},
 pdfsubject={A natural-language reconstruction of the proposed proof},
 pdfauthor={A mathematical reconstruction}}
\titleformat{\section}{\Large\bfseries\color{ink}}{\thesection}{.65em}{}
\titlespacing*{\section}{0pt}{.8em}{.65em}
\setlength{\parindent}{1.2em}
\linespread{1.045}
\setlength{\parskip}{.22em plus .05em minus .03em}
\setlength{\emergencystretch}{2em}
\setcounter{topnumber}{2}
\setcounter{tocdepth}{3}\setcounter{secnumdepth}{3}
\widowpenalty=4500
\clubpenalty=4500
\pagestyle{fancy}\fancyhf{}
\fancyhead[L]{\small\color{muted}Dirichlet GRH \textperiodcentered\ Volume I}
\fancyhead[R]{\small\color{muted}Dirichlet zeros and power intervals}
\fancyfoot[C]{\small\thepage}
\renewcommand{\headrulewidth}{.3pt}
\newcommand{\lead}[1]{\par\smallskip\noindent{\bfseries\color{ink}#1.}\enspace}
\newcommand{\sour}[2]{\unskip\penalty150\hspace{.35em}{\small\hyperlink{ref#1}{[#2]}}}
\titleformat{\subsection}{\large\bfseries}{\thesubsection}{.6em}{}
\titlespacing*{\subsection}{0pt}{1.15em plus .15em}{.4em}
\setlength{\abovedisplayskip}{9pt plus 2pt minus 3pt}
\setlength{\belowdisplayskip}{9pt plus 2pt minus 3pt}
\newcommand{\RR}{\mathbb R}\newcommand{\CC}{\mathbb C}\newcommand{\QQ}{\mathbb Q}
\newcommand{\ZZ}{\mathbb Z}\newcommand{\ord}{\operatorname{ord}}
\newcommand{\Frob}{\operatorname{Frob}}\newcommand{\Gal}{\operatorname{Gal}}
\newcommand{\len}{\operatorname{length}}\newcommand{\Res}{\operatorname*{Res}}
\newcommand{\wmean}{\operatorname{av}}\newcommand{\Ker}{\operatorname{ker}}
\newcommand{\Lh}{\mathcal L}\newcommand{\Gg}{\mathcal G}
\newcommand{\real}{\operatorname{Re}}\newcommand{\imag}{\operatorname{Im}}
\newcommand{\abs}[1]{\left|#1\right|}\newcommand{\eps}{\varepsilon}
\newcommand{\anchoritem}[2]{\par\Needspace{2.8\baselineskip}\hypertarget{ref#1}{}\noindent{\bfseries\color{ink}#2.}\enspace}
\newcommand{\srcloc}[2]{{\small\nolinkurl{#1.lean}}, line~#2}
\numberwithin{equation}{section}
\raggedbottom

\newcommand{\Sel}{\mathcal S}
\newcommand{\Cl}{\operatorname{Cl}}
\newcommand{\Aut}{\operatorname{Aut}}
\newcommand{\GL}{\operatorname{GL}}
\newcommand{\OO}{\mathcal O}
\hypersetup{bookmarksdepth=3}


\addto\captionsenglish{\renewcommand{\partname}{Volume}}


\makeatletter
\expandafter\def\csname r@cor:cost128\endcsname{{100.17}{79}{}{proposition.100.17}{Dirichlet-GRH-v930-Volume-XIX.pdf}}
\expandafter\def\csname r@cor:count111\endcsname{{9.15}{119}{}{proposition.9.15}{Dirichlet-GRH-v930-Volume-I.pdf}}
\expandafter\def\csname r@cor:interval56\endcsname{{65.23}{41}{}{proposition.65.23}{Dirichlet-GRH-v930-Volume-XIII.pdf}}
\expandafter\def\csname r@cor:normaldegree347\endcsname{{127.2}{27}{}{proposition.127.2}{Dirichlet-GRH-v930-Volume-XXIV.pdf}}
\expandafter\def\csname r@cor:original200\endcsname{{83.7}{57}{}{proposition.83.7}{Dirichlet-GRH-v930-Volume-XVI.pdf}}
\expandafter\def\csname r@cor:rank111\endcsname{{9.14}{119}{}{proposition.9.14}{Dirichlet-GRH-v930-Volume-I.pdf}}
\expandafter\def\csname r@cor:regular207\endcsname{{98.17}{17}{}{proposition.98.17}{Dirichlet-GRH-v930-Volume-XIX.pdf}}
\expandafter\def\csname r@cor:sameprimes517\endcsname{{9.34}{134}{}{proposition.9.34}{Dirichlet-GRH-v930-Volume-I.pdf}}
\expandafter\def\csname r@cor:trace235\endcsname{{110.7}{43}{}{proposition.110.7}{Dirichlet-GRH-v930-Volume-XXI.pdf}}
\expandafter\def\csname r@eq:action29\endcsname{{121.2}{45}{}{equation.121.2}{Dirichlet-GRH-v930-Volume-XXIII.pdf}}
\expandafter\def\csname r@eq:array81\endcsname{{66.24}{68}{}{equation.66.24}{Dirichlet-GRH-v930-Volume-XIII.pdf}}
\expandafter\def\csname r@eq:avg81\endcsname{{66.25}{69}{}{equation.66.25}{Dirichlet-GRH-v930-Volume-XIII.pdf}}
\expandafter\def\csname r@eq:common29\endcsname{{121.3}{45}{}{equation.121.3}{Dirichlet-GRH-v930-Volume-XXIII.pdf}}
\expandafter\def\csname r@eq:factor83\endcsname{{66.26}{70}{}{equation.66.26}{Dirichlet-GRH-v930-Volume-XIII.pdf}}
\expandafter\def\csname r@eq:gap\endcsname{{6.3}{26}{}{equation.6.3}{Dirichlet-GRH-v930-Volume-I.pdf}}
\expandafter\def\csname r@eq:geom163\endcsname{{86.13}{18}{}{equation.86.13}{Dirichlet-GRH-v930-Volume-XVII.pdf}}
\expandafter\def\csname r@eq:iter\endcsname{{70.6}{8}{}{equation.70.6}{Dirichlet-GRH-v930-Volume-XIV.pdf}}
\expandafter\def\csname r@eq:letters111\endcsname{{9.1}{117}{}{equation.9.1}{Dirichlet-GRH-v930-Volume-I.pdf}}
\expandafter\def\csname r@eq:level85\endcsname{{66.28}{73}{}{equation.66.28}{Dirichlet-GRH-v930-Volume-XIII.pdf}}
\expandafter\def\csname r@eq:mass205\endcsname{{98.4}{9}{}{equation.98.4}{Dirichlet-GRH-v930-Volume-XIX.pdf}}
\expandafter\def\csname r@eq:norm83\endcsname{{66.27}{71}{}{equation.66.27}{Dirichlet-GRH-v930-Volume-XIII.pdf}}
\expandafter\def\csname r@eq:op81\endcsname{{66.23}{67}{}{equation.66.23}{Dirichlet-GRH-v930-Volume-XIII.pdf}}
\expandafter\def\csname r@eq:quot207\endcsname{{98.7}{14}{}{equation.98.7}{Dirichlet-GRH-v930-Volume-XIX.pdf}}
\expandafter\def\csname r@eq:rd207\endcsname{{98.13}{18}{}{equation.98.13}{Dirichlet-GRH-v930-Volume-XIX.pdf}}
\expandafter\def\csname r@eq:region198\endcsname{{20.12}{33}{}{equation.20.12}{Dirichlet-GRH-v930-Volume-IV.pdf}}
\expandafter\def\csname r@eq:regular207\endcsname{{98.12}{18}{}{equation.98.12}{Dirichlet-GRH-v930-Volume-XIX.pdf}}
\expandafter\def\csname r@eq:scale207\endcsname{{98.10}{17}{}{equation.98.10}{Dirichlet-GRH-v930-Volume-XIX.pdf}}
\expandafter\def\csname r@eq:src109\endcsname{{87.4}{37}{}{equation.87.4}{Dirichlet-GRH-v930-Volume-XVII.pdf}}
\expandafter\def\csname r@eq:sub81\endcsname{{66.22}{67}{}{equation.66.22}{Dirichlet-GRH-v930-Volume-XIII.pdf}}
\expandafter\def\csname r@eq:tail114\endcsname{{19.15}{17}{}{equation.19.15}{Dirichlet-GRH-v930-Volume-IV.pdf}}
\expandafter\def\csname r@eq:window\endcsname{{{W}}{2}{}{AMS.2}{Dirichlet-GRH-v930-Volume-I.pdf}}
\expandafter\def\csname r@lem:angle149\endcsname{{86.10}{10}{}{proposition.86.10}{Dirichlet-GRH-v930-Volume-XVII.pdf}}
\expandafter\def\csname r@lem:basis63\endcsname{{65.14}{34}{}{proposition.65.14}{Dirichlet-GRH-v930-Volume-XIII.pdf}}
\expandafter\def\csname r@lem:cas96\endcsname{{121.25}{55}{}{proposition.121.25}{Dirichlet-GRH-v930-Volume-XXIII.pdf}}
\expandafter\def\csname r@lem:deep163\endcsname{{86.23}{19}{}{proposition.86.23}{Dirichlet-GRH-v930-Volume-XVII.pdf}}
\expandafter\def\csname r@lem:degree181\endcsname{{104.6}{6}{}{proposition.104.6}{Dirichlet-GRH-v930-Volume-XX.pdf}}
\expandafter\def\csname r@lem:delay197\endcsname{{19.12}{10}{}{proposition.19.12}{Dirichlet-GRH-v930-Volume-IV.pdf}}
\expandafter\def\csname r@lem:flag215\endcsname{{98.30}{29}{}{proposition.98.30}{Dirichlet-GRH-v930-Volume-XIX.pdf}}
\expandafter\def\csname r@lem:forms203\endcsname{{110.10}{46}{}{proposition.110.10}{Dirichlet-GRH-v930-Volume-XXI.pdf}}
\expandafter\def\csname r@lem:free32\endcsname{{121.16}{47}{}{proposition.121.16}{Dirichlet-GRH-v930-Volume-XXIII.pdf}}
\expandafter\def\csname r@lem:mass32\endcsname{{121.17}{48}{}{proposition.121.17}{Dirichlet-GRH-v930-Volume-XXIII.pdf}}
\expandafter\def\csname r@lem:move80\endcsname{{7.8}{94}{}{proposition.7.8}{Dirichlet-GRH-v930-Volume-I.pdf}}
\expandafter\def\csname r@lem:orbit110\endcsname{{86.30}{24}{}{proposition.86.30}{Dirichlet-GRH-v930-Volume-XVII.pdf}}
\expandafter\def\csname r@lem:ray88\endcsname{{7.10}{96}{}{proposition.7.10}{Dirichlet-GRH-v930-Volume-I.pdf}}
\expandafter\def\csname r@lem:realflag265\endcsname{{105.1}{32}{}{proposition.105.1}{Dirichlet-GRH-v930-Volume-XX.pdf}}
\expandafter\def\csname r@lem:root94\endcsname{{121.20}{51}{}{proposition.121.20}{Dirichlet-GRH-v930-Volume-XXIII.pdf}}
\expandafter\def\csname r@lem:shuffle107\endcsname{{9.10}{115}{}{proposition.9.10}{Dirichlet-GRH-v930-Volume-I.pdf}}
\expandafter\def\csname r@lem:size97\endcsname{{87.1}{33}{}{proposition.87.1}{Dirichlet-GRH-v930-Volume-XVII.pdf}}
\expandafter\def\csname r@lem:sphericalcore252\endcsname{{75.1}{112}{}{proposition.75.1}{Dirichlet-GRH-v930-Volume-XIV.pdf}}
\expandafter\def\csname r@lem:split29\endcsname{{121.13}{44}{}{proposition.121.13}{Dirichlet-GRH-v930-Volume-XXIII.pdf}}
\expandafter\def\csname r@lem:topchar266\endcsname{{88.10}{65}{}{proposition.88.10}{Dirichlet-GRH-v930-Volume-XVII.pdf}}
\expandafter\def\csname r@nav:I-28\endcsname{{6.6}{32}{}{subsection.6.6}{Dirichlet-GRH-v930-Volume-I.pdf}}
\expandafter\def\csname r@nav:I-29\endcsname{{6.6.1}{56}{}{subsubsection.6.6.1}{Dirichlet-GRH-v930-Volume-I.pdf}}
\expandafter\def\csname r@nav:III-13\endcsname{{25.2}{26}{}{subsection.25.2}{Dirichlet-GRH-v930-Volume-V.pdf}}
\expandafter\def\csname r@nav:V-18\endcsname{{66.2}{49}{}{subsection.66.2}{Dirichlet-GRH-v930-Volume-XIII.pdf}}
\expandafter\def\csname r@nav:VI-08\endcsname{{121.16}{51}{}{subsection.121.16}{Dirichlet-GRH-v930-Volume-XXIII.pdf}}
\expandafter\def\csname r@prop:average163\endcsname{{86.25}{21}{}{proposition.86.25}{Dirichlet-GRH-v930-Volume-XVII.pdf}}
\expandafter\def\csname r@prop:balance\endcsname{{6.3}{26}{}{proposition.6.3}{Dirichlet-GRH-v930-Volume-I.pdf}}
\expandafter\def\csname r@prop:cancelsquare256\endcsname{{74.9}{111}{}{proposition.74.9}{Dirichlet-GRH-v930-Volume-XIV.pdf}}
\expandafter\def\csname r@prop:center78\endcsname{{8.5}{104}{}{proposition.8.5}{Dirichlet-GRH-v930-Volume-I.pdf}}
\expandafter\def\csname r@prop:cofactor286\endcsname{{105.12}{43}{}{proposition.105.12}{Dirichlet-GRH-v930-Volume-XX.pdf}}
\expandafter\def\csname r@prop:column110\endcsname{{86.29}{23}{}{proposition.86.29}{Dirichlet-GRH-v930-Volume-XVII.pdf}}
\expandafter\def\csname r@prop:corecusp251\endcsname{{71.8}{22}{}{proposition.71.8}{Dirichlet-GRH-v930-Volume-XIV.pdf}}
\expandafter\def\csname r@prop:corehecke272\endcsname{{88.6}{60}{}{proposition.88.6}{Dirichlet-GRH-v930-Volume-XVII.pdf}}
\expandafter\def\csname r@prop:cost177\endcsname{{72.23}{65}{}{proposition.72.23}{Dirichlet-GRH-v930-Volume-XIV.pdf}}
\expandafter\def\csname r@prop:denom235\endcsname{{110.8}{44}{}{proposition.110.8}{Dirichlet-GRH-v930-Volume-XXI.pdf}}
\expandafter\def\csname r@prop:dilation179\endcsname{{104.3}{4}{}{proposition.104.3}{Dirichlet-GRH-v930-Volume-XX.pdf}}
\expandafter\def\csname r@prop:fiber274\endcsname{{90.3}{85}{}{proposition.90.3}{Dirichlet-GRH-v930-Volume-XVII.pdf}}
\expandafter\def\csname r@prop:field12\endcsname{{6.4}{27}{}{proposition.6.4}{Dirichlet-GRH-v930-Volume-I.pdf}}
\expandafter\def\csname r@prop:fourier232\endcsname{{110.19}{54}{}{proposition.110.19}{Dirichlet-GRH-v930-Volume-XXI.pdf}}
\expandafter\def\csname r@prop:good205\endcsname{{98.7}{8}{}{proposition.98.7}{Dirichlet-GRH-v930-Volume-XIX.pdf}}
\expandafter\def\csname r@prop:index159\endcsname{{86.17}{14}{}{proposition.86.17}{Dirichlet-GRH-v930-Volume-XVII.pdf}}
\expandafter\def\csname r@prop:lattices277\endcsname{{90.5}{87}{}{proposition.90.5}{Dirichlet-GRH-v930-Volume-XVII.pdf}}
\expandafter\def\csname r@prop:linkedline265\endcsname{{74.4}{104}{}{proposition.74.4}{Dirichlet-GRH-v930-Volume-XIV.pdf}}
\expandafter\def\csname r@prop:op81\endcsname{{66.28}{67}{}{proposition.66.28}{Dirichlet-GRH-v930-Volume-XIII.pdf}}
\expandafter\def\csname r@prop:others77\endcsname{{66.27}{66}{}{proposition.66.27}{Dirichlet-GRH-v930-Volume-XIII.pdf}}
\expandafter\def\csname r@prop:parity262\endcsname{{71.18}{35}{}{proposition.71.18}{Dirichlet-GRH-v930-Volume-XIV.pdf}}
\expandafter\def\csname r@prop:positions264\endcsname{{121.4}{33}{}{proposition.121.4}{Dirichlet-GRH-v930-Volume-XXIII.pdf}}
\expandafter\def\csname r@prop:rank11\endcsname{{6.9}{30}{}{proposition.6.9}{Dirichlet-GRH-v930-Volume-I.pdf}}
\expandafter\def\csname r@prop:ray50\endcsname{{6.30}{63}{}{proposition.6.30}{Dirichlet-GRH-v930-Volume-I.pdf}}
\expandafter\def\csname r@prop:res80\endcsname{{7.7}{94}{}{proposition.7.7}{Dirichlet-GRH-v930-Volume-I.pdf}}
\expandafter\def\csname r@prop:rowfinite262\endcsname{{88.15}{70}{}{proposition.88.15}{Dirichlet-GRH-v930-Volume-XVII.pdf}}
\expandafter\def\csname r@prop:shift203\endcsname{{110.11}{47}{}{proposition.110.11}{Dirichlet-GRH-v930-Volume-XXI.pdf}}
\expandafter\def\csname r@prop:tracelat253\endcsname{{122.7}{63}{}{proposition.122.7}{Dirichlet-GRH-v930-Volume-XXIII.pdf}}
\expandafter\def\csname r@prop:transfer286\endcsname{{89.6}{80}{}{proposition.89.6}{Dirichlet-GRH-v930-Volume-XVII.pdf}}
\expandafter\def\csname r@prop:translate109\endcsname{{87.9}{38}{}{proposition.87.9}{Dirichlet-GRH-v930-Volume-XVII.pdf}}
\expandafter\def\csname r@sec:absolute163\endcsname{{86.7}{15}{}{subsection.86.7}{Dirichlet-GRH-v930-Volume-XVII.pdf}}
\expandafter\def\csname r@sec:bessel29\endcsname{{121.13}{45}{}{subsection.121.13}{Dirichlet-GRH-v930-Volume-XXIII.pdf}}
\expandafter\def\csname r@sec:chain81\endcsname{{66.10}{67}{}{subsection.66.10}{Dirichlet-GRH-v930-Volume-XIII.pdf}}
\expandafter\def\csname r@sec:char94\endcsname{{121.16}{51}{}{subsection.121.16}{Dirichlet-GRH-v930-Volume-XXIII.pdf}}
\expandafter\def\csname r@sec:depth163\endcsname{{86.8}{19}{}{subsection.86.8}{Dirichlet-GRH-v930-Volume-XVII.pdf}}
\expandafter\def\csname r@sec:fischer83\endcsname{{66.11}{69}{}{subsection.66.11}{Dirichlet-GRH-v930-Volume-XIII.pdf}}
\expandafter\def\csname r@sec:frames207\endcsname{{98.5}{16}{}{subsection.98.5}{Dirichlet-GRH-v930-Volume-XIX.pdf}}
\expandafter\def\csname r@sec:global110\endcsname{{86.10}{23}{}{subsection.86.10}{Dirichlet-GRH-v930-Volume-XVII.pdf}}
\expandafter\def\csname r@sec:grid85\endcsname{{66.12}{72}{}{subsection.66.12}{Dirichlet-GRH-v930-Volume-XIII.pdf}}
\expandafter\def\csname r@sec:middle131\endcsname{{86.1}{3}{}{subsection.86.1}{Dirichlet-GRH-v930-Volume-XVII.pdf}}
\expandafter\def\csname r@sec:missing\endcsname{{6.6}{32}{}{subsection.6.6}{Dirichlet-GRH-v930-Volume-I.pdf}}
\expandafter\def\csname r@seriesguide\endcsname{{1.1}{4}{}{subsection.1.1}{Dirichlet-GRH-v930-Volume-I.pdf}}
\expandafter\def\csname r@thm:bestransfer265\endcsname{{88.16}{72}{}{proposition.88.16}{Dirichlet-GRH-v930-Volume-XVII.pdf}}
\expandafter\def\csname r@thm:bottom271\endcsname{{89.8}{82}{}{proposition.89.8}{Dirichlet-GRH-v930-Volume-XVII.pdf}}
\expandafter\def\csname r@thm:center\endcsname{{70.7}{9}{}{proposition.70.7}{Dirichlet-GRH-v930-Volume-XIV.pdf}}
\expandafter\def\csname r@thm:common111\endcsname{{9.13}{118}{}{proposition.9.13}{Dirichlet-GRH-v930-Volume-I.pdf}}
\expandafter\def\csname r@thm:common133\endcsname{{98.2}{4}{}{proposition.98.2}{Dirichlet-GRH-v930-Volume-XIX.pdf}}
\expandafter\def\csname r@thm:commonnorm128\endcsname{{100.16}{78}{}{proposition.100.16}{Dirichlet-GRH-v930-Volume-XIX.pdf}}
\expandafter\def\csname r@thm:compose81\endcsname{{66.29}{68}{}{proposition.66.29}{Dirichlet-GRH-v930-Volume-XIII.pdf}}
\expandafter\def\csname r@thm:conductor232\endcsname{{19.15}{12}{}{proposition.19.15}{Dirichlet-GRH-v930-Volume-IV.pdf}}
\expandafter\def\csname r@thm:corecopy251\endcsname{{71.1}{13}{}{proposition.71.1}{Dirichlet-GRH-v930-Volume-XIV.pdf}}
\expandafter\def\csname r@thm:coreeuler252\endcsname{{71.2}{14}{}{proposition.71.2}{Dirichlet-GRH-v930-Volume-XIV.pdf}}
\expandafter\def\csname r@thm:cross132\endcsname{{86.3}{4}{}{proposition.86.3}{Dirichlet-GRH-v930-Volume-XVII.pdf}}
\expandafter\def\csname r@thm:deep163\endcsname{{86.24}{20}{}{proposition.86.24}{Dirichlet-GRH-v930-Volume-XVII.pdf}}
\expandafter\def\csname r@thm:direction249\endcsname{{71.6}{19}{}{proposition.71.6}{Dirichlet-GRH-v930-Volume-XIV.pdf}}
\expandafter\def\csname r@thm:downcore254\endcsname{{71.13}{28}{}{proposition.71.13}{Dirichlet-GRH-v930-Volume-XIV.pdf}}
\expandafter\def\csname r@thm:edge242\endcsname{{67.9}{82}{}{proposition.67.9}{Dirichlet-GRH-v930-Volume-XIII.pdf}}
\expandafter\def\csname r@thm:factor83\endcsname{{66.30}{70}{}{proposition.66.30}{Dirichlet-GRH-v930-Volume-XIII.pdf}}
\expandafter\def\csname r@thm:first250\endcsname{{71.10}{25}{}{proposition.71.10}{Dirichlet-GRH-v930-Volume-XIV.pdf}}
\expandafter\def\csname r@thm:firstsquare255\endcsname{{71.3}{16}{}{proposition.71.3}{Dirichlet-GRH-v930-Volume-XIV.pdf}}
\expandafter\def\csname r@thm:frame55\endcsname{{65.1}{21}{}{proposition.65.1}{Dirichlet-GRH-v930-Volume-XIII.pdf}}
\expandafter\def\csname r@thm:initial247\endcsname{{89.2}{74}{}{proposition.89.2}{Dirichlet-GRH-v930-Volume-XVII.pdf}}
\expandafter\def\csname r@thm:invar148\endcsname{{20.11}{31}{}{proposition.20.11}{Dirichlet-GRH-v930-Volume-IV.pdf}}
\expandafter\def\csname r@thm:jactransfer266\endcsname{{88.14}{69}{}{proposition.88.14}{Dirichlet-GRH-v930-Volume-XVII.pdf}}
\expandafter\def\csname r@thm:lattice32\endcsname{{121.18}{49}{}{proposition.121.18}{Dirichlet-GRH-v930-Volume-XXIII.pdf}}
\expandafter\def\csname r@thm:levels88\endcsname{{7.11}{97}{}{proposition.7.11}{Dirichlet-GRH-v930-Volume-I.pdf}}
\expandafter\def\csname r@thm:localmass156\endcsname{{87.24}{48}{}{proposition.87.24}{Dirichlet-GRH-v930-Volume-XVII.pdf}}
\expandafter\def\csname r@thm:markedcore253\endcsname{{102.7}{95}{}{proposition.102.7}{Dirichlet-GRH-v930-Volume-XIX.pdf}}
\expandafter\def\csname r@thm:markedranks623\endcsname{{117.46}{85}{}{proposition.117.46}{Dirichlet-GRH-v930-Volume-XXII.pdf}}
\expandafter\def\csname r@thm:mean56\endcsname{{65.22}{40}{}{proposition.65.22}{Dirichlet-GRH-v930-Volume-XIII.pdf}}
\expandafter\def\csname r@thm:midlevel131\endcsname{{86.1}{3}{}{proposition.86.1}{Dirichlet-GRH-v930-Volume-XVII.pdf}}
\expandafter\def\csname r@thm:midnorm131\endcsname{{86.2}{4}{}{proposition.86.2}{Dirichlet-GRH-v930-Volume-XVII.pdf}}
\expandafter\def\csname r@thm:mix277\endcsname{{90.8}{91}{}{proposition.90.8}{Dirichlet-GRH-v930-Volume-XVII.pdf}}
\expandafter\def\csname r@thm:neighborindex356\endcsname{{90.15}{96}{}{proposition.90.15}{Dirichlet-GRH-v930-Volume-XVII.pdf}}
\expandafter\def\csname r@thm:norm211\endcsname{{98.21}{21}{}{proposition.98.21}{Dirichlet-GRH-v930-Volume-XIX.pdf}}
\expandafter\def\csname r@thm:norm212\endcsname{{98.41}{36}{}{proposition.98.41}{Dirichlet-GRH-v930-Volume-XIX.pdf}}
\expandafter\def\csname r@thm:norm58\endcsname{{65.25}{43}{}{proposition.65.25}{Dirichlet-GRH-v930-Volume-XIII.pdf}}
\expandafter\def\csname r@thm:ordinary163\endcsname{{86.22}{18}{}{proposition.86.22}{Dirichlet-GRH-v930-Volume-XVII.pdf}}
\expandafter\def\csname r@thm:orientation46\endcsname{{6.23}{57}{}{proposition.6.23}{Dirichlet-GRH-v930-Volume-I.pdf}}
\expandafter\def\csname r@thm:pol49\endcsname{{64.10}{12}{}{proposition.64.10}{Dirichlet-GRH-v930-Volume-XIII.pdf}}
\expandafter\def\csname r@thm:proj\endcsname{{70.1}{3}{}{proposition.70.1}{Dirichlet-GRH-v930-Volume-XIV.pdf}}
\expandafter\def\csname r@thm:quot207\endcsname{{98.13}{14}{}{proposition.98.13}{Dirichlet-GRH-v930-Volume-XIX.pdf}}
\expandafter\def\csname r@thm:quot208\endcsname{{98.37}{34}{}{proposition.98.37}{Dirichlet-GRH-v930-Volume-XIX.pdf}}
\expandafter\def\csname r@thm:quot80\endcsname{{7.9}{95}{}{proposition.7.9}{Dirichlet-GRH-v930-Volume-I.pdf}}
\expandafter\def\csname r@thm:rangenew50\endcsname{{6.29}{62}{}{proposition.6.29}{Dirichlet-GRH-v930-Volume-I.pdf}}
\expandafter\def\csname r@thm:ratio248\endcsname{{89.3}{76}{}{proposition.89.3}{Dirichlet-GRH-v930-Volume-XVII.pdf}}
\expandafter\def\csname r@thm:raw97\endcsname{{87.2}{33}{}{proposition.87.2}{Dirichlet-GRH-v930-Volume-XVII.pdf}}
\expandafter\def\csname r@thm:rd207\endcsname{{98.18}{18}{}{proposition.98.18}{Dirichlet-GRH-v930-Volume-XIX.pdf}}
\expandafter\def\csname r@thm:rd208\endcsname{{98.39}{35}{}{proposition.98.39}{Dirichlet-GRH-v930-Volume-XIX.pdf}}
\expandafter\def\csname r@thm:real108\endcsname{{87.5}{36}{}{proposition.87.5}{Dirichlet-GRH-v930-Volume-XVII.pdf}}
\expandafter\def\csname r@thm:realcore253\endcsname{{71.16}{32}{}{proposition.71.16}{Dirichlet-GRH-v930-Volume-XIV.pdf}}
\expandafter\def\csname r@thm:realrow257\endcsname{{74.1}{100}{}{proposition.74.1}{Dirichlet-GRH-v930-Volume-XIV.pdf}}
\expandafter\def\csname r@thm:region198\endcsname{{20.13}{33}{}{proposition.20.13}{Dirichlet-GRH-v930-Volume-IV.pdf}}
\expandafter\def\csname r@thm:rest96\endcsname{{121.24}{54}{}{proposition.121.24}{Dirichlet-GRH-v930-Volume-XXIII.pdf}}
\expandafter\def\csname r@thm:scale207\endcsname{{98.16}{17}{}{proposition.98.16}{Dirichlet-GRH-v930-Volume-XIX.pdf}}
\expandafter\def\csname r@thm:scale232\endcsname{{110.18}{53}{}{proposition.110.18}{Dirichlet-GRH-v930-Volume-XXI.pdf}}
\expandafter\def\csname r@thm:secondrow256\endcsname{{74.5}{105}{}{proposition.74.5}{Dirichlet-GRH-v930-Volume-XIV.pdf}}
\expandafter\def\csname r@thm:section110\endcsname{{86.31}{24}{}{proposition.86.31}{Dirichlet-GRH-v930-Volume-XVII.pdf}}
\expandafter\def\csname r@thm:signal77\endcsname{{66.26}{65}{}{proposition.66.26}{Dirichlet-GRH-v930-Volume-XIII.pdf}}
\expandafter\def\csname r@thm:src109\endcsname{{87.7}{37}{}{proposition.87.7}{Dirichlet-GRH-v930-Volume-XVII.pdf}}
\expandafter\def\csname r@thm:suffix237\endcsname{{110.9}{45}{}{proposition.110.9}{Dirichlet-GRH-v930-Volume-XXI.pdf}}
\expandafter\def\csname r@thm:transverse250\endcsname{{106.2}{53}{}{proposition.106.2}{Dirichlet-GRH-v930-Volume-XX.pdf}}
\expandafter\def\csname r@thm:twists\endcsname{{5.1}{19}{}{proposition.5.1}{Dirichlet-GRH-v930-Volume-I.pdf}}
\expandafter\def\csname r@thm:twofreq160\endcsname{{87.29}{51}{}{proposition.87.29}{Dirichlet-GRH-v930-Volume-XVII.pdf}}
\expandafter\def\csname r@thm:types276\endcsname{{90.7}{89}{}{proposition.90.7}{Dirichlet-GRH-v930-Volume-XVII.pdf}}
\expandafter\def\csname r@thm:upcore254\endcsname{{71.9}{23}{}{proposition.71.9}{Dirichlet-GRH-v930-Volume-XIV.pdf}}
\expandafter\def\csname r@thm:word275\endcsname{{83.10}{59}{}{proposition.83.10}{Dirichlet-GRH-v930-Volume-XVI.pdf}}
\expandafter\def\csname r@thm:zero106\endcsname{{86.27}{22}{}{proposition.86.27}{Dirichlet-GRH-v930-Volume-XVII.pdf}}
\makeatother
\begin{document}
\setcounter{part}{17}
\fancyhead[L]{\small\color{muted}Dirichlet GRH \textperiodcentered\ Volume XVIII}
\fancyhead[R]{\small\color{muted}Terminal characters and measures}
\thispagestyle{plain}\begin{center}
{\LARGE\bfseries The generalized Riemann hypothesis\\[2pt]for Dirichlet \(L\)-functions}\\[6pt]
{\small 27 November 2026 \quad\textperiodcentered\quad Edition 930}\end{center}
\part{Terminal characters, measures and Whittaker normalizations}\label{part:XVIII}
\begin{abstract}\noindent We fix the marked characters, measures and integral highest lines at the terminal level. The compact projections and the Whittaker pairing determine the normalizations on residue flags and their restriction subspaces.\end{abstract}
\setcounter{section}{91}
\noindent The \hyperref[seriesguide]{guide in Volume~I} lists all twenty-seven volumes.
\begingroup\small\setlength{\parskip}{0pt}\tableofcontents\endgroup
\clearpage

\section{Weights, terminal lines and finite levels}\label{nav:g-VIII-0}\label{nav:volume-XVIII}

\subsection{The prescribed ranks and the final Bessel subgroup}\label{nav:V-06}
\label{sec:ranks47}

At each \(x\) on the intervals of Volume~I take \(E=E_x\), \(F_+=E^+\),
and \(n=[E:\QQ]/4\). Thus \([F_+:\QQ]=2n\).
The ranks in \srcloc{Formalism/LocalRingCardinality/Params}{6} satisfy
\(m,k\) even, \(t\) odd, \(m+3\le k,\ 3m<t,\ k+m+4\le t\),
\(m\le 2n\le k\), and \(m^2\le k+m(t+1)\). We write \(t_{\rm r}\)
for this odd rank parameter to distinguish it from a Galois translate.
The two towers have different ranks, specified in
\srcloc{Formalism/LocalRingCardinality/DataTheta}{627}. Their signatures
and embeddings are set out in
\srcloc{Development/AbelianVarietyTheory/Definitions/AutomorphicMembership/AdelicDataRowMeasureScaleTrivialTuple}{260}.

\begin{proposition}[Simultaneous ranks and signatures]\label{prop:ranks47}
Put \(m=2\lfloor\sqrt{x}/2\rfloor,\ k=2n,\ r=k+m-1\),
\(t_{\rm r}=k+m+5,\ d_0=t_{\rm r}+1\). For all sufficiently large \(x\)
these integers obey the rank conditions above and the size inequalities
of Proposition \ref{prop:balance}, with the fields and translates of
Proposition~\ref{prop:field12}. At step \(0\le j\le m\) there are Hermitian spaces
\(X_j,Y_j\) over \(E/F_+\) with dimensions and numbers of negative
directions at each real place given by
\[
  N_j=k+jd_0,\qquad M_j=r+jt_{\rm r},\qquad
  a(X_j)=0,\quad a(Y_j)=\lfloor(m-j)/2\rfloor .
\]
There are isometric inclusions \(X_j\subset Y_j\) for \(j<m\) and
\(Y_m\subset X_m\) of codimension one at the endpoint.
\end{proposition}
\begin{proof}
We have \(n\asymp_\psi x^{2+8\eps}\), \(m\sim\sqrt{x}\) and \(t_{\rm r}\)
odd. Thus \(k\ge m+3,\ t_{\rm r}>3m\), \(m\le2n=k\), and \(m^2\le k\)
for large \(x\). The stronger inequality \(k+m+4\le t_{\rm r}\)
holds by construction. The bounds \(xm\le k\le n^2x^{1/2-\eps}\)
follow directly from the inequalities for \(n,m\). With \(h\ge n/2\)
and \(a_E=xb\), all the other size inequalities apply whenever the
chosen distinguished set has density at least \(x^{-\eps}\).
Also \(4t_{\rm r}r+1<\min(q,d)\) on sufficiently distant common
intervals. Indeed \(x\) there is bounded by a fixed power of the
initial selection scale, whereas the lower bounds for \(q,d\) in
Theorem \ref{thm:twists}, grow faster than every fixed power of that scale.

Take diagonal forms with the indicated numbers of coefficients \(1,-1\).
For \(j<m\) the complement of \(X_j\) has dimension \(m-j-1\).
If \(m-j=2v\), its signature is \((v-1,v)\), and if \(m-j=2v+1\)
it is \((v,v)\). These are nonnegative pairs, so \(X_j\) embeds as
the first \(N_j\) positive directions. At \(m\) the identity
\(M_m=N_m-1\) gives the asserted positive coordinate hyperplane.
For the definite chain of dual pairs one can multiply alternate
Hermitian forms by the same nonzero purely imaginary scalar to obtain
skew-Hermitian spaces. This operation preserves their unitary groups.
\end{proof}

Let \(P,Q_0\ge5\) be the two fixed distinct auxiliary rational primes
of Lemma~\ref{lem:split29}. At \(v\mid P\) write \(q_v=|\OO_{F_+}/v|\),
\(R=\prod_{v\mid P}(\OO_{F_+}/v)^\times\), and \(U=\widehat R\).
Use the definite chain \(U(X_0),\ldots,U(X_m)\), the splitting of
\eqref{eq:action29} and one fixed initial function with one central
character and full integral invariance at \(P\). Take the common
vector \(\phi^0_{N_j,N_{j+1}}\) of \eqref{eq:common29} on every edge.
The other Schwartz functions are fixed across all twists. Project the
restriction to \(H=U(Y_m)\) to a fixed right-Iwahori-invariant test
space, as in Section~\ref{sec:bessel29}.

\begin{theorem}[Residue characters at the prescribed endpoint]\label{thm:block47}
There is a choice of \(|R|^{m-2}\) tuples of global determinant twists
with one common output central character. Their restrictions to \(H\)
have distinct right residue characters. In particular their projected
columns are pairwise orthogonal for every such common test space.
Put \(\delta_i=\chi_i^{(-1)^i}|_R\in U\) for \(1\le i<m\) and
\(v_i=\delta_i\cdots\delta_{m-1}\).
In the \(N_m-1\) endpoint coordinates the character is constant on
successive blocks. The first block has length \(N_1\) and value \(v_1\).
Next come the values \(v_2,\ldots,v_{m-1}\), each on \(d_0\)
coordinates, and the final block has length \(d_0-1\) and value one.
\end{theorem}
\begin{proof}
At a place \(v\mid P\), expand each common vector by the exact finite
Fourier formula \eqref{eq:common29}. In the first edge all \(k\)
residue labels are trivial in every surviving term, since integration
against the initial function is unchanged by independent diagonal units.
At vertex \(i\), the incoming edge contributes the inverse of its
label on the first \(N_{i-1}\) coordinates and one on the remaining
\(d_0\). The internal determinant twist contributes \(\delta_i\) in
every coordinate. Invariance of integration under the Iwahori forces
the outgoing labels to be the incoming labels times \(\delta_i^{-1}\)
in the first \(N_{i-1}\) coordinates and \(\delta_i^{-1}\) in the new
ones. Failure at any coordinate makes that term vanish by translating
by a diagonal unit on which its character differs from one.
This argument works independently at every \(v\mid P\).
On the final output one inverts the last edge labels and appends \(d_0\)
trivial coordinates. Restricting to \(Y_m\) removes the very last
coordinate. Iteration gives the blocks claimed, whose consecutive ratios
recover all the \(\delta_i\), with \(v_m=1\).

To choose the global twists, write \(g=\gcd(N_1,\ldots,N_{m-1})\)
and let \(w_i=N_i/g\). The vector \(w\) is primitive in \(\ZZ^{m-1}\).
Elementary integer column operations take it to \((0,\ldots,0,1)\):
the Euclidean algorithm replaces each successive pair of nonzero
entries by its greatest common divisor and zero. The same operations
act as automorphisms on tuples in any abelian character group, using
integer powers and multiplication. In particular the kernel of
\(\delta\mapsto\prod\delta_i^{w_i}\) on \(U^{m-1}\) is a copy of
\(U^{m-2}\). Each tuple in this kernel lifts to a tuple of global
finite-order characters with the primitive weighted product one.
Indeed, in the transformed coordinates lift the first \(m-2\)
coordinates by Lemma~\ref{lem:split29}, take the last coordinate
trivial, and undo the integer operations. Taking \((-1)^i\)-th powers
to return to the \(\chi_i\) gives
\(\prod\chi_i^{(-1)^iN_i}=1\). The central-character formula of
Theorem~\ref{thm:center} applies at these ranks. Distinct endpoint
characters lie in orthogonal subspaces, as Haar measure is translation
invariant. The common orthogonal projection commutes with the Iwahori
action and preserves this orthogonality.
\end{proof}

\begin{proposition}[Local norms and denominators at these ranks]\label{prop:integral47}
Let \(A_P=R^{N_m-1}\), let \(\Xi\subset\widehat A_P\) be the family
of endpoint characters in Theorem~\ref{thm:block47}, and
\(Z=\bigcap_{\gamma\in\Xi}\ker\gamma\). Put
\[
 S_1=mk+d_0m(m-1)/2,\qquad
 S_2=\sum_{j=0}^{m-1}\bigl(N_jN_{j+1}-N_j(N_j-1)/2\bigr).
\]
The product of the common local vectors has squared norm
\(P^{-[F_+:\QQ]S_2}\) with the integral matrix lattices of volume one.
Its full Fourier expansion has common scalar denominator
\(|R|^{S_1}\). A single residue character projector is integral after
multiplication by \(|R|^{N_m-1}\). The projector onto functions on
one level orbit constant on the cosets of \(Z\) has exact denominator
\[
                         |Z|=|R|^{N_m-m+1}=|R|^{k+mt_{\rm r}+1}
\]
on integral scalar functions on that orbit. At the ranks of
Proposition~\ref{prop:ranks47}, the logarithms of these denominators
and the negative logarithm of the squared norm are \(o(b/x)\).
\end{proposition}
\begin{proof}
For an \(N_{j+1}\)-by-\(N_j\) residue matrix in the support of
\(\phi^0\), exactly \(N_j(N_j-1)/2\) entries are free. All the others
are prescribed, including the diagonal ones. Multiply the resulting
volumes at all edges and places and use \(\prod_{v\mid P}q_v=P^{[F_+:\QQ]}\).
Formula \eqref{eq:common29} contributes
\(\prod_{v\mid P}(q_v-1)^{-N_j}\) at edge \(j\). Summing the \(N_j\)
gives \(S_1\). The output projectors are averages over \(A_P\) and
\(Z\), respectively. The character family is the injective image of
the primitive kernel in the preceding proof, hence is a subgroup of
order \(|R|^{m-2}\). Finite character duality gives
\([A_P:Z]=|R|^{m-2}\). The principal level at \(Q_0\) makes the finite
level action free at any rank by the positive-lattice stabilizer
argument of Lemma~\ref{lem:free32}. On one free orbit the \(Z\)-average
of a point mass is \(1/|Z|\) on its coset. This proves exactness of
that denominator. Finally \(\log|R|\le 2n\log P\),
\(N_j=O((j+1)n)\), \(S_1=O(m^2n)\), and \(S_2=O(m^3n^2)\).
All logarithms in question are \(O_P(m^3n^3)\), a fixed power of \(x\)
for fixed \(\psi,\eps\), whereas \(b=\exp(x^\eps)\).
\end{proof}

For example, on a single free level orbit and one unit archimedean
coordinate the squared norm of the projected column is still
\(|p_\gamma(y)|^2/h\), with \(h\) the number of orbits at these new
ranks. This follows by summing its coordinates
\(p_\gamma(y)\gamma(a)\) with the Haar probability calculation of
Theorem~\ref{thm:lattice32}. No lower bound for this scalar follows
from character orthogonality. The representations of the row tower
\(U(Y_j)\), which can have indefinite signatures before the endpoint,
enter the prescribed Bessel period as well. In the supplied source
\srcloc{Formalism/SharpRestriction/Divisibility/FunctionModelRallisFromInputs}{534}
specifies nonvanishing for vectors in a given chain, and line 548
specifies its continuation into the next representation. The global
normalization and comparison of those row and column integrals remain
necessary in Section \ref{nav:I-28}.

\subsection{Integral highest weight projections in the definite chain}\label{nav:VI-05}
\label{sec:hooks86}

We compute the norm of the corner minors that occur in the
archimedean model and the value of their normalized differential
on the polynomial lattice of Section~\ref{sec:chain81}. The choice
of lattice enters the arithmetic assertion.

For a partition \(\lambda_1\ge\cdots\ge\lambda_r\ge0\), put
\(\lambda_{r+1}=0\). Write \(\Delta_i(X)\) for the leading
\(i\)-rowed minor and set
\[
 p_\lambda(X)=\prod_{i=1}^r\Delta_i(X)^{\lambda_i-\lambda_{i+1}},
 \qquad H(\lambda)=
 \frac{\prod_{i=1}^r(\lambda_i+r-i)!}
      {\prod_{1\le i<j\le r}(\lambda_i-\lambda_j+j-i)} .
\]
Use the empty product for \(r=0\). Extra rows and columns of \(X\)
are allowed. The Fischer inner product makes monomials orthogonal
with \(\|X^\alpha\|_{\rm F}^2=\alpha!\).

\begin{lemma}[The norm of the two corner factors]\label{lem:hook86}
One has \(\|p_\lambda\|_{\rm F}^2=H(\lambda)\).
In an \(M\)-by-\(n\) matrix, use the first \(p\) rows and first
columns for a partition \(\lambda^+\) of length at most \(p\).
For a partition \(\lambda^-\) of length at most \(q\), form its
\(i\)-rowed minor from rows \(p+1,\ldots,p+i\) and the last \(i\)
columns in their natural order. If \(p+q\le M\) and
\(\max(p,q)\le n\), the squared norm of the product is
\(H(\lambda^+)H(\lambda^-)\). These squared norms multiply over
independent matrices at different embeddings.
\end{lemma}
\begin{proof}
It suffices first to use a square \(r\)-by-\(r\) matrix \(G\)
of independent standard complex Gaussians. By monomial integration
\(\|p_\lambda\|_{\rm F}^2=\mathbb E|p_\lambda(G)|^2\).
Write \(G=UT\) by complex Gram--Schmidt with positive diagonal
on \(T\). Left unitary invariance makes the conditional law of
\(U\) given \(T\) Haar on \(U(r)\). Indeed left multiplication
changes only \(U\), and the invariant probability measure is unique.
Every leading \(i\)-rowed minor is multiplied by
\(\prod_{j\le i}T_{jj}\) under right multiplication by \(T\).
As in the proof of Theorem~\ref{thm:factor83}, the moments of the
squared residual lengths give
\[
 \mathbb E |p_\lambda(G)|^2=
 \left(\int_{U(r)}|p_\lambda(U)|^2\,dU\right)
           \prod_{i=1}^r(r-i+1)_{\lambda_i}.
\]
To evaluate the integral take the unit tensor made of
\(\lambda_i-\lambda_{i+1}\) copies of
\(e_1\wedge\cdots\wedge e_i\) for each \(i\). The matrix
coefficient of this vector with itself is \(p_\lambda(U)\).
The weight \(\lambda\) is the highest weight in this tensor
product and has multiplicity one. Each factor weight has every
initial partial sum bounded by that of its highest weight.
Equality in all the sums for the tensor forces the highest
weight in each factor. Complete reducibility and the highest
weight theorem therefore place the vector in the irreducible
module of weight \(\lambda\). Its dimension is
\[
 d_r(\lambda)=\prod_{i<j}
               \frac{\lambda_i-\lambda_j+j-i}{j-i}
\]
by the Weyl dimension formula.~\cite[Thm.~5.22.1]{XVIII-XIII-IV-etingof86}
Averaging the orthogonal rank-one projection onto the tensor
gives a commuting operator on that module. Its trace is one, so
Schur's lemma makes it \(1/d_r(\lambda)\) times the identity.
Pair with the same vector to obtain the Haar integral. Substituting
the dimension formula and cancelling
\(\prod_i(r-i)!=\prod_{i<j}(j-i)\) proves the first assertion.
Reversing the order of the last columns changes the second corner
polynomial into a leading-minor polynomial up to sign. The two
blocks have disjoint rows, hence use independent Gaussian entries
even when their sets of columns intersect. Independence also
applies across embeddings.
\end{proof}

\begin{theorem}[A normalized differential on the integral lattice]
\label{thm:projection86}
Let \(A\) be an integral domain of characteristic zero. Take an
\(s\)-by-\(r\) matrix \(X\), \(s\ge r\), and let \(J\) have
the identity in its first \(r\) rows and zeros below. Suppose
\(Q\in A[X]\) has column degrees \(\lambda_i\) and is invariant
under adding any multiple of an earlier column to a later one, as
a polynomial identity. Its component of row degrees
\((\lambda_1,\ldots,\lambda_r,0,\ldots,0)\) is
\(Q(J)p_\lambda(X)\). Consequently
\[
       \frac{p_\lambda(\partial)Q(0)}{H(\lambda)}=Q(J)\in A .
\]
The orthogonal projection onto the indicated row weight preserves
this lattice. The functional \(Q\mapsto Q(J)\) on its complex
span has Fischer operator norm \(H(\lambda)^{-1/2}\).
\end{theorem}
\begin{proof}
Keeping only the indicated row degrees commutes with the column
identities because a column addition preserves the row of every
variable. Call the component \(Q'\). Only the diagonal monomial
\(\prod_i X_{ii}^{\lambda_i}\) can contribute at \(J\), so
\(Q'(J)=Q(J)\). Since \(Q'\) uses only the top \(r\) rows we
may now use square matrices. On a lower triangular matrix every
surviving monomial with these row and column degrees is diagonal.
Indeed \(\sum_{ij}(i-j)\alpha_{ij}=0\) and every summand is
nonnegative there. In particular \(Q'(L)=Q(J)\) on lower
unitriangular matrices. On the dense open set where the leading
minors do not vanish, factor \(X=LT\) with \(L\) lower unitriangular
and \(T\) upper triangular. The column identities give
\(Q'(X)=Q(J)\prod_i T_{ii}^{\lambda_i}=Q(J)p_\lambda(X)\).
The identity holds over the fraction field by density and hence
holds in \(A[X]\). The polynomial \(p_\lambda\) has real integral
coefficients and only the indicated row degrees. Applying its
differential at zero pairs it with \(Q'\) in the Fischer form.
Lemma~\ref{lem:hook86} gives the displayed formula. Over \(\CC\)
Cauchy--Schwarz gives the norm bound, with equality on \(p_\lambda\).
\end{proof}

\begin{proposition}[The exact denominator on an ambient weight space]
\label{prop:lcm86}
On integral polynomials in a \(2\)-by-\(2\) matrix with all row
and column degrees \(b\ge0\), the least positive integer clearing
the values of \((\det\partial)^b/H((b,b))\) at zero is
\(\operatorname{lcm}(1,\ldots,b+1)\).
\end{proposition}
\begin{proof}
The monomial basis is \((X_{11}X_{22})^{b-a}(X_{12}X_{21})^a\)
for \(0\le a\le b\). On this monomial the value equals
\((-1)^a/((b+1)\binom ba)\), by the binomial theorem and
\(H((b,b))=b!(b+1)!\). Put \(n=b+1\). For a prime \(\ell\le n\)
let \(h=\lfloor\log_\ell n\rfloor\) and \(e=v_\ell(n)\).
The factorial valuation formula counts \(v_\ell\binom{n-1}{a}\)
as the number of borrows when subtracting \(a\) from \(n-1\)
in base \(\ell\). The lowest \(e\) digits of \(n-1\) are
\(\ell-1\), and no borrow can leave the highest position.
There are at most \(h-e\) borrows. At \(a=\ell^h-1\) the identity
\[
    n\binom{n-1}{\ell^h-1}=\ell^h\binom n{\ell^h}
\]
attains valuation \(h\), since the degree \(\ell^h\) coefficient
of \((1+z)^n\bmod\ell\) is the nonzero leading base-\(\ell\)
digit of \(n\). This follows by multiplying the digit factors
\((1+z^{\ell^j})^{n_j}\). No prime above \(n\) occurs. These
are exactly the prime powers in the claimed least common multiple.
\end{proof}

\begin{theorem}[The definite chain with intermediate projections]
\label{thm:projectchain86}
In \eqref{eq:array81} start with \(Q_0=1\). At every step keep
the highest row weight of the output at each embedding, using
Theorem~\ref{thm:projection86}. Use the same masses and terminal
identity frame. The simultaneous nonvanishing, stabilizer denominator
and finite-place bound of Theorem~\ref{thm:compose81} hold for the
projected values. The torsion bound \eqref{eq:level85} also holds.
The polynomial acting on the translated product section now has
squared Fischer norm
\begin{equation}
 \|P^{\rm pr}\|_{\rm F}^2=\prod_{j,\tau}H(v_{j,\tau})
   =\|P\|_{\rm F}^2\bigg/\prod_{j,\tau}d_{R_j}(v_{j,\tau}),
                                                        \label{eq:prnorm86}
\end{equation}
where \(v_{j,\tau}\) is the unpadded partition in
Section~\ref{sec:fischer83}. The field-norm upper estimate
\eqref{eq:norm83} holds with this polynomial and the projected value.
\end{theorem}
\begin{proof}
After projection the output at class \(h\) is a scalar times
\(\prod_\tau p_{\lambda_{j,\tau}}(Y_\tau)\), by
Theorem~\ref{thm:projection86} applied separately to each
matrix. The scalar is its value at \(I\). At the next step set
the new frame equal to \(I\) in \eqref{eq:sub81}. This replaces
the previous frame by the transpose of the leading square block
of the derivative matrix. Principal minors are unchanged under
transposition. Multiplication by the determinant power adds the
same integer to every part. It follows by induction that the
scalar is obtained by applying the single polynomial
\[
       P^{\rm pr}(D)=\prod_{j,\tau}
          p_{v_{j,\tau}}(D_{j,\tau,[1:R_j,\,1:R_j]})
\]
in independent derivative variables to the very same translated
product section \(\mathcal G_h\) as in the proof of
Theorem~\ref{thm:compose81}. This polynomial takes value one on
the rectangular inclusions. Its block degrees \(b_j\) agree with
those in that theorem. The positive vacuum overlap gives
\(\mathcal G_h\ne0\). The product-torus norm identity
\eqref{eq:avg81} and the sampling proof in Section~\ref{sec:grid85}
use only nonvanishing of the polynomial and these block degrees.
They give the asserted common torsion inputs, with algebraic
values for the corrected translations.

Projection keeps a subset of the monomials, so after division by
each step's previously established prime power the coefficients
are still integral at every selected place. Apply Proposition
~\ref{prop:op81} inductively, clearing the same masses.
Independent square blocks and Lemma~\ref{lem:hook86} give the
first equality in \eqref{eq:prnorm86}. Its second equality is
the one-block norm calculation divided into \eqref{eq:factor83}.
All coefficients of \(P^{\rm pr}\) are integers, and the section,
degrees and Gaussian envelope in the proof of \eqref{eq:norm83}
are identical. That proof therefore gives the stated field-norm
estimate.
\end{proof}

The two corners used for the indefinite archimedean norm appear in
\srcloc{Formalism/LocalToGlobal/IndefiniteReal/LocalNewformFunction}{293}.
The subtraction involving the indefinite row weights is specified in
\srcloc{Formalism/LocalRingCardinality/KernelPlaceMeanBound}{25}.
Theorem~\ref{thm:projectchain86} computes what the integral
projection supplies on the definite polynomials in this manuscript.

\subsection{Elementary divisors of the finite map at the terminal prime}\label{nav:VI-06}
\label{sec:fourier89}

Let \(T\) be any terminal prime furnished by
Theorem \ref{thm:levels88}.
There are \(f=[E^+:\QQ]\) split places of \(E^+\) above \(T\),
each with residue field \(\mathbb F_T\). At one place put
\(V=\mathbb F_T^k\). On functions on \(V^*\setminus\{0\}\)
use the operator
\begin{equation}
       (Dh)(x)=\sum_{\xi(x)=1}h(\xi),\qquad 0\ne x\in V .
                                                        \label{eq:inc89}
\end{equation}
The same definition is used in
\srcloc{Formalism/LogVandermondeBounds/FinrankFiberSumZero}{16}.
We compute all its elementary divisors over the \(T\)-adic integers
on the summand with zero sum on each scalar orbit. In particular
the computation measures both the best common divisor of its image
and the largest power needed to invert it.

Let \(\omega:\mathbb F_T^\times\longrightarrow\ZZ_T^\times\)
be the Teichm\"uller character, which lifts each residue to the
unique root of unity of order dividing \(T-1\) in its residue
class. For \(1\le j\le T-2\) write \(L_j(V^*)\) for functions
with \(h(c\xi)=\omega(c)^{-j}h(\xi)\). The analogous lattice on
\(V\) is denoted by \(L_j(V)\). Each can be extended by zero at
the origin. The multiplicative group has order invertible in
\(\ZZ_T\), and its character table has inverse given by
orthogonality over \(\ZZ_T\). Thus the common kernel of the
orbit sums is exactly the direct sum of the \(L_j\) with \(j\)
in this range.

\begin{lemma}[Integral Fourier bases in one variable]\label{lem:gauss89}
Choose a primitive \(T\)-th root \(\zeta\), put
\(\varpi=\zeta-1\) and normalize its valuation to one in
\(\ZZ_T[\zeta]\). Put \(e_0=\mathbf1_{\{0\}}\),
\(e_{T-1}=\mathbf1_{\mathbb F_T}\), and, for \(1\le d\le T-2\),
let \(e_d(x)=\omega(x)^{-d}\) for \(x\ne0\) and \(e_d(0)=0\).
These functions form a \(\ZZ_T\)-basis. For the unnormalized
Fourier operator \((\mathcal Fh)(y)=\sum_x\zeta^{xy}h(x)\),
\[
      \mathcal Fe_d=\gamma_d e_{T-1-d},\qquad
      v_\varpi(\gamma_d)=d .
\]
Here \(\gamma_0=1,\ \gamma_{T-1}=T\). For \(1\le d\le T-2\),
\(\gamma_d=\sum_{a=1}^{T-1}\omega(a)^{-d}\zeta^a\) has squared
complex absolute value \(T\) at every complex embedding.
\end{lemma}
\begin{proof}
Character orthogonality gives a basis on \(\mathbb F_T^\times\).
Its trivial character extended by zero equals \(e_{T-1}-e_0\),
so adjoining \(e_0\) proves the basis assertion. Fourier
transformation interchanges the two endpoint functions with the
displayed factors. For the others its value at zero vanishes by
orthogonality. For \(y\ne0\) change \(x\) to \(a/y\) to obtain
\(\gamma_d\omega(y)^d\).

The shifted cyclotomic polynomial
\(((1+X)^T-1)/X\) is Eisenstein, so \(v_\varpi(T)=T-1\).
Expand \(\zeta^a=(1+\varpi)^a\). The coefficient of
\(\varpi^l\) in the resulting sum is in \(\ZZ_T\) and reduces
modulo \(T\) to \(\sum_{a\ne0}a^{-d}\binom al\).
For \(0\le l<d\) this vanishes by summing the monomials over
\(\mathbb F_T^\times\). For \(l=d\) the same sum is
\(-1/d!\bmod T\), since the leading coefficient is \(1/d!\).
The lower coefficients therefore have \(\varpi\)-valuation
at least \(T-1\), and higher powers have order greater than
\(d\). Thus \(v_\varpi(\gamma_d)=d\). At any complex embedding
the additive and multiplicative characters in this Gauss sum
remain nontrivial. In the squared absolute value write \(a=cb\).
The inner sum over \(b\ne0\) is \(T-1\) for \(c=1\) and
\(-1\) otherwise. Summing the character of \(c\) gives \(T\).
\end{proof}

\begin{theorem}[All elementary divisors in each scalar character]
\label{thm:smith89}
On \(L_j(V^*)\) the map \eqref{eq:inc89} takes values in
\(L_{T-1-j}(V)\). Up to units its elementary divisors over
\(\ZZ_T\), with multiplicities, are
\begin{equation}
   T^{(|d|-j)/(T-1)},\qquad
   d\in\{0,1,\ldots,T-1\}^k,\quad |d|\equiv j\pmod{T-1},
   \qquad |d|=\sum_i d_i.                                  \label{eq:smith89}
\end{equation}
The minimum exponent is zero and the maximum is \(k-1\), both
in every nontrivial character. At any complex embedding the
squared norm of every vector in such a character is multiplied
by \(T^{k-1}\) for the counting inner products.
\end{theorem}
\begin{proof}
Use dual coordinate bases for \(V,V^*\). The tensor products
\(e_d=\bigotimes_i e_{d_i}\) form integral bases on the whole
spaces. Scalar multiplication acts on a tensor by
\(\omega^{-|d|}\). Keeping the displayed congruence gives bases
of precisely the character lattices, and their elements vanish
at the origin. The Fourier transform in \(k\) variables sends
\(e_d\) to \(\prod_i\gamma_{d_i}\) times the tensor indexed by
\((T-1-d_i)_i\).

For an input \(h\) of character \(\omega^{-j}\), its sum over
each nonzero line is zero. The terms with \(\xi(x)=0\) in the
Fourier sum vanish by this cancellation and \(h(0)=0\). Write
each remaining covector uniquely as \(c\xi_0\) with
\(\xi_0(x)=1\). Summing first over \(c\ne0\) gives
\(\mathcal Fh=\gamma_j Dh\) on \(x\ne0\). At the origin both
sides extend by zero. Hence \(D\) sends a basis tensor to a
single target basis tensor with coefficient
\(\prod_i\gamma_{d_i}/\gamma_j\). This coefficient lies in
\(\ZZ_T\), since \eqref{eq:inc89} and both bases are integral
over \(\ZZ_T\). Its \(T\)-valuation is \((|d|-j)/(T-1)\)
by Lemma~\ref{lem:gauss89}. The target permutation is bijective.
Independent unit scalings and ordering the diagonal entries
give exactly \eqref{eq:smith89}. Take
\(d=(j,0,\ldots,0)\) and \(d=(j,T-1,\ldots,T-1)\) for the two
extremes. Every admissible total lies between their totals.
Finally \(\mathcal F^*\mathcal F=T^k I\) by orthogonality of
the additive characters and \(|\gamma_j|^2=T\). This yields
the asserted complex norm.
\end{proof}

\begin{corollary}[The joint finite map on the definite class fibres]
\label{cor:fibres89}
Give the initial definite rank-\(k\) Hermitian genus the
upper unitriangular levels at the two primes of Theorem \ref{thm:levels88}. At each place above \(T\), refine the level by a
nonzero residue covector. Let \(P\) be the number of classes
before this refinement and let the free coefficient lattice
over \(\ZZ_T\) have rank \(a_1\) per class. One may equally
extend scalars to the integers of a coefficient field at \(T\).
In the space with zero sum on every individual scalar orbit the
rank is
\begin{equation}
      Pa_1\left(\frac{T^k-1}{T-1}(T-2)\right)^f .
                                                        \label{eq:rank89}
\end{equation}
For fixed local characters \(1\le j_v\le T-2\), the elementary
divisors of the joint incidence map are the products in
\eqref{eq:smith89} over the \(f\) places, each repeated \(Pa_1\)
times before collecting equal powers. Its image contains
\(T^{f(k-1)}\) times the target lattice in those characters.
It has an image vector with a unit coordinate. The exponent
\(f(k-1)\) is the least common exponent for the inclusion.
\end{corollary}
\begin{proof}
An automorphism of a definite integral lattice has finite order.
Indeed its integral matrices form a discrete subset of the
compact stabilizer of the positive trace form. At \(g_1\) the
chosen level forces unipotent reduction of each first split
component. The other component is its inverse transpose in
dual bases. The underlying rational representation therefore
has unipotent reduction at \(g_1\). The kernel of reduction in
\(\GL_{2fk}(\ZZ_{g_1})\) has no torsion. To see this, for a
matrix \(I+g_1^b A\) with \(b\ge1\) and a unit entry of \(A\),
raising to a positive exponent prime to \(g_1\) preserves \(b\).
Raising to \(g_1\) increases it by exactly one by the binomial
formula, since \(g_1\) is odd. A nonidentity finite-order
stabilizer element would consequently have \(g_1\)-power order.
A power would have order \(g_1\), so its rational characteristic
polynomial would contain the cyclotomic polynomial of degree
\(g_1-1>2fk\). The inequality holds since
\(g_1>(kb)^8\) and \(f,k\) are polynomially bounded on the
common intervals. This contradiction proves freeness.

Above each base class the additional local choices thus form
the full product of the nonzero residue covector sets. Indeed
a split integral basis identifies each local module with
\(\mathbb F_T^k\), and the only identifications within the fibre
would be made by base-class automorphisms. Choosing other
dual bases conjugates the incidence maps by integral
permutations. Each scalar character on a single factor occurs
once on each of its \((T^k-1)/(T-1)\) scalar orbits. Summing
over the \(T-2\) nontrivial characters independently at each
place proves \eqref{eq:rank89}. Taking tensor products of the
integral diagonal bases and repeating for each coefficient
coordinate and each class proves the assertions on elementary
divisors. The extremal tuples in the one-place theorem can be
chosen simultaneously at all places.
\end{proof}

The normalized squared complex norm of this finite incidence
map grows by \(T^{f(k-1)}\) on the joint character space (with
the same counting norm on both sides). Over \(\ZZ_T\) it has a
unit elementary divisor even when all the local characters are
nontrivial. The formula concerns the full initial class fibres.
The valuation of a subsequent global period on a specified
vector also depends on the remaining lifts and on the terminal
linear functional.

\subsection{Primitive negative lines at the terminal level}\label{nav:VI-07}
\label{sec:line92}

Use the arithmetic row lattice \(\Lambda=\OO_E^r\) of
Theorem \ref{thm:norm91}.
Its diagonal Hermitian form \(h\) has \(p_+=k+m/2-1\) entries
of sign \(+\) and \(p_-=m/2\) of sign \(-\). At large \(x\)
these supply at least one positive and two negative coordinates.
Take the form to be linear in its second variable. Write
\(M=N(g_1g_2)^2\) and fix a prime from Corollary \ref{cor:real91},
with \(T=\pi\bar\pi\) and \(\pi\equiv1\bmod M\OO_{K_d}\).
The two principal ideals over \(T\) in \(K_d\) are distinct.
Choose three coordinate vectors \(e_+,e_-,e\), with the first
positive and the last two negative, and set \(u=e_++e_-\).
This isotropic vector is orthogonal to \(e\).

\begin{theorem}[The norm, complement and metric of a negative line]
\label{thm:line92}
For each integer \(c\) with \(1\le c<T\) put
\[
                    z_c=\pi e+Mc\,u .
\]
Then \(h(z_c,z_c)=-T\), \(z_c\equiv e\bmod M\Lambda\), and
\(\{t\in E:t z_c\in\Lambda\}=\OO_E\). Its orthogonal
complement \(L_c=z_c^\perp\cap\Lambda\) has signature
\((p_+,p_--1)\) everywhere at infinity. As \(\OO_E\)-modules,
\[
 \Lambda/(\OO_E z_c\oplus L_c)\simeq\OO_E/T\OO_E,
 \qquad L_c^\#/L_c\simeq\OO_E/T\OO_E ,
\]
where the dual is taken for the restricted Hermitian form.
For the positive norm obtained by replacing every diagonal
sign by \(+\), at every archimedean embedding one has
\[
          \|z_c/\pi\|^2=1+2M^2c^2/T .
\]
\end{theorem}
\begin{proof}
The orthogonal and isotropic relations give the Hermitian norm,
and the prescribed congruence of \(\pi\) gives the level
condition. No prime of \(E\) can divide both \(Mc\) and \(\pi\),
since \(T\nmid Mc\). The coordinates thus generate the unit
ideal. Multiplication by a scalar is integral on \(z_c\) precisely
when that scalar is integral. The functional
\(\ell(y)=h(z_c,y)\) is onto \(\OO_E\) by the same coordinate
ideal, or by self-duality of \(\Lambda\). Its kernel is \(L_c\)
and \(\ell(z_c)=-T\). This proves the first quotient formula
and the signature follows by orthogonal decomposition over
the real places.

The projection \(y\mapsto y+\ell(y)z_c/T\) from \(E^r\) onto
\(z_c^\perp\) sends \(\Lambda\) into \(L_c^\#\). It maps onto
that dual. Indeed any integral linear functional on \(L_c\)
extends to \(\Lambda\), since \(\Lambda/L_c\simeq\OO_E\)
is free. Represent the extension by an element of the self-dual
\(\Lambda\) using the first variable of \(h\). Projecting this
element leaves its pairing with \(L_c\) unchanged. The kernel
on \(\Lambda\) is its intersection with \(E z_c\), which is
\(\OO_Ez_c\) by primitivity. Projection is the identity on
\(L_c\), so passing to the quotient gives the second formula.
Finally every embedding of \(K_d\) sends \(\pi\) to a complex
number of squared absolute value \(T\), and the three coordinate
vectors are orthogonal in the positive norm. This proves the
metric formula.
\end{proof}

An integral polynomial \(P\) homogeneous of degree \(A\) in a
terminal vector restricts to a binary form
\[
        P(su+te)=\sum_{i=0}^{A} b_i s^i t^{A-i}.
\]
The same notation applies to a multihomogeneous polynomial of
total degree \(A\) on several copies of this plane on inserting
the same \(s,t\) in every copy. When the copies come from a CM
type above a fixed embedding of \(K_d\), the element \(\pi\)
is the same scalar at all of them. Coefficients may be in a
number field \(J\) containing \(E\).

\begin{theorem}[Integral evaluation in the normalized negative line]
\label{thm:jets92}
Let \(0\le A<T-1\) and assume the \(b_i\) are integral at a
place \(w\) above \((\pi)\). The normalized evaluations
\(P(z_c/\pi)\) are integral at \(w\) for all \(1\le c\le A+1\)
if and only if
\begin{equation}
           b_i\in \pi^i\OO_{J,w}\quad(0\le i\le A).     \label{eq:jets92}
\end{equation}
In fact, with \(v_w(T)=1\) and \(v_w(0)=+\infty\),
\[
     \min_{1\le c\le A+1}v_w(P(z_c/\pi))
            =\min_{0\le i\le A}\bigl(v_w(b_i)-i\bigr).
\]
At places above \((\bar\pi)\) the normalization is by a unit.
For every nonzero evaluation \(y=P(z_c)\) in a normal coefficient
field, replacing \(y\) by \(P(z_c/\pi)\) decreases the normalized
logarithm of the absolute field norm of \(y\bar y\) by exactly
\(A\log T\). For integral values before division,
\eqref{eq:jets92} at all \(w\mid(\pi)\) prevents any new finite
denominator in any of these evaluations.
\end{theorem}
\begin{proof}
Homogeneity gives \(P(z_c/\pi)=\sum_i b_i(M/\pi)^i c^i\).
The matrix \((c^i)_{1\le c\le A+1,\ 0\le i\le A}\) has
Vandermonde determinant \(\prod_{c<d}(d-c)\), a unit at \(w\)
because all the positive differences are less than \(T\).
Both this matrix and its inverse preserve the coefficient
lattice. Thus its output coordinates generate the same
fractional ideal in \(\OO_{J,w}\) as the input coordinates.
Now \(M\) and \(\bar\pi\) are units at \(w\), and
\(v_w(\pi)=v_w(T)=1\). Taking valuations proves the criterion
and the equality, also if the restricted form is identically
zero. At all finite places away from \((\pi)\), division by
\(\pi^A\) preserves integrality. Finally \(P(z_c/\pi)=
P(z_c)/\pi^A\). Forming its product with its complex conjugate
divides by the rational integer \(T^A\), giving the asserted
change in the logarithmic norm.
\end{proof}

\begin{corollary}[Metric cost at all the high orders]\label{cor:cost92}
Let \(a=\lceil xb\rceil\) and let \(A=ha\) be the sum of the
orders at the high embeddings of Theorem \ref{thm:single82}. Then
\(A+1<T\) eventually on the common intervals. Uniformly for
\(1\le c\le A+1\) and every archimedean embedding,
\[
       0\le A\log\|z_c/\pi\|^2
                  \ll \exp(-c_0 m x^\eps)
\]
for an absolute \(c_0>0\) and all sufficiently large \(x\)
in the given family.
\end{corollary}
\begin{proof}
Here \(h\le f\) is polynomially bounded,
\(\log M=O_{\psi,\eps}(x^\eps)\) and
\(\log(A+1)=O_{\psi,\eps}(x^\eps)\). Corollary \ref{cor:real91}
gives \(\log T\ge(16-o(1))m x^\eps\), where \(m\to\infty\).
The bound now follows from
\(A\log(1+2M^2c^2/T)\le 2A M^2(A+1)^2/T\).
\end{proof}

The norm and primitivity requirements for the Bessel line are
visible in
\srcloc{Development/AbelianVarietyTheory/Proofs/DecoratedLatticeProofs/BesselLevelLineVectorPrimitiveRowLatticeNormStatement}{12}.
The factor \(T^{-A}\) in Theorem~\ref{thm:jets92} has an exact
finite integrality condition in this family of normalized lines.
For instance \(P(X)=X_+^A\), \(A>0\), has \(b_A=1\), and every
tested normalized evaluation has valuation \(-A\) at \(w\).
For the prescribed global period, its restriction to these lines
and the required divisibilities of its coefficients still have
to be proved. A net saving in its height also requires a bound
for the undivided value on \(z_c\), whose positive norm grows
as \(\sqrt T\) over the tested range. The example exhibits the
role of the coefficients even for primitive lines of the correct
norm and congruence.

\subsection{Terminal levels for the refined column periods}\label{nav:VI-12}
\label{sec:refined112}

Fix a field, a prime \(T=\varpi\bar\varpi\), and a family of column
types as in Theorem \ref{thm:common111} and Corollary \ref{cor:count111}. At every \(v\mid T\) of \(F_+\)
choose a nonzero residue covector \(\ell_v\) on the \(k\)-dimensional
standard column lattice. Write \(K_\pi(\ell)\) for its stabilizer
at every such place, with the previous column level elsewhere.
Write \(X_\pi(\ell)=[H]/K_\pi(\ell)\). The scalar orbits in each
refinement fibre and the multiplicities on their zero-sum subspace
are those of Corollary~\ref{cor:fibres89}.
At the unrefined row level of Section~\ref{sec:global110} the
Bessel coefficient is invariant under the unrefined column level.
Hence its pairing with every such zero-sum column vanishes,
for every frequency. Indeed the coefficient is constant on each
refinement fibre, and averaging a zero-sum column over the
covectors in any one place gives zero. A finite translate of
the row changes the stabilizer of the coefficient. We give one
translate that realizes the full refined column space at both
frequencies and preserves the volume of the row level.

\begin{lemma}[A split row lattice inducing a covector stabilizer]
\leavevmode\label{lem:level112}
There is a finite adelic element \(a_f\in G(\mathbb A_{F_+,f})\),
equal to the identity away from \(T\), with the following properties.
Put \(K_\sigma'=a_fK_\sigma a_f^{-1}\), \(K_B'=B(\mathbb A_{F_+,f})
\cap K_\sigma'\), and let \(K_\Lambda\) be the full original
lattice stabilizer. Then
\[
 K_B'=(N(\mathbb A_{F_+,f})\cap K_\sigma')\rtimes K_\pi(\ell),
 \qquad \chi_\beta(K_B')=1\quad(\beta\in\OO_E).
\]
At each \(v\mid T\) the index of \(K'_{\sigma,v}\cap K_{\Lambda,v}\)
in \(K_{\Lambda,v}\) equals the number of subspaces of codimension
\(s+1\) in \(\mathbb F_T^r\), where \(r=2s+k+1\). The finite
volume of \(K_\sigma'\) is that of \(K_\sigma\).
\end{lemma}
\begin{proof}
At \(v\mid T\), take \(O=\ZZ_T\) and write the split Hermitian
space in its two factors \(V^+,V^-\). The perfect pairing
identifies the second factor with the dual of the first. The
split unitary group acts by an arbitrary \(g\in\GL(V^+)\) on
the first and by its contragredient on the second. Use the Witt
basis of Section~\ref{sec:char94}. It is an integral basis
locally, since \(T\) is odd. Use a primitive row vector lifting
\(\ell_v\), also denoted \(\ell_v\), on the \(W_0^+\) coordinates.
In the positive split factor take the lattice
\[
 L^+=\bigoplus_{i=1}^s TOu_i^+
       \ \oplus\ M_\ell\
       \oplus\bigoplus_{i=1}^s Ov_i^+,\qquad
 M_\ell=\{(w,z)\in O^k\oplus O:z\equiv\ell_v(w)\bmod T\}.
\]
Here \(z\) is the coordinate on \(Ee\), and superscripts on
basis vectors denote their split components. Take \(L^-\) to
be the lattice dual to \(L^+\). Multiplying each \(u_i^+\)
by \(T\), acting on the middle coordinates by
\((w,z)\mapsto(w,\ell_v(w)+Tz)\), and fixing each \(v_i^+\)
takes the original positive-factor lattice onto \(L^+\).
Together with the dual action this defines \(a_v\). Taking the
product of these elements gives \(a_f\). At \(v\) the conjugate
level is precisely the stabilizer of \(L^+\) and its dual.

Every \(nh\in B\) preserves the Witt flag. On the graded middle
quotient \(n\) acts trivially and \(h\) acts by \((w,z)\mapsto
(hw,z)\). The induced middle lattice is exactly \(M_\ell\).
Its projection onto \(w\) is \(O^k\). Preserving it in both
directions therefore requires \(h\in\GL_k(O)\) and
\(\ell_vh=\ell_v\bmod T\). Conversely every such \(h\), embedded
as the identity on the remaining coordinates, preserves \(L^+\)
and the dual lattice. Thus for any \(nh\) in the stabilizer
one has \(h\) there, and then \(n\) there as well. This proves
the local semidirect product formula.

It remains to check the character on this unipotent stabilizer.
Inside the flag the \(O\)-coefficients of \(u_i^+\) in \(L^+\)
form \(TO\), for every \(i\). In the dual lattice the corresponding
coefficients of \(u_i^-\) form \(O\) because \(u_i^-\) pairs
with \(v_i^+\). Applying \(n-1\) to \(Tu_{i+1}^+\) and to
\(u_{i+1}^-\) proves integrality of both split components of
\(d_i(n)\) for \(i<s\). Moreover \(Te^+\in L^+\), and
\(e^-\in L^-\) because the \(z\)-coordinate of \(M_\ell\) is
integral. Applying \(n-1\) to these two vectors proves the
same assertion for \(d_s(n)\). Thus every summand in the
finite additive character defining \(\chi_\beta\) is integral
if \(\beta\in\OO_E\). The character is trivial on \(H\) by its
definition. Away from \(T\), Lemma~\ref{lem:orbit110} applies.

The lattice \(L^+\) contains \(T\) times the original lattice
and is contained in the original one. Modulo \(T\), the flag
conditions give \(s\) independent zero coordinates and the
middle graph gives one more independent equation. The stabilizer
inside \(\GL_r(O)\) is the inverse image of the stabilizer of
that codimension-\((s+1)\) subspace. Reduction onto \(\GL_r(\mathbb F_T)\)
is surjective by lifting a matrix with invertible determinant.
The latter group is transitive on subspaces of the given dimension
by extending bases. This proves the index. The volume identity
follows from inner invariance of Haar measure on the unimodular
group \(G(\mathbb A_{F_+,f})\).
\end{proof}

\begin{theorem}[Prescribing the refined periods by one finite translate]
\label{thm:onto112}
Let \(\mathcal C_\ell\) be the sum of the full isotypic spaces
of any finite set of archimedean column types on \(X_\pi(\ell)\).
Let \(J\) be any finite set of distinct nonzero integral
frequencies in \(E\), and let \(\eta>0\). There is an
\(H_\infty\)-equivariant operator from \(\mathcal C_\ell^J\)
to compactly supported smooth, \(K_\infty\)-finite functions
on the original-level quotient \([G]/K_\sigma\), of \(L^2\)
operator norm at most \(\eta\). Its output \(F\) for the tuple
\((\phi_\alpha)_{\alpha\in J}\) satisfies
\[
         {\cal F}_\alpha(R(a_f)F)=\phi_\alpha
                          \quad\hbox{on }X_\pi(\ell)\quad(\alpha\in J).
\]
The same \(a_f\) works for all columns, all frequencies and all
inputs. For \(J=\{1,\varpi\}\), pairing with an orthonormal
basis of the refined zero-sum spaces gives the identity matrix
at each matching frequency and the zero matrix across frequencies.
\end{theorem}
\begin{proof}
Use the quotient with row level \(K_\sigma'\) first. The compact
period orbit \(B(F_+)\backslash B(\mathbb A_{F_+})/K_B'\) embeds
there by the argument of Lemma~\ref{lem:orbit110}: a rational
identification of two points lies in \(B\) at infinity, hence in
\(B(F_+)\). Its image is an embedded submanifold by compactness
and the inclusion of the real Lie groups. Lemma~\ref{lem:level112}
ensures that \(\sum_\alpha\chi_\alpha(n)\phi_\alpha(h)\) descends
to this orbit. Fourier orthogonality on \([N]\) gives precisely
the prescribed restrictions. The proof of Theorem~\ref{thm:section110}
now applies. Extending along a small normal tube gives the
equivariant extension of arbitrarily small \(L^2\) norm.
The finite compact-group projections followed by inversion of
the induced endomorphism preserve the prescribed coefficients
exactly. Here finite dimensionality of the complete column
isotypic spaces follows as in Proposition~\ref{prop:column110},
since the refined definite genus still has finitely many classes.

If \(F'\) denotes this output at level \(K_\sigma'\), put
\(F=R(a_f^{-1})F'\). This is at level \(K_\sigma\), and right
translation is unitary for the adelic quotient measure. It
commutes with the archimedean actions. Using the entire column
isotypic spaces in the construction allows us afterwards to
restrict the input to the zero-sum subspaces. Choosing one
basis element in one frequency component proves the matrix
assertion by orthonormality.
\end{proof}

\begin{corollary}[The analytic cost on a prescribed row type]
\label{cor:cost112}
For the parameters and terminal primes above, the bound
\[
 |{\cal P}_{\varpi}(R(a_f)F,\phi)|
 \ \le\ \exp(-c\sqrt T/r)\,\|F\|_{L^2([G])}\|\phi\|_{L^2([H])}
\]
holds for any row form \(F\) at the original level satisfying
the archimedean conditions of Theorem \ref{thm:raw97}, and for its bounded real translates as
in Proposition~\ref{prop:translate109}, with an absolute \(c>0\).
Thus a unit column receiving the scalar value one from a row
in this class forces \(\|F\|_{L^2([G])}\ge\exp(c\sqrt T/r)\).
\end{corollary}
\begin{proof}
Restrict the translated row to the common open subgroup
\(K_\sigma'\cap K_\Lambda\). At each \(v\mid T\) the index
from Lemma~\ref{lem:level112} is at most
\(|\GL_r(\mathbb F_T)|\le T^{r^2}\).
The added factor in the restriction norm of Theorem~\ref{thm:rest96}
is at most \(T^{fr^2/2}\). Its logarithm is polynomially bounded
in \(x\) by Lemma \ref{lem:size97}. The finite translation commutes
with all the real differentiations and preserves the global
\(L^2\) norm. In Theorem \ref{thm:raw97}, the choice of
differentiation order in the \(\varpi\)-frequency consequently
gives the same bound with just this extra factor. Since
\(\log T\gg m x^\eps\) and \(r\) is polynomial in \(x\),
the additional logarithm is \(o(\sqrt T/r)\). It is absorbed
by decreasing \(c\). The real-translation derivative estimate
of Proposition \ref{prop:translate109}, commutes with the finite action
as well. Substituting a scalar value one gives the last claim.
\end{proof}

The finite change here preserves the row volume in the comparison
with the refined multiplicities in Corollary \ref{cor:rank111}.
Theorem~\ref{thm:onto112} leaves the decomposition of its smooth
row functions into automorphic representations undetermined.
To deduce the arithmetic estimate one
still needs a nonzero coefficient on the iterated integral theta
image, with its arithmetic and analytic normalizations compared.

\subsection{The Whittaker character at the two auxiliary primes}\label{nav:VI-13}
\label{sec:whit125}

The unfolded Gram matrix for the projected Poincar\'e series has a
finite factor indexed by changes of basis at \(g_1,g_2\). We compute
each entry, including its sign and its denominator. Write
\(F_+=E^+\), \(f=[F_+:\QQ]\), and retain \(E,x,k,b,g_1,g_2,T\) from
Section \ref{nav:I-28}.
Both \(g_i\) split completely in \(E\). At each place \(v\) of
\(F_+\) over \(g_1g_2\) write \(q_v=\#\kappa(v)\). Thus \(g_1\)
and \(g_2\) each occur \(f\) times in the list of \(q_v\).

Over \(\mathbb F_q\), let \(G_n=\GL_n(\mathbb F_q)\) and let \(U_n\)
be its upper unitriangular subgroup. Fix a nontrivial additive
character \(\psi\) of \(\mathbb F_q\). The character
\(\theta_n(u)=\psi(\sum_{i=1}^{n-1}u_{i,i+1})\) is multiplicative
on \(U_n\). Let \(\gamma_n\) be the character of the complex
representation \(\operatorname{Ind}_{U_n}^{G_n}\theta_n\).
The formula for an induced character gives
\begin{equation}
 \gamma_n(g)=\frac{1}{|U_n|}
 \sum_{\substack{a\in G_n\\aga^{-1}\in U_n}}\theta_n(aga^{-1}).
 \label{eq:ind125}
\end{equation}
Indeed, on the lines indexed by left cosets \(U_n a\), right
translation permutes the lines. Each fixed line contributes the
indicated character to the trace, and has \(|U_n|\) representatives.
In particular \(\gamma_n\) vanishes outside the unipotent elements.
For \(I=\prod_{v\mid g_1g_2}\GL_k(\kappa(v))\), the normalized
matrix in question is
\begin{equation}
 A_{ij}=\prod_{v\mid g_1g_2}
 \frac{\overline{\gamma_{k,v}(j_v^{-1}i_v)}}{\gamma_{k,v}(1)}
 \qquad(i,j\in I).
 \label{eq:actual125}
\end{equation}
Here one may choose the nontrivial additive character separately at
each place. The source gives the conjugation sum in
\srcloc{Formalism/LocalToGlobal/SiegelLatticeTheory/ClosedLocalOrbits/GGMirabolic}{28}
and the relative basis change \(j^{-1}i\) in
\srcloc{Formalism/LocalToGlobal/SiegelLatticeTheory/ClosedLocalOrbits/TermOpStmt}{140}.
We use \(\gamma_n(1)\) to make the diagonal equal to one.

\begin{lemma}[Averages on a block parabolic]\label{lem:parabolic125}
If \(n=a+b\), \(a,b>0\), and \(A_0\in G_a,\ C_0\in G_b\), then
\begin{equation}
 q^{-ab}\sum_{X\in\operatorname{Mat}_{a,b}(\mathbb F_q)}
 \gamma_n\!\begin{pmatrix}A_0&X\\0&C_0\end{pmatrix}
     =\gamma_a(A_0)\gamma_b(C_0).
 \label{eq:parabolic125}
\end{equation}
\end{lemma}
\begin{proof}
Let \(P\) stabilize the first \(a\) coordinate directions and let
\(V\) be the normal subgroup with identity diagonal blocks. Averaging
over \(V\) projects any complex representation onto its \(V\)-fixed
vectors. This projection commutes with block diagonal matrices.
The left side of \eqref{eq:parabolic125} is therefore the trace
of \((A_0,C_0)\) on these fixed vectors.

Realize the induced representation by functions \(h\) on \(G_n\)
with \(h(ug)=\theta_n(u)h(g)\), acting by right translation.
Partition their supports into \(U_n w P\). The cells are indexed by
subsets \(J\subseteq\{1,\ldots,n\}\) of size \(a\). Indeed a subspace
in the Grassmannian has a basis in column echelon form with leading
(last nonzero) positions \(J\). Upper unitriangular row operations
take it to the corresponding coordinate subspace. Take \(w\) to
send the first block to \(J\) and the second to its complement,
in increasing order within each. If \(i\in J,\ i+1\notin J\),
then \(1+tE_{i,i+1}\) belongs to \(U_n\cap wVw^{-1}\).
It acts on the left by \(\psi(t)\). Invariance on the right by
\(V\), normal in \(P\), forces a fixed function to vanish on
\(wP\) as soon as \(\psi(t)\ne1\). The only remaining subset
is the final segment of size \(a\).

For this final segment the intersection \(w^{-1}U_nw\cap P\)
consists of block diagonal elements in \(U_a\times U_b\).
The character there is \(\theta_a\theta_b\), since consecutive
coordinates within either block are still consecutive after the
permutation. On the surviving cell a right \(V\)-fixed function is
thus determined by its restriction to \(w(G_a\times G_b)\), with
precisely these two left transformation laws. Conversely such a
function extends by the two transformation laws and right
\(V\)-invariance to the cell, and by zero elsewhere. This identifies
the fixed space with the tensor product of the two induced
representations. Taking traces proves the formula.
\end{proof}

\begin{theorem}[Value at every relative basis change]\leavevmode\label{thm:char125}
Set \(D_q(d)=\prod_{j=1}^{d}(q^j-1)\), with \(D_q(0)=1\).
For a unipotent \(g\in G_n\) let \(d=\dim\ker(g-1)\). Then
\[
                       \gamma_n(g)=(-1)^{n-d}D_q(d).
\]
\end{theorem}
\begin{proof}
Define \(B_n\) by the displayed formula on unipotents and by zero
elsewhere. We first prove that the functions \(B_n\) satisfy
\eqref{eq:parabolic125}. If either diagonal block is not unipotent,
its characteristic polynomial gives a factor of the block triangular
characteristic polynomial, and both sides vanish. Otherwise put
\(\alpha=\dim\ker(A_0-1)\), \(\beta=\dim\ker(C_0-1)\). A uniformly
chosen \(X\) induces a uniformly chosen linear map
\[
 Y:\ker(C_0-1)\longrightarrow \mathbb F_q^a/\operatorname{im}(A_0-1).
\]
Restriction followed by quotient is surjective on these spaces of
linear maps, by extending bases. Solving the lower equation first
and then the upper one gives kernel dimension
\(\alpha+\beta-\operatorname{rank}Y\) for the block matrix minus
the identity.

Expose the columns of \(Y\) one at a time. After \(j\) columns of
rank \(h\), write \(l=\alpha+j-h\). The next column belongs to
their span with probability \(q^{h-\alpha}\). The conditional
expectation of \((-1)^{\operatorname{rank}Y}D_q(\alpha+j+1-
\operatorname{rank}Y)\) after that column is
\[
 (-1)^h\{q^{h-\alpha}D_q(l+1)-(1-q^{h-\alpha})D_q(l)\}
        =(-1)^h D_q(l)(q^{j+1}-1).
\]
Starting with \(D_q(\alpha)\) and taking \(\beta\) successive
expectations gives \(D_q(\alpha)D_q(\beta)\). Multiplying by
\((-1)^{n-\alpha-\beta}\) proves the block formula for \(B_n\).

We prove \(B_n=\gamma_n\) by induction. For \(n=1\) both are the
character of the regular representation of \(\mathbb F_q^\times\).
At the identity in every rank the induced character has degree
\(|G_n|/|U_n|=D_q(n)\), by counting ordered independent columns
and the free entries in \(U_n\). Let \(h=\gamma_n-B_n\). Induction
and Lemma~\ref{lem:parabolic125} give zero for the average of \(h\)
over the upper right block with any fixed proper diagonal blocks.
We have \(h(1)=0\), and \(h\) is supported on unipotents.

For a partition \(\lambda=(\lambda_1,\ldots,\lambda_m)\) of \(n\)
with decreasing parts, let \(c_\lambda(g)\) count flags with these
successive dimensions such that \((g-1)F_i\subseteq F_{i-1}\) for
all \(i\), where \(F_0=0\). Then \(\sum_g h(g)c_\lambda(g)=0\).
For \(m=1\) this is \(h(1)=0\). Otherwise fix one flag and choose
a basis adapted to it. In the sum over matrices lowering the flag
after subtraction of the identity, fix the first diagonal block
and all entries in the remaining diagonal block. The upper right
rectangle ranges freely, and its sum is zero by the block average.
Summing also over flags proves the assertion.

Write \(\mu\) for the partition of Jordan block sizes of \(g-1\),
and \(\mu'\) for its conjugate partition. A counted flag has
\(F_i\subseteq\ker((g-1)^i)\), so
\(\sum_{j\le i}\lambda_j\le\sum_{j\le i}\mu'_j\) for every \(i\),
with missing parts read as zero.
When \(\lambda=\mu'\) the flag is uniquely
\(F_i=\ker((g-1)^i)\). Order the partitions by any total order
extending these partial-sum inequalities. The matrix of flag
counts on unipotent classes is triangular with ones on the
diagonal when columns are indexed by \(\mu'\).
It is invertible. The vector of class sizes times the class
values of \(h\) is in its kernel. Every such value vanishes.
\end{proof}

\begin{corollary}[Content on the integral basis and pointwise decay]
\label{cor:content125}
An entry of \eqref{eq:actual125} vanishes if a local relative
matrix is not unipotent. Otherwise set
\(t_v=\operatorname{rank}(j_v^{-1}i_v-1)\). Then
\begin{equation}
 A_{ij}=\prod_{v\mid g_1g_2}
     \frac{(-1)^{t_v}}
          {\displaystyle\prod_{h=k-t_v+1}^{k}(q_v^h-1)} .
 \label{eq:entry125}
\end{equation}
On the lattice with basis \(I\), the values of this pairing generate
the fractional ideal \(D^{-1}\ZZ\), where
\begin{equation}
             D=\prod_{v\mid g_1g_2}\prod_{h=2}^{k}(q_v^h-1).
 \label{eq:den125}
\end{equation}
For \(k\ge2\), the maximum of \(|A_{ij}|\) over \(i\ne j\) equals
\((\min(g_1,g_2)^k-1)^{-1}\).
\end{corollary}
\begin{proof}
Theorem~\ref{thm:char125} divided by its value at the identity
gives \eqref{eq:entry125}, including the complex conjugation.
A unipotent in positive dimension has a fixed vector, so
\(t_v\le k-1\). Every denominator divides \(D\).
Fix any \(j\) and at each place let \(j_v^{-1}i_v\) have one
Jordan block with eigenvalue one. The entry is then \(\pm1/D\).
This proves the ideal assertion and the minimality of the common
denominator. For distinct indices any nonzero entry has a factor
with \(t_v\ge1\), giving the stated upper bound. A single
transvection at one place of smallest residue size and identities
elsewhere give equality.
\end{proof}

Thus one and the same pair of basis changes gives a primitive entry
after clearing the denominator for every word and terminal prime
of the prescribed family. The integral basis in this assertion
consists of the changes of residue basis at the two auxiliary
primes. Forming a coefficient of the projected theta lift also
involves summing the global orbital terms with their archimedean
weights.

\subsection{The exact second moment on the residue bases}\label{nav:VI-14}
\label{sec:moment126}

Let \(A\) be the normalized finite matrix of \eqref{eq:actual125} on
\(\CC^I\) with its orthonormal basis indexed by \(I\). Put \(Q=\min(g_1,g_2)\).
The number of entries in a row is too large to bound its squared
norm by multiplying it by the square of the largest entry. We
compute the whole sum by counting the intertwining operators of
the finite induced representation.

\begin{lemma}[Intertwining operators at one place]\label{lem:intertwine126}
For \(\Gamma_n=\operatorname{Ind}_{U_n}^{G_n}\theta_n\) over
\(\mathbb F_q\) one has
\[
            \dim_{\CC}\operatorname{End}_{G_n}(\Gamma_n)
                         =(q-1)q^{n-1}.
\]
\end{lemma}
\begin{proof}
An equivariant endomorphism of the induced representation is
determined by the image of the function supported on \(U_n\)
with value one at the identity. Its translates are a basis after
one representative per coset has been chosen. Their relations for
other representatives impose on the image the right transformation
law by \(\theta_n\), in addition to its defining left transformation
law. Conversely every function with both laws defines an
endomorphism on those translates. Such a function's value at \(y\)
can be chosen freely on its double coset precisely when
\(\theta_n(u)=\theta_n(yuy^{-1})\) for all \(u\in U_n\) with
\(yuy^{-1}\in U_n\). A failure of this equality forces the value
to be zero. Hence the dimension is the number of compatible double
cosets.

The Bruhat decomposition writes these cosets uniquely as \(U_n t w U_n\)
with \(w\) a permutation matrix and \(t\) invertible diagonal.
Indeed the upper triangular double cosets are indexed by the
permutation of pivots under elimination. Nonzero diagonal factors
on either side can be moved through the permutation and combined.
If two diagonal factors represent the same \(U_n\)-double coset
at \(w\), their ratio is the diagonal of a conjugate by \(w\) of
a unitriangular matrix, hence is the identity.

For \(i<j\) with \(w(i)<w(j)\), conjugation by \(tw\) sends
\(1+zE_{ij}\) to \(1+z(t_{w(i)}/t_{w(j)})E_{w(i),w(j)}\).
These elementary subgroups generate the intersection in the
compatibility condition. Its allowed entries are precisely the
pairs increasing in both orders. One can clear them in increasing
order of \(j-i\), since concatenating two allowed pairs again gives
an allowed pair. Testing every \(z\) detects equality of the
coefficients, as a nonzero multiple of \(z\) runs over the field.
Every adjacent ascent of \(w\) must therefore increase by one.
Its increasing runs are disjoint intervals of consecutive values.
They occur in decreasing order of intervals, since a descent
separates consecutive runs. Conversely these permutations have
the required matching of adjacent pairs in both orders. The
coefficients then agree exactly when \(t\) is constant on the
image of each run. For \(l\) runs there are
\(\binom{n-1}{l-1}\) choices of lengths and \((q-1)^l\) choices
of scalars. Summing over \(l\) gives \((q-1)q^{n-1}\).
\end{proof}

\begin{theorem}[Exact squared deviation and common denominator]
\label{thm:moment126}
Every row of \(A\) has the same squared distance from the
corresponding coordinate vector. Write this squared distance
as \(\delta\). With \(D\) the exact denominator of
Corollary~\ref{cor:content125}, one has
\begin{align}
 1+\delta&=\prod_{v\mid g_1g_2}\prod_{j=2}^k(1-q_v^{-j})^{-1},
 &\|A-\mathrm{Id}\|_{\mathrm{HS}}^2&=|I|\delta, \label{eq:moment126}\\
 D(1+\delta)&=(g_1g_2)^{\,f(k(k+1)/2-1)}.&& \label{eq:massden126}
\end{align}
On the fields of Section~\ref{sec:whit125} and throughout the
specified intervals,
\[
 0\le\delta\le \exp\!\big(2f(g_1^{-2}+g_2^{-2})\big)-1
             =O\!\left(\frac{f}{(kb)^{16}}\right).
\]
All these statements hold for every word using the two auxiliary
primes, independently of the choice of terminal prime.
\end{theorem}
\begin{proof}
At a single place the average of conjugation by \(G_k\) on
\(\operatorname{End}_{\CC}(\Gamma_k)\) projects onto its invariant
endomorphisms. Conjugation by \(g\) has trace
\(\gamma_k(g)\overline{\gamma_k(g)}\), by taking
\(\Gamma_k\otimes\Gamma_k^\vee\) and a unitary metric. Hence
\[
   \sum_{g\in G_k}|\gamma_k(g)|^2=|G_k|(q-1)q^{k-1}
\]
by Lemma~\ref{lem:intertwine126}. Divide by
\(\gamma_k(1)^2=D_q(k)^2\). Since
\(|G_k|=q^{k(k-1)/2}D_q(k)\), cancellation of \(q-1\)
leaves \(\prod_{j=2}^k(1-q^{-j})^{-1}\). In each row of
\eqref{eq:actual125} the relative elements run over the entire
product group. The sum of the squared moduli of the row factors
as the product of the local sums. Its diagonal entry is one.
Subtracting this entry's squared modulus proves the first line.
Multiplication by \(D\) replaces each factor \(q_v^j-1\) by \(q_v^j\),
proving the second line.

For \(q\ge3\) one has
\(\sum_{j=2}^\infty-\log(1-q^{-j})
\le q^{-2}/((1-q^{-1})(1-q^{-2}))\le2q^{-2}\).
Each residue size occurs \(f\) times. Furthermore
\(g_i\ge(kb)^8,\ b=\exp(x^\eps)\) by
Proposition \ref{prop:ray50} and its choice of the two auxiliary primes.
Here \(f\) grows at most polynomially in \(x\), so the exponent
in the bound tends uniformly to zero as the left endpoint tends
to infinity. The inequality \(e^t-1\le2t\) for sufficiently small
nonnegative \(t\) proves the stated order.
\end{proof}

\begin{proposition}[Quadratic estimates for the orbital amplitudes]
\label{prop:amplitudes126}
For vectors \(z_i\) in a complex Hilbert space set
\[
   S=\sum_i\|z_i\|^2,\qquad
   B=\sum_{i,j}A_{ij}\langle z_i,z_j\rangle .
\]
In addition to positivity, \(B\) has the uniform upper bound
\begin{equation}
 0\le B\le
 \prod_{v\mid g_1g_2}\prod_{j=1}^{k-1}(1-q_v^{-j})^{-1}\, S
 \le (1+3/g_1)^f(1+3/g_2)^f S .
 \label{eq:spectral126}
\end{equation}
If \(s\) of the \(z_i\) are nonzero, then
\begin{equation}
 |B-S|\le\frac{(\sum_i\|z_i\|)^2-S}{Q^k-1}
                  \le\frac{s-1}{Q^k-1} S \qquad(s\ge1).
 \label{eq:sparse126}
\end{equation}
In particular the quadratic form is positive definite on any set
of fewer than \(Q^k\) indices. On the full index set and for
\(k\ge2\) it annihilates each array constant in any one local index.
\end{proposition}
\begin{proof}
The single-place matrix is a Gram matrix of the operators
\(\Gamma_k(i)\) in the Hilbert--Schmidt inner product divided by
\(\gamma_k(1)\), since
\(\operatorname{tr}(\Gamma_k(j)^*\Gamma_k(i))=\gamma_k(j^{-1}i)\).
The character is real by Theorem~\ref{thm:char125}. Products of
these matrices are tensor products of positive operators.

To bound the norm at one place, use the regular representation
on \(\CC^{G_k}\). The matrix is convolution by the central function
\(\overline{\gamma_k}/\gamma_k(1)\). Its centrality gives the same
operator in the left and right convolution conventions. On each irreducible
constituent \(\pi\) of dimension \(d_\pi\), Schur's lemma and
the trace give the scalar \(|U_k|m_\pi/d_\pi\), where
\(m_\pi=\dim\operatorname{Hom}_{U_k}(\theta_k,\pi)\).
Indeed the character inner product with \(\gamma_k\) equals this
multiplicity by the induced function model, or directly by
\eqref{eq:ind125} and the character projection on \(U_k\).

For completeness, \(d_\pi\ge m_\pi\prod_{j=1}^{k-1}(q^j-1)\).
This inequality in rank \(n\) holds for any representation of
the subgroup \(P_n^0\) fixing the last coordinate covector.
For \(n=1\) it is equality. In rank \(n\) its subgroup of last
column translations is the additive group \(\mathbb F_q^{n-1}\).
Its characters are \(\psi(\ell(x))\) for row covectors \(\ell\).
Distinct covectors give distinct characters, and their count is
the order of the additive group. The block \(G_{n-1}\) acts
transitively on the \(q^{n-1}-1\) nonzero covectors, so their
character spaces have the same dimension. The space for the last
coordinate covector is stable under \(P_{n-1}^0\). Inside it the
\(\theta_{n-1}\)-space has dimension \(m_\pi\), since adjoining
the last column translations generates \(U_n\) with the character
\(\theta_n\). Apply the induction hypothesis there and count all
the nonzero covectors. Restricting \(\pi\) to \(P_k^0\) gives the
asserted bound. Since \(|U_k|=q^{k(k-1)/2}\), the bound for the
scalar is \(\prod_{j=1}^{k-1}(1-q^{-j})^{-1}\).

For \(q\ge3\), the product of \(1-q^{-j}\) is at least
\(1-\sum_{j\ge1}q^{-j}=1-1/(q-1)\), by multiplying the factors
successively. Its reciprocal is at most \(1+1/(q-2)\le1+3/q\).
Tensoring the local bounds, also with the identity on the Hilbert
space of amplitudes, proves \eqref{eq:spectral126}. Alternatively,
in the off-diagonal part take absolute values and use the sharp
entry bound in Corollary~\ref{cor:content125}. This gives the first
inequality of \eqref{eq:sparse126}. Cauchy--Schwarz on the support
gives the second. Finally the local row sum is zero for \(k\ge2\).
In \eqref{eq:ind125} its sum over \(g\) contains
\(\sum_{u\in U_k}\theta_k(u)=0\). Tensoring with the other factors
proves the last assertion.
\end{proof}

In the source expansion, group the orbital amplitudes by their
relative residue bases \(i\), and let \(z_i\) be the sum in group
\(i\). After division by the value of the character at the identity,
the finite quadratic expression on these amplitudes is \(B\).
The grouping and its normalization appear in
\srcloc{Formalism/HermitianCurvature/StatementPB}{86} and
\srcloc{Formalism/LocalToGlobal/SiegelLatticeTheory/ClosedLocalOrbits/TermOpStmt}{135}.
The estimates above control the effect of the finite matrix on
these vectors. In particular a lower bound for \(S\) alone allows
a zero value for \(B\) when the amplitudes are constant in a local
index. For the actual projected operator \(P_u\), the source uses
an estimate for
\(\|c_u^{-1}P_u^*P_u-\mathrm{Id}\|_{\mathrm{HS}}^2\) on its
column space, with \(c_u>0\) the unfolding constant. It combines
the finite matrix with estimates for the degenerate and corank-one
orbital terms in
\srcloc{Development/AbelianVarietyTheory/Proofs/DecoratedLatticeProofs/DiagonalTheoremOfFourPointBounds}{49}.
Bounding those global sums in this manuscript and comparing the
resulting period with the prescribed integral theta pairing remain
necessary for the argument in Volume~I.

\subsection{The identity double coset at the common terminal level}\label{nav:VI-15}
\label{sec:identity127}

We compute the global identity term itself, before estimating the
other terms of the Gram matrix. Use \(B=N\rtimes H\subset G\) of
Section~\ref{sec:char94} over \(F_+=E^+\). The archimedean group
\(H_\infty\) is compact. The finite row level \(K\) is either
\(K_\sigma\) or \(K_\sigma'=a_fK_\sigma a_f^{-1}\) of
Section~\ref{sec:refined112}. Put \(K_N=K\cap N_f,\ K_H=K\cap H_f\),
where a subscript \(f\) denotes finite adeles. We have
\(B_f\cap K=K_N\rtimes K_H\). The characters \(\chi_1,\chi_\varpi\)
of Section~\ref{sec:char94} are trivial on \(K_N\) and on \(H\).
At the translated level the column classes for \(K_H\) include
the full refinement by a nonzero covector at every place over \(T\).

Use product Haar measures on \(N(\mathbb A_{F_+})\), and fix Haar
measures on \(G_f,G_\infty,H_\infty\). Write \(v_N\) for the volume
of \(N(F_+)\backslash N(\mathbb A_{F_+})\). This volume is finite
and positive by the compactness of the adelic unipotent quotient.
We use the quotient measure \(d\dot g\) on \(B_\infty\backslash
G_\infty\) defined by integration first over \(N_\infty,H_\infty\).
To see that this is an invariant quotient measure, \(N_\infty\)
is unipotent and \(H_\infty\) is compact, so both \(B_\infty\) and
the reductive group \(G_\infty\) are unimodular.

\begin{lemma}[A common unipotent covolume for all column classes]
\label{lem:covol127}
Let \(h_i\) run through representatives of
\(H(F_+)\backslash H_f/K_H\). The finite support \(B_fK\)
has these same representatives for the action of \(B(F_+)\) on
its right \(K\)-cosets. The rational stabilizer of \(h_iK\) is
\[
       \Lambda_i=N(F_+)\cap h_i K_Nh_i^{-1},
       \qquad
       \operatorname{vol}(\Lambda_i\backslash N_\infty)
              =\frac{v_N}{\operatorname{vol}_f(K_N)} .
\]
\end{lemma}
\begin{proof}
The group \(N(F_+)\) is dense in \(N_f\). Indeed the finite
exponential and logarithm polynomials for strictly upper block
triangular matrices identify this unitary unipotent group with
its Lie algebra as varieties in characteristic zero. The logarithm
of an isometry is skew adjoint, since taking adjoints commutes
with the logarithm and the logarithm of the inverse is its negative.
The same polynomials identify their finite adelic spaces, since they
are integral outside finitely many primes. Density now follows from
additive approximation in a vector space. In each coordinate one
may clear the finitely many prescribed denominators and solve the
resulting finitely many congruences in a fractional ideal by the
Chinese remainder theorem, remaining integral at all other places.

Given \(n hK\), first translate \(h\) by \(H(F_+)\) into \(h_iK_H\).
Then density supplies a rational unipotent element in the open
coset of \(h_iK_Nh_i^{-1}\) needed to remove the remaining
unipotent component. Conversely a rational identification of
two resulting cosets gives a rational identification of the
column classes on projection to \(H\), so their indices coincide.
Each column class has trivial rational stabilizer in \(H\).
For the unrefined levels this is Proposition~\ref{prop:column110},
and passage to a subgroup retains triviality. It follows that its
stabilizer in \(B(F_+)\) is exactly the subgroup displayed above.

For any compact open \(C\subset N_f\), density also gives
\(N_f=N(F_+)C\). A fundamental region for \(N(F_+)\) in the
adelic group is given, up to measure zero, by a real fundamental
region for \(N(F_+)\cap C\) times \(C\). This follows by reducing
the finite coordinate first and the real coordinate second.
Hence \(v_N=\operatorname{vol}((N(F_+)\cap C)\backslash N_\infty)
\operatorname{vol}_f(C)\). Take \(C=h_iK_Nh_i^{-1}\).
Conjugation by \(H_f\) preserves the Haar measure on \(N_f\).
In the Witt basis the only independent Lie algebra coordinates
on which \(h\in H\) acts nontrivially may be taken to be the
block from \(W_0\) to \(E^s\). The opposite block is its negative
adjoint. On the independent block \(h\) acts by precomposition
with \(h^{-1}\). Its determinant over \(F_+\) is
\(\operatorname{N}_{E/F_+}(\det_E h)^{-s}=1\). All other
independent coordinates are fixed. This computes the modulus
at each place and proves the common covolume formula.
\end{proof}

Let \(\rho\) be one of the prescribed irreducible unitary column
types on \(H_\infty\), with representation space \(V_\rho\) of
dimension \(d_\rho\). A column in class coordinates is a covector
\(\ell_i\in V_\rho^\vee\) for each \(h_i\), with the ordinary sum
inner product. Equivalently, write
\(\ell(\gamma_f h_i k)=\ell_i\circ\rho(\gamma_\infty)^{-1}\)
for \(\gamma\in H(F_+)\), \(k\in K_H\). Triviality of the
stabilizers makes this well defined.
For \(\beta\in\{1,\varpi\}\) take a smooth real-place function
\(W:G_\infty\to V_\rho\) satisfying
\begin{equation}
 W(nhg)=\chi_{\beta,\infty}(n)\rho(h)W(g),\qquad
 \int_{B_\infty\backslash G_\infty}\|W(g)\|^2\,d\dot g=1 .
 \label{eq:real127}
\end{equation}
The calculation applies to every such \(W\). For existence with
these two properties alone, choose a compactly supported smooth
section on a coordinate ball of the quotient where the bundle
with fibre \(V_\rho\) is trivial. Extend it by zero and normalize
its norm. For a specified automorphic representation the choice of
its real-place function is additional data.

Define the seed \(\Phi_{\beta,\rho,\ell}\) on \(G(\mathbb A_{F_+})\)
by zero outside \(G_\infty\times B_fK\) and by
\begin{equation}
     \Phi_{\beta,\rho,\ell}(g_\infty,n_fh_fk)
            =\chi_{\beta,f}(n_f)\ell(h_f)(W(g_\infty))
            \quad(k\in K)
 \label{eq:seed127}
\end{equation}
on its support. Independence of the finite decomposition follows
from \(B_f\cap K=K_N\rtimes K_H\), triviality on \(K_N\), and
\(H_f\)-invariance of the character of \(N_f\). Under left
translation by a rational element of \(H\) the two occurrences
of \(\rho\) cancel. Under a rational element of \(N\) the finite
and infinite characters multiply to one. Thus the seed descends
to \(B(F_+)\backslash G(\mathbb A_{F_+})\). Give a column type
and a frequency one fixed \(W\) when comparing different columns.

\begin{theorem}[The normalized global identity term]
\label{thm:identity127}
In the integral of \(\Phi_{\beta,\rho,\ell}\,
\overline{\Phi_{\beta',\rho',\ell'}}\) over
\(B(F_+)\backslash G(\mathbb A_{F_+})\) the answer is zero
if \(\beta\ne\beta'\) or \(\rho\not\simeq\rho'\). For matching
data the answer is
\begin{equation}
 c_{\rho,K}\sum_i(\ell_i,\ell_i'),\qquad
 c_{\rho,K}=
 \frac{\operatorname{vol}_f(K)\,v_N\,\operatorname{vol}(H_\infty)}
      {\operatorname{vol}_f(K_N)\,d_\rho}.
 \label{eq:main127}
\end{equation}
Here the Hermitian pairing on covectors is linear in its first
entry. After division by \(c_{\rho,K}>0\) the operator is the
identity also on the refined subspace with zero sum on each
scalar orbit. For the common words of Theorem~\ref{thm:common111},
distinct words have zero mutual identity term.
\end{theorem}
\begin{proof}
Reduce the finite coordinates using Lemma~\ref{lem:covol127}.
On the coset \(h_iK\) the product of seeds is independent of the
right coordinate \(k\). Its finite integration contributes
\(\operatorname{vol}_f(K)\) and its real integral is over
\(\Lambda_i\backslash G_\infty\). Indeed one may choose a
fundamental region there in the real variable and the full coset
in the finite variable. The residual rational stabilizer acts by
measure-preserving permutations of that coset.

Disintegrate the real integral in the order
\(\Lambda_i\backslash N_\infty,\ H_\infty,\
B_\infty\backslash G_\infty\), using the quotient integration
formula twice. If \(\beta\ne\beta'\), the ratio of their characters
descends to \(\Lambda_i\backslash N_\infty\). For a rational point
of \(\Lambda_i\) its finite value is one by \(H_f\)-invariance and
triviality on \(K_N\), and the product of all local values is one.
The ratio is nontrivial on \(N_\infty\) since its terminal trace
coefficient is \(\beta-\beta'\ne0\). Translating the integration
variable on the right by an element with nontrivial ratio shows
the integral is zero. For equal frequencies the unipotent integral
is the covolume in Lemma~\ref{lem:covol127}.

For fixed \(\dot g\), averaging the resulting rank-one operator
over \(H_\infty\) gives an intertwiner from \(\rho'\) to \(\rho\).
It vanishes for inequivalent irreducibles by Schur's lemma.
For equal data it is scalar, and its trace is
\(\operatorname{vol}(H_\infty)\|W(g)\|^2\). Its scalar is this
trace divided by \(d_\rho\). Integrating over the last quotient
in \eqref{eq:real127} and summing over \(i\) proves
\eqref{eq:main127}. These integrations are absolutely convergent.
For the square of one seed this follows by nonnegative integration
and the displayed finite factors. For a mixed pair use the two
square integrals and Cauchy--Schwarz. Orthogonal compression of
the identity to any subspace is its identity. Finally the words
have distinct column highest weights, because any differing
interior bit changes the corresponding weight by \(\Delta>0\).
Their column types are consequently inequivalent.
\end{proof}

The same identity-term formula holds with a common unitary
character \(\kappa\) on the finite row level whenever
\(\kappa|_{K_N}=\chi_{\beta,f}|_{K_N}\) for the frequencies in
use. In that case include \(\kappa(k)\) in \eqref{eq:seed127} and
the factor \(\kappa(k_H)\) in the right transformation law of
\(\ell\). These factors make the two decompositions agree, and at
each \(h_iK\) the common factor and its conjugate multiply to one.

\begin{proposition}[The change of normalization at \(T\)]
\label{prop:vol127}
For \(r=2s+k+1\) the actual levels of Section~\ref{sec:refined112}
satisfy
\begin{equation}
 [N_f\cap K_\sigma:N_f\cap K_\sigma']=T^{\,fs(r-s)},\qquad
          c_{\rho,K_\sigma'}=T^{\,fs(r-s)}c_{\rho,K_\sigma}.
 \label{eq:unipindex127}
\end{equation}
Both identities apply at every selected \(x,T,E\) and to every
chosen system of nonzero terminal covectors.
\end{proposition}
\begin{proof}
Away from \(T\) the two levels agree. In the first split factor at
\(v\mid T\) the completed Witt flag has block dimensions
\((1,\ldots,1,k+1,1,\ldots,1)\), with \(s\) lines on each side.
The chosen split model identifies \(N_v\) with the group of upper
block unitriangular matrices for this list. In fact the flag and
its Hermitian orthogonal flag specify the two strings of lines,
and the action on all the successive quotients is the identity.
The action on the other split factor is then determined by duality.
At the original level its upper blocks are arbitrary integral
matrices. The level change \(a_v\) multiplies the first \(s\)
lines by \(T\), acts on the middle block of size \(d=k+1\)
by the matrix \(M:(w,z)\mapsto(w,\ell_v(w)+Tz)\), and fixes the
last \(s\) lines. In particular \(\det M=T\), and both \(M\)
and \(TM^{-1}\) are integral.

Conjugating by \(a_v\) multiplies an upper block \(X_{ij}\) by
\(a_i\) on the left and by \(a_j^{-1}\) on the right. The
preceding integrality shows that every resulting block is still
integral. Product additive measure in the independent upper
blocks is Haar measure, since group translation is triangular
in increasing block distance with identity diagonal. For a block
of size \(d_i\times d_j\), the determinant of its linear change
is \((\det a_i)^{d_j}(\det a_j)^{-d_i}\). The blocks with row in
the first string and column in the middle contribute valuation
\(s(d-1)\). The blocks with row in the middle and column in the
last string contribute \(s\), and the positions between the two
strings of lines contribute \(s^2\). Within either string the
valuation is zero. The total is \(s(d+s)=s(r-s)\). With
\(\lvert T\rvert_v=T^{-1}\) this gives the local subgroup index
\(T^{s(r-s)}\). All \(f\) places have this residue size and their
indices multiply. Finally the finite row levels are conjugate in
the unimodular group \(G_f\), so their volumes agree. Substitution
into \eqref{eq:main127} gives the second formula.
\end{proof}

In the Gram matrix of the Poincar\'e sums over
\(B(F_+)\backslash G(F_+)\), the terms with identical rational
cosets unfold to the integral in Theorem~\ref{thm:identity127}.
For a square this identity contribution unfolds by nonnegative
integration, by tiling from a fundamental region for \(G(F_+)\).
Its mixed version follows by polarization, or by the absolute
integrability of the two matching-seed products. This identifies
the main term even when estimates for interchanging the other
orbital sums have still to be proved. The seed and the normalization
occur in \srcloc{Formalism/LocalToGlobal/SiegelLatticeTheory/HaarTransversals/PoincareInHH}{163}
and line~257 of the same file. The identity term alone gives no
lower bound for the full squared norm. Terms from distinct rational
cosets can have either sign in its real expansion. Their control
for the prescribed real theta vectors is still required.

\section{The common terminal prime and its characters}\label{nav:g-VIII-1}

\subsection{Integral theta test functions at the common split prime}\label{nav:VII-10}
\label{sec:split161}

Use the prime \(T\), the lattices \(\Lambda_1\subset\Lambda_0\)
and the row levels of Proposition~\ref{prop:index159}. We work on
the first pair \(Y_0,Y_1\) of the indefinite theta tower. Its ranks
are \(r\) and \(w=r+t_{\rm r}\), respectively, by
\srcloc{Formalism/LocalRingCardinality/Params}{40}. Fix a
place \(v\mid T\), put \(O=\ZZ_T\), and let \(\mathbf k=\mathbf F_T\).
The split Hermitian spaces have two factors, identified with a
vector space and its dual by the Hermitian pairing. On the first
row factor the rational element \(a=a_\varpi\) carries
\(\Lambda_0\) onto \(\Lambda_1\). Its determinant has valuation
\(s\). On the second factor it acts by the contragredient with
determinant of valuation \(-s\).

Choose one factor at each \(v\) and write \(\epsilon_v=1\) for
the first and \(\epsilon_v=-1\) for the second. In the associated
matrix model a position is a map \(Y\) from the chosen
\(r\)-dimensional row space to the \(w\)-dimensional partner.
Choose an integral partner frame, identifying its lattice with
\(O^w\). On the chosen row factor write \(g_\epsilon\) for the
action of \(g\). With the matrix splitting the row acts by
\begin{equation}
   \omega_v(g,1)\phi(Y)
       =|\det g_\epsilon|_T^{\,w/2}\phi(Yg_\epsilon).
   \label{eq:weil161}
\end{equation}
Indeed the split symplectic space is the sum of this matrix
space and its trace dual. Right multiplication by \(g_\epsilon\)
on \(w\) rows has Jacobian \(|\det g_\epsilon|_T^w\), giving the
factor for a unitary pullback. This is the split Weil action of
Kudla, Chapter III, Remark 6.3~\cite{XVIII-XIII-IV-kudla29}. Another splitting
multiplies the formula by a unitary character \(\eta_v(g)\),
since the two operators lift the same symplectic action. We
retain this character when passing to the global kernel.

Let \(\mathcal S_v\) be the set of \((r-s)\)-planes in
\(\Lambda_0/T\Lambda_0\). For \(S\in\mathcal S_v\) let
\(L_S\) be its inverse image in \(\Lambda_0\) and put
\(D_S=L_S\) in the first factor, \(D_S=L_S^\vee\) in the second.
Choose an ordered basis \(b_S\) of \(D_S\). Two kinds of
Schwartz functions occur in the argument. The first is the
lattice vacuum
\[
    v_S(Y)=\mathbf1_{\{Y(D_S)\subseteq O^w\}}.
\]
In this position model, invariance under translations by
\(\operatorname{Hom}(D_S,O^w)\) forces the vacuum to be constant
on its cosets. Invariance under multiplication by characters
from its dual lattice forces it to vanish off that lattice,
since the additive character has conductor \(O\). Its value
one at zero therefore gives the stated function. For the second
kind fix a character \(\chi:\mathbf k^\times\to\CC^\times\).
On integral matrices in the frame \(b_S\) write
\[
     \overline{Yb_S}=\binom{C}{D},\qquad
        C\in M_r(\mathbf k),\quad D\in M_{w-r,r}(\mathbf k).
\]
Define \(p_S(Y)=\chi(\det C)\) if \(Yb_S\) is integral,
\(D=0\) and \(C\) is invertible, and zero otherwise. These
are the three residue conditions on a local factor at a
place in the extra level set \(Q\) in
\srcloc{Formalism/LocalToGlobal/IndefiniteReal/DatumVec}{66}.
The power or inverse of the character there is included in
\(\chi\). Changing \(b_S\) multiplies \(p_S\) by a root of
unity. The vacuum and this residue function are treated in
their indicated frames.

\begin{theorem}[An integral left inverse on the orbit of lattices]
\label{thm:eval161}
Choose either kind of function at each \(v\mid T\), keeping
that choice throughout its orbit. Write \(\phi_{\boldsymbol S}\)
for their product and \(A\) for the ring of integers of a
number field containing all the residue character values.
The map from the free \(A\)-module on \(\prod_{v\mid T}\mathcal S_v\)
to the compactly supported locally constant \(A\)-valued functions
given by \([\boldsymbol S]\mapsto\phi_{\boldsymbol S}\) has an
\(A\)-linear left inverse. The ideal generated by the values
of \(\sum_{\boldsymbol S}c_{\boldsymbol S}\phi_{\boldsymbol S}\)
is exactly \((c_{\boldsymbol S}:\boldsymbol S)\). At a residue
place the supports belonging to distinct \(S\) are disjoint.
\end{theorem}
\begin{proof}
Take \(Y_S\) to send the ordered basis \(b_S\) to the first
\(r\) unit vectors of \(O^w\). The inequality \(w\ge r\)
gives such a linear map, and \(Y_S^{-1}(O^w)=D_S\).
All the \(D_{S'}\) have the same covolume, also when duals
are used. If \(Y_S(D_{S'})\subseteq O^w\), containment
\(D_{S'}\subseteq D_S\) and equality of covolumes imply
\(D_{S'}=D_S\), hence \(S'=S\). Both test functions at \(S\)
take the value one at \(Y_S\). Evaluation at the product
of these matrices therefore recovers exactly the corresponding
coefficient. Conversely every function value is an \(A\)-linear
combination of the coefficients, proving the assertion on ideals.
For a point of the support of \(p_S\), the top square block of
\(Yb_S\) lies in \(\mathrm{GL}_r(O)\). Its inverse shows directly
that \(Yy\in O^w\) implies \(y\in D_S\). Thus on that support
\(Y^{-1}(O^w)=D_S\), proving disjointness.
\end{proof}

\begin{theorem}[The Gram matrix of the finite theta vectors]
\label{thm:gram161}
Give each test vector its unit norm in the complex \(L^2\)
space of positions with additive Haar measure. If \(S,S'\) at a
vacuum place meet in codimension \(j\) in \(S\), their inner
product equals \(T^{-jw}\). At a residue place the Gram matrix
is the identity. Before normalization the ratio of the squared
norm of \(p_S\) to that of \(v_S\) is
\begin{equation}
      T^{-r(w-r)}\prod_{h=1}^{r}(1-T^{-h}).
      \label{eq:prvol161}
\end{equation}
For the product of all \(f=[E^+:\QQ]\) places, the Gram matrix
\(\mathcal G_T\) of the normalized vectors satisfies
\[
    \|\mathcal G_T-1\|_{\rm op}
      \le (1+32T^{-(t_{\rm r}+1)})^{f}-1.
\]
\end{theorem}
\begin{proof}
The intersection of two vacuum supports consists of the maps
integral on \(D_S+D_{S'}\). This sum contains \(D_S\) with
index \(T^j\). For the first choice of factor this follows by
reducing modulo \(T\), and for the dual choice it follows from
\(L_S^\vee+L_{S'}^\vee=(L_S\cap L_{S'})^\vee\). The dual of
the sum on the left is \(L_S\cap L_{S'}\), so the identity
follows by biduality for lattices over \(O\). In each of the
\(w\) rows, integrality on the larger lattice cuts the volume
by \(T^{-j}\). The vacuum norms are equal along the orbit,
which proves the inner product formula. Theorem~\ref{thm:eval161}
proves orthogonality at a residue place. On one vacuum support
the reduction of the matrix is uniform. The bottom
\((w-r)r\) entries must vanish and the first \(r\) columns of
the square block must be linearly independent. Counting their
successive choices gives \eqref{eq:prvol161}.

For fixed \(S\), the number of planes at distance \(j\) is
\[
       {r-s\brack j}_T{s\brack j}_T T^{j^2}.
\]
Choose the intersection in \(S\), the \(j\)-dimensional image
in the quotient by \(S\), and then a linear map from that image
to the \(j\)-dimensional quotient of \(S\) by the intersection.
Its graph determines the remaining choice. The bound
\({a\brack j}_T\le4T^{j(a-j)}\) follows by removing the leading
powers in the product of Proposition~\ref{prop:index159}, since
the resulting denominator is at least
\((1-1/2)(1-\sum_{h\ge2}2^{-h})=1/4\).
The off-diagonal row sum is thus bounded by
\[
     16\sum_{j\ge1}T^{-j(w-r+j)}\le32T^{-(t_{\rm r}+1)}.
\]
Absolute row sums multiply for the product of local Gram
matrices. Removing the diagonal leaves row and column sums
bounded by the claimed expression. Cauchy--Schwarz with these
two bounds gives the operator norm bound.
\end{proof}

\begin{proposition}[The exact denominator after finite averaging]
\label{prop:exact161}
Let \(\boldsymbol S_*\) correspond to \(\Lambda_1\) at all places.
On \(\prod_{v\mid T}K_{0,v}\), average the translates of
\(\phi_{\boldsymbol S_*}\) in the character by which its
stabilizer acts, extended by the determinant and the chosen
splitting. The least positive integer making this function
\(A\)-valued is \(I={r\brack s}_T^f\). One has \(I\equiv1\bmod T\).
\end{proposition}
\begin{proof}
Write \(\chi=1\) at a vacuum place. The local stabilizer is
taken here to be the stabilizer of the label \(S_*\). It acts on
the vector by \(\eta_v(k)\chi(\overline{\det k_\epsilon})\).
Indeed its determinant is a unit, and changing the domain
basis multiplies the top determinant by that unit modulo
\(T\). This gives the required character on all of \(K_{0,v}\).
Multiply each translate in the average by the inverse of this
character. The corrected translate depends only on the coset
of the stabilizer. There is one term for every tuple of planes,
because reduction acts transitively, with exactly \(I\) cosets
by Proposition~\ref{prop:index159}. Each term is its indicated
test function times a root of unity in \(A\) after cancellation
of the splitting factor. Hence multiplication by \(I\) gives integral
values. At a tuple of the matrices in Theorem~\ref{thm:eval161}
the value is an \(A\)-unit divided by \(I\). If an integer \(b\)
clears this denominator, \(b/I\in A\cap\QQ=\ZZ\). Finally all
factors \(T^a-1\) in the Gaussian product are units congruent
to \(-1\bmod T\). This proves the stated congruence.
\end{proof}

\subsection{Integral theta tests for the depth filtration}\label{nav:VII-13}
\label{sec:integraldepth165}

In the original-level formula \eqref{eq:geom163} one may also
discard every rational position stabilizing the full isotropic
space \(A_s=\sum_{i=1}^s Eu_i\). Indeed let \(Q_s\) be its
stabilizer. The sum of all terms indexed by \(Q_s(F)\) is the
pairing of the two intermediate sums over \(B(F)\backslash
Q_s(F)\). This follows by the same unfolding as in
Theorem~\ref{thm:ordinary163}, since the integral of the
absolute majorant for this subseries is bounded by that for
the full series. At \(K^0\) Theorem~\ref{thm:midnorm131}
evaluates this pairing as the seed pairing, including on the
refined columns. Thus the entire deviation at the original
level comes from pairs with distinct isotropic spaces. For a
relevant such pair, Lemma~\ref{lem:angle149} also imposes
the lower bound there on the product of the two frame heights.

We now compute the averaging index in
Proposition~\ref{prop:average163} at the actual terminal lattice.
This also gives a realization of its finite subgroup by matrix
coefficients on the first theta pair. At \(v\mid T\) take the
split factor where \(\varpi\) has valuation one. The residue size
is \(T\), and \(O=\ZZ_T\). Choose the basis of the level lattice
\(\Lambda=O^s\oplus M_v\oplus O^s\) in Section~\ref{sec:middle131},
so \(K=K^0_v=\GL_r(O)\) in these coordinates. Write \(J_j\) for
its principal congruence subgroup of depth \(j\ge1\), and put
\[
    I_j=[K:(B_v\cap K)J_j],\qquad
    D_B=\frac{r(r+1)-k(k-1)}2 .
\]
Using \(\Lambda\) together with its dual in the other split
factor gives the same congruence subgroup in the underlying
\(F_v\)-linear action of Section~\ref{sec:depth163}.

\begin{theorem}[The exact terminal index]\label{thm:index165}
For \(j\ge1\),
\[
    I_j=T^{(j-1)D_B}(T^k-1)\prod_{h=k+1}^{r}(T^h-1).
\]
In particular \(v_T(I_j)=(j-1)D_B\), and
\(0\le jD_B\log T-\log I_j\le 4T^{-k}\).
\end{theorem}
\begin{proof}
In the split factor the parabolic has \(s\) line blocks on either
side of a middle block of size \(k+1\). Its unitriangular radical
in the adapted basis has independently chosen entries above the
diagonal blocks. There are
\[
               a=\frac{r^2-(k+1)^2-2s}{2}
\]
such entries. Reduction of its integral subgroup modulo \(T^j\)
therefore has order \(T^{ja}\). Theorem~\ref{thm:midlevel131}
identifies the rest of \(B_v\cap K\) with the integral
\(k\)-by-\(k\) matrices \(h\in\GL_k(O)\) satisfying
\(\ell h=\ell\bmod T\), for the primitive row \(\ell=\ell_v\).
In middle coordinates \((w,z)\mapsto(w,\ell(w)+Tz)\) their
matrices are
\[
                  \begin{pmatrix}h&0\\(\ell-\ell h)/T&1\end{pmatrix}.
\]
To reduce to the identity at depth \(j\) therefore requires
\(h=1+T^j X\) and \(\ell X=0\bmod T\). The map \(1+T^jX\mapsto
\ell X\bmod T\) is a surjective homomorphism on the principal
congruence subgroup of \(\GL_k(O)\), since \(j\ge1\) and
\(\bar\ell\ne0\). Its kernel has index \(T^k\). The original
condition on \(h\) has index \(T^k-1\) in \(\GL_k(O)\), by
transitivity on nonzero residue rows. The number of reduced
middle matrices is thus
\[
         |\GL_k(O/T^j)|\,T^k/(T^k-1).
\]
The radical and middle images intersect trivially by their
diagonal blocks and give the whole image. Finally
\[
 |\GL_n(O/T^j)|=T^{(j-1)n^2+n(n-1)/2}
                           \prod_{h=1}^n(T^h-1).
\]
One obtains this by counting ordered independent columns over
the residue field and lifting all entries freely. Dividing
the order for \(n=r\) by the two image orders gives the formula since
\(r(r-1)/2-k(k-1)/2=a+k\) and \(r^2-k^2-a=D_B\).
Also \(D_B=k+\sum_{h=k+1}^r h\). The logarithmic difference
is \(-\sum_{h=k}^r\log(1-T^{-h})\), bounded by
\(2\sum_{h\ge k}T^{-h}\le4T^{-k}\).
\end{proof}

Use a partner lattice basis on the first theta edge, whose
partner rank \(w=r+t_{\rm r}\) is at least \(r\). Write a
position matrix as \(Y=\binom{X}{Z}\), with \(r\)-by-\(r\)
top block \(X\). Let \(G_j=\GL_r(O/T^j)\), let \(B_j\) be the
image of \(B_v\cap K\) there, and let \(a\) range over the left
cosets \(B_j\backslash G_j\). The coarse residue support is
\[
          X\in K,\qquad Z\in T M_{w-r,r}(O).
\]
This is the residue condition of Section~\ref{sec:split161} in
a basis of \(\Lambda\). Let \(\kappa\) be the desired smooth
unitary level character and let \(\zeta\) be the splitting
character in the matrix realization \eqref{eq:weil161}.
On the coarse support put \(\theta(X)=\kappa(X)/\zeta(X)\)
and extend by zero outside to obtain a test \(\phi_*\).
Without either character this is the characteristic function.
The fine test \(\phi_{a,j}\) is its restriction to
\(X\bmod T^j\in a\). Write \(\phi_{1,j}\) for the identity
coset and let \(A\) be the ring of integers of a number field
containing the finite character values. The values on
\(K\) of each smooth character are roots of unity because
the open kernel there has finite index.

\begin{theorem}[A matrix coefficient and an integral trace]
\label{thm:tests165}
The \(\phi_{a,j}\) have disjoint supports, have equal squared
norm \(\|\phi_*\|^2/I_j\), and span a free \(A\)-module of rank
\(I_j\) with an \(A\)-linear left inverse given by evaluations.
With inner product linear in the first variable,
\[
 \frac{\langle\omega(g)\phi_{1,j},\phi_{1,j}\rangle}
      {\|\phi_{1,j}\|^2}
    =\begin{cases}
       \kappa(g),&g\in (B_v\cap K)J_j,\\
       0,&\text{otherwise}
     \end{cases}
          \qquad (g\in\GL_r(F_v)).
\]
On these tests the \(\kappa\)-average satisfies
\(I_j e_K\phi_{1,j}=\phi_*\). Its exact denominator over \(A\)
is \(I_j\). The same \(\phi_*\) occurs at every depth.
\end{theorem}
\begin{proof}
Choose a lift \(x_a\in K\) of one member of each coset. Evaluation
at \(\binom{x_a}{0}\), followed by multiplication by
\(\theta(x_a)^{-1}\in A^\times\), extracts precisely the
coefficient at \(a\). Formula \eqref{eq:weil161} gives
\(\omega(h)\phi_{a,j}=\kappa(h)\phi_{a\bar h^{-1},j}\) for
\(h\in K\), where \(\bar h\) denotes reduction modulo \(T^j\).
These unitary translations permute all the supports transitively.
Their disjoint union is the entire coarse support, proving the
norm formula. If both \(Y\) and \(Yg\) are in the identity
support, their top blocks \(X\) and \(Xg\) are in \(K\).
Hence \(g\in K\). The two tests then have disjoint supports
unless \(\bar g\in B_j\), and in the latter case the displayed
transformation law gives the coefficient. This proves the
formula everywhere. Averaging with \(\overline{\kappa(h)}\)
gives the sum of all \(I_j\) translates with equal coefficient
\(I_j^{-1}\). At the matrix for the identity coset its value
is a unit divided by \(I_j\). A positive integer \(m\) clears
this value only if \(m/I_j\in A\cap\QQ=\ZZ\). Conversely
multiplication by \(I_j\) restores the integral coarse test.
\end{proof}

\begin{corollary}[A common coarse theta period at all depths]
\label{cor:trace165}
For each \(v\mid T\) choose a depth \(j_v\ge0\), taking the
coarse test at depth zero and \(\phi_{1,j_v}\) otherwise.
Let \(I=\prod_{v:j_v>0}I_{j_v}\). Keep all the other theta
data fixed. Suppose the theta rows \(F_{\mathbf j}\) and \(F_*\)
of the fine and coarse tests are given by convergent lifts in
automorphic \(L^2\). Let \(S_{\mathbf j}\) be the Bessel seed
periodization of Section~\ref{sec:depth163} with the additional
indicators at all the places of positive depth. For the product
\(e\) of their level-character averages one has
\begin{align}
    I e F_{\mathbf j}=F_* ,\qquad I e S_{\mathbf j}=S_0,\qquad
    \langle F_*,S_0\xi\rangle
            =I^2\langle eF_{\mathbf j},eS_{\mathbf j}\xi\rangle .
    \label{eq:pair165}
\end{align}
Moreover, for fixed arithmetic data,
\[
  \big\| I S_{\mathbf j}^* S_{\mathbf j}-1\big\|_{\rm op}
                \le C_x T^{-\eta\sum_v j_v},
               \qquad \eta=L-4fr^2 .
\]
\end{corollary}
\begin{proof}
Right translation of a theta row by a finite row element \(h\)
replaces the test by \(\omega(h)\) applied to it. This follows
by substituting \(\omega(g)\omega(h)=\omega(gh)\) in the
rational theta sum, which commutes with the independent
partner action. The average on each finite orbit is a finite
sum after character correction. Linearity of the convergent
lift and Theorem~\ref{thm:tests165} prove the first equality.
For Bessel seeds use Proposition~\ref{prop:average163} at
each place. The averages commute since they act on different
factors. Taking the inner product gives \eqref{eq:pair165}.
The same calculation holds after the real row projection,
which commutes with the finite averages.

For the last estimate repeat the norm calculation of
Lemma~\ref{lem:deep163} on its single nonzero affine frame
coordinate \(z\in F\). At every place \(v\) of positive depth
the same \(z\) has valuation at least \(j_v-O_x(1)\). Clearing
one common denominator and taking its field norm gives
\(T^{\sum_v j_v}\le C_x V(g)V(\gamma_\infty g)\) for every
nonidentity coset meeting both supports. All constants may be
taken over the full finite set of places above \(T\), independent
of the depths. Disintegration over the finite class cells gives
volume fraction \(I^{-1}\) on the identity term. The dyadic
tail bound and the integrable first height moment in the proof
of Theorem~\ref{thm:deep163} give the displayed operator bound.
\end{proof}

For the averaging factor in \eqref{eq:pair165} one has
\(\log I^2\le 2 D_B\log T\sum_v j_v\). If the total depth is
bounded by a fixed power of \(x\), this cost is \(o(b/x^A)\)
for every fixed \(A>0\) on the intervals of Volume~I, since
\(D_B,\log T\) are polynomially bounded and \(b=\exp(x^\eps)\).
The theorem gives the denominator of the finite test and the
exact trace on convergent theta lifts. It leaves the arithmetic
normalization and nonzero value of \(\langle F_*,S_0\xi\rangle\)
on the prescribed full chain to be established.

\subsection{The size of the common terminal prime}\label{nav:VII-14}
\label{sec:terminal209}

Write \(m=2d,\ r=k+m-1,\ b=\exp(x^\eps)\), and let
\(Q=g_1g_2\). Use the actual words, spacings and common
terminal prime from Theorem \ref{thm:common111}. Thus \(m\asymp x^{1/2-\eps/16}\),
\(k\asymp x^{9/2+10\eps}\), \(\log Q=(16+o(1))\log b\),
and \(f=[E^+:\QQ]\). The field \(E\) is abelian and has conductor
dividing \(N\) with \(\log(2N)=o(\log b)\) throughout the selected
intervals. Indeed the conductor bound in
Lemma \ref{lem:delay197}
is \(O_{\psi,\eps}(\log(2x)+x^{\eps/5}\log^2(2x))\) already
with initial delay one. All limits in this section are uniform in
the words and translates.

\begin{lemma}[The mass in two adjacent ranges of ranks]\label{lem:mass209}
Let \(\mathfrak m_n\) denote the mass of the standard positive
definite Hermitian genus over \(E/F\), \(F=E^+\), in rank \(n\).
If \(\eta\) is the quadratic character of \(E/F\), then
\begin{equation}
 \frac{\mathfrak m_r}{\mathfrak m_k}
    =2^{-f(r-k)}\prod_{j=k+1}^{r}|L_F(1-j,\eta^j)|,\qquad
 \frac{1}{fk}\log\frac{\mathfrak m_r}{\mathfrak m_k}
                           =o(m\log b).                   \label{eq:mass209}
\end{equation}
At an odd prime \(p\) the valuation of the ratio is also
\[
 \sum_{j=k+1}^r
 \left(v_p\big|K_{2j-2}(\mathcal O_E)_{\rm odd}^{\,c=(-1)^j}\big|
                   -v_p w_j(J_j)\right),
\]
where \(c\) is complex conjugation, \(J_j=F\) for even \(j\)
and \(J_j=E\) for odd \(j\), and
\(w_j(J)=|H^0(J,(\QQ/\ZZ)(j))|\).
\end{lemma}
\begin{proof}
Complex conjugation acts nontrivially on the two imaginary quadratic
subfields of relatively prime discriminants \(-q_0,-d_0\).
At every finite place inertia fixes one of these subfields. Thus
\(E/F\) is unramified at all finite places. The identity Hermitian
lattice is maximal for its integral norm condition. In fact if
an integral lattice contains it and \(y\) belongs to that lattice,
polarizing the norm against the multiples of a standard basis
vector shows that \(\operatorname{Tr}_{E/F}(a y_i)\) is integral
for all \(a\in\mathcal O_E\). The relative different is one,
so \(y_i\in\mathcal O_E\) for every coordinate, also locally
above two.

Apply the mass formula of Gan, Hanke and Yu for maximal definite
Hermitian lattices, Propositions 4.4 and 4.5 and (4.2)--(4.3)
of \cite{XVIII-XIII-IV-ghy209}. Its factor at a split place is one. At an
unramified nonsplit place the signed determinant of the standard
form is a unit, hence a norm. Indeed the norm on finite residue
fields is surjective and one corrects successive unit digits
by the surjective residue trace in the unramified quadratic
extension. Consequently all finite correction factors are one,
in both parities of the rank. The formula is
\(\mathfrak m_n=2^{1-fn}\prod_{j=1}^n |L_F(1-j,\eta^j)|\).
Its mass convention gives the maximal compact group volume one
at infinity and each standard lattice stabilizer volume one at
the finite places, as in the definition of \(\mathfrak m_n\).
Division proves the product in \eqref{eq:mass209}.

For even \(j\) the factor is \(\zeta_F(1-j)\), and for odd \(j\)
it is the value of \(\zeta_E/\zeta_F\) at \(1-j\), taken by
continuation. In either case it is a product of \(f\) primitive
Dirichlet values of the parity of \(j\) with conductors at most
\(N\), by factorization for this abelian extension. Each has
absolute value at most \(C(j-1)!\,N^j\). This follows from
the Bernoulli polynomial expression at \(1-j\) and the bound
\(|B_j(y)|\le Cj!\) for \(0\le y\le1\), obtained by Cauchy's
estimate on the unit circle for
\(t e^{yt}/(e^t-1)\). These are nonzero values by the primitive
functional equation at \(j>1\). Thus the log of the mass ratio
is bounded above by \(C f(r-k)r\log(2rN)\).
For a lower bound, append \(r-k\) ordered orthonormal vectors
to each rank-\(k\) lattice. A standard-genus lattice of rank
\(r\) has at most \(rw_E\) choices for a norm-one vector, with
\(w_E=|\mu(E)|\). Indeed two such vectors are orthogonal or
generate the same line with root-of-unity ratio, by integrality
and Cauchy--Schwarz at all embeddings. Removing the ordered
frame recovers its orthogonal complement. Counting with reciprocal
automorphism orders gives
\(\mathfrak m_k\le(rw_E)^{r-k}\mathfrak m_r\), as in
Lemma \ref{lem:mass90}.
Here \(w_E\le8f^2\) by the elementary totient bound. Now
\(\log r=O(\log x)\), \(r/k\to1\) and \(\log N=o(\log b)\).
Dividing the two bounds by \(fk\) proves the second assertion.
For the valuation formula, apply the critical-sign higher
class-number formula at weight \(j\) to each factor of the
first identity, as proved in
Section \ref{nav:III-13}.
\end{proof}

\begin{lemma}[The contribution of the extra row parameters]
\label{lem:weyl209}
For the word parameters \(u_i\) and row parameters \(\lambda_i\)
at every real place,
\[
       \frac{1}{k}\log\frac{W_r(\lambda)}{W_k(u)}
                             =(1+o(1))\,m\log b .
\]
\end{lemma}
\begin{proof}
Among the row entries are \(k-1\) numbers
\(v_i=4\Delta(k-1-i)\), each between \(u_i\) and \(u_{i+1}\).
The other \(m\) entries are \(-b_j\) with
\(b_j=B_\tau+2(j-1)H_0\). In the notation of the word
construction \(X'=\exp(x^{\eps/3})\) and the smallest base
weight is at least \(b/(nX')\), where \(n\) is bounded by a
fixed power of \(x\). The largest base weight is \(O(xb)\).
The two spacings \(\Delta,H_0\), the ranks and the parity
adjustment all have polynomial size. Hence uniformly for every
pair \(i,j\),
\[
      \log(v_i+b_j)=(1+o(1))\log b.
\]
There are exactly \(m(k-1)\) such factors. The other new factors,
between pairs of the \(-b_j\), are between one and \(2mH_0\).

Write \(\Delta(v)\), \(\Delta(u)\) for the positive absolute
Vandermonde products. The inequality
\(v_i-v_j\le u_i-u_{j+1}\) for \(i<j\) matches each factor
of \(\Delta(v)\) with a distinct factor of \(\Delta(u)\).
All unused factors are at least one. Conversely, for
\(2\le i<j\le k-1\), \(u_i-u_j\le v_{i-1}-v_j\).
Only \(2k-3\) factors of \(\Delta(u)\) involve an endpoint.
They are bounded by the full column span, at most
\(3\lceil xb\rceil\). Thus
\[
 -(2k-3)\log(3\lceil xb\rceil)\le
           \log(\Delta(v)/\Delta(u))\le0 .
\]
Finally the quotient of the Weyl factorial denominators contributes
\(-\sum_{i=k}^{r-1}\log(i!)\), of absolute value \(O(km\log k)\).
After division by \(k\), the contributions from the last three
estimates are \(O(m^2\log(2mH_0)/k+\log b+m\log k)\).
This is \(o(m\log b)\), since \(k\gg m^2\), \(m\to\infty\)
and \(\log x=o(\log b)\). The \(m(k-1)\) mixed factors give
the stated leading term.
\end{proof}

\begin{theorem}[Asymptotic for the actual common terminal prime]
\label{thm:prime209}
For all the selected words over the selected fields and intervals,
the prime in the original finite theta and Bessel tests satisfies
\begin{equation}
            \log T=(17+o(1))m\log b.                    \label{eq:prime209}
\end{equation}
In particular, if \(R=|\Delta_E|^{1/(2f)}\) and
\(\omega=2\pi/(RQ)\), then
\[
             \log(T\,\omega^{\,2(d-1)})
                                      =(1+o(1))m\log b .
\]
\end{theorem}
\begin{proof}
For a selected word \(c\), the prime construction gives
\(0<fk\log(T/Y_c)<1\) for large \(x\). By the arithmetic
volume and formal-degree identities of Volume~I, the normalizing
target satisfies
\[
 fk\log Y_c
   =\log(I_\sigma/I_\pi)+\log(\mathfrak m_r/\mathfrak m_k)
     +\sum_\tau\log\{W_r(\lambda_\tau)/W_k(u_\tau)\}
     -2\log\Gamma_c .
\]
For a split level prime \(g\), the index ratio is
\((g^k\prod_{i=k+1}^r(g^i-1))^f\): it follows by counting
independent columns of a general linear matrix and dividing by
the sizes of the two prescribed unipotent reductions.
The two such primes have \((kb)^8\le g\le 2(kb)^8\).
Consequently the first term after division by \(fk\) is
\((16+o(1))m\log b\), using \(r=k+m-1\) and \(m/k\to0\).
Lemma~\ref{lem:mass209} bounds the second term and
Lemma~\ref{lem:weyl209} evaluates the third. The size bound
\(0\le\log\Gamma_c\le C f m^5\log b\) gives an error
\(o(m\log b)\) after division, since \(m^4/k\to0\).
This proves \eqref{eq:prime209} for the same prime \(T\).
The conductor--discriminant formula gives \(\log R\le\log N\).
Now substitute \(\log Q=(16+o(1))\log b\), \(\log R=o(\log b)\)
and \(2(d-1)=m-2\) into the second displayed logarithm.
\end{proof}

\subsection{Integral parameters for the prescribed real character}\label{nav:VII-15}
\label{sec:match209}

The frame change with \(t=1\) in \eqref{eq:rd207} gives
two different scales for the block entries and the character.
We now allow the terminal norm in the real model to change and
reindex the rational matrices. The finite adelic integral is
still the one at its original prime \(T\) and level \(K^0\).
Use indices \(1,\ldots,d\), and put
\[
 \alpha=2d-3,\qquad T^\#=Q^\alpha T,\qquad
 L_Q=\operatorname{diag}(Q^{i-1})_{i=1}^d,\qquad
 J_0=\operatorname{diag}((-1)^{d-j})_{j=1}^d .
\]

\begin{lemma}[An integral cross matrix for each supported label]
\label{lem:match209}
A supported relevant pair with \(\mathcal G\ne0\) has a representative
cross matrix \(Z_0\) with entries in \(\mathcal O_E\).
Put \(Z^\#=Q^{-1}L_Q Z_0 J_0 L_Q\).
This matrix is integral and has constant antidiagonals in its
monomial normal form. Writing \(\beta_\nu\) for its block scalars
and \(D_j^\#\) for its leading minors gives
\begin{align}
 \beta_\nu&=(-1)^{d-j_\nu}Q^{\alpha-e_\nu}\zeta_\nu,\qquad
 |D_j^\#|=Q^{j(j-2)}|D_j|, \label{eq:label209}\\
 \mathcal G^\#&=Q^{2d(d-2)}\mathcal G,\qquad
 w(\beta)\ge Q^{-4f} w_\gamma .                         \notag
\end{align}
The Gram determinant on the left uses \(T^\#\) on both sides.
The map from rational orbit labels to these integral block
points is injective, including when the first block is a singleton.
\end{lemma}
\begin{proof}
For every finite place, support gives a representative in
\(K_v^0\) after two \(B_v\) actions. The resulting cross matrix
is integral there. The initial rational matrix is integral away
from finitely many places. Upper unitriangular changes in its
two frames suffice to pass locally to the integral representative.
Choose two rational upper unitriangular changes sufficiently close
to those local changes at the finitely many places in question
and integral everywhere else. For each of their entries this is
additive approximation in \(E\). Indeed the Chinese remainder theorem in
a fractional ideal \(\prod_{\mathfrak p\in S}\mathfrak p^{-n}\)
allows arbitrary prescribed residues at \(S\) to the required
precision once \(n\) is sufficiently large. Apply it at both
primes of \(E\) over a split place of \(F\). Integrality of the
resulting matrix at the prescribed places is an open condition
on the changes since matrix multiplication and conjugation are
continuous. This gives \(Z_0\). Its entry \(Z_{0,11}=D_1\)
lies in \(Q\mathcal O_E\) by Proposition~\ref{prop:shift203}.

The power of \(Q\) in entry \(i,j\) of \(Z^\#\) is \(i+j-3\).
Only the top left entry has negative exponent, and its divisibility
was just proved. Diagonal conjugation preserves lower and upper
unitriangular groups, so the monomial form of \(Z^\#\) is
\(Q^{-1}L_Q M J_0 L_Q\), with \(M\) the previous monomial form.
On an antidiagonal \(i+j=j_{\nu-1}+j_\nu+1\), this exponent
is \(\alpha-e_\nu\). The alternating signs down the block cancel
against those in \(J_0\), giving the first formula. Leading
determinants of the diagonal factors acquire \(Q^{j(j-2)}\)
and the sign of the first \(j\) entries of \(J_0\).
Since \(\alpha+(d-1)(d-3)=d(d-2)\), the two Gram terms
acquire the same factor. The new Gram invariant is nonzero
and integral. All endpoint exponents \(j(j-2)\) are nonnegative
except at \(j=1\). At that endpoint
\(1+|D_1/Q|^2\ge Q^{-2}(1+|D_1|^2)\). Squaring and taking
the \(f\) places gives the weight bound. The block lengths and
the scalars \(\zeta_\nu\) are recovered from \(\beta_\nu\) by
\eqref{eq:label209}. Lemma~\ref{lem:forms203} then recovers the
rational double coset. An integral matrix in the cell and its
nonzero integral Gram invariant are precisely the arithmetic
conditions in
\srcloc{Formalism/SharpRestriction/CuspSpanNormBounds/LeadMinorK}{150}.
\end{proof}

\begin{theorem}[Matching the character and the rotated block point]
\label{thm:match209}
Set \(t=Q^{\alpha/2}\) and \(\kappa=1/(RQ)\). At each real
place scale the \(i\)-th isotropic vector by
\(t(-\mathrm i\kappa)^{d-i}\) and its dual by the inverse
conjugate, and replace each terminal vector by \(tz_i\).
Use \(\gamma^\#=A^{-1}\gamma A\), where \(A\) is this isometry
on the hyperbolic planes and the identity on the middle space.
The transported section has character coefficient
\(c^\#=-2\mathrm i\,\omega\), with \(\omega=2\pi/(RQ)\).
In monomial coordinates its cross matrix is the constant
antidiagonal block point \(\beta\), embedded at that place and
multiplied blockwise by \(\mathrm i^{e_\nu}R^{-e_\nu}\).
In particular these are exactly the rotated local point, the
character coefficient, and the two terminal norms \(T^\#,T^\#\)
in the source's real orbital integral. They satisfy
\begin{align*}
      \log T^\#&=(33+o(1))m\log b,\\
      \log\{T^\#\omega^{\,2(d-1)}\}&=(17+o(1))m\log b .
\end{align*}
\end{theorem}
\begin{proof}
The character transformation in Theorem~\ref{thm:scale207}
applies to a complex scale \(p\kappa\) of phase \(p\), since the
simple entries of the transformed unipotent action are divided
by \(p\kappa\), so the coefficient is multiplied by \(p\kappa\).
Take \(p=-\mathrm i\) and \(c=4\pi\) to obtain \(c^\#\).
In the cross matrix the old \((i,j)\) entry is multiplied by
\(t^2\kappa^{2d-i-j}\mathrm i^{j-i}\).
In a monomial block \(2d-i-j=e_\nu\) and
\[
                  \mathrm i^{j-i}
                    =\mathrm i^{2d-i-j}(-1)^{d-j}.
\]
This is just the \(J_0\) sign in \eqref{eq:label209} followed by
the rotation \(\mathrm i^{e_\nu}\).
The real block factor is \(R^{-e_\nu}\) after absorbing
\(Q^{\alpha-e_\nu}\) into \(\beta_\nu\). In zero-based indices
the rotation in the source is
\((-1)^{d-1}(-\mathrm i)^{(i-1)+(j-1)}=\mathrm i^{2d-i-j}\).
One may take group representatives with exactly monomial cross
matrix. Each upper unitriangular change of the standard decorated
frame is induced by an element of \(N\). To see this, prescribe its action
on the span of the frame, which preserves the self pairing, and
extend over the dual isotropic space. Explicitly the upper
isotropic matrix gives a Levi unitriangular action, and a map
from the middle space to the isotropic space can send the
nonzero terminal vector to any prescribed column. Its negative
adjoint and the usual half of its quadratic product complete
it to a unipotent isometry, by multiplication in the Witt basis.
These changes preserve the simple character and within-block
Haar measure as in Section~\ref{sec:frames207}. They multiply
the kernels on either side by unitary column factors, which
can equivalently be included in the finite contractions.

For the source conventions see
\srcloc{Formalism/FormalDegreeBessel/StatementG}{55},
\srcloc{Formalism/SharpRestriction/CuspSpanNormBounds/FrobeniusCoefficientMemLp/FramesPt}{52}
and \srcloc{Formalism/SharpRestriction/CuspSpanNormBounds/RotatedWeightedBound}{11}.
Finally \(\log T^\#=\log T+(m-3)\log Q\). Apply
Theorem~\ref{thm:prime209}, and then subtract
\((m-2)\log(RQ/(2\pi))\), to get the two asymptotics. These
verify the terminal size and positivity requirements in lines
39 and 73 of
\srcloc{Formalism/SharpRestriction/CuspSpanNormBounds/LeadMinorK}{35}
with a fixed sufficiently large upper size constant.
\end{proof}

\begin{corollary}[The arithmetic sum in the matching coordinates]
\label{cor:match209}
Let \(P^\#_\infty\) be the product of point densities in
Theorem~\ref{thm:match209}. Write
\(\widehat K_{\gamma^\#}=t^{-2kf}K_{\infty,\gamma^\#}^{\rm root}\)
for the kernel of the transported sections including this scalar.
With consistent monomial representatives on the two sides,
\begin{align}
 c_E^a P_\infty^{1/2}K_{\infty,\gamma}^{\rm area}
    &=R^{-fS}Q^{\,f(d(d-2)-a)}
                           (P^\#_\infty)^{1/2}\widehat K_{\gamma^\#},
                                                        \label{eq:match209}\\
 \sum_c |J_{\gamma,c}|
    &\le B_\ell B_{\ell'}\,m_C
       \Bigl(\prod_{v\in S(\gamma)}L_{\gamma,v}\Bigr)
       R^{-fS}Q^{-f(dk+a)}
             (P^\#_\infty)^{1/2}\|\widehat K_{\gamma^\#}\|_{\rm op}.
                                                        \notag
\end{align}
The squared regular norm of the renormalized local kernel
\(t^{-2k}K_{\gamma^\#}^{\rm root}\) equals that of the
original local kernel. The equality of these regular integrals
allows arbitrary Borel truncations carried by the same map.
\end{corollary}
\begin{proof}
All Jacobians in Theorem~\ref{thm:scale207} depend only on the
absolute values of the diagonal entries. The added phase has
real Jacobian one on every complex entry, leaves the compact
column action fixed, and its root phases have absolute value one.
Therefore \eqref{eq:scale207} and \eqref{eq:regular207} hold for
the transported kernel, with the rotation included in the
change of the regular variables. Passing from area to root
measure and multiplying by \(c_E^a\) gives \(R^{fa}\) before
the inverse frame change. After writing the new kernel in terms
of \(\widehat K\), the scalar becomes
\[
 R^{fa}t^{2df}(RQ)^{-f(S+a)}
      =R^{-fS}Q^{\,f(d(d-2)-a)} .
\]
This proves \eqref{eq:match209}. Apply
\(P_\infty^{1/2}\ge Q^{fd(k+d-2)}\) and the finite covolume
formula as in Theorem~\ref{thm:rd207}. The exponent after
division is \(d(d-2)-a-d(k+d-2)=-dk-a\). Finally the scalar
\(t^{-2k}\) cancels the factor \(t^{4k}\) in the squared norm
from Corollary~\ref{cor:regular207}.
\end{proof}

\begin{proposition}[The spectral parameters of the matched sections]
\label{prop:hyp209}
Fix any positive gap constant \(C_g\). At all sufficiently large
\(x\), the transported two sections of a fixed translate meet
the standard local parameter inequalities for the source's averaged
sum at \(T^\#\), with real derivative parameter \(\omega\), upper
size constant \(20\), and size parameter \(b_{\rm sp}=b^2\).
For pure row inputs, the kernels \(\widehat K_{\gamma^\#}\) at
all the rational labels come from a single compact-finite coefficient
at each real place. Its full squared regular norm is the original
local regular integral.
\end{proposition}
\begin{proof}
The relevant inequalities are in
\srcloc{Formalism/SharpRestriction/CuspSpanNormBounds/LeadMinorK}{35}
and lines 68--73 there. The sorted row parameters are integers
in odd rank and the columns are half-integers in even rank.
The interlacing inequalities \(u_i>v_i>u_{i+1}\) and \(u_k>-b_1\)
hold for both words. All absolute parameters are bounded by
\(b^2\) for large \(x\), since the endpoints and negative base
weights in the formulas of Lemma~\ref{lem:weyl209} are \(O(xb)\)
and all spacings have polynomial size. The signed chamber stays
the same since the representation still acts by right translation.
The quantities \(\omega,\omega^{-1}\) are between \(b_{\rm sp}^{-20}\)
and \(b_{\rm sp}^{20}\), using \(\log R=o(\log b)\) and the formula
for \(Q\). The new terminal norms obey
\(1\le T^\#\le b_{\rm sp}^{20m}\) and
\(T^\#\omega^{2(d-1)}\ge1\) by Theorem~\ref{thm:match209}.

The entries strictly between two column entries are exactly the
\(v_i\). They lie in \([0,B_\tau/8]\) and the others are in
\([-2B_\tau,-B_\tau]\) for large \(x\), since \(B_\tau\) grows
faster than every fixed power of \(x\). This also gives
\(B_\tau\ge C_g r^7\). Every row gap is at least
\(\min(4\Delta,B_\tau,2H_0)\ge C_g m^3\) by the spacing choice
of Section~\ref{sec:absolute163}. For two entries with opposite
chamber signs the distance is at least \(\min(B_\tau,2H_0)\), since
all the \(v_i\) have the same sign in the chamber. Since
\(H_0\ge r^{19} f^5 x^\eps\), this exceeds
\(C_g m^{9/2}(3dr^2+r^3)\).
For equal columns \(u_1-u_2\ge4k\Delta\) exceeds
\(2C_gd^2+1\), and the entries \(v_1,v_2\) lie on opposite
sides of \(u_2\). Both are between column entries.
For different words an interior coordinate differs by \(\Delta\).
Inside the interval between the negatives of the two coordinates
leave distance two at each end. The remaining closed interval has
length \(\Delta-4\ge C_gd^2\) at large \(x\). These prove the
separating-gap and cut clauses. The two derivative parameters in
the local size test now both equal \(\omega\), so their ratio
in its terminal norms is one. The fixed nonzero Haar measure
on \(N\) is a positive finite scalar times its exponential
coordinate Haar measure, by uniqueness.

By Theorem~\ref{thm:quot207}, applied in the new model, one
may choose compact-finite test vectors \(u'_i,\xi'_i\) with
\(\ell'_i\ne0\). Use the coefficient
\[
 F^\#(g)=
 t^{-2k}\frac{\langle v_1,v_2\rangle}{\overline{\ell'_1}\ell'_2}
       \langle\sigma(g)u'_2,u'_1\rangle\,\xi'_1(\xi'_2)^* .
\]
The isometric transported maps are the maps of that theorem
with a choice of phase. The integral of \(F^\#\) on the group
pair quotient equals the local renormalized kernel for every
label by the same identity. It is a coefficient \(P\sigma(g)Q_0\)
with \(P^*\) and \(Q_0\) having compact-finite image, since those
images are spanned by \(u'_1,u'_2\), respectively. In particular
these test vectors are independent of the orbit.

There is a single smooth choice of middle isometry in the regular
integral. In the plane with signed unit basis \(w_0,e\) and
\(e=z/\sqrt{T^\#}\) for the new terminal vector \(z\), put
\(\lambda=-q_d/T^\#\), \(h=(|\lambda|^2-1)^{1/2}\). Take the matrix
\(\left(\begin{smallmatrix}\bar\lambda&h\\h&\lambda\end{smallmatrix}\right)\)
on this plane and the identity on its positive complement.
Its columns have norms \(1,-1\) and their pairing is zero.
It gives the required terminal pairing \(q_d\) and varies
smoothly for \(|q_d|>T^\#\). Outside this domain the regular
density vanishes almost everywhere. The singleton-block instance
of \eqref{eq:quot207} identifies the regular orbital transform of
\(F^\#\) with the local renormalized kernel. Its squared integral
is therefore the one in Corollary~\ref{cor:match209}. This claim
also holds in the extended nonnegative sense.
\end{proof}

Put \(M_\gamma=m_C\prod_{v\in S(\gamma)}L_{\gamma,v}\).
Applying Cauchy--Schwarz on any finite set \(\mathcal A\) of the
relevant rational labels gives
\begin{align*}
 \sum_{\gamma\in\mathcal A,c}|J_{\gamma,c}|
 &\le B_\ell B_{\ell'}R^{-fS}Q^{-fdk+2f}\\
 &\quad{}\times
 \left(\sum_{\gamma\in\mathcal A}
   w_\gamma^{-1}Q^{-2fa} M_\gamma^2\right)^{1/2}
 \left(\sum_{\gamma\in\mathcal A} w(\beta_\gamma)P^\#_\infty
         \|\widehat K_{\gamma^\#}\|_{\rm HS}^2\right)^{1/2}.
\end{align*}
Indeed insert \(w(\beta)^{1/2}w(\beta)^{-1/2}\) in the
previous bound and use \eqref{eq:label209}. The second sum runs
over distinct integral points with a common pair of terminal norms.
For pure row inputs its summands are bounded above by the
corresponding terms of the source's nonnegative sum at \(T^\#\)
and \(b_{\rm sp}\). In fact Proposition~\ref{prop:hyp209} supplies
one compact-finite coefficient at each place and the regular
section. The supremum in
\srcloc{Formalism/SharpRestriction/CuspSpanNormBounds/LeadMinorK}{103}
includes the monomial representative, its normalized root Haar
and a disintegration of the group pair measure along that Haar.
Such a disintegration exists by the quotient integration formula,
using a measurable section of the finite-dimensional homogeneous
space. The orbit is relevant by transport. Taking the product
over places and then using injectivity in Lemma~\ref{lem:match209}
proves the comparison with the sum in
\srcloc{Formalism/SharpRestriction/CuspSpanNormBounds/RotatedWeightedBound}{15}.
An upper estimate for this sum in terms of its regular norm, its
normalization against the finite masses, and the arithmetic lower
bound for the designated period still require proof here.

\begin{theorem}[One definite mass for all rational complements]
\label{thm:complement209}
The finite Haar measures on \(U(C)\) can be chosen in
\eqref{eq:mass205} so that for every relevant supported label
\[
 m_C=\mathfrak m_{k-1},\qquad
 \frac{m_C}{\mathfrak m_k}
   =|\Delta_F|^{1/2-k}
       \left(\frac{(2\pi)^k}{(k-1)!}\right)^f\zeta_F(k)^{-1}.
\]
All the ordinary factors \(L_{\gamma,v}=1\) retain that value.
In the preceding sum one may therefore take the same
\(\mathfrak m_{k-1}\) outside the square root for all labels,
leaving the local factors \(L_{\gamma,v}\) inside.
\end{theorem}
\begin{proof}
The positive definite complement \(C\) has odd rank \(n=k-1\).
In any \(E\)-basis let \(\delta\in F^\times\) be its Hermitian
determinant. It is totally positive. Scaling the form by
\(\delta^{-1}\) preserves its unitary group and gives determinant
\(\delta^{1-n}\), the norm of the element
\(\delta^{(1-n)/2}\in F^\times\subset E^\times\).
The scaled space is isometric to the standard form at every
finite place. At a split place this follows by identifying the
two factors as mutually dual vector spaces. At an unramified
field place diagonalize the form and remove unit factors by
surjectivity of the norm on units. Scaling basis vectors reduces
each diagonal entry to one or a uniformizer of \(F_v\).
The number of uniformizers is even by the determinant condition.
A pair of equal diagonal entries gives a hyperbolic plane since
\(-1\) is a unit norm. Indeed it has an isotropic vector, to
which one adjoins a vector pairing to one and subtracts half
the latter's norm times the isotropic vector to make the new
vector isotropic as well. The same construction applies to
two diagonal ones, so every pair of uniformizers can be removed.

Choose a self-dual lattice in each of these finite scaled spaces.
At an ordinary place start with the self-dual lattice for the
original form in Proposition~\ref{prop:good205}. There
\(v(\delta)\) is even at a nonsplit place, by expressing
the determinant in a self-dual basis. A scalar rescaling of
the lattice then makes it self-dual for the scaled form, with
the same stabilizer. At a split place keep one of its two
lattices and take the dual for the new pairing in the second
factor, again leaving its stabilizer unchanged. These choices
agree with integral lattices in fixed global coordinates
at almost all places. Their intersection with the diagonally
embedded global \(E\)-vector space \(C\)
gives a global lattice with the chosen localizations. Indeed
it is trapped between two fractional multiples of a basis
lattice, and additive approximation realizes every local residue
class to arbitrary precision. This lattice is maximal by the
unramified trace-dual argument in Lemma~\ref{lem:mass209}.

Give each of its stabilizers finite volume one. The mass of
this positive definite genus is \(m_C\), by decomposition of
the adelic quotient into its compact open classes. All finite
norm-discriminant classes of the scaled form have unit representative,
so the same mass formula as in Lemma~\ref{lem:mass209} gives
\(m_C=\mathfrak m_{k-1}\).
At the exceptional places changing a Haar measure on \(U(C_v)\)
multiplies \(m_C\) by that scalar and \(L_{\gamma,v}\) by its
inverse. Thus these normalizations give the original cell
measure, including at the marked levels.

In the consecutive-rank mass quotient, only
\(|L_F(1-k,\eta^k)|=|\zeta_F(1-k)|\) occurs because \(k\) is even.
The completed zeta function for totally real \(F\) has factor
\(|\Delta_F|^{s/2}(\pi^{-s/2}\Gamma(s/2))^f\) and is invariant
under \(s\mapsto1-s\). Evaluation at \(k\ge2\), followed by the
gamma reflection and duplication identities, therefore gives
\[
 |\zeta_F(1-k)|=|\Delta_F|^{k-1/2}
       \left(\frac{2(k-1)!}{(2\pi)^k}\right)^f\zeta_F(k).
\]
Division of the two masses is \(2^f/|\zeta_F(1-k)|\), which
gives the claimed value.
\end{proof}

\subsection{The rational change of the full Bessel character}\label{nav:VI-24}
Retain the terminal prime \(T=\varpi\bar\varpi\) and the group of
Theorem~\ref{thm:norm91}. In
Section \ref{nav:VI-08} the
signed row space is written as \(s=m/2-1\) hyperbolic planes with
Witt bases \(u_i,v_i\), and a middle space \(W=W_0\perp Ee\).
Write \({\cal P}_\beta(F,\phi)\) for its Bessel period with last
vector \(\beta e\), normalized quotient measures, and coefficient
one on each preceding simple root. We include any chosen finite
right translate in the row form \(F\). The root coordinates of the
character are specified in
\srcloc{Formalism/LocalToGlobal/BesselHermitianTheta/TiedBessel}{47}.

\begin{theorem}[Transport of the character and the level]\label{thm:full95}
Let \(d\) multiply every \(u_i\) by \(\varpi^{-1}\), every \(v_i\)
by \(\bar\varpi\), and act identically on \(W\). Then \(d\in G(F_+)\).
With \(R(d)F(y)=F(yd)\), one has
\begin{equation}
            {\cal P}_{\varpi}(R(d)F,\phi)={\cal P}_1(F,\phi).
            \label{eq:full95}
\end{equation}
Among diagonal maps scaling the indicated Witt vectors and fixing
\(W\) pointwise, this is the unique one for which conjugation
\(n\mapsto d n d^{-1}\) carries the character with last vector
\(\varpi e\) to the one with last vector \(e\) by pullback.

Away from \(T\), \(d\) preserves the original lattice and is the
identity modulo \(M\). At \(T\) the orbit of \(d\widehat\Lambda\)
under the finite stabilizer of \(\widehat\Lambda\), also with the
levels at \(g_1,g_2\) imposed, has cardinality
\begin{equation}
       \left(\prod_{j=0}^{s-1}\frac{T^{r-j}-1}{T^{s-j}-1}\right)^f .
       \label{eq:grass95}
\end{equation}
\end{theorem}
\begin{proof}
The \(u_i,v_i\) retain their pairings, and the other summands are
fixed. This gives the claimed rational isometry. For a diagonal
unitary \(d_0\) with \(d_0u_i=a_i u_i\) and \(d_0|_W=1\),
conjugation multiplies the \((i,i+1)\)-coordinate by
\(a_i/a_{i+1}\), and the last coordinate, from \(e\) to \(u_s\),
by \(a_s\). Thus the required root coefficients are equivalent to
\[
              a_i/a_{i+1}=1\quad(i<s),\qquad \varpi a_s=1 .
\]
For necessity one may vary each coordinate alone. For \(i<s\)
send \(u_{i+1}\) to \(u_{i+1}+t u_i\) and \(v_i\) to
\(v_i-\bar t v_{i+1}\), fixing the other basis vectors. For the
last coordinate use the element \(n(t)\) of Lemma \ref{lem:root94}.
At an archimedean place, varying the complex scalar \(t\) detects
any nonzero difference of coefficients by the real trace pairing.
The equations give exactly \(d_0=d\).

The map \(d\) commutes with \(H=U(W_0)\). In the integral defining
the left side of \eqref{eq:full95}, automorphy under left
\(d^{-1}\) replaces \(F(nh_0d)\) by \(F(d^{-1}nd\,h_0)\).
Set \(n'=d^{-1}nd\). The map preserves \(N(F_+)\). The induced
automorphism on \([N]\) carries its invariant probability measure to
itself. The old character at \(dn'd^{-1}\) now has last coefficient
one and all its earlier coefficients one. This proves the identity,
including its normalization.

Apply the two-dimensional matrix calculation of
Theorem~\ref{thm:neighbour93} to each hyperbolic plane to obtain
the congruence away from \(T\). At a place above \(T\) choose the
first split component in which \(\varpi\) has valuation one.
The enlarged lattice is the inverse image of an \(s\)-dimensional
residue subspace in \(T^{-1}\Lambda_w/\Lambda_w\). The second
component is its dual. Every change of residue basis lifts to the
integral general linear group. Counting ordered independent
\(s\)-tuples in an \(r\)-space and dividing by the number of bases
of an \(s\)-space gives the factor in \eqref{eq:grass95}.
The \(f\) places contribute independently. Right translation by
\(d\) takes a vector fixed by a level \(K\) to one fixed by
\(dKd^{-1}\). Consequently applying \eqref{eq:full95} to a form
already at the original level gives
\({\cal P}_{\varpi}(F,\phi)={\cal P}_1(R(d^{-1})F,\phi)\),
with the inverse change of level on the right.
\end{proof}

\begin{proposition}[An archimedean factor for the undivided period]
\label{prop:der95}
At each real place \(\tau\) of \(F_+\) fix its chosen complex
embedding of \(E\), and set
\(\eta_\tau=\overline{\varpi_\tau}/\sqrt T\). Write
\(X_{\tau,\xi}\) for differentiation of right translation along
\(b\mapsto n_\tau(b\xi)\), \(b\in\RR\). For integers
\(a_\tau\ge0\), put
\[
 A=\sum_\tau a_\tau,\qquad
 D_{\varpi}=\prod_\tau X_{\tau,\eta_\tau}^{\,a_\tau},\qquad
 D_1=\prod_\tau X_{\tau,1}^{\,a_\tau}.
\]
For every smooth row form and column form as in the theorem,
\begin{align}
 {\cal P}_{\varpi}(D_{\varpi}F,\phi)
       &=(4\pi i\sqrt T)^A{\cal P}_{\varpi}(F,\phi),
          \label{eq:dperiod95}\\
 |{\cal P}_{\varpi}(F,\phi)|^2
       &\le \frac{\|\phi\|_{L^2([H])}^2}{(16\pi^2 T)^A}
              \int_{[H]}\int_{[N]} |(D_{\varpi}F)(nh_0)|^2\,dn\,dh_0 .
          \label{eq:dbound95}
\end{align}
The transported form satisfies
\begin{equation}
              D_{\varpi}R(d)F=T^{A/2}R(d)D_1F .   \label{eq:dcost95}
\end{equation}
\end{proposition}
\begin{proof}
On each real line of parameters, the displayed \(n_\tau(b\xi)\)
form a one-parameter subgroup. Multiplication of the matrices in
Lemma \ref{lem:root94} gives the sum of the two real parameters,
with the quadratic term of their sum. At the place \(\tau\),
\(\Psi_E(\varpi_\tau b\eta_\tau)=\exp(4\pi i b\sqrt T)\),
since the complex trace is twice the real part. Apply the global
translation law of that lemma and differentiate \(a_\tau\) times.
The differentiations under the integral are valid for smooth
forms on the compact domain. More explicitly, one may take compact
sets of representatives for \([H]\) and \([N]\). Images of their
products under translations by a bounded interval of parameters
are compact, and all the differentiated integrands are bounded
there. Directions from distinct places commute. This proves
\eqref{eq:dperiod95}, including \(A=0\).
Apply Cauchy--Schwarz on the two probability quotients to its left
side. The character has absolute value one, yielding
\eqref{eq:dbound95}.

Finally \(d^{-1} n(t)d=n(\varpi t)\), as one verifies on \(u_s,e,v_s\).
Since \(R\) is a representation, differentiating after translating
by \(d\) replaces the direction \(\eta_\tau\) by
\(\varpi_\tau\eta_\tau=\sqrt T\). This gives
\eqref{eq:dcost95} after taking all the derivatives.
\end{proof}

In particular \eqref{eq:dbound95} applies with order \(a=\lceil xb\rceil\)
at the \(h\) high embeddings and zero at the other embeddings on
the very intervals of Corollary~\ref{cor:real91}. Its \(T^{-ha}\)
factor is obtained on the undivided global integral. An application
to the arithmetic period must estimate the differentiated vector
in its specified normalization. Formula \eqref{eq:dcost95} gives
the exact cost when this vector itself undergoes the global change
of character. The integral lattice comparison and the strict
cohomology estimate of Proposition~\ref{prop:res80} remain open.

\subsection{The terminal prime and the parameters of the real integral}\label{nav:VII-31}
We continue on the same intervals as Theorem~\ref{thm:region198},
with exactly the fields, words, level and terminal prime of
Theorem~\ref{thm:common111}. Put \(F=E^+\), \(f=[F:\QQ]\),
\(m=2d\), \(r=k+m-1\), \(Q=g_1g_2\), \(b=\exp(x^\eps)\),
and \(R=|\Delta_E|^{1/(2f)}\). The conductor of \(E\) divides
a number \(N\) with \(\log(2N)=o(\log b)\) on these intervals,
by the initial delay in Theorem~\ref{thm:region198}.

\begin{theorem}[The terminal prime at the original level]
\label{thm:prime210}
Uniformly in the selected words and translates,
\[
   \log T=(17+o(1))m\log b,\qquad
   \log\{T(2\pi/(RQ))^{m-2}\}=(1+o(1))m\log b .
\]
\end{theorem}
\begin{proof}
Use the defining target \(Y_c\) in Theorem~\ref{thm:common111}.
Theorem~\ref{thm:norm91} expresses \(fk\log Y_c\) as the sum
of the logarithms of the level index ratio, the mass ratio
\(\mathfrak m_r/\mathfrak m_k\), and the product of the local
Weyl dimension ratios, minus \(2\log\Gamma_c\).
Proposition~\ref{prop:index90} gives
\((16+o(1))m\log b\) for the normalized level contribution.
For these fields \(E/F\) is unramified at every finite place.
The standard lattice is self-dual with unit determinant, and
every unit is a norm at the unramified nonsplit places. Thus the
definite unitary mass formula gives
\[
 \mathfrak m_r/\mathfrak m_k
   =2^{-f(r-k)}\prod_{j=k+1}^{r}|L_F(1-j,\eta^j)|,
\]
where \(\eta\) is the quadratic character of \(E/F\).
Each value is a product of \(f\) primitive Dirichlet values
of the parity of \(j\). Their Bernoulli polynomial expressions
bound the logarithm of each Dirichlet value above by
\(O(j\log(2jN))\). Together with the lower mass estimate of
Lemma~\ref{lem:mass90}, this gives
\(\log(\mathfrak m_r/\mathfrak m_k)/(fk)=o(m\log b)\).
The local hypotheses and exact odd-primary relation to the
critical cohomological sign are proved in
Lemma \ref{lem:mass209}.

There are \(m(k-1)\) differences \(v_i+b_j\) between the two
parts of the row \eqref{eq:src109}. Their logarithms are
\((1+o(1))\log b\) uniformly, since the smallest \(b_j\)
is at least \(b/(n\exp(x^{\eps/3}))\), the largest is
\(O(xb)\), and the \(v_i\) have polynomial size. In the quotient
of the shared Vandermonde by the column Vandermonde, interlacing
matches each shared difference to a larger column difference.
In the reverse direction it matches each interior column
difference to a larger shared difference. Only the \(2k-3\)
endpoint factors can cost \(O(\log b)\) each. The factorials
cost \(O(km\log k)\), and differences within the negative row
cost \(O(m^2\log x)\). Thus the Weyl terms divided by \(fk\)
contribute \((1+o(1))m\log b\). Finally
\(\log\Gamma_c/(fk)=o(m\log b)\) and \(fk|\log(T/Y_c)|<1\).
This proves the first asymptotic. The other follows from
\(\log Q=(16+o(1))\log b\) and \(\log R\le\log N=o(\log b)\).
\end{proof}

\begin{theorem}[Integral labels at the real additive character]
\label{thm:match210}
Set \(\alpha=2d-3\), \(T^\#=Q^\alpha T\) and \(\omega=2\pi/(RQ)\).
The relevant nondegenerate rational labels at \(K^0\) inject into
the source's integral block points for terminal norms
\((T^\#,T^\#)\). At every real place, the original normalized
sections transport to sections with character coefficient
\(-2\mathrm i\omega\) whose cross matrices represent the rotated
local points of these very labels.
The new terminal parameters satisfy
\(\log T^\#=(33+o(1))m\log b\) and
\(\log(T^\#\omega^{m-2})=(17+o(1))m\log b\).
The finite test still has its original prime \(T\).
\end{theorem}
\begin{proof}
Choose an integral cross matrix for the original pair by rational
unitriangular changes. This is possible because support supplies
an integral representative at all finite places, and additive
approximation in the finitely many nonintegral frame entries
realizes those open conditions simultaneously.
In entry \(i,j\) multiply this matrix by
\((-1)^{d-j}Q^{i+j-3}\), numbering from one.
The only negative exponent occurs at \((1,1)\), where the
original entry \(D_1\) belongs to \(Q\mathcal O_E\).
The new matrix is integral. The transformation is a diagonal
scaling on each side and changes the alternating antidiagonal
in each monomial block to a constant one. At an endpoint \(j\)
the absolute minor acquires \(Q^{j(j-2)}\). With the new terminal
norms the Gram determinant acquires \(Q^{2d(d-2)}\), so it is
again nonzero and integral. Each old block scalar can be recovered
from its new scalar and length. This proves injectivity.

To match the real variable, scale the \(i\)-th isotropic vector
by \(Q^{\alpha/2}(-\mathrm i/(RQ))^{d-i}\), its dual by the
reciprocal conjugate, and the terminal vector by \(Q^{\alpha/2}\).
The associated unitary conjugation sends the coefficient \(4\pi\)
to \(-4\pi\mathrm i/(RQ)\). In a block with endpoints
\(j_{\nu-1},j_\nu\), put \(e_\nu=2d-j_{\nu-1}-j_\nu-1\).
On the antidiagonal \(e_\nu=2d-i-j\), and the phase on entry
\(i,j\) is \(\mathrm i^{j-i}=\mathrm i^{e_\nu}(-1)^{d-j}\).
After the rational sign and power of \(Q\) already included in
the new point, this leaves precisely the source multiplier
\(\mathrm i^{e_\nu}R^{-e_\nu}\). This proves the simultaneous
matching. The size assertions follow by adding \(\alpha\log Q\)
to the two estimates of Theorem~\ref{thm:prime210}.
Section \ref{nav:VII-15}
gives the group representatives and checks the remaining real
parameter inequalities for each pair of words.
\end{proof}

\begin{corollary}[Finite cell contributions in the matching coordinates]
\label{cor:match210}
Let \(\beta_\gamma\) be the new point, \(P_\infty^\#\) the
product of its local point densities, \(t=Q^{\alpha/2}\), and
\(\widehat K\) the transported product quotient kernel multiplied
by \(t^{-2kf}\). Use \(S=d(d-1)\), the block dimension \(a\),
and \(M_\gamma=m_C\prod_{v\in S(\gamma)}L_{\gamma,v}\) as in
Theorem~\ref{thm:rd208}. Then
\[
 \sum_c|J_{\gamma,c}|
   \le B_\ell B_{\ell'}\,M_\gamma R^{-fS}Q^{-f(dk+a)}
                 (P_\infty^\#)^{1/2}\|\widehat K\|_{\rm op}.
\]
For a pure input the renormalized kernel is the orbital integral
of one compact-finite coefficient at each real place, simultaneously
for every label. Its regular squared norm at each place equals
the original regular integral.
\end{corollary}
\begin{proof}
Apply the general two-parameter identity of
Theorem \ref{thm:scale207}
with \(\kappa=1/(RQ)\) and the given \(t\). Its real Jacobians
are unchanged by the added unit-modulus phase. Converting the
arithmetic unipotent area to root measure combines with the
unipotent covolume to give \(R^{fa}\). After the indicated
regular norm scalar the multiplier of
\((P_\infty^\#)^{1/2}\widehat K\) is
\[
          R^{fa}t^{2df}(RQ)^{-f(S+a)}
             = R^{-fS}Q^{f(d(d-2)-a)} .
\]
Divide by the physical density lower bound
\(P_\infty^{1/2}\ge Q^{fd(k+d-2)}\) from the proof of
Theorem~\ref{thm:rd208} and contract with the finite covectors.
This gives the asserted exponent. Equality of the squared
regular norms follows by the change of variables on the regular
cell, including on every transported truncation. The compact-finite
coefficient comes from Theorem~\ref{thm:quot208} in the new
isometric model. The choices of the two test vectors depend on
the words and the place alone. Full details and the comparison of
the nonnegative sum over all labels are in
Corollary \ref{cor:match209} and Proposition \ref{prop:hyp209}.
\end{proof}

The transformed parameter \(T^\#\) is used only inside the real
orbit sum. The source estimate needed there bounds the weighted
sum of squared orbital integrals by a uniform multiple of their
regular squared norm. Even with that estimate, the finite masses
and norm normalizations in Corollary~\ref{cor:match210} must be
combined for the original integral row. The lower bound for its
designated projected theta pairing in Corollary~\ref{cor:original200}
and the complementary cohomology orders in \eqref{eq:region198}
remain open on the same intervals.

\begin{proposition}[The definite mass of the complement]
\label{prop:complement210}
At every label in Corollary~\ref{cor:match210}, use hyperspecial
stabilizers of volume one in the definite complement at finite places.
Then
\[
       M_\gamma=\mathfrak m_{k-1}\prod_{v\in S(\gamma)}L_{\gamma,v},
 \qquad
       \frac{\mathfrak m_{k-1}}{\mathfrak m_k}
       =\frac{|\Delta_F|^{1/2-k}}{\zeta_F(k)}
                   \left(\frac{(2\pi)^k}{(k-1)!}\right)^f .
\]
\end{proposition}
\begin{proof}
The complement has odd dimension \(k-1\) and totally positive
determinant \(\delta\). Multiplying its form by \(\delta^{-1}\)
preserves its unitary group and makes the determinant a norm.
Over an unramified quadratic local extension, diagonalization and
unit norms reduce this form to ones and an even number of
uniformizers. A pair of uniformizers and a pair of ones both
give a hyperbolic plane. The split case is a pair of dual spaces.
Thus one can choose self-dual local lattices for this form, with
the old stabilizer at ordinary places. They give a global maximal
lattice for which every finite factor of the mass formula is one.
Its mass is \(\mathfrak m_{k-1}\). Scaling the exceptional
Haar measures changes their local factors inversely, as required
in the cell formula. The consecutive-rank quotient is
\(2^f/|\zeta_F(1-k)|\), since the rank \(k\) is even. The
functional equation for the completed Dedekind zeta function
at \(k\) gives the displayed expression.
Theorem \ref{thm:complement209}
gives the measure and lattice details.
\end{proof}

\section{Measures, formal degrees and integral weights}\label{nav:g-VIII-2}

\subsection{A single iterated scalar for each translate}\label{nav:VI-18}
On the fields of Theorem~\ref{thm:rangenew50}, use its definite
tower of ranks \(N_j=k+j(t_{\rm r}+1)\), \(0\le j\le m\).
Write \(R_j=N_{j-1}\), \(s_j=N_j\), \(R_\Sigma=\sum_{j=1}^m R_j\).
These ranks are all even. Keep \(a=\lceil xb\rceil\), the sets
\(T\subset\Phi_0\), \(T_2+u\subset\Phi_1\) in the notation of
Proposition~\ref{prop:center78}, and the thresholds \(w_\mu=B'\)
on the translated set and \(w_\mu=W\) elsewhere in \(\Phi_1\).
We take \(a,W,B'\) to be the integers chosen in
Theorem~\ref{thm:orientation46}. We start with the constant one
at every initial class. The following choice of determinant powers
keeps total degrees under control as the output polynomials pass
through the successive substitutions.

\begin{lemma}[Rounding the accumulated column degrees]\label{lem:balance82}
There are nonnegative integer powers \(n_{j,\tau}\) such that the
total degrees \(B_{j,\tau}\) in the \(j\)-th derivative matrix
in equation \ref{eq:sub81}, satisfy
\begin{align}
 R_j a\le B_{j,t}<R_j(a+1) &\quad (t\in T),\notag\\
 R_j(w_\mu-1)<B_{j,\mu}\le R_jw_\mu &\quad(\mu\in\Phi_1),\qquad
 B_{j,\tau}=0\quad(\tau\in\Phi_0\setminus T).             \label{eq:balance82}
\end{align}
\end{lemma}
\begin{proof}
The column-weight rule in
Proposition \ref{prop:op81}
gives \(B_{0,\tau}=0\) and
\(B_{j,\tau}=B_{j-1,\tau}+R_j n_{j,\tau}\).
At a high embedding choose the ceiling of
\((R_j a-B_{j-1,t})/R_j\). At the first stage the choice is \(a\).
At every later stage the preceding upper bound gives
\(B_{j-1,t}< R_{j-1}(a+1)\le R_j(a+1)\). The ceiling is
therefore nonnegative, and its rounding error gives the two
high inequalities. At a low embedding choose the floor of
\((R_jw_\mu-B_{j-1,\mu})/R_j\). The preceding upper bound
and monotonicity of the ranks show that this is nonnegative.
Rounding down gives the two low inequalities. Use zero at all
remaining embeddings and apply induction.
\end{proof}

Let \(S\) consist of the unramified primes \(5\le p\le a/W\)
inert in \(K_d\) and split in \(K_q\). Let \(\mathcal P_x\)
be the primes of Lemma~\ref{lem:move80}. Take one fixed split
support prime for torsion as in Section \ref{nav:V-18}, and the
positive split representatives there. At each step choose its
one basis integral at every prime in \(S\cup\mathcal P_x\).
Fix any terminal class and the identity frame. Put
\[
 D_x=\prod_{p\in\mathcal P_x}
       p^{\,2\lceil h/2\rceil\sum_{j=1}^m d_p(R_j a,R_j s_j)},
 \qquad d_l(v,s)=\max(0,\lceil(v-s(l-1))/l\rceil).
\]
The product, depending on the fields and ranks, is independent of \(u\).

\begin{theorem}[The centered count from one composed theta scalar]
\label{thm:single82}
There is one tuple of rational torsion inputs, one at each step,
such that the terminal value for the lists of Lemma~\ref{lem:balance82}
is nonzero for every \(u\). For each \(u\) multiply this value by
the common stabilizer denominator \(a_{\rm cl}\) of Theorem \ref{thm:compose81}, and call the result \(y_u\). In one normal coefficient
field put \(Z_u=y_u\overline{y_u}/D_x\), using complex conjugation
at the given embedding. These algebraic numbers have nonnegative
valuation everywhere above \(S\cup\mathcal P_x\), and
\begin{equation}
  \Delta_u\ \le\
     \frac{m}{R_\Sigma}\mathcal H_S(Z_u)
         -hm M_{\rm sp}+O_{\psi,\eps}(hmb/x).                \label{eq:single82}
\end{equation}
Here \(\mathcal H_S\) and \(M_{\rm sp}\) have the meanings in
Proposition~\ref{prop:center78}. The common divisor satisfies
\(\log D_x/R_\Sigma\ge(1/14-o_{\psi,\eps}(1))ha\).
\end{theorem}
\begin{proof}
The ranks \(R_j,s_j\) give the required even products. The bases
exist by Lemma \ref{lem:basis63}. The split changes of representatives
and the torsion support occur away from all the primes under
consideration. Apply
Theorem \ref{thm:compose81}
to the finitely many translates at this \(x\). It proves the
nonvanishing through all intermediate class sums, algebraicity
over a common coefficient field, and the addition of the
valuation bounds along this same composition.

At \(p\in S\) the permutation \(\pi=\sigma_pc_E\) interchanges
\(\Phi_0,\Phi_1\). Take the halves in its cycles contained in
\(\Phi_0\). For each \(t\in T\) and each odd \(k\ge1\), the
second floor bound in equation \ref{eq:op81}, and
\eqref{eq:balance82} give a valuation contribution at stage \(j\)
of at least
\[
 [d_{p^k}(B_{j,t},R_j s_j)-B_{j,\pi^kt}]_+
       \ \ge\ R_j[a/p^k-w_{\pi^k t}-s_j]_+ .
\]
Indeed \(d_l(v,s)\ge v/l-s\). The proof of the cited coefficient
bound adds these contributions over all \(t,k\) before the
next class sum. Let \(s_*=\max_j s_j\). At every place above
\(p\in S\), division by \(R_\Sigma\) gives
\[
 \frac{v_p(y_u\overline {y_u})}{R_\Sigma}
       \ \ge\ 2\sum_{t\in T}\sum_{k\ {\rm odd}}
                         [a/p^k-w_{\pi^k t}-s_*]_+ .
\]
Conjugation permutes the places above \(p\), accounting for the
factor two. All exponents \(k\) in this proof are positive.

At \(p\in\mathcal P_x\) every cycle is even and \(\pi T\) lies
in \(\Phi_0\setminus T\). Choose in each cycle the half meeting
\(T\) in more points, with either choice in a tie. These halves
contain at least \(\lceil h/2\rceil\) points of \(T\). For each
of them the \(k=1\) bound at stage \(j\) is at least
\(d_p(R_j a,R_j s_j)\), since the accumulated degree at its
neighbor is zero. The product \(y_u\overline{y_u}\) is therefore
divisible at all these places by the stated rational factor.
The sets \(S,\mathcal P_x\) are disjoint, so the quotient \(Z_u\)
retains the estimate over \(S\). Moreover
\[
 \frac{\log D_x}{R_\Sigma}\ \ge\
    2\lceil h/2\rceil\sum_{p\in\mathcal P_x}(a/p-s_*)\log p .
\]
Now use \(p\le7x/4\), \(\sum_{\mathcal P_x}\log p=x/8+o(x)\)
and \(s_*x=o(a)\), as in Theorem~\ref{thm:quot80}.

At \(p\in S\) the desired summand of the prime count for \(t\)
is \(2\log p\sum_{k\ {\rm odd}}[a/p^k-w_{\pi^k t}]_+\).
Writing \(X=a/W\), all contributing prime powers satisfy \(p^k\le X\).
The total loss from \(s_*\) in the positive parts is at most
\(2hs_*X\log(2X)\), since the sum of \(\log p\) over prime powers
up to \(X\) is at most \(X\log(2X)\). This is \(o(hb/x)\) on
the common intervals, where \(X\asymp x\sqrt b\) and \(s_*\) is
polynomially bounded. The missing ramified and small primes
contribute \(O_{\psi,\eps}(hb/x)\) by the inertia-average argument
of Theorem \ref{thm:signal77}. For each nonzero field element
\(\mathcal H_S\) bounds above the normalized sum of its signed
valuations at \(S\), by the product formula. We obtain
\[
 \sum_{t\in T} C_{\rm sp}(t,a,w)
       \le \mathcal H_S(Z_u)/R_\Sigma+O_{\psi,\eps}(hb/x).
\]
The thresholds in this count are the same for every step.
By Proposition \ref{prop:others77}, the omitted quadratic sign
differs from its main term by \(O_{\psi,\eps}(b/x)\) for each
\(t\). The main term for the sign retained here is \(M_{\rm sp}\)
by Proposition~\ref{prop:center78}. Subtract the main terms and
multiply by \(m\) to obtain \eqref{eq:single82}.
\end{proof}

Thus an upper estimate
\[
  \frac{m}{R_\Sigma}\mathcal H_S(Z_u)
       \le hm M_{\rm sp}+\tfrac12 x^{5\eps}\mathcal B
\]
on the common intervals would give \eqref{eq:gap}, since
\(hmb/x=o(x^{5\eps}\mathcal B)\). This step requires an upper
bound including the denominators outside \(S\). The lower
divisibilities and the positive translation mean already apply
to the same scalar. The source's automorphic period also has a
terminal Bessel functional, level constraints, and additional
frame and character factors. The transfer to that functional and
the uniform bound on \(\mathcal R_{**}\) in
Proposition~\ref{prop:res80} are unresolved.

\subsection{The size of the archimedean factorial in the same comparison}\label{nav:VI-19}
Keep the scalar of Theorem~\ref{thm:single82} and write \(P_u\)
for the polynomial which computes its terminal derivative. The
full factorial in the product translation mean is known exactly
by Theorem \ref{thm:factor83}.
Here we compute its asymptotic on the prescribed tower. Put
\(D=t_{\rm r}+1\), so \(R_j=k+(j-1)D\).

\begin{theorem}[Asymptotic for the factorial on the definite tower]
\label{thm:asym84}
For the particular integer powers in the proof of
Lemma~\ref{lem:balance82}, uniformly in \(u\),
\begin{equation}
 \frac{\log\|P_u\|_{\rm F}^2}{hR_\Sigma}
     =\log(a!)+(1-\gamma+o_{\psi,\eps}(1))a.             \label{eq:asym84}
\end{equation}
\end{theorem}
\begin{proof}
Let \(v_{j,i}\) be the column weights at a fixed high embedding.
They are the same at every high embedding. Write
\(B_j=R_ja+\delta_j\), \(0\le\delta_j<R_j\), and \(\eta=k/D\).
The ranks in Theorem~\ref{thm:rangenew50} give \(0<\eta\le1\)
for large \(x\). At stage one \(n_1=a\). For \(j\ge2\),
the recurrence in Lemma~\ref{lem:balance82} gives
\[
 \frac{n_j}{a}=\frac1{j-1+\eta}
                  +\frac{\delta_j-\delta_{j-1}}{aR_j},
 \qquad |\delta_j-\delta_{j-1}|<R_j .
\]
There are \(k\) first-stage columns and \(D\) new columns at
each stage \(l\ge2\). A column first present at \(l\) has
\(v_{j,i}=\sum_{s=l}^j n_s\). For \(2\le l\le j\) this yields
\[
 \frac{v_{j,i}}a=\sum_{s=l}^j\frac1{s-1+\eta}+O(m/a).
\]
The first-stage columns satisfy \(v_{j,i}/a\ll 1+\log(2j)\),
and the error \(m/a\) tends to zero uniformly in \(j\).

Put \(F(z)=z\log z\) for \(z>0\), \(F(0)=0\). We claim
that for \(x\) tending to infinity and any \(j\le m(x)\)
tending to infinity,
\begin{equation}
  \frac1{R_j}\sum_{i=1}^{R_j} F(v_{j,i}/a)
    \longrightarrow \int_0^1 F(-\log t)\,dt=1-\gamma.    \label{eq:dist84}
\end{equation}
Indeed the \(k\) initial terms have proportion
\(\eta/(j-1+\eta)\) and absolute contribution
\(O((1+\log(2j))^2/j)\). Each remaining group has proportion
\(1/(j-1+\eta)\). For any fixed \(0<\sigma<1\) and
\(\sigma j\le l\le j\), comparison of the harmonic sum with
its integral gives
\(v_{j,i}/a=-\log(l/j)+O(1/(\sigma j)+m/a)\), uniformly in
\(0<\eta\le1\). Uniform continuity of \(F\) on the resulting
compact interval and Riemann summation give the integral over
\([\sigma,1]\). For \(2\le l<\sigma j\) the absolute summand is
bounded by \(C(1+\log(j/(l-1)))^2\), since \(F\) is bounded
on \([0,1]\) and \(z\log z\le z^2\) for \(z\ge1\). The decreasing
function \((1+\log(1/t))^2\) is integrable on \((0,1]\).
Comparison with its integral over \((0,\sigma]\) shows that the
omitted mean tends to zero as \(\sigma\downarrow0\), uniformly
in the indicated limits. Finally \(t=e^{-z}\) gives the
integral \(\int_0^\infty z\log z\,e^{-z}dz=\Gamma'(2)=1-\gamma\),
by the gamma recurrence and \(\Gamma'(1)=-\gamma\).

Use \(\log(v!)=v\log v-v+O(\log(v+2))\), also valid at zero
with the preceding convention. Since \(v_{j,i}\ll a(1+\log(2m))\),
\(\sum_i v_{j,i}=R_ja+O(R_j)\), and \(m\) is polynomial in \(x\),
\[
 \frac{\sum_i\log(v_{j,i}!)-R_j\log(a!)}{aR_j}
       =\frac1{R_j}\sum_i F(v_{j,i}/a)+O_\eps(\log(2a)/a).
\]
The finite initial stages carry a proportion tending to zero in
the \(R_j/R_\Sigma\)-weighted average since \(m\to\infty\).
The tail convergence in \eqref{eq:dist84} holds along every such
sequence \(j\to\infty\), and the bounds in its proof give a
uniformly bounded mean over all stages. The weighted mean
therefore tends to \(1-\gamma\).

To pass to the exact product in equation \ref{eq:factor83}, use
\((R_j-i+1)_v=v!\binom{v+R_j-i}{R_j-i}\).
The total logarithm of the binomial factors is between zero and
\[
 |\Phi|\sum_j R_j^2\log(2+s_*+R_j(a+B'+1)),\qquad s_*=\max_j s_j.
\]
Indeed each column weight is bounded by the total at its embedding.
This expression is polynomially bounded in \(x\), hence
\(o(haR_\Sigma)\). At the low embeddings \(\sum_i v_{j,\mu,i}
\le R_j w_\mu\). Their factorials have total logarithm divided
by \(hR_\Sigma\) at most
\[
 \frac{\sum_{\mu\in\Phi_1} w_\mu}{h}
                  \log(2+s_*(a+B'+1))
                  \ll_{\psi,\eps} b x^{9\eps/8}=o(a).
\]
Here \(\sum w_\mu\le n B'\), \(h\ge nx^{-\eps/8}/2\),
and the logarithm is \(O_\eps(x^\eps)\) for large \(x\).
The remaining embeddings have zero weights. Combining the
estimates for all embeddings proves \eqref{eq:asym84}.
\end{proof}

\begin{proposition}[Size of the explicit divisor and of the available upper estimate]
\label{prop:loss84}
For \(D_x\) of Theorem~\ref{thm:single82},
\[
   \frac{\log D_x}{hR_\Sigma}
       =\bigl(\tfrac14\log(7/5)+o_{\psi,\eps}(1)\bigr)a .
\]
Let \(U_u\) be the right side of equation \ref{eq:norm83},
less \(\log D_x\). It is an upper bound for the normalized
logarithmic field norm of \(Z_u\). Its distance above the
centered main term satisfies
\begin{equation}
 \frac{U_u}{hR_\Sigma}-M_{\rm sp}
     \ \ge\ \bigl(\tfrac34+\tfrac12\omega-o_{\psi,\eps}(1)\bigr)
                         a x^\eps,\qquad \omega=h_2/(2n).       \label{eq:loss84}
\end{equation}
\end{proposition}
\begin{proof}
The secant estimate in the proof of Lemma~\ref{lem:move80} applies
uniformly for \(5x/4\le y\le7x/4\). Indeed on this compact interior
interval it uses the same bound on the remainder over \([x,2x]\)
and a common choice of the relative secant displacement. The
prime powers and the excluded residue classes have total weight
\(o(x)\) on the full interval. Hence, writing \(L=5x/4\),
the sum over \(p\in\mathcal P_x\), \(p\le y\), of \(\log p\) is
\((y-L)/4+o(x)\) uniformly there. Partial summation against
\(1/y\) gives \(\sum_{\mathcal P_x}\log p/p=\log(7/5)/4+o(1)\)
because the endpoint and integral errors are \(o(x)/L=o(1)\).
For \(s\ge1\) the floor bound satisfies
\(v/p-s\le d_p(v,s)\le v/p\). Use this in the definition of
\(D_x\). The aggregate error after division by \(hR_\Sigma\)
is \(O(s_*x)=o(a)\), and \(2\lceil h/2\rceil/h\to1\).
This gives the first formula.

By \eqref{eq:asym84} and \(\log a=x^\eps+\log x+o(1)\),
the contribution of \(\log\|P_u\|_{\rm F}^2/(hR_\Sigma)\) to
\(U_u/(hR_\Sigma)\) is \((1+o(1))ax^\eps\).
The product \(a_{\rm cl}M\) in the same estimate is at
least one. At each stage it includes a sum of positive integers
\(A_{N_{j-1}}m_c\) over a nonempty class set. The terms
\(d_j\log E_j\) and \(e(b_j,d_j)\) there are nonnegative.
The absolute values of the two elliptic means \(L_d,V_d\)
are \(O_{\psi,\eps}(x^{\eps/4})\) by Theorem~\ref{thm:mean56},
Corollary~\ref{cor:interval56}, and Theorem~\ref{thm:norm58}.
Also \(\sum_j b_j/(hR_\Sigma)=O(a)\) by \eqref{eq:balance82}
and \(nB'/h\ll b x^{\eps/8}\), while all \(d_j\) and \(m\)
are polynomially bounded. Thus any negative contribution from
those means is \(o(ax^\eps)\) after this normalization.
The divisor just computed also has size \(o(ax^\eps)\) here.

Finally \(\log(a/W)=x^\eps/2+\log x+o(1)\) and
\(\log(B'/W)=x^\eps/2+o(1)\). Each of the three logarithmic
derivatives in \(M_{\rm sp}\) is \(O(\log x)\) by
Theorem~\ref{thm:twists}. Proposition~\ref{prop:center78} gives
\(M_{\rm sp}=(1/4-\omega/2+o(1))a x^\eps\), uniformly since
\(0\le\omega\le1/2\). Subtraction proves \eqref{eq:loss84}.
\end{proof}

Inequality \eqref{eq:loss84} measures the excess in the explicit
Gaussian upper estimate. The actual period can have a smaller
field norm. The allowed error in \eqref{eq:single82}, divided
by \(hm\), is \(o(a)\) by Theorem~\ref{thm:quot80} and the
size of \(x^{5\eps}\mathcal B\) recorded immediately after it.
Consequently controlling the error of the Gaussian integration
alone leaves a discrepancy much larger than the permitted one.
An upper estimate for the actual period at the required precision
and a bound for its negative valuations outside \(S\) are still
missing on these same fields.

\subsection{The size of the integral highest weight projections}\label{nav:VI-20}
Keep the fields, intervals, ranks, balanced orders and divisor
\(D_x\) of Theorem~\ref{thm:single82}. We now project onto the
highest row weight after every definite step. At the terminal
class use the identity frame. Denote the new scalar, with its
stabilizer denominator cleared, by \(\widetilde y_u\) and put
\(\widetilde Z_u=\widetilde y_u\overline{\widetilde y_u}/D_x\).
The degree partition on the \(R_j\) input columns is \(v_{j,\tau}\)
from Theorem~\ref{thm:asym84}. Define
\[
     C_u=\sum_{j,\tau}\log d_{R_j}(v_{j,\tau}),\qquad
     d_R(v)=\prod_{i<l}\frac{v_i-v_l+l-i}{l-i}.
\]

\begin{theorem}[The projected scalar on the common intervals]
\label{thm:project87}
One can choose common rational torsion inputs for which all the
\(\widetilde y_u\) are nonzero, and \eqref{eq:single82} holds
with \(\widetilde Z_u\) in place of \(Z_u\). The available
Gaussian upper estimate for its normalized logarithmic field norm
is \(\widetilde U_u=U_u-C_u\), with \(U_u\) the explicit bound
in Proposition~\ref{prop:loss84}. Uniformly in \(u\),
\begin{gather}
 0\le C_u=o_{\psi,\eps}(b/x), \qquad
 \frac{\log\|P^{\rm pr}_u\|_{\rm F}^2}{hR_\Sigma}
       =\log(a!)+(1-\gamma+o_{\psi,\eps}(1))a, \notag\\
 \frac{\widetilde U_u}{hR_\Sigma}-M_{\rm sp}
       \ge \bigl(\tfrac34+\tfrac12\omega-o_{\psi,\eps}(1)\bigr)
                     a x^\eps.                             \label{eq:proj87}
\end{gather}
\end{theorem}
\begin{proof}
The integral projection, its nonvanishing on the translated
product torus, and the coefficientwise prime powers are proved in
Theorem \ref{thm:projectchain86}.
In particular every valuation bound used in the proof of
Theorem~\ref{thm:single82} holds for \(\widetilde y_u\). Its
proof uses only those bounds, the nonzero value, and the product
formula. This proves the asserted inequality on the same fields
and yields the same divisor \(D_x\).

The field-norm calculation for the new polynomial uses the same
degree \(b_j\) in each matrix block, the same representatives
and masses, and the same metric envelopes, uniformly in the
torsion choice. The squared Fischer norm is divided by the
product of the dimensions by equation \ref{eq:prnorm86}.
Consequently the right hand side of the field-norm estimate
decreases by exactly \(C_u\).

For \(i<l\), the factor in \(d_R(v)\) is at least one and is
at most \(1+v_1\). Here \(v_1\le R_j(a+B'+1)\) by the balanced
total degrees. Consequently, writing \(s_*=\max_j N_j\),
\[
    0\le C_u\le |\Phi|\sum_j R_j^2
                           \log(2+s_*(a+B'+1)).
\]
The ranks, the number of stages, and \(|\Phi|\) are bounded by
fixed powers of \(x\), while the logarithm is
\(O_{\psi,\eps}(x^\eps+\log x)\). The whole upper bound is
polynomial in \(x\), and \(b=\exp(x^\eps)\). Hence \(C_u=o(b/x)\).
Subtraction in Theorem~\ref{thm:asym84} and in
Proposition~\ref{prop:loss84} proves both formulas.
\end{proof}

There is a separate integral issue if one applies a corner
differential directly to an arbitrary polynomial of the indicated
row and column degrees. In rank two, Proposition \ref{prop:lcm86}
gives the exact common denominator \(\operatorname{lcm}(1,\ldots,a+1)\)
for shape \((a,a)\). Its logarithm is
\(\sum_{n\le a+1}\Lambda(n)=(1+o(1))a\) by the ordinary prime
number theorem. The normalized differential is integral on the
column-invariant lattice by Theorem \ref{thm:projection86}. The field-norm
calculation above uses that lattice throughout. The additional
hook factor in the source uses a second, indefinite row space
with its own weights and integral period pairing. Controlling
that pairing on these fields and bounding
\(\mathcal R_{**}(x)\) of Proposition~\ref{prop:res80} remain
necessary for the proposed argument.

\subsection{A lower bound for the terminal scale from levels, masses and weights}\label{nav:VI-21}
Write \(r=k+m-1\). At either of the two split primes \(g\)
from Theorem~\ref{thm:levels88}, impose the full upper
unitriangular group \(U_k(g)\) on the column reduction. On the
row reduction in \(\mathrm{GL}_r(\mathbb F_g)\) take
\(H_r(g)\) upper unitriangular with \(k\) specified entries
immediately above the diagonal in one column set to zero.
These are the consecutive entries ending just above that
column's diagonal entry. The column is the last of the middle
block of length \(k+1\). This is the level defined in
\srcloc{Formalism/LocalToGlobal/Columns/GenericUnipotentCharacter}{526}.
Let \(I_\pi,I_\sigma\) be the products of the respective
indices at the \(f\) split places for each \(g\).

\begin{proposition}[The finite level contribution]
\label{prop:index90}
Uniformly on the common intervals,
\begin{align}
 \frac{I_\sigma}{I_\pi}
 &=\prod_{g\in\{g_1,g_2\}}
            \left(g^k\prod_{i=k+1}^{r}(g^i-1)\right)^f,
                                                     \label{eq:idx90}\\
 \frac{\log(I_\sigma/I_\pi)}{fk}
 &=(16+o_{\psi,\eps}(1))\,m x^\eps .                 \notag
\end{align}
\end{proposition}
\begin{proof}
Choosing independent columns successively gives
\(|\mathrm{GL}_s(\mathbb F_g)|=
g^{s(s-1)/2}\prod_{i=1}^s(g^i-1)\). The order of the
column subgroup is \(g^{k(k-1)/2}\). In the row subgroup
all entries above the diagonal except the \(k\) specified ones
are free, so its order is \(g^{r(r-1)/2-k}\).
The consecutive zero entries are preserved by upper unitriangular
multiplication and inversion, as is also seen by preserving
the last basis vector in the middle subquotient. Division gives
the one-place ratio and hence \eqref{eq:idx90}. Its logarithm
divided by \(fk\) is
\[
 \left(m+\frac{m(m-1)}{2k}\right)(\log g_1+\log g_2)
       +\frac1k\sum_{g\in\{g_1,g_2\}}\sum_{i=k+1}^{r}\log(1-g^{-i}).
\]
The double sum is \(O(\min(g_1,g_2)^{-k})\) by a geometric
series. Both \(g\) lie between \((kb)^8\) and \(2(kb)^8\).
Now \(\log b=x^\eps,\ \log k=O(\log x)\), and \(m/k\to0\).
\end{proof}

Let \(\mathfrak m_s\) be the sum of reciprocal automorphism
orders over the isometry classes of positive integral Hermitian
lattices in the genus of the standard lattice in \(E^s\).
These are finite sums by finiteness of the definite class set.
Write \(w_E\) for the number of roots of unity in \(E\). For a
strictly decreasing real list \(s=(s_1,\ldots,s_l)\) put
\[
  W_l(s)=\frac{\prod_{i<j}(s_i-s_j)}
                    {\prod_{i=0}^{l-1}i!}.
\]
For integral differences at least one this is the Weyl
dimension for the associated compact highest weight after the
usual shift and a central translation.

\begin{lemma}[Bounds for the definite masses and interlacing weights]
\label{lem:mass90}
For \(r\ge k\) one has
\(\mathfrak m_r/\mathfrak m_k\ge (w_E r)^{-(r-k)}\).
Let \(u_1>\cdots>u_k\), \(k\ge2\), and let
\(u_i\ge v_i\ge u_{i+1}\) for \(1\le i<k\). Suppose a list
\(\lambda\) of length \(r=k+m-1\) is obtained by adjoining
\(m\) numbers to the \(v_i\) and arranging the result in
decreasing order. Suppose all entries of \(\lambda\) differ
pairwise by at least one in absolute value. For
\(A=\max(1,u_1-u_k)\),
\[
     \frac{W_r(\lambda)}{W_k(u)}
          \ \ge\ A^{-(2k-3)}\left(\prod_{i=k}^{r-1}i!\right)^{-1}.
\]
\end{lemma}
\begin{proof}
In a positive integral Hermitian lattice, take vectors \(v,w\)
of Hermitian norm one. Then \(\alpha=(v,w)\) is an algebraic
integer and every conjugate has absolute value at most one
by Cauchy--Schwarz. If \(\alpha\ne0\), its absolute norm is
a positive integer, so all these absolute values equal one.
It is a root of unity. Indeed the monic integral polynomials
whose roots are the conjugates of \(\alpha^j\) have bounded
degree and bounded coefficients as \(j\) varies. Their finite
union of roots contains all powers of \(\alpha\). Equality
in Cauchy--Schwarz gives \(w=v\alpha\) in the convention
linear in the second argument. Thus there are at most \(r\)
mutually orthogonal norm-one lines and at most \(rw_E\) choices
for a norm-one vector.

Count rank-\(r\) lattices with an ordered orthonormal frame
of length \(r-k\), weighting each framed class by the reciprocal
of its stabilizer. Over a single unframed lattice the
orbit--stabilizer formula gives the number of frames divided
by its automorphism order, so the full framed mass is at most
\((rw_E)^{r-k}\mathfrak m_r\). Every rank-\(k\) lattice \(L\)
in the standard genus gives the framed lattice
\(L\perp\mathcal O_E^{r-k}\) with the standard ordered frame.
It belongs to the rank-\(r\) genus at every place. The orthogonal
complement recovers \(L\), and fixing the frame leaves precisely
\(\operatorname{Aut}(L)\). These classes already contribute
\(\mathfrak m_k\). This proves the mass bound.

For \(2\le i<j\le k-1\), match the column factor \(u_i-u_j\)
with \(v_{i-1}-v_j\), which is at least as large by interlacing.
These are distinct factors of the absolute row Vandermonde.
All its other factors are at least one. The column factors
involving the first or the last entry number \(2k-3\), and
each is at most \(A\). Division by the two factorial products
proves the inequality, including \(k=2\).
\end{proof}

At every embedding choose such regular interlacing lists
\(u_\tau,\lambda_\tau\) with \(u_{\tau,1}-u_{\tau,k}\le 3a\),
\(a=\lceil xb\rceil\). The bound below is uniform in the lists
and in the translate. As a concrete choice, take the common
row entries \(v_i=2(k-1-i)\), \(1\le i<k\), and append the
distinct negative integers \(-j(t_{\rm r}+4)\), \(1\le j\le m\).
One can take \(u_i=v_i+1/2\) for \(i<k\),
\(u_k=v_{k-1}-1/2\), and increase \(u_1\) by \(a\) at the
prescribed high embeddings. These give regular integral or
half-integral parameters, strict interlacing, and the stated
span for large \(x\). On the lowest compact types these explicit lists give zero Bessel period by Theorem \ref{thm:zero106}. The replacement lists of Theorem~\ref{thm:real108} admit all of the column types and their real Bessel models.

\begin{theorem}[A terminal prime at the compact dimension scale]
\label{thm:target90}
With the two levels in \eqref{eq:idx90}, put
\[
 d_{u,*}=\mathfrak m_k I_\pi\prod_{\tau}W_k(u_\tau),\qquad
 D_*^{\rm c}=\mathfrak m_r I_\sigma\prod_{\tau}W_r(\lambda_\tau).
\]
Let \(\Gamma\ge1\) satisfy \(\log\Gamma\le C f m^5 x^\eps\)
for a fixed \(C\). Then the positive target
\[
     Y_{\rm c}=\left(\frac{D_*^{\rm c}}
                                  {d_{u,*}\Gamma^2}\right)^{1/(fk)}
     \quad\hbox{satisfies}\quad
       \log Y_{\rm c}\ge(16-o_{\psi,\eps,C}(1))\,m x^\eps .
\]
Consequently there are \(Y_{\rm c}^{1-o(1)}\) primes as in
Theorem~\ref{thm:levels88} with
\(Y_{\rm c}<T\le(1+1/(fk))Y_{\rm c}\).
The same conclusion with a positive constant in the logarithmic
lower bound holds if the numerator is multiplied by a factor
of logarithm at least \(-c_1 fkm x^\eps\) for a fixed \(c_1<16\).
\end{theorem}
\begin{proof}
Take logarithms and divide by \(fk\). The finite index term
was evaluated in Proposition~\ref{prop:index90}. For the mass
term apply Lemma~\ref{lem:mass90}. The roots of unity generate
a cyclotomic subfield, whence \(\varphi(w_E)\le 2f\).
The prime-power formula for \(\varphi\) gives
\(\varphi(l)\ge\sqrt{l/2}\) for all \(l\), since each odd
prime-power factor is at least its square root and the power
of two loses at most a factor \(\sqrt2\). Thus \(w_E\le8f^2\).
The loss in the logarithm of the mass ratio before division
is at most \(O(m\log x)\).

For the weight ratio, Lemma~\ref{lem:mass90} bounds the loss
after division by
\[
     2\log(3a)+\frac1k\sum_{i=k}^{r-1}\log(i!)
                   =O_{\psi,\eps}(x^\eps+m\log x).
\]
Here \(r\le2k\), \(\log(i!)\le i\log i\), and there are
\(m-1\) terms. Both losses are \(o(m x^\eps)\) after the
prescribed normalization. Finally
\(2\log\Gamma/(fk)\le 2Cm^5 x^\eps/k=o(m x^\eps)\) since
\(m^4/k\to0\) at the ranks of Theorem~\ref{thm:rangenew50}.
The lower bound follows. Apply Theorem~\ref{thm:levels88}
to this target. Multiplying the numerator as indicated loses
at most \(c_1m x^\eps\), leaving a positive fixed constant.
\end{proof}

At each of these primes the joint incidence map in
Theorem \ref{thm:smith89} and Corollary \ref{cor:fibres89} has elementary exponents ranging from zero
to \(f(k-1)\), with these extremes attained in each fixed tuple
of nontrivial scalar characters. In particular a primitive
integral input cannot be sent into \(T^{f(k-1)+1}\) times the
full target lattice. To see this, use its coordinate with a
unit coefficient in the integral diagonal basis. Its image
coefficient has valuation at most \(f(k-1)\). The sum \(ha\)
of the high orders in Theorem~\ref{thm:single82} exceeds this
bound by an unbounded factor because \(f,k\) grow polynomially
and \(a\ge x\exp(x^\eps)\). A subsequent scalar functional
can change the valuation by cancellation.

The term \(D_*^{\rm c}\) uses the compact Weyl dimensions.
The source uses formal degrees of specified real representations
and an arithmetic volume, with their normalizations tied to the
indefinite row in \srcloc{Formalism/LocalRingCardinality/DataTies}{17}
and \srcloc{Formalism/SharpRestriction/GlobalSchurBounds/GeneralLinearActions/GapAtLeast}{85}.
Theorem~\ref{thm:norm91} below gives the comparison of the numerator for the specified arithmetic unitary group and its discrete series.

\subsection{The indefinite arithmetic volume and the actual formal degrees}\label{nav:VI-22}
We now identify the two factors in the preceding compact calculation
for global unitary groups on the same fields. Put \(F_+=E^+\).
Fix \(p_++p_-=r\), taking \(p_+=k+m/2-1\), \(p_-=m/2\) for
the prescribed signature. In \(E^r\) use the lattice
\(\Lambda=\mathcal O_E^r\) with Hermitian matrix
\(\operatorname{diag}(I_{p_+},-I_{p_-})\). Write \(G\) for its
unitary group over \(F_+\) and \(G_{\rm c}\) for the unitary
group of the positive identity matrix in rank \(r\).

At finite places give the local stabilizer of \(\Lambda\) volume
one. At a real place use the Haar density on \(U(p_+,p_-)\)
corresponding at the identity to Haar probability on \(U(r)\).
In the Cartan decompositions the comparison is the identity on
\(\mathfrak k\) and multiplication by \(i\) from \(\mathfrak p\)
to \(i\mathfrak p\). Equivalently take absolute values of the
same invariant complex top form on the two real forms and
normalize its compact volume to one. The equivalence follows
by evaluating the top form on bases of \(\mathfrak k\) and
\(\mathfrak p\), since the inserted powers of \(i\) have
absolute value one. This is the compact-dual normalization
described in \srcloc{Formalism/LocalRingCardinality/DataCD}{57}.
Let \(\mathcal V(G)\) denote the adelic quotient volume for the
product of these local measures, before imposing the two finite
levels in \eqref{eq:idx90}.

\begin{theorem}[Two normalization identities on the common fields]
\label{thm:norm91}
For the group and measures just defined, \(\mathcal V(G)=\mathfrak m_r\).
At any real place, let \(\sigma\) be a discrete series of
\(U(p_+,p_-)\) with regular integral infinitesimal character,
written as a decreasing Harish--Chandra parameter \(\lambda\)
after permutation. Its formal degree for the indicated Haar
density is \(W_r(\lambda)\). In this assertion integral means
that the shifted parameter is the highest weight of a
finite-dimensional representation, with its specified central
character.
\end{theorem}
\begin{proof}
The fields \(K_q,K_d\) have the distinct prime discriminants
\(-q,-d\). At a finite prime of \(E\), an inertia element
restricts to inertia on both quadratic subfields. At least one
of those fields is unramified at its underlying rational prime.
Complex conjugation acts nontrivially on both, so it lies in no
finite inertia group. Inertia in \(E/F_+\) is the intersection
with \(\Gal(E/F_+)\), whose only nonidentity element is
conjugation. Thus \(E/F_+\) is split or unramified quadratic
at every finite place.

In either case \(-1\) is the norm of a local integral unit.
At a split place take \((-1,1)\). For an unramified quadratic
extension, the norm surjects on nonzero residue classes.
Starting with a lift of a solution, one corrects modulo
successive powers of a common uniformizer \(\pi\) by multiplying
by \(1+\pi^j z\). Its norm is
\(1+\pi^j\operatorname{Tr}(z)\bmod\pi^{j+1}\) for \(j\ge1\).
The trace of the separable quadratic residue extension is onto,
also in characteristic two. Completeness gives a unit with the
desired norm. Scaling each negative basis vector by such a unit
exhibits an integral isometry between the two local Hermitian
lattices. The finite stabilizers therefore correspond.

For clarity, compare the product measures through Tamagawa
measure. The two groups are inner forms over \(F_+\). Indeed
after splitting \(E/F_+\) both are \(\mathrm{GL}_r\), and
changing the Hermitian matrix conjugates their descent
involutions by an inner automorphism. Transport an invariant
rational top form from \(G_{\rm c}\) to \(G\) by an inner
twisting isomorphism over an algebraic closure. The transported
form descends to \(F_+\), since inner automorphisms preserve
every such form. For \(\mathrm{GL}_r\) this last assertion
also follows directly from determinant one of conjugation on
its matrix algebra. The ratio of the local absolute-form
density to our local Haar density is thus the same at each
finite place by the integral isometries, and at each real
place by the comparison at the identity above. The convergence
factors in the Tamagawa measure are the same as well. They
depend on the Galois module of algebraic characters, which in
both cases is generated by the determinant with conjugation
acting by inversion. The global discriminant and dimension
factors are identical. Consequently changing our measure to
Tamagawa measure multiplies the two quotient volumes by the
same scalar. Neither group has a nontrivial \(F_+\)-rational
character, so these quotients are the full adelic quotients.
Kottwitz's invariance of Tamagawa numbers under inner twisting
now gives equality of their volumes.~\cite{XVIII-XIII-I-kott91}

For \(G_{\rm c}\) the archimedean factors are compact of
volume one. The finite double cosets with the given stabilizer
parametrize the genus in the standard positive Hermitian space.
Indeed a finite adelic element transports the completed lattice,
and intersecting with \(E^r\) recovers a global lattice with
these completions. Conversely the defining local isometries
of any lattice in the genus give such a tuple, integral almost
everywhere. The rational stabilizer of a tuple up to its compact
level is precisely the automorphism group of that global lattice.
Each double coset contributes the reciprocal of that finite order
to the adelic volume. Summing gives \(\mathfrak m_r\), as claimed.

For the real-place assertion first use \(SU(p_+,p_-)\) with
both indices positive. This is a connected linear semisimple
group without compact factors, with compact Cartan. Fix a
symmetric-space metric on \(X=SU(p_+,p_-)/S(U(p_+)\times U(p_-))\)
and the dual metric on the compact symmetric space \(X^{\rm c}\).
For the Haar measure given by Riemannian measure on \(X\) and
Haar probability on its maximal compact subgroup, the discrete
series formal degree is
\(\dim(\tau)/\operatorname{vol}(X^{\rm c})\), where \(\tau\)
has the same infinitesimal character.~\cite[Cor.~5.2]{XVIII-XIII-I-olbrich91}
Here one can take the representation of the complexified group
of highest weight \(\lambda-\rho_r\), restricted to the derived
group, where \((\rho_r)_i=(r+1-2i)/2\). The hypotheses on
\(\lambda\) are precisely those giving this highest weight.
The Weyl dimension formula gives \(\dim(\tau)=W_r(\lambda)\).
On the compact dual the product of its Riemannian quotient
measure with Haar probability on the same maximal compact
subgroup has total mass \(\operatorname{vol}(X^{\rm c})\).
Dividing by that volume on both real forms gives the measure
specified in the theorem. A formal degree scales by the inverse
of the factor multiplying the Haar measure, by the matrix
coefficient orthogonality integral. This gives \(W_r(\lambda)\).

The map \(U(1)\times SU(p_+,p_-)\to U(p_+,p_-)\) given by
multiplication is onto, by taking an \(r\)-th root of the
determinant, and its kernel is finite and central. Use Haar
probability on \(U(1)\). The product Haar measure pushes forward
with the indicated compact-dual normalization. Its compact
counterpart pushes forward to Haar probability on
\(U(r)\), and the differential of multiplication compares the
same central line and the same semisimple tangent densities
on both sides. Pullback of the representation is a central
character times an irreducible representation of the derived
group, since the center acts by scalars. The squared absolute
matrix coefficients integrate to the same values under this
pushforward. This extends the degree formula to the required
unitary group. For a compact signature the formula follows
directly from compact Schur orthogonality.
\end{proof}

\begin{corollary}[Terminal primes for the indefinite row scale]
\label{cor:real91}
Use the regular integral lists from Theorem~\ref{thm:target90}
at the real places and take any discrete series in their
indicated packets. Such series exist for these parameters by
the discrete-series classification for a compact Cartan,
including the explicit lists above since \(r\) is odd. Put
\[
       D_*^{\rm row}=\mathcal V(G) I_\sigma
                    \prod_{\tau}\deg(\sigma_\tau),\qquad
       T_*=\left(D_*^{\rm row}/(d_{u,*}\Gamma^2)\right)^{1/(fk)} .
\]
With \(\Gamma\) in the range of Theorem~\ref{thm:target90},
\(\log T_*\ge(16-o(1))m x^\eps\). There are \(T_*^{1-o(1)}\)
primes in \((T_*,(1+1/(fk))T_*]\) with a generator congruent
to one modulo \(N(g_1g_2)^2\mathcal O_{K_d}\). All the finite
incidence conclusions of Section \ref{nav:VI-06} apply to them.
\end{corollary}
\begin{proof}
Theorem~\ref{thm:norm91} identifies the numerator exactly with
\(D_*^{\rm c}\). Apply Theorem~\ref{thm:target90} and then the
finite-place calculation at each of the chosen primes.
\end{proof}

These are volume and formal-degree identities for specified
global groups and local representations. A period evaluated on
specified automorphic vectors requires an additional estimate.
In particular the negative term from the terminal level in
the height of the projected Bessel period, and the uniform
bound for the residual length \(\mathcal R_{**}(x)\) in
Proposition~\ref{prop:res80}, remain open here.

\subsection{A self-dual lattice for the normalized terminal lines}\label{nav:VI-23}
Let \(K=K_d,\ M=N(g_1g_2)^2\), and use the terminal scale of
Corollary~\ref{cor:real91}. We may choose its prime with
\(\pi\equiv1\bmod 2M\mathcal O_K\). Indeed in
Lemma~\ref{lem:ray88} replacing \(M\) by \(2M\) changes
\(\log(4dM^2)\) by \(\log4\). All interval estimates in
Theorem~\ref{thm:levels88} still apply.
Write \(T=\pi\bar\pi\). In the signed lattice \(\Lambda\) of
Theorem~\ref{thm:norm91} take pairwise distinct coordinate
vectors \(e_+,e_-,e\) of signs \(+,-,-\). Put
\[
     u=e_++e_-,\qquad v=(e_+-e_-)/2,\qquad z_c=\pi e+Mc\,u .
\]
The Hermitian form pairs \(u\) and \(v\) to one and vanishes
on each of them. The lattice calculations for \(z_c\), including
its norm \(-T\) and the local evaluation criterion, are in
Section \ref{nav:VI-07}.

\begin{theorem}[One rational isometry for all the normalized lines]
\label{thm:neighbour93}
Define \(g\) on the Hermitian space by
\[
          gu=\pi^{-1}u,\qquad gv=\bar\pi v,\qquad
          g=1\quad\text{on }\langle u,v\rangle^\perp .
\]
Then \(g\in G(F_+)\) and \(\Lambda'=g\Lambda\) is self-dual.
At every finite place away from \(T\) it is the original lattice,
and \(g\) acts as the identity modulo \(M\) there. All the vectors
\(z_c/\pi\) for \(1\le c<T\) split off orthogonally as integral
lines of norm \(-1\) in \(\Lambda'\).

The orbit of \(\widehat{\Lambda}'\) under the finite adelic
stabilizer of \(\widehat\Lambda\) has cardinality
\begin{equation}
                \left(\frac{T^r-1}{T-1}\right)^f .           \label{eq:orbit93}
\end{equation}
The same count holds if the auxiliary levels at \(g_1,g_2\)
are imposed on the stabilizer. At every archimedean place the
largest and smallest singular values of \(g\), for the positive
coordinate norm, are \(\sqrt T\) and \(1/\sqrt T\).
\end{theorem}
\begin{proof}
The cross pairing after scaling is
\(\overline{\pi^{-1}}\bar\pi h(u,v)=1\), so the map is a
unitary \(E\)-linear map defined over \(F_+\) as a point of
the unitary group. Isometry preserves self-duality. To check
the integral assertion at all places away from \(T\), including
those above two, use the original ordered pair \(e_+,e_-\).
The diagonal entries of its block are
\((\pi^{-1}+\bar\pi)/2\) and both off-diagonal entries are
\((\pi^{-1}-\bar\pi)/2\). Away from \(T\), both scalars
\(\pi^{-1},\bar\pi\) are units congruent to one modulo \(2M\).
Thus the block and its inverse are integral and equal to the
identity modulo \(M\). Finally \(g^{-1}(z_c/\pi)=e+Mc\,u\)
is integral of norm \(-1\). A vector of unit Hermitian norm
gives an integral orthogonal projection onto its line by
dividing the pairing by that norm. This proves the splitting.

At a place of \(F_+\) above \(T\), identify the split unitary
group with \(\mathrm{GL}_r\) by choosing the factor of \(E\)
in which \(\pi\) has valuation one. Its residue field is
\(\mathbb F_T\). The hyperbolic basis above is integral and
invertible there since \(T\) is odd. In the chosen factor the
lattice is enlarged by one inverse uniformizer on the line
of \(u\), and all other directions keep their integral spans.
In the conjugate factor the lattice is its dual, determined
uniquely by the first. The possible enlargements are indexed
by the residue lines in \(T^{-1}\Lambda_w/\Lambda_w\), where
\(\Lambda_w\) denotes the lattice in the chosen factor. Each
line has its full inverse image as the enlargement.
The reduction map
\(\mathrm{GL}_r(\mathcal O_{F_+,v})\to\mathrm{GL}_r(\mathbb F_T)\)
is onto by lifting entries of an invertible matrix. Its action
on lines is transitive. There are \((T^r-1)/(T-1)\) lines,
independent at the \(f\) places. All other places contribute
one orbit element, including the two auxiliary levels. This
gives \eqref{eq:orbit93}. In the positive metric \(u,v\) are orthogonal
and orthogonal to the remaining coordinates. Their eigenvalues
have absolute values \(T^{-1/2},T^{1/2}\). This proves the final
assertion.
\end{proof}

\begin{corollary}[The two polynomial lattices on the evaluation plane]
\label{cor:plane93}
Let \(P\) be homogeneous of degree \(A<T-1\), with restricted
binary form \(P(su+te)=\sum_{i=0}^{A}b_i s^it^{A-i}\) integral
on the original coordinate plane over \(\mathcal O_{J,w}\),
where \(w\mid(\pi)\) in a coefficient field containing \(E\).
The following are equivalent.
\begin{enumerate}
\item The polynomial on the plane of \(\Lambda'\) has integral
coefficients in the basis \(\pi^{-1}u,e\).
\item Every \(P(z_c/\pi)\) for \(1\le c\le A+1\) is integral at \(w\).
\item \(b_i\in\pi^i\mathcal O_{J,w}\) for every \(i\).
\end{enumerate}
The operator \(P\mapsto P\circ g\) on all complex homogeneous
polynomials of degree \(A\), with the positive Fischer norm,
has norm \(T^{A/2}\). On polynomial functions of the plane
\(\langle u,e\rangle\) with its induced positive norm, the
same pullback has norm one.
\end{corollary}
\begin{proof}
At \(w\) the original plane lattice has basis \(u,e\) and the
new one has basis \(\pi^{-1}u,e\), by the integral hyperbolic
basis in the theorem. In that new basis the coefficients are
\(\pi^{-i}b_i\). This proves the equivalence of the first and
third assertions. Theorem \ref{thm:jets92} proves equivalence
with the second one by a Vandermonde determinant that is a
unit at \(w\). For the norm statement choose an orthonormal
positive basis in the directions of \(u,v\) and the remaining
coordinates. The normalized coordinate monomials form an
orthonormal Fischer basis. Pullback scales a monomial by
the product of the corresponding eigenvalues, one for each
factor of its degree. The largest modulus on the full space
is \(T^{A/2}\), attained in the direction dual to \(v\).
On the plane only the eigenvalues \(\pi^{-1},1\) occur and
the largest modulus is one, attained in the coordinate dual to \(e\).
\end{proof}

This identifies the precise lattice change underlying the
factor \(T^{-A}\) in the squared terminal evaluation. A
polynomial integral on \(\Lambda'\) passes the local test on
the plane. For a global theta period one must compare the
specified row polynomial and test vector at these two levels,
including the normalization in their archimedean pairing.
The rational isometry transports a lattice and a vector
together. Keeping the prescribed period functional in that
transport is a further arithmetic assertion. The uniform
height of this period and the residual cohomology bound of
Proposition~\ref{prop:res80} are still open.

\subsection{The mass and the real terminal parameters}\label{nav:VI-29}
Use the common prime of Theorem~\ref{thm:common111}, with the
original split level primes \(g_1,g_2\) and \(Q=g_1g_2\). In
this calculation \(F=E^+\), \(f=[F:\QQ]\), \(m=2d\),
\(s=m-1\) and \(r=k+s\). Define
\[
 R=|\Delta_E|^{1/(2f)},\qquad \omega=2\pi/(RQ),\qquad
 L_j=L_F(j,\eta^j),
\]
where \(\eta\) is the quadratic character of \(E/F\) and
\(\Delta_E\) the absolute field discriminant. For a list
\(v\) of distinct real entries write
\(\mathcal V(v)=\prod_{i<j}|v_i-v_j|\). For each selected word
\(c\), use the shifted highest weights \(u_\tau(c)\) and
\(\lambda_\tau\) that occur in its two Weyl dimensions. Put
\[
     A_c=\left(\prod_{\tau:F\hookrightarrow\RR}
       \frac{\mathcal V(\lambda_\tau)}{\mathcal V(u_\tau(c))}
                    \right)^{1/(fk)}.
\]
The entries here are the parameters in the Weyl product, including
the shifts in Lemma~\ref{lem:weyl209}.

\begin{lemma}[The factorials and discriminants in the adjacent masses]
\label{lem:mass232}
On the period fields one has
\[
 \frac{\mathfrak m_{k+s}}{\mathfrak m_k}
     =\prod_{j=k+1}^{k+s}
       \left(\frac{(j-1)!\,R^{j-1/2}}{(2\pi)^j}\right)^f L_j .
\]
For \(k\ge1\) the product of the \(L_j\) is positive and the
absolute value of its logarithm is \(O(f2^{-k})\), with an
absolute constant independent of \(s\) and the field.
\end{lemma}
\begin{proof}
Inertia at each finite prime of \(E/\QQ\) fixes at least one of
the two imaginary quadratic subfields of relatively prime
discriminants \(-q,-d_0\). Hence it cannot contain complex
conjugation. This proves that \(E/F\) is unramified at the finite
places, and \(|\Delta_E|=|\Delta_F|^2\).
For even \(j\) the relevant abelian factorization is that of
\(\zeta_F\), and for odd \(j\) it is that of \(\zeta_E/\zeta_F\).
In each case it consists of \(f\) primitive Dirichlet functions
with parity \((-1)^j\). The product of their conductors is
\(|\Delta_F|\) in the even case, by the conductor--discriminant
formula, and \(|\Delta_E|/|\Delta_F|=|\Delta_F|\) in the odd case.

For a primitive factor of conductor \(h\) and this parity, the
absolute-value functional equation at \(1-j\) supplies a factor
\[
   h^{j-1/2}\,\frac{2(j-1)!}{(2\pi)^j}
\]
times the absolute value of the conjugate character's function
at \(j\). Indeed the quotient of the two archimedean factors is
\( |\Gamma_{\RR}(j+a)/\Gamma_{\RR}(1-j+a)|\),
where \(a\in\{0,1\}\) is \(j\bmod2\) and
\(\Gamma_{\RR}(t)=\pi^{-t/2}\Gamma(t/2)\). The reflection
and duplication formulas give exactly the displayed factor
without \(h^{j-1/2}\). Taking the product gives
\(R^{f(j-1/2)}\) from the conductors.
Conjugation permutes the primitive factors. The Euler product
at \(j\ge2\) is positive, since the quadratic character takes
real values. Substitute in the mass formula of Lemma~\ref{lem:mass209}.
Its \(2^{-fs}\) cancels the factors \(2^{fs}\).
Finally \(|\log L_j|\le f\log\zeta(j)\le3f2^{-j}\).
The first inequality follows prime by prime from the positive
absolute-value series for the logarithm, as there are at most
\(f\) primes of \(F\) above a rational prime. Summing the
geometric series proves the bound.
\end{proof}

\begin{theorem}[Multiplicative estimates for the two real terminal parameters]
\label{thm:terminal232}
For the selected words on the fields of Theorem~\ref{thm:conductor232},
let \(T\) be their actual common terminal prime and
\(T^\#=TQ^{m-3}\) its value after the change of frame.
Uniformly for every selected \(c\) and every allowed choice of the
numbers \(\Gamma_c\) in Theorem~\ref{thm:common111},
\[
 \begin{aligned}
 T\,\omega^{m-2}
     &=\frac{RQ^2}{2\pi}\,A_c\bigl(1+O_{\psi,\varepsilon}(x^{-2})\bigr),\\
 T^\#\,\omega^{m-2}
     &=\frac{RQ^{m-1}}{2\pi}\,A_c\bigl(1+O_{\psi,\varepsilon}(x^{-2})\bigr).
 \end{aligned}
\]
These are the parameters of the original and transported real
integrals for the same finite tests.
\end{theorem}
\begin{proof}
Let \(S=\sum_{j=k+1}^{k+s}j=ks+s(s+1)/2\) and put
\[
  Z=\prod_{j=k+1}^{k+s}\left(L_j
              \prod_{g\in\{g_1,g_2\}}(1-g^{-j})^f\right).
\]
Then \(|\log Z|\ll f2^{-k}\) by the lemma and
\(|\log(1-g^{-j})|\le2g^{-j}\).
At a split level prime the index ratio, per real place, is
\(g^k\prod_{j=k+1}^{k+s}(g^j-1)\), by the column count in
Theorem~\ref{thm:prime209}.
Thus its product contributes \(Q^{f(k+S)}\) and the displayed
local corrections. In rank \(n\) the Weyl dimension is the
positive Vandermonde of the shifted entries divided by
\(\prod_{i=1}^{n-1} i!\). In the ratio between ranks \(k+s\)
and \(k\), the new denominators are \(\prod_{j=k+1}^{k+s}(j-1)!\)
at each real place. They cancel the factorials in
Lemma~\ref{lem:mass232} exactly. In the definition of \(Y_c\)
in Theorem~\ref{thm:common111}, this leaves the identity
\[
 Y_c\omega^s=Q A_c\,\Gamma_c^{-2/(fk)}Z^{1/(fk)}
      \left(\frac{Q}{2\pi}\right)^{s(s+1)/(2k)}
       R^{s^2/(2k)} .
\]
Indeed the mass gives \(R^{f(S-s/2)}(2\pi)^{-fS}\).
Multiplying the \(fk\)-th root by \(\omega^s\) gives exactly
the two powers on the right.

The choice of prime gives \(0<fk\log(T/Y_c)\le3e^{-x^{1/4}}/8\).
Also \(0\le\log\Gamma_c\ll fm^5 x^\varepsilon\),
\(\log Q=O(x^\varepsilon)\) and
\(\log R\le \log N_x=O_{\psi,\varepsilon}(\log(2x)
 (1+\log\log(3x))^2)\) under either conductor bound.
It follows that
\[
 \left|\log\frac{T\omega^s}{Q A_c}\right|
 \ll_{\psi,\varepsilon}
       \frac{m^5x^\varepsilon+m^2\log(2RQ)+1}{k}
        = O_{\psi,\varepsilon}(x^{-2}).
\]
For the last equality use \(m\asymp x^{1/2-\varepsilon/16}\)
and \(k\asymp x^{9/2+10\varepsilon}\). Exponentiation preserves
this relative error. Dividing by \(\omega\) gives the first
assertion, and multiplying by \(Q^{m-3}\) gives the second,
using the exact frame change in Theorem~\ref{thm:match209}.
\end{proof}

The factor \(A_c\) is determined by the word and the row.
The integral against the designated projected theta vector
still requires an estimate for its signed value at these
parameters.

\subsection{Changing the global theta character at the prescribed period}\label{nav:VI-30}
On the fields of Theorem~\ref{thm:conductor232}, use the
notation of Theorem~\ref{thm:scale232}. Thus \(F=E^+\),
\(d=[E:\QQ]=2f\), \(\beta'=(d_0\gamma)^{-1}\), and \(u=\beta'/\beta\)
is totally positive. Both normalizations use the same skew form.
The fractional ideal of \(u\) is supported on split places. Write
\(v\) and \(w\) for normalized additive valuations in \(F\) and
\(E\), respectively, and write \(Nv,Nw\) for residue field sizes.

\begin{theorem}[A norm representative with prescribed split components]
\label{thm:norm233}
Let \(S\) consist of finitely many split places of \(F\), including
those where \(v(u)\ne0\). Above each \(v\in S\) choose a prime
\(w_v\) of \(E\) and a positive integer \(a_v\). There are
\(\xi\in E^\times\) and one prime \(P\) of \(E\) such that
\[
  \xi\bar\xi=u,\qquad \xi_{w_v}\equiv1\pmod{w_v^{a_v}},\qquad
  (\xi)=\left(\prod_{v\in S}\bar w_v^{\,v(u)}\right)P\bar P^{-1}.
\]
The rational prime below \(P\) splits completely in \(E\) and can
avoid any given finite set \(U\) of rational primes. Enlarge \(U\)
to include the rational primes under \(S\) and those ramified in
\(E\). Put \(\mathfrak m=\prod_{v\in S}(w_v\bar w_v)^{a_v}\). Then
\[
 \log NP\ll 1+\log|D_E|+\log N_{E/\QQ}\mathfrak m
              +d\sum_{l\in U}\log(2dl)+d\log(2d).
\]
For the split test places and prime levels in Proposition~\ref{prop:cost177}, the logarithm of \(NP\) and
\(\sum_{w<\infty}|w(\xi)|\log Nw\) are
\(O_{\psi,\eps}((1+\max_{v\in S}a_v)x^C)\) for a fixed \(C\).
\end{theorem}
\begin{proof}
At a split place every element is a norm. At an inert place \(u\)
is a unit and \(E/F\) is unramified. The residue field norm is onto.
Lift an element of the required residue norm to a unit upstairs.
The quotient of \(u\) by its norm is a principal unit and has a square
root by Hensel's lemma, since all dyadic places split. This square
is itself the norm of an element of \(F_v\). At a real place
positivity is exactly the local norm condition. Applying the
Hasse--Minkowski theorem to \(X^2+d_0Y^2-u Z^2\) now gives a
nonzero rational zero over \(F\). Here \(Z\ne0\) by total reality and
\(d_0>0\). Division by \(Z\) gives \(\xi_0\in F(\sqrt{-d_0})=E\)
with norm \(u\).

Multiply \(\xi_0\) by \(t/\bar t\) to achieve the required congruences.
Indeed weak approximation in \(E\) lets \(t\) be arbitrarily close
to \(\xi_{0,w_v}^{-1}\) at \(w_v\) and to one at \(\bar w_v\),
simultaneously for all \(v\in S\). Call the result \(\xi_1\) and put
\(J=\prod_{v\in S}\bar w_v^{v(u)}\). The ideal
\((\xi_1)J^{-1}\) is trivial over \(S\), and its conjugate is its
inverse. It can therefore be written \(A/\bar A\) for an integral
ideal \(A\) coprime to \(\mathfrak m\). Take precisely the positive
exponents of the fractional ideal. They occur in opposite pairs
away from the conjugation-fixed primes.

Take the ray class of \(A\) modulo \(\mathfrak m\). For each
\(l\in U\) let \(b_l\) be the least positive integer with
\(\varphi(l^{b_l})>d\), and put \(M=\prod_{l\in U}l^{b_l}\).
Then \(\log M\le\sum_{l\in U}\log(2dl)\), by minimality.
In the compositum over \(E\) of its ray class field and \(E(\zeta_M)\),
lift the chosen ray class to an element of the abelian Galois group.
Every prime above \(l\in U\) ramifies in this extension, since the
absolute ramification index of \(\QQ_l(\zeta_{l^{b_l}})\) exceeds
\(d\). The finite conductor norm of each character is at most
\((N_{E/\QQ}\mathfrak m)M^d\). In fact a character extends from
the subgroup of the product of the ray and cyclotomic groups to
that product. The ray factor has conductor dividing \(\mathfrak m\).
The cyclotomic factor comes by norm from \(\QQ\), and the determinant
of multiplication by \(1+l^{b_l}y\) is one modulo \(l^{b_l}\) for
integral \(y\). This proves the asserted bound on its conductor.

The abelian case of
\href{https://arxiv.org/pdf/1604.01750}{Thorner and Zaman, Theorem~1.1},
applied over \(E\), gives an unramified prime \(P\) of absolute
residue degree one with this Frobenius and with the claimed bound
on \(\log NP\). It avoids \(U\). Its rational prime therefore also
has absolute ramification index one in the Galois field \(E\), so
it splits completely. The equality of ray classes gives
\(P=(s)A\) with \(s\equiv1\bmod\mathfrak m\). Taking
\(\xi=\xi_1s/\bar s\) proves the three identities.

Finally the disjoint supports give the exact finite sum
\[
  \sum_{w<\infty}|w(\xi)|\log Nw
      =\sum_{v\in S}|v(u)|\log Nv+2\log NP.
\]
The first sum is polynomially bounded by Proposition~\ref{prop:fourier232}. For the stated test places the
rational primes in \(U\) may be taken from the conductor, the two
original prime ideals and the fixed levels \(g_1,g_2,T\). Their
logarithms have polynomial sum by Proposition~\ref{prop:cost177}
and Lemma~\ref{lem:size97}. Also
\(\sum_{v\in S}\log Nv\le f\sum_{l\in U}\log l\),
\(\log |D_E|\le d\log N_x\), and \(d\) is polynomially bounded.
These estimates prove the final assertion. At infinity the absolute
values are fixed exactly by \(|\sigma\xi|^2=\tau u\) for
\(\sigma\mid\tau\).
\end{proof}

\begin{theorem}[The transported kernel at the same Bessel character]
\label{thm:transport233}
Start with the global splittings, Schwartz functions and Haar measures
of a given convergent theta chain for \(\psi_\beta\). Use the same
groups and skew form, and use \(\psi_{\beta'}\) in every oscillator.
There is an explicit transport of the splittings and Schwartz functions
for which each theta kernel agrees pointwise with its original kernel.
In particular the iterated lift, its prescribed terminal Bessel
coefficient and every fixed linear projection of that coefficient
agree, including finite signed linear combinations using those
transported data. The external character and all group measures in
the period have their original values.
\end{theorem}
\begin{proof}
On a tensor space write \(B\) for its alternating trace pairing.
Multiplication by \(\xi\) commutes with both unitary groups and
satisfies \(B(\xi v,\xi w)=uB(v,w)\). For the Heisenberg group,
whose central term in the multiplication is \(B(v,w)/2\), the map
\((v,t)\mapsto(\xi v,ut)\) is consequently an automorphism defined
over \(F\). Pulling the old central character back gives exactly
\(\psi_\beta(ut)=\psi_{\beta'}(t)\).

Take a rational polarization of the original space with position
space \(X\). Pull both halves back by \(\xi^{-1}\), writing
\(X'=\xi^{-1}X\), and transport the position measure. In the
Schr\"odinger models the map from old functions to new functions is
\(D\phi(y)=\phi(\xi y)\), for \(y\in X'(\mathbb A_F)\).
Translation by \(v\) in the new coordinates is translation by
\(\xi v\) in the old ones. Multiplication by a character from the
other half gives the same identity by the pairing formula. Thus
\(D\) intertwines the two Heisenberg actions with the stated
automorphism. Conjugating each old unitary Weil operator by \(D\)
gives a lift of the same group element for the new central character.
These operators multiply as before, giving the transported splitting.
For every pair \((g,h)\) the theta distribution now gives
\[
  \sum_{y\in X'(F)}\big(D\omega_\beta(g,h)\phi\big)(y)
     =\sum_{z\in X(F)}\big(\omega_\beta(g,h)\phi\big)(z).
\]
This is a bijective change of rational summation variables. In
particular it respects the rational splittings and automorphy of the
kernel. Applying the identity on each edge gives identical functions
before each successive integration. Integration against the original
Bessel character, and then applying the same projection, preserves
equality. If the character is written \(\psi_\beta(\ell(n))\) for a
linear functional on the unipotent quotient, its expression with
the new additive character is \(\psi_{\beta'}(\ell(n)/u)\).
This specifies the frequency in the new notation even for a
character with several nonzero coordinates.
\end{proof}

\begin{proposition}[The finite test at the one additional split place]
\label{prop:split233}
Use the norm representative in Theorem~\ref{thm:norm233} and the
split position model at each split place, taking its position half
at the chosen prime of \(E\). Protect every prescribed nonspherical
split test by choosing \(a_v\) so scalar units congruent to one to
that depth fix its position function. All such position functions
can be retained under transport. At every other finite place except
\(v_*=P\cap F\) the standard lattice vacuum can also be retained.
At \(v_*\) the only change, up to unitary phase in the local model,
is a central translate of the spherical test. On an input with
one unitary central character at \(v_*\), removing these translates
on all edges changes the resulting lift by a scalar of absolute
value one. The pointwise identity for a signed sum of different
central characters uses the translated tests of the preceding theorem.
\end{proposition}
\begin{proof}
Locally use the same position Haar measure for both split models
and the corresponding dual measures in Fourier inversion. On the
position half of dimension \(n\) the unitary transport can then
be taken to be
\[
       \phi(x)\longmapsto |\xi_w|_v^{n/2}\phi(\xi_w x).
\]
On translations this gives the coordinate change by \(\xi_w\).
On characters the other coordinate scales by
\(\xi_{\bar w}=u/\xi_w\), giving the new additive character. The
positive factor is fixed by the change of variable in the squared
norm. Unitary identifications with other polarized models differ
by a phase for the same Heisenberg intertwining relation, by
irreducibility. At \(v\in S\) the congruence on \(\xi_w\) makes
the displayed operator fix the test. Every Schwartz function admits
such a depth, by compact support and local constancy. At another
split place with \(w(\xi)=0\), scalar multiplication preserves any
full integral lattice. At an inert place \(\xi\) is a unit in the
quadratic completion. It carries the inverse image of the same
self-dual lattice in the Heisenberg group onto itself and pulls
back its inducing character to the new one. In the lattice model
this carries the vacuum normalized to have value one at the
identity to the correspondingly normalized vacuum.

At \(v_*\) take \(w=P\), so \(w(\xi)=1\). In the split unitary
group the central norm-one pair \((\xi_w,\xi_w^{-1})\), or its
inverse according to the pullback convention, has this very action
on the position half. Its Weil operator for the transported splitting
is the displayed unitary dilation times a scalar of absolute value
one. This follows directly from its linear action and its inverse
dual action in the split polarization. Translating the Schwartz
vector by that central element translates the source argument of
the kernel. Changing variables in the source integral therefore
extracts a value of the unitary central character or its inverse.
At every next edge the lift still has a single central character
there. Indeed a scalar on the source paired with the inverse scalar
on the target acts trivially on the tensor space, hence by a scalar
in the Weil representation. Moving one scalar through the integral
gives the asserted target character. Iteration proves the claim.

For the rational models in Theorem~\ref{thm:transport233} the
summation formula fixes the global phase. In a separate local model
the phase must travel with a summand if coefficients are added
before taking an absolute value. At a real place \(\xi\) acts by a
positive dilation \(\sqrt{\tau u}\) followed by a norm-one scalar.
The transported real vectors must therefore also be included in
that theorem. The size of the fixed-coordinate real test is given
by Theorem~\ref{thm:scale232}.
\end{proof}

\begin{theorem}[The original finite tests at both endpoints of an even chain]
\label{thm:even233}
Use the tensor convention and the intermediate determinant factors
of equation~\eqref{eq:iter}. Suppose the number \(m\) of
edges is even, as in Theorem~\ref{thm:rangenew50}. Fix
one splitting and one set of unitary model identifications per edge.
Allow a finite set of columns \(c\), with homogeneous real Fock
polynomials of degrees \(d_{i,\tau,c}\) on edge \(i\), and with
finite-order intermediate determinant characters \(\chi_{i,c}\).
Use the inert lattice vacua and the good split lattice indicators of
Proposition~\ref{prop:split233}. Choose the norm representative to
protect the other split test functions of all columns. Write \(A_c\)
for the original operator and
\(A'_c\) for the operator at \(\psi_{\beta'}\) with the transported
splittings, the same finite test shapes and the same real polynomials
in fixed coordinates. Use unit-norm real vacua, the split position
measures of Proposition~\ref{prop:split233}, and the tensor product
of unitary model identifications on each side. Then there is a
single \(C\in\CC\) of modulus one, independent of the columns,
such that, on every input for which the chain converges,
\[
  A_c=C\,a_c\zeta_c A'_c,\qquad
  a_c=\prod_{\tau}(\tau u)^{D_{\tau,c}/2},\qquad
  \zeta_c=\prod_{\substack{1\le j<m\\ j\ {\rm odd}}}
                       \chi_{j,c}(z)^{-n_j},\qquad
  D_{\tau,c}=\sum_{i=1}^m d_{i,\tau,c}.
\]
Here \(z\) is a norm-one idele specified below, \(n_j\) is the rank
of the intermediate group and every positive square root is taken
in the chosen real embedding. This is an equality of functions on
the original target. It consequently applies to the original
terminal Bessel character and projection for arbitrary source vectors.
\end{theorem}
\begin{proof}
Use the convention that a linear map \(L\) on a position space
acts on functions by pullback by \(L^{-1}\), with its unitary
modulus factor. At \(v_*\) let the first component of \(z\) be
\(\xi_P^{-1}\) and its other component \(\xi_P\), and at each
real place set \(z_\tau=(\sigma\xi/|\sigma\xi|)^{-1}\) in the
chosen complex embedding \(\sigma\mid\tau\). Take \(z_v=1\) at
the other finite places. The finite comparison is now exactly the
central translation in Proposition~\ref{prop:split233}, up to a
phase independent of the test.

At infinity take a compatible complex structure invariant under the
compact scalar circle, as in the Fock model of Proposition~\ref{prop:dilation179}. Write multiplication by
\(\sigma\xi\) as a rotation \(U_\tau\) in this circle times
\(\sqrt{\tau u}\). Pullback sends a homogeneous polynomial \(P\) of
degree \(d'\) to \((\tau u)^{d'/2}P(U_\tau y)\). It is the unitary
map between the two Fock norms since the Gaussian parameter is
multiplied by \(\tau u\) and the normalization of the vacuum
cancels the change of real volume. Pullback by the rotation is the
action of \(z_\tau\) in the source center up to its vacuum phase.
This phase concerns the representation and is independent of \(P\).
These unitary local intertwiners and the rational transport in
Theorem~\ref{thm:transport233} satisfy the same Heisenberg
intertwining relation. Irreducibility gives a single phase at each
place, and at almost all places the
identifications agree on the unit vacuum. Thus the transported
kernel on edge \(i\) equals its fixed-coordinate kernel at
\((zg_{i-1},g_i)\), multiplied by the positive degree factors on
that edge and a phase independent of \(c\).

For each odd edge \(i\) replace the arguments \((zg_{i-1},g_i)\)
by \((g_{i-1},zg_i)\). The scalar introduced is independent of the
column since \((z,z^{-1})\) acts trivially on the symplectic tensor
space and its lift is central. This is the kernel identity of
Theorem~\ref{thm:center}. On each even edge the translation
was already on its odd-indexed source. The translated arguments are
now precisely all occurrences of \(g_j\) with odd \(j\). Since
\(m\) is even, they are all integration variables. Change variables
to \(g'_j=zg_j\) there. The quotient measures are invariant. The
intermediate factor becomes
\(\chi_{j,c}(\det g'_j)\chi_{j,c}(z)^{-n_j}\).
Neither endpoint is translated. Multiplication gives the displayed
identity and the same overall phase \(C\) for every column.
\end{proof}

\begin{corollary}[Valuations of the projected Gram minors]
\label{cor:gram233}
Apply the same complex-linear terminal map, for example the specified
Bessel coefficient followed by its projection, to both operators
in Theorem~\ref{thm:even233}. For any finite set of columns and
chosen inputs let \(G,G'\) be the Gram matrices of the resulting
vectors, using an inner product linear in the second argument.
Put \(D=\operatorname{diag}(a_c\zeta_c)\). Then \(G=D^*G'D\).
If the entries of one matrix are algebraic, so are those of the
other. Choose a number field containing \(F\), these entries, the
scalars and their conjugates. Let \(v_{\mathfrak l}\) be a finite
valuation of this field.
For every nonzero minor with row and column sets \(I,J\) of equal
size one has
\[
 v_{\mathfrak l}(\det G_{I,J})-v_{\mathfrak l}(\det G'_{I,J})
   =\tfrac12\sum_\tau\left(\sum_{c\in I}D_{\tau,c}+
             \sum_{c\in J}D_{\tau,c}\right)v_{\mathfrak l}(\tau u).
\]
The two minors vanish together. Away from the rational primes
supporting \((u)\), the corresponding minors have identical
valuations. In particular the additional split prime below \(P\)
has zero correction at every place above it.
\end{corollary}
\begin{proof}
The common phase cancels in the inner product. Each \(a_c\) is
positive algebraic, since its square is the product of the indicated
algebraic conjugates of \(u\) to integer powers. Each \(\zeta_c\) is
a root of unity. Taking a minor of the matrix identity multiplies
it by \(\prod_{c\in I}\overline{a_c\zeta_c}\prod_{c\in J}a_c\zeta_c\).
Its valuation gives the formula. All conjugates of \(u\) are units
away from the rational support of its ideal, which was excluded in
the choice of \(P\).
\end{proof}

For example, suppose the determinant tuples range over a finite
group of twists and each \(D_{\tau,c}\) is independent of \(c\).
Then \(a_c=a\) is constant and \(c\mapsto\zeta_c\) is a character
\(\zeta\) of that group. In a Fourier projection weighted by
\(\bar\eta(c)\), the identity reads \(A_\eta=C a A'_{\eta\bar\zeta}\),
since \(\bar\eta(c)\zeta_c=\overline{(\eta\bar\zeta)(c)}\).
For a general signed sum the individual weights \(a_c\zeta_c\)
must be kept. The minors and the norm transport above give an
exact change of normalization. The required arithmetic lower bound
for the projected sum and its comparison with the full global
orbital contribution remain open in Section~\ref{sec:missing}.

\subsection{Integral restrictions on the relative scalar}\label{nav:VI-31}
\begin{proposition}[The scalar congruence when the isotropic flags meet]
\leavevmode\label{prop:finite237}
For a relevant, finitely supported pair in
Theorem~\ref{thm:suffix237}, put \(a=D_t/Q^t\), taking
\(D_0=1\). Then
\[
 a\in\mathcal O_E\setminus\{0\},\qquad
 a\lambda\in\mathcal O_E,\qquad a(1-\lambda)\in Q\mathcal O_E .
\]
If \(\lambda\ne1\), necessarily
\(|N_{E/\QQ}a|^{1/(2f)}\ge Q/2\). At \(t=0\) the only such
label is the identity double coset for \(Q>2\).
For every \(t\), the ideal \(\mathfrak d\) in
Proposition~\ref{prop:denom235} satisfies
\(N_{E/\QQ}\mathfrak d\le |N_{E/\QQ}a|\).
\end{proposition}
\begin{proof}
Choose local integral representatives of the two decorated
frames at a finite place, as in Lemma~\ref{lem:flag215}.
Their prefixes are primitive. Above \(Q\) choose their
reductions equal by upper unitriangular changes, as in
Proposition~\ref{prop:shift203}. All prefix minors are
unchanged by such choices. When \(t>0\) the integral
matrix \(B=(h(u_{0,i},u_{1,j})/Q)_{i,j\le t}\) is nonsingular
with determinant \(a\). The entries of
\(y=(h(u_{0,i},u_{1,t+1})/Q)_{i\le t}\) are integral for
the same reason: above \(Q\) both isotropic lists have the
same reduction and at the other places \(Q\) is a unit.
Let \(z=B^{-1}y\), so that \(az\) is integral by the adjugate
identity. The vector
\[
 v=u_{1,t+1}-\sum_{j\le t}z_j u_{1,j}
\]
is the unique vector in \(I\cap X_{1,t+1}\) of leading
coefficient one in the second list. Indeed it is perpendicular
to \(X_{0,t}\); the prefix pairing is perfect and its kernel
on \(X_{1,t+1}\) is the one-dimensional intersection with \(I\).
Its leading coefficient in the first list is \(\lambda\) by
Theorem~\ref{thm:suffix237}. At \(t=0\) take the empty system,
\(a=1\) and \(v=u_{1,1}\).

The vector \(av\) is integral and belongs to \(X_{0,t+1}\).
By primitivity its coordinates in the first prefix are integral,
so its last coefficient \(a\lambda\) is integral. Above \(Q\)
the expression in the second list, reduced modulo the prime,
has last coefficient \(a\) in the first reduced prefix, since
every \(az_j\) is integral. Its reduced vectors are linearly
independent; hence \(a\lambda\equiv a\) there. This local
argument is made in both factors at split places. It proves
the global inclusions. If \(\lambda\ne1\), take the norm of
\(a(1-\lambda)\), whose absolute value is at most
\(2^{2f}|N_{E/\QQ}a|\) and at least \(Q^{2f}\). In particular
for \(t=0\) one has \(\lambda=1\); the two whole decorated
lists then agree by Theorem~\ref{thm:suffix237}.

In the integral map of Proposition~\ref{prop:denom235} the
images of the first \(t\) vectors on either side span sublattices
of full rank in \(U_i=(\Lambda\cap X_i)/(\Lambda\cap I)\).
In these sublattices the determinant of the pairing divided
entrywise by \(Q\) equals \(\bar a\) in the order of slots
used there. Passing from \(U_0,U_1\) to the sublattices
adds two nonnegative local determinant valuations (with
conjugation on the first slot). Thus the determinant valuation
for the full lattices is no larger. The largest elementary
divisor is bounded by this determinant valuation at every
prime. Taking norms gives the last assertion.
\end{proof}

\begin{proposition}[An orbit attaining the scale on every initial column]
\leavevmode\label{prop:sharp237}
For \(s\ge2\) the original finite test on every refined
initial column admits one and the same relevant rational label
with \(t=1\), \(\lambda=1\) and
\[
                w_\gamma^{-1}=(1+16 Q^2 q)^{-2f}.
\]
Here \(q\) is the auxiliary rational prime with
\(\sqrt{-q}\in E\). There is a compatible cell whose two
oriented class coordinates agree and whose rational isometry
in Corollary~\ref{cor:trace235} is the identity.
\end{proposition}
\begin{proof}
Use the first coordinate plane of Section~\ref{sec:char94}
with \(u=e_++e_-\), \(v=(e_+-e_-)/2\), \(h(u,v)=1\).
Set \(c=4Q\sqrt{-q}\) and define
\(\gamma w=w+c\,v h(v,w)\).
Since \(h(v,v)=0\) and \(c+\bar c=0\), expansion of the
pairing shows \(\gamma\) is an isometry with inverse obtained
by negating \(c\). It sends \(u\) to \(u+cv\) and fixes its
orthogonal coordinate-plane complement. In the signed integral
coordinate basis both denominators from \(v\) divide \(4\)
together, so \(\gamma-1\) has entries in \(Q\mathcal O_E\).
It therefore belongs to the marked level at \(g_1,g_2\) and
preserves the standard lattice everywhere else. At \(T\) the
element \(a^0\) defining \(K^0\) acts only on the middle
space of the Witt decomposition and commutes with \(\gamma\)
by Theorem~\ref{thm:midlevel131}. Thus \(\gamma_f\in K^0\).

The common isotropic subspace is spanned by \(u_2,\ldots,u_s\).
The two terminal vectors coincide and the only nonzero leading
minor is \(h(u,u+cv)=c\). Both radicals fix their own first
vector pointwise. An element of their intersection fixes
both \(u,u+cv\), hence their nondegenerate plane pointwise,
and preserves its perpendicular. The remaining two lists on
that perpendicular coincide. Thus every simple character
coordinate agrees: the first is zero on both sides and the
remaining ones are read on the same vectors.
The label is relevant and has the asserted weight since
\(|\tau c|^2=16Q^2q\) at every embedding.

Represent any refined class by \(h\in H(\mathbb A_{F,f})\),
where \(H=U(W_0)\) fixes the first coordinate plane.
Use \(h\) as the finite cell representative. Both finite
coordinates are in the required support since
\(\gamma_f^{-1}h=h\gamma_f^{-1}\) and \(\gamma_f\in K^0\).
They give exactly the same column representative. The common
positive complement may be chosen to be \(W_0\); both
projections are the identity there since \(\gamma\) fixes it.
Thus the class factors may both be taken to be one, and the
rational isometry in the trace is the identity. The assertion
concerns finite support and character compatibility. The
archimedean coefficient on this cell is still part of the
period integral.
\end{proof}

\subsection{The size of the complete column space}\label{nav:VI-32}
\begin{lemma}[Column dimension at the original period level]
\label{lem:dimension239}
Fix a translate on an interval of Theorem~\ref{thm:region198}
and use the covector refinements at its original prime \(T\).
Let \(D_c\) be the full finite multiplicity space for word \(c\)
and \(D_{\rm all}=\bigoplus_c D_c\) over the selected words.
Let \(D_*\) denote the real comparison dimension called \(D\) in
Theorem~\ref{thm:common111}. For sufficiently large \(x\),
uniformly over the translates and selected data of that theorem,
\[
          \log(2+\dim D_{\rm all}+D_*)\le x^{40}.
\]
\end{lemma}
\begin{proof}
Write \(F=E^+\), \(f=[F:\QQ]\), \(Q=g_1g_2\), \(b=\exp(x^\eps)\),
and let \(N\) be the cyclotomic modulus containing \(E\).
By Corollary~\ref{cor:rank111}, the dimension for one word is
\(\mathfrak m_k I_\pi(T^k-1)^f\prod_\tau W_k(u_\tau(c))\).
Here \(\mathfrak m_k\) is the mass of the definite genus,
\(I_\pi\) the column index at the two auxiliary primes, and
\(W_k\) the compact Weyl dimension. The mass is bounded by the
number of full-level classes because each term is the reciprocal
of a positive integer. The full-level count in the proof of
Lemma~\ref{lem:mass32}, with underlying integral rank \(2fk\),
gives \(\log^+\mathfrak m_k\ll(2fk)^4\log(4fkN)\).
Counting all residue matrices gives \(I_\pi\le Q^{fk^2}\).
At every real place the shifted weights have spread at most
\(3a\) in the construction of Theorem~\ref{thm:common111},
where \(a=O(xb)\). The product formula for \(W_k\) therefore
gives \(\log W_k\ll k^2\log(2+k+xb)\).
Together these bounds give
\[
 \log(2+\dim D_c)\ll
  1+(2fk)^4\log(4fkN)+fk^2\log Q+fk\log T
                    +fk^2\log(2+k+xb).
\]
The ranges \(0<\eps<1/52\), \(f\ll_\psi x^{2+8\eps}\),
\(k\ll x^{9/2+10\eps}\) and \(m\ll\sqrt x\) give
\(f\ll x^3,\ k\ll x^5\). On the same intervals
\(\log N=o(x^\eps)\), \(\log Q=(16+o(1))x^\eps\) and
\(\log T=(17+o(1))m x^\eps\) by Theorem~\ref{thm:prime210}.
The last display is \(O_{\psi,\eps}(x^{33})\).
There are at most \(2^{k-2}\) words. Finally
\(\Gamma_c^2\dim D_c/D_*=1+O(\exp(-x^{1/4}))\) by
Corollary~\ref{cor:rank111}, with
\(\log\Gamma_c\ll f m^5 x^\eps\).
These bounds imply the assertion after enlarging the initial scale.
\end{proof}

\subsection{Denominators for compact unitary integration}\label{nav:VI-33}
In the compact period of Theorem~\ref{thm:edge242} the Haar
average acts on explicit polynomials in unitary matrix entries.
The following denominator is independent of their coefficients and
of the number of finite classes. Write \((h)_b=h(h+1)\cdots(h+b-1)\)
with \((h)_0=1\).

\begin{proposition}[An integral multiple of every compact moment]
\label{prop:moments242}
Let \(p(g,\bar g)\) be a polynomial in the entries and conjugate
entries on \(U(n)\), of degree at most \(L\) in each set of
variables. Let \(I\) be the fractional ideal generated by its
coefficients in a number field. With Haar probability \(dg\), set
\[
 \mathcal D(n,L)=\prod_{j=1}^{\min(n,L)}(n-j+1)_{L-j+1}.
 \qquad
 \mathcal D(n,L)\int_{U(n)}p(g,\bar g)\,dg\ \in I .
\]
The same assertion on a product of unitary groups uses the product
of the corresponding integers. Their prime factors on a given group
are at most \(n+L-1\).
\end{proposition}
\begin{proof}
It suffices to integrate a monomial. Right multiplication by
diagonal unit complex numbers makes its integral zero unless, in
each column \(j\), the numbers of ordinary and conjugate entries
are equal, say \(k_j\). Right permutation of columns puts the
\(s\) occupied columns first. Thus \(s\le\min(n,L)\),
\(k_j\ge1\) for \(j\le s\), and \(\sum k_j\le L\).

Conditioned on the first \(j-1\) columns \(v_a\) of a Haar
unitary matrix, column \(j\) is a uniform unit vector \(v\)
in their orthogonal complement. This follows from unitary
invariance and transitivity of the stabilizer on that sphere.
Set \(h=n-j+1\) and \(Q=1-\sum_{a<j}v_a v_a^*\). If \(j\)
is the last column left in a monomial and occurs in \(b\)
matched pairs, then its conditional integral is
\[
 \mathbb E\!\left[\prod_{\ell=1}^b v_{i_\ell}
           \overline{v_{t_\ell}}\ \middle|\ v_1,\ldots,v_{j-1}\right]
   =\frac{1}{(h)_b}\sum_{\sigma\in S_b}
                  \prod_{\ell=1}^b Q_{i_\ell,t_{\sigma(\ell)}}.
\]
To prove the identity take a standard complex Gaussian in this
\(h\)-dimensional subspace. Its squared radius is a sum of \(h\)
independent mean-one exponentials, with \(b\)-th moment
\((h)_b\), by its generating function \((1-t)^{-h}\).
Its angle is independent and uniform. The mixed moments of the
Gaussian are the sum of products on the right before division.
Indeed the coefficient obtained by differentiating its generating
function \(\exp(\sum_{ij}x_i Q_{ij} y_j)\) once for each
labeled occurrence is exactly the sum over the matchings.
Dividing by the radial moment proves the formula.

Integrate the occupied columns from \(s\) down to \(1\).
Expanding a product of \(b\) entries of \(Q\) creates at most
\(b\) matched pairs in preceding columns, with integer
coefficients. The earlier original factors are still present.
At step \(j\), at least \(j-1\) original pairs remain in the
first \(j-1\) columns, so the number \(b\) of pairs at column
\(j\) satisfies \(b\le L-j+1\). Its denominator divides
\((n-j+1)_{L-j+1}\). Multiplication of these common
denominators proves the assertion, including \(L=0\).
Successive integration proves the product assertion. The
largest factor in any rising factorial is at most \(n+L-1\).
\end{proof}

For example, on \(U(n)\) with \(n\ge2\) the degree-two bound is
\(n(n-1)(n+1)\). For local computations the exact cost of the
displayed denominator at a rational prime \(p\) is
\[
 v_p\mathcal D(n,L)=\sum_{j=1}^{\min(n,L)}\ \sum_{a\ge1}
 \left(\left\lfloor\frac{n+L-2j+1}{p^a}\right\rfloor
             -\left\lfloor\frac{n-j}{p^a}\right\rfloor\right).
\]
This follows by counting multiples of \(p^a\) in each interval
of factors. An alternative evaluation of unitary moments uses
the character theory of the symmetric group.
\href{https://arxiv.org/abs/math-ph/0402073}{B.~Collins and
P.~\'Sniady, \emph{Integration with respect to Haar measure}}.

\subsection{The finite change of level on a theta lift}\label{nav:VII-28}
Keep the fields and intervals of Theorem~\ref{thm:invar148} and
the levels and terminal characters of Theorem~\ref{thm:twofreq160}.
On the first edge of the indefinite tower write \(w=r+t_{\rm r}\)
for the rank of \(Y_1\), so that \(Y_0\) has rank \(r\).
At the \(f\) places of \(E^+\) above \(T\), choose either of the
two matrix realizations of this split pair. Use the sign
\(\epsilon_v=1\) when the row domain is the factor in which
\(\varpi\) has valuation one and \(-1\) when it is its dual.
Position matrices \(Y\) have \(w\) rows and \(r\) columns. Put
\[
    h_\theta=\frac{sw}{2}\sum_{v\mid T}\epsilon_v,\qquad
    h_B=\frac{fs(r-s)}{2},\qquad d=T^{h_B+h_\theta}.
\]

\begin{theorem}[Two projected periods evaluated on theta rows]
\label{thm:thetarow162}
Fix a Schwartz test \(\Phi_0\) on the first theta pair, with
arbitrary real vector and fixed splitting. Form \(\Phi_1\) at
each \(v\mid T\) by the unscaled substitution
\(Y\mapsto Y a_{v,\epsilon_v}\) in \(\Phi_0\), where \(a_{v,+}\)
is the action of \(a_\varpi\) on the first split row space and
\(a_{v,-}\) is its inverse transpose on the dual. At finite
places away from \(T\) translate the test by the row Weil
action of \(a_f\), and keep its real component fixed. These
operations act linearly on the tensor factors. At \(v\mid T\)
the substitution gives a bijection of coefficient-valued
Schwartz functions, preserving their ideals of values.

Let \(F_i\) be the theta rows formed with these two tests and
the same automorphic input on \(U(Y_1)\). Whenever the lifts
are realized by convergent integrals and the rows are in
automorphic \(L^2\), for every common refined class tuple
\(\ell\), every selected word and matching real Bessel vectors,
\begin{equation}
   \left|\big\langle F_1,\mathcal P_{1,1}(\ell)\big\rangle\right|
      = d\,
   \left|\big\langle F_0,\mathcal P_{\varpi,0}(\ell)\big\rangle\right|.
   \label{eq:thetarow162}
\end{equation}
This also applies when the input on \(U(Y_1)\) is produced by
earlier integrations in a convergent iterated lift. The number
\(d\) is positive rational. For every fixed \(A>0\) it satisfies
on these intervals
\[
     \log\big(\operatorname{num}(d)\operatorname{den}(d)\big)
         \le \tfrac12 fs(w+r-s)\log T=o(b/x^A).
\]
For the same choice of first factor at all places, \(d\) is
the integer \(T^{fs(w+r-s)/2}\).
\end{theorem}
\begin{proof}
In the stated matrix model the determinant of \(a_{v,\epsilon_v}\)
has valuation \(s\epsilon_v\). The unitarily normalized row Weil
action therefore multiplies the unscaled substitution by
\(\eta_v(a_v)T^{-sw\epsilon_v/2}\), with \(|\eta_v(a_v)|=1\),
by equation \ref{eq:weil161}.
The unscaled substitution is a bijective linear change of
variables on positions, giving equality of value sets and
preservation of local constancy and compact support in both
directions. Put \(z=\prod_{v\mid T}\eta_v(a_v)\). In the
global Weil representation these definitions give
\[
          \omega(a_f,1)\Phi_0=zT^{-h_\theta}\Phi_1 .
\]
To see the effect on the rows, write the theta kernel as the
rational theta functional applied to \(\omega(g,h)\Phi\).
In a global position model this functional is the sum over the
rational position vectors. For fixed \(g,h\) the finite
support conditions bound their denominators and the real
Schwartz function decreases faster than any inverse power.
This proves convergence of the sum. Changing the realization
at finitely many places conjugates the action by its local
intertwiners in this identity. Since
\(\omega(g,h)\omega(a_f,1)=\omega(g a_f,h)\), the two
theta kernels agree pointwise after the same translation and
scalar. Integrating against the given input yields
\(R(a_f)F_0=zT^{-h_\theta}F_1\). This argument applies with
any fixed input for which the displayed theta integrals exist.

The first identity of Theorem~\ref{thm:twofreq160} gives
\(\mathcal P_{1,1}=T^{h_B}R(a_f)\mathcal P_{\varpi,0}\) on
the same ordinary column inputs. Taking the indicated inner
products and using unitarity of right translation gives
\eqref{eq:thetarow162}, including zero values on both sides.
By Theorem~\ref{thm:commonnorm128}, \(f\) is even. The sum of
the signs therefore is even as well, so both exponents are
integers. The sum of the logarithms of the numerator and
denominator of \(d\) in lowest terms is
\(|h_B+h_\theta|\log T\). The ranks, \(f\) and \(\log T\)
are bounded by fixed powers of \(x\) as in
Corollary~\ref{cor:cost128}. This proves the claimed bound
because \(b=\exp(x^\eps)\).
\end{proof}

\begin{corollary}[Integral first-edge tests on the period fields]
\label{cor:finite162}
At each \(v\mid T\) take either the lattice vacuum or the
residue function with invertible top \(r\)-by-\(r\) block and
zero bottom \(t_{\rm r}\)-by-\(r\) block. Use the lattice and
its dual of Theorem~\ref{thm:twofreq160} and form the orbit of
the level-\(K_1\) tests under \(K_0\) at these places. In the
ring of integers containing the residue character values,
their \(I={r\brack s}_T^f\) coordinates can all be recovered
by evaluation up to roots of unity. The finite average with
its determinant character has exact denominator \(I\), which
is a unit at \(T\). For unit vectors in the finite Weil
representation, their Gram matrix satisfies
\[
      \|\mathcal G_T-1\|_{\rm op}
           \le 64fT^{-(t_{\rm r}+1)}=o(e^{-x^{1/4}}).
\]
If every place uses the residue function, \(\mathcal G_T=1\).
The negative logarithm of the ratio of the product squared
norm to that for the unnormalized vacua is \(o(b/x^A)\) for
every fixed \(A>0\).
\end{corollary}
\begin{proof}
Theorem \ref{thm:eval161} constructs at each lattice a matrix
that takes its basis to the first \(r\) coordinate vectors of
the partner. Evaluation on products of those matrices gives
the claimed recovery. Proposition~\ref{prop:exact161} computes the
denominator by the same evaluations. The product estimate of
Theorem~\ref{thm:gram161} gives the stated norm bound by
\((1+u)^f-1\le e^{fu}-1\le2fu\) for \(fu\) sufficiently
small, with \(u=32T^{-(t_{\rm r}+1)}\).
Here Lemma~\ref{lem:size97} and the prime range of
Theorem~\ref{thm:common111} give
\(\log T\ge(16-o(1))m x^\eps\), with
\(m\asymp x^{1/2-\eps/16}\) and \(f\) polynomial in \(x\).
This yields the asserted decay. At a residue place the test
matrix determines its own input lattice as the preimage of
the integral partner lattice, which accounts for exact
orthogonality. Equation \ref{eq:prvol161} gives a cost at
most \(rt_{\rm r}\log T+\log 4\) per residue place by
\(\prod_{h=1}^r(1-T^{-h})\ge1/4\). Multiplication over the
places and the polynomial bounds used above give the final
estimate.
\end{proof}

The equality in Theorem~\ref{thm:thetarow162} concerns two
values of actual theta rows when the tests are transported
as stated. Corollary~\ref{cor:finite162} recovers the integral
coefficients before global summation. A sum over rational
positions and projection can annihilate a nonzero test.
To deduce \eqref{eq:window} from these results one still needs
the nonvanishing, arithmetic normalization and uniform estimates
for the designated projected theta forms on the full chain,
together with the remaining strict orders in \eqref{eq:tail114}.

\subsection{Changing the terminal frequency by a rational isometry}\label{nav:VI-45}
Use the terminal space \(S_m=C_m\perp U\) of
Lemma~\ref{lem:topchar266}, with \(H=U(C_m)\) and
\(G=U(S_m)\). The staircase \(U\) consists of \(s=m/2-1\)
hyperbolic pairs \(b_j,b_j^\vee\) and a middle vector \(z\) of norm
\(-1\). Let \(N\) and \(c_j\) denote the flag radical and its
additive coordinates in that lemma. Normalize each of \([N]\)
and \([H]\) to have mass one. For a character \(\chi\) trivial on
\(N(F)\) put
\[
 F_\chi B(h)=\int_{[N]} B(nh)\overline{\chi(n)}\,dn,
 \qquad R(g)B(y)=B(yg).
\]
Use the additive character \(\Psi_E\) of Section~\ref{sec:char94}.
Write the character \(\chi_+\) of the raw period as
\(\Psi_E(\sigma\sum_jc_j)\), with \(\sigma\in E^\times\).
Indeed the additive characters of \(\mathbb A_E/E\) are precisely
the rational multiples of \(\Psi_E\), by adelic additive duality.
This absorbs both the choice of sign and the scalar \(\xi_0\)
in Lemma~\ref{lem:topchar266}.

\begin{lemma}[A change of frequency on the whole terminal quotient]\leavevmode
\label{lem:flag272}
At each point of the common intervals one can choose the common
prime \(T=\varpi\bar\varpi\) of Theorem~\ref{thm:common111} outside
all the marked primes of the horizontal grid, for fixed grid data.
On \(G\) set \(\chi_\varpi=\Psi_E(\sum_{j<s}c_j+\varpi c_s)\).
There is \(t\in G(F)\) centralizing \(H\) for which
\[
                   F_{\chi_\varpi}(R(t)B)=F_{\chi_+}B
\]
as functions on \([H]\). The last vector defining \(\chi_\varpi\)
is \(\varpi z\) and has norm \(-T\).
For the horizontal lifts of Theorem~\ref{thm:bottom271}, the common
function is fixed by \(K_v=\operatorname{Stab}_{H(F_v)}
(\Lambda_{C_m,v})\) at each \(v\mid T\).
\end{lemma}
\begin{proof}
The marked set lies above the two original level primes, the two
type primes and a single auxiliary rational prime. This follows by
taking the union of the primes of \(g_1g_2p_1p_2\) with the one
chosen place above \(p_0\) in the product defining \(\mathfrak m\).
\srcloc{Development/AbelianVarietyTheory/Definitions/Divisors/LeastSplitAuxiliaryPrimeThetaChainSpec}{250--267}.
\srcloc{Development/AbelianVarietyTheory/Definitions/Divisors/BasePlaceNotMemChiGridBadPlaces}{70--82}.
The number of admissible common primes in
Theorem~\ref{thm:common111} tends to infinity. Discarding at
most five preserves its conclusion. Use this one choice for
the comparisons at the point \(x\).

Put \(a_s=\sigma/\varpi\) and \(a_j=\sigma a_{j+1}\) for \(j<s\).
Let \(t\) send \(b_j\) to \(a_j b_j\) and \(b_j^\vee\) to
\(\bar a_j^{-1}b_j^\vee\), fixing \(C_m\) and \(z\). It preserves
all the pairings of these summands, so is a rational isometry.
Conjugating \(n\) by \(t\) multiplies \(c_j(n)\) by
\(a_j/a_{j+1}\) for \(j<s\), and \(c_s(n)\) by \(a_s\).
It follows that \(\chi_\varpi(tnt^{-1})=\chi_+(n)\).
By left rational invariance and commutation with \(H\),
\(B(nht)=B(t^{-1}nt\,h)\). Conjugation gives a bijection of
the rational unipotent quotient and preserves its probability
measure, since it carries invariant measure to a multiple of
itself. Changing variables proves the identity. Linearity in the
last vector gives the assertion about its norm and character.

At the selected places the last horizontal test is the lattice
vacuum by Proposition~\ref{prop:rowfinite262}. It is fixed under
the full target lattice stabilizer with the prescribed splitting.
The terminal lattice is the orthogonal sum of the core lattice
and the staircase ideals, so \(K_v\) embeds into this stabilizer.
Right translation by elements of \(H(\mathbb A_F)\) commutes with
the integral \(F_{\chi_+}\) by its definition. This proves the
last assertion. In particular a nontrivial character of a finite
subgroup of \(\prod_{v\mid T}K_v\) has zero projection on this
function before changing its finite test.
\end{proof}

\subsection{Embedding the original graph and its Bessel character}\label{nav:VI-46}
Take the common fields of Theorem~\ref{thm:conductor232} and
use the unit CM ideal in Theorem~\ref{thm:pol49}. Fix the
horizontal grid data and choose the common \(T=\varpi\bar\varpi\)
as in Lemma~\ref{lem:flag272}. Use the raw group \(G=U(S_m)\),
its core \(H=U(C_m)\) and the simple coefficients \(c_j\).
Write \(\chi_+=\Psi_E(\sigma\sum_jc_j)\). In the absolute-trace
normalization \(\sigma=\xi_0\) up to the sign of the base
character. It is a unit at \(T\), since
\(\xi_0=-1/\sqrt{-d_0}\) for the chosen CM ideal and \(T\)
splits completely in \(E\). The staircase lines at \(T\) have
unit ideals. Indeed \(\mathfrak m\) is a unit there and the
ideal \(\mathfrak A\) of Proposition~\ref{prop:tracelat253}
is supported only on the relative different for this choice.
The staircase ideals are the products of these ideals and
their conjugates in Proposition~\ref{prop:positions264}.

\begin{theorem}[A rational embedding with the prescribed terminal lattices]\leavevmode
\label{thm:graph274}
There is a rational isometric embedding of the original row
space \(S_0\) of Section~\ref{sec:char94}, with definite
column space \(W_0\), into \(S_m\). It maps \(W_0\) into
\(C_m\). On the embedded small flag radical \(N_0\), the
restriction of \(\chi_+\) is exactly \(\chi_\varpi\).
At each \(v\mid T\) one can choose the core graph change
\(a_v\) of Proposition~\ref{prop:corehecke272} so that the
stabilizer of \(a_v\Lambda_{S_m,v}\) intersects \(U(S_0)(F_v)\)
in the full stabilizer of the prescribed \(\Lambda_{\ell_{0,v}}\)
of Theorem~\ref{thm:midlevel131}. The subgroup \(Z_v\) is central
in \(U(W_0)(F_v)\). It fixes the perpendicular of \(W_0\) and
permutes the graph lattices through their residue covectors.
\end{theorem}
\begin{proof}
Inside the positive core embed the definite space
\(W_0\perp E v_0\) with \(h(v_0,v_0)=T-1\). To construct it,
complete this space to the dimension of \(C_m\) by a positive
diagonal space of determinant
\(\det C_m/((T-1)\det W_0)\), using positive Gram representatives.
There is room for this last coordinate. The new space and \(C_m\)
have the same determinant class and all the same signatures.
Over a nonarchimedean place these invariants and the dimension
classify a Hermitian space, and at the split places the dimension
suffices. The local-to-global theorem identifies the two global
spaces. \href{https://people.math.harvard.edu/~gross/preprints/incoherent.pdf}
{B.~Gross, Incoherent definite spaces and Shimura varieties, Section 3}.

Move the embedding by a rational element of \(H\) so that at
every \(v\mid T\) the prescribed lattice of \(W_0\) and the
unit-norm-ideal line \(Ov_0\) are orthogonal unimodular
summands of the core lattice. Such embeddings exist locally.
At a split place a unitary isometry is an arbitrary invertible
map on one factor and its inverse dual on the other. The
Gram matrix on the desired \(k+1\) coordinates is invertible
over \(O=\mathbb Z_T\), so coordinate columns for a direct
summand and its corresponding dual rows realize this Gram
matrix. Completing bases gives a local isometry carrying the
rational embedding to this local one. The required integral
summand condition is open. Weak approximation holds for \(H\)
at these places by its rational Cayley chart
\(X\mapsto(1-X)(1+X)^{-1}\) on skew-adjoint matrices. In each
split local group the open chart is dense, since the vanishing
of a nonzero determinant polynomial has empty interior.
Approximation of the chart coordinates in the underlying
\(F\)-vector space proves the simultaneous assertion.

Let \(z\) be the raw middle vector of norm \(-1\). Use the
Hermitian convention linear in the second argument. Take
\[
       z'=\bar\varpi z+v_0,\qquad
       u'_j=\sigma^{\,j-s-1}b_j\quad (1\le j\le s).
\]
The norm of \(z'\) is \(-T+(T-1)=-1\), and
\(h(z',z)=-\varpi\). Thus \(z=\varpi z'+z^\perp\)
with \(z^\perp\) perpendicular to the small row. Map its
middle vector \(e\) to \(z'\), its flag generators to \(u'_j\)
and their duals to the inverse conjugate multiples of the
raw duals. This is an isometry. For \(j<s\) the small
simple coordinate equals \(\sigma c_j\). For \(n\in N_0\)
the extra perpendicular is fixed, so the coefficient of
\(b_s\) in \((n-1)z\) is \(\varpi\) times that in
\((n-1)z'\). Since \(u'_s=\sigma^{-1}b_s\), the last term
\(\sigma c_s\) is \(\varpi\) times the small last coordinate.
This proves the character identity.

At \(v\mid T\) use the first split factor, where \(\varpi=0\bmod T\).
The flag generators and their duals have been multiplied by
units. The integral line \(Oz'\) splits off orthogonally with
unit norm. On the perpendicular of \(W_0\) in the core use
the primitive coordinate \(y(X)=-h(v_0,X)\) in the first
factor, extended by zero off the core. It is primitive since
\(h(v_0,v_0)=T-1\) is a unit. Choose the other coordinates in
its kernel and use the given covector \(\ell_{0,v}\) on \(W_0\)
in the graph change \(a_v\). On \(z'\) this coordinate is
\(1-T\equiv1\bmod T\). On the integral perpendicular of
the whole small row it vanishes modulo \(T\), because
\(h(z',X)=-\varpi X_z+h(v_0,X)=0\) there. Consequently the
new lattice defined by \(y=\ell_{0,v}(w)\bmod T\) is precisely
the direct sum of the required small graph lattice and the
old perpendicular lattice. The conjugate split factor is
their dual. Intersecting stabilizers gives the claim in both
directions. The scalar action has the asserted block form.
\end{proof}

\subsection{The raw period with the larger auxiliary spacing}\label{nav:VI-47}
Use Proposition~\ref{prop:spacing276} on the original Hermitian
spaces, with the unit CM ideal and the terminal prime of
Theorem~\ref{thm:graph274}. The bottom vector in the following
theorem is any unit lowest-type vector in a given old cuspidal
realization as in Theorem~\ref{thm:first250}.

\begin{theorem}[Nonvanishing through the modified auxiliary core]\leavevmode
\label{thm:raw276}
There is a compact-finite horizontal lift \(B\) of the given
vector and an integral class coefficient \(e\) of the modified
terminal core type such that
\[
       P_{\rm raw}=\int_{[H]}e(h)F_{\chi_+}B(h)\,dh\ne0 .
\]
At every marked place each horizontal test can be a single
coset indicator at the original depth. The last test is
spherical at every \(v\mid T\). The assertion holds on each
sufficiently distant common interval for the larger spacing \(g\).
\end{theorem}
\begin{proof}
First construct a partner for the bottom vector. The combined
signed list of the two real parameters starts with
\(a_0',\ldots,a_m'\), continues with
\(v_1,v_1-1/2,\ldots,v_{k-1},v_{k-1}-1/2\), and ends with
\(-b_1,\ldots,-b_m\). The signs alternate in the first and last
parts, and the middle consists of positive entries alternating
between the two groups. The endpoints have exactly the same
signs as in Theorem~\ref{thm:first250}. Hence the real
corank-one branching law gives the required discrete summand.
The least internal gap in each parameter exceeds \(24n^2+2\)
as in that theorem. The added top gaps have only increased.
The bound \(T_0\) used in Theorems~\ref{thm:transverse250}
and~\ref{thm:direction249} is \(10^4n^4\) times one plus an
upper bound on the absolute entries. Proposition~\ref{prop:spacing276}
and the original bounds on the \(b_j\) give \(\log T_0=O(\log b)\).
Those two estimates give a vector \(v\) for the branching
embedding \(j\) such that
\(d_S d_\sigma\mathcal T(v)/|\langle v,jw\rangle|^2
\le\exp(Cfn^5\log b)\). Here \(w\) is the chosen bottom
vector in the real model and the factor \(d_S\) from compact
projection satisfies \(\log d_S=O(fn^2\log T_0)\) by the Weyl
formula. This is the ratio of the two terms in
Theorem~\ref{thm:ratio248} before their level factors. The old auxiliary
depth has \(\log N\mathfrak m>C_{\rm dep} f^2n^6\log b\)
by Theorem~\ref{thm:initial247}. It dominates this cost and
the lattice return factor \(C_n^f\), for which
\(\log C_n<6n+9\). Consequently the nonidentity tail is less
than half the identity term at a point above half the supremum
of the given bottom vector, as in Theorem~\ref{thm:first250}.
The resulting Poincar\'e series is bounded in a cuspidal
realization of the new core series, by
Proposition~\ref{prop:corecusp251} and Theorem~\ref{thm:corecopy251}.
The first period of its compact translates is continuous on
the real maximal compact group and is nonzero at the identity.
Indeed the row is integrable on its finite-volume quotient
by Cauchy--Schwarz and the compact translates are uniformly
bounded. Some compact matrix coefficient has nonzero integral
against this function by Peter--Weyl. Convolution with that
coefficient therefore gives a compact-finite partner at the
same finite level.

We check the lift through the even core for these data.
On edge \(i\) center at \(c_{m-i}'\). The old entries are at
distances at least \(\min(g/2-1,g-(d-1)/2)\), exceeding
\((d+3)/2\). Their signs and ordering relative to the center
are those in Theorem~\ref{thm:realcore253}. The complete new
plus string in the reverse calculation has the \(d\) centered
entries from \((d-1)/2\) to \(-(d-1)/2\). Thus the two real
first-occurrence tests in Theorems~\ref{thm:realcore253}
and~\ref{thm:downcore254} have identical counts and give the
successive new parameters in the two directions. In particular
any nonzero lift in either direction at these first occurrences
is cuspidal. The local real Rallis pairings are nonzero on
polynomial tests by the real doubling result used there.

At unmarked finite places the old spherical bottom list is
twisted by a unitary character. Its factor at \((d+1)/2\)
is nonzero by the absolute logarithmic bound of
Lemma~\ref{lem:sphericalcore252}. Each earlier added string
gives the factors \(L_E^{S_E}(a,\kappa_h/\kappa_i)\),
\(1\le a\le d\), of Theorem~\ref{thm:coreeuler252},
where \(\kappa_i\) is the inverse horizontal character on edge
\(i\). Their ratios have nonzero angular exponent when \(h\ne i\),
since consecutive new centers differ by \(g\). At one the
nontrivial unitary Hecke value is nonzero, and at larger integers
the Euler product converges absolutely. The denominators in
doubling are at integers \(d+j>1\). The marked local integral
is the positive scalar of Theorem~\ref{thm:markedcore253},
since the characters agree on finite units and the same
rectangular coset is used. These verifications give exactly the
local factors of the global inner product in
Theorem~\ref{thm:upcore254}. Its argument therefore gives for
every nonzero bottom core vector a chain of polynomial tests
with nonzero ascent, including vectors in a finite sum of
cuspidal constituents by its mixed inner-product identity.

Adjunction now implies that the descents span each compact
type at the bottom. In fact a vector orthogonal to all the
descents would have zero inner product with its own ascent
after pulling the kernels to the other side. The resulting
quantity is the positive squared norm of a nonzero ascent.
Projection to a bottom compact type preserves the span by
equivariance of the first polynomial test. That type has
finite multiplicity at the fixed cuspidal level. This gives
the span assertion of Theorem~\ref{thm:downcore254}. Apply
the nonzero first period to the span. On the terminal compact
quotient the integral class coefficients span its prescribed
real type, by Proposition~\ref{prop:terminalcore255}, for
arbitrary integral highest weights. Linearity selects one
coefficient \(e\) and one chain of descents with nonzero
first period.

It remains to transport this same coefficient to the horizontal
row. In its quotient coordinates all old entries at edge \(i\)
lie below \(c_{m-i}'\) by more than \((r+7)/2\). The cloud lies
below \(c_m'\), and the preceding centers are separated by \(g\).
The signed count in Theorem~\ref{thm:realrow257} gives first
occurrence and cuspidality for each nonzero horizontal output.
In the first two seesaws the same horizontal character is
used on the even core and the horizontal row of each edge.
The splitting quotient is therefore the prescribed parity
character or its inverse, just as in
Proposition~\ref{prop:cancelsquare256}. All finite test lattices
and the Fourier expansion of the negative line are the same.
The two kernel identities in Theorems~\ref{thm:firstsquare255}
and~\ref{thm:secondrow256} express the nonzero first period
through the second restriction. Expansion in real Hermites
converges in Schwartz space, and the marked expansion is finite.
It follows that one selection of tests gives a nonzero second
restriction with the given \(e\).

At each ensuing odd index, Theorem~\ref{thm:bestransfer265}
takes the nonzero Jacobi period to a nonzero Bessel period.
At an even index Theorem~\ref{thm:jactransfer266} takes it
back to a Jacobi period on the next space. Their orbital
formulas use the same lattices, additive character, rational
frames and measures. On the residual line the splitting
ratio is \((\gamma_i\eta^{-1})/\gamma_i=\eta^{-1}\)
for the new horizontal character \(\gamma_i\), as verified
in Proposition~\ref{prop:linkedline265}. The nonzero value
at the base of the orbit persists on a real neighborhood.
The supported real test and the finite coset expansion in
those theorems thus give one nonzero term at each step.
Rapid decrease on cuspidal inputs justifies the integrations
as in their proofs. This finite induction reaches \(m\)
with coefficient \(e\). The rational sign change of
Lemma~\ref{lem:topchar266} gives \(P_{\rm raw}\).
The tests are unchanged off the marked set up to the indicated
unramified splittings. In particular the terminal prime is
outside that set and its lattice vacuum is spherical.
\end{proof}

\subsection{Compact branching and determinant characters}\label{nav:V-42}
At a real place of \(F=E^+\) the embedding of the definite
column in the large core is conjugate to the coordinate inclusion
\(U(k)\subset U(N)\). Write \(\lambda_1\ge\cdots\ge\lambda_N\)
and \(\mu_1\ge\cdots\ge\mu_k\) for integral highest weights.
For the norm-one torus \(D=U(1)\) of \(E/F\), put
\(w_E=|\mu(E)|\), the number of roots of unity in \(E\).

\begin{proposition}[Branching and twists unramified at every finite place]\leavevmode
\label{prop:branch275}
The type \(\mu\) occurs in the restriction of type \(\lambda\),
followed by multiplication by \(\det^j\), precisely when
\[
       \lambda_{i+N-k}\ \le\ \mu_i-j\ \le\ \lambda_i
                               \qquad(1\le i\le k).
\]
For any integers \(j_\tau\in w_E\mathbb Z\), one for each
real place, the twists \(\det^{j_\tau}\) are the real parts
of one adelic determinant character of \(G_0\). This
character is unitary, trivial on rational points and
unramified at every finite place.
\end{proposition}
\begin{proof}
For adjacent ranks the branching rule follows from the Schur
character formula. First shift the highest weight by a
determinant power to make it a partition. Removing all boxes
of the largest letter from a semistandard tableau leaves a
horizontal strip as its complement. The remaining shapes
interlace the original one and each of their tableaux extends
uniquely for the given strip. Setting the last diagonal
entry to one and comparing characters gives the branching
rule, including the shift back by the determinant.
Iterating the inequalities proves necessity for arbitrary
ranks. For sufficiency write \(b_i=\mu_i-j\). If \(N>k\),
make a row of length \(N-1\) by putting
\(l_i=\max(\lambda_{i+1},b_i)\) for \(i\le k\) and
\(l_i=\lambda_{i+1}\) otherwise. It is decreasing and lies
between adjacent entries of \(\lambda\). It satisfies
\(l_i\ge b_i\) and \(l_{i+N-1-k}\le b_i\). Indeed both the
candidate \(\lambda_{i+N-k}\) and any candidate
\(b_{i+N-1-k}\) at that position are at most \(b_i\).
Induction ends at a row of length \(k\) forced to equal \(b\).

Write \(K_D\) for the product of the maximal finite compact
subgroups of \(D\). At split places the valuation identifies
the quotient by the compact subgroup with \(\mathbb Z\).
At nonsplit places the norm-one group is compact. Hence
\(D(\mathbb A_f)/K_D\) is the group of fractional ideals
of the form \(J/\bar J\). Each principal \(J=(a)\) gives
the global point \(a/\bar a\). The quotient by global
points is therefore finite, being an image of the ideal
class group of \(E\). The rational points with finite
components in \(K_D\) are exactly the norm-one units.
All their conjugates have absolute value one. Their powers
have bounded integral characteristic-polynomial coefficients,
which makes these units roots of unity.

On \(D(F)D(F_\infty)K_D\) prescribe trivial values on
\(D(F)K_D\) and the character
\((z_\tau)\mapsto\prod_\tau z_\tau^{j_\tau}\) at infinity.
The unit calculation makes the prescription well defined.
It is continuous because \(D(F_\infty)K_D\) is open and the
character is continuous there. Extend it through the finite
abelian quotient one cyclic factor at a time. For a new
generator whose \(l\)-th power is already in the subgroup,
choose an \(l\)-th root of the prescribed value in the unit
circle. This defines a continuous unitary extension.
Composition with determinant on \(G_0\) gives the assertion.
In particular it is trivial on the determinants of every
finite compact subgroup of \(G_0\), including the terminal
graph stabilizers and the scalars \(Z\).
\end{proof}

\subsection{A larger spacing for the auxiliary compact type}\label{nav:V-43}
Use the fields, bottom row, words and terminal prime of
Theorem~\ref{thm:word275}. Only the intermediate core weights
and their splitting characters will vary. Write \(g_*\) for
the least auxiliary gap, \(n=r+1=k+m\), \(d=n+4\) and
\(N=k+m(r+6)\). Define integers
\[
 A=1+a+b_*+8k\Delta+2k,\qquad
 g=g_*+2w_E\left\lceil\frac{8(A+N+w_E)}{2w_E}\right\rceil .
\]
In the formulas of Theorem~\ref{thm:first250} use
\(a_l'=v_1+1/2+(m+1-l)g\) and
\(c_l'=\lfloor(a_{l-1}'+a_l')/2\rfloor\).
Keep the original \(u_i,v_i,b_j,e_j\), with their place
dependence. The source distinguishes the inequalities on
the gap from the choice of their least solution.
\srcloc{Development/AbelianVarietyTheory/Definitions/Divisors/LeastSplitAuxiliaryPrimeThetaChainSpec}{121--152}.

\begin{proposition}[A common spacing for all real column weights]\leavevmode
\label{prop:spacing276}
The enlarged auxiliary parameters admit horizontal splitting
characters with the same restrictions to all finite unit
groups as before. They leave the bottom-row parameter and
the norm-\(T\) terminal frequency intact. After one determinant
twist unramified at every finite place, the compact terminal
type with the new parameters contains every prescribed word
type on restriction to \(\prod_\tau U(k)\), including at
the high places. The same holds for its dual with the twist
chosen in the dual convention. If \(d_\rho\) is the dimension
of either large type and \(d_\mu\) the dimension of a word
type, then on the common intervals
\[
       \log(d_\rho d_\mu)\ \ll_{\psi,\varepsilon}\ fN^2\log b,
              \qquad b=\exp(x^\varepsilon).
\]
\end{proposition}
\begin{proof}
The difference of the two centers at \(l\) is
\((m+3/2-l)(g-g_*)\), an integer divisible by \(w_E\).
By the torus construction in Proposition~\ref{prop:branch275} is a norm-one character with these differences as real
exponents and trivial values on every finite compact subgroup.
Pull it back from the norm-one torus by \(z\mapsto z/\bar z\)
and multiply the old horizontal Hecke character by this
pullback. Its angular exponent changes by twice the stated
difference, its restriction to the real subfield is trivial,
and every local finite conductor and character on units
is preserved. Use its inverse wherever the old inverse was
used. The parity character of Proposition~\ref{prop:parity262}
also works for these choices. The bottom row uses only the
cloud and the entries \(-b_j\), and the Fourier character
uses the rational Hermitian flag and the fixed additive
character. Neither depends on \(g\).

Let \(\lambda\) be the highest weight of the new terminal
compact type at one real place. The same string count as in
Theorem~\ref{thm:word275} gives
\(\lambda_k-\lambda_{N-k+1}\ge (m-1)g-N\), also in dual
coordinates. The interval with these endpoints has half-length
greater than \(A+w_E\) for \(m\ge4\). Every desired highest
weight \(\mu_i\) has absolute value less than \(A\) by
\eqref{eq:letters111}, including both endpoints at every high
place. At each place choose a multiple \(j_\tau\) of \(w_E\)
nearest the negative midpoint of the interval. Then all
\(\mu_i-j_\tau\) belong to it. Proposition~\ref{prop:branch275}
gives the branching assertion and the global character with
these exponents. The same choice works for every word with
the given orientation.

The old ranks, \(g_*,w_E,k\Delta\) and \(f\) are polynomially
bounded in \(x\). Since \(b_*=O(xb)\), we have \(g=O(xb)\)
on all sufficiently distant intervals. Each new core
parameter has absolute value at most
\(v_1+2+(m+1)g+d\). Thus the logarithm of two plus the
largest absolute entry, including the half-root shifts,
is \(O_{\psi,\varepsilon}(\log b)\) on both groups.
In the Weyl dimension formula each factor for \(U(l)\) is
\(1+(\lambda_i-\lambda_j)/(j-i)\), \(i<j\). Applying the
logarithmic bound to these factors and summing over all
real places proves the estimate.
\end{proof}

\subsection{Finite levels for the restricted row}\label{nav:V-44}
Take the fields, intervals, enlarged auxiliary weights and raw
period of Theorem~\ref{thm:raw276}. Use the finite translate
of Proposition~\ref{prop:lattices277}. Let \(S_{\rm m}\)
be the marked set of the horizontal tests, disjoint from the
places above \(T\). Let \(a_v\) be their prescribed depths,
extended by zero outside \(S_{\rm m}\). Choose \(v_*\) above
the marked odd split prime \(g_1\). For a lattice \(L\) write
\(K_L(0)=U(L)\), and let \(K_L(b)\) for \(b>0\) be the
kernel of reduction modulo \(\pi_v^b\) on \(U(L)\).

On \(H=U(C)\) take the level \(J\) of the translated core
lattices with depths \(a_v\) at the marked places, one at
every \(v\mid T\), and zero elsewhere. Its finite class number
is \(h_H\). Every stabilizer is trivial since it is finite
and embeds at \(v_*\) in a conjugate of a torsion-free odd
principal congruence group. All class masses are consequently
\(1/h_H\), as in Proposition~\ref{prop:terminalcore255}.
Apply Theorem~\ref{thm:mix277} with \(\omega=1/h_H\) and
\(d=d_\rho\) for the real type of the conjugate raw column.
For several such types on the same finite data use their
largest dimension in the formula for \(j\). The bound for
\(c_Q^j\) then holds for each type separately.
Write \(j_v=j\) at \(v=v_*\) and zero elsewhere.

\begin{theorem}[A restricted row at explicit finite levels]\leavevmode
\label{thm:row277}
For every prescribed word \(\mu\) and terminal character
\(\lambda\) of Theorem~\ref{thm:types276} is a restricted
row \(A\) on the original group, fixed by the original
terminal graph groups, and fixed at \(v\nmid T\) by
\(K_{\Lambda_{0,v}}(a_v+2e_v+2j_v)\).
At those same places it is also fixed by the embedded column
group \(K_{\Lambda_{W,v}}(a_v+2j_v)\).
For \(p(h)=F_0A_\lambda(h,1_f)\) its real type is \(\mu\) and
\[
 \|p\|_2\ \ge\
 {\sqrt{|Z|}(1-c_T)^{f/2}\over 2\sqrt{2}}
       \sqrt{\frac{m_\nu}{d_\rho d_\mu}}\,M_B .
\]
Here \(M_B=|P_{\rm raw}|/\|e\|>0\) is calculated from the
original raw tests for the enlarged spacing. The depths in
the assertion are common to the choices of bottom vector,
word and terminal character for the fixed grid and type.
\end{theorem}
\begin{proof}
Translate the large row \(B\) and its raw column \(e\) on the
right by the same finite \(u\). Equivariance of the large
Fourier map under \(H\) and right invariance of Haar measure
preserve the raw period and the norm of its column. Denote by
\(\widetilde B\) the projection of \(R(u)B\) to real type
\(\rho\) under the definite group. This retains the raw pairing.
The unramified vacua and marked coset tests of
Theorem~\ref{thm:raw276} are fixed by their full lattice
groups and principal congruence groups, respectively, by
Proposition~\ref{prop:rowfinite262} and matrix substitution
on an integral marked coset. After translation these are the
groups of \(L_v\) in Proposition~\ref{prop:lattices277}.
The terminal vacua are the original ones.

Set \(P=D_\lambda F_{\chi_+}\widetilde B\), with the graph
change \(a\) and the unnormalized finite character sum of
Proposition~\ref{prop:corehecke272}. At \(v\mid T\) each
translated core lattice is the inverse image of a hyperplane
modulo \(T\). Its stabilizer contains the principal subgroup
of depth one. Thus \(P\) is fixed by \(J\). The same
compression identity and real nontriviality as in
Theorem~\ref{thm:types276} give
\(\|P\|_2\ge\sqrt{|Z|}(1-c_T)^{f/2} M_B\).
Let \(\nu\) be the occurrence supplied for \(\mu\) by
Proposition~\ref{prop:spacing276}. Theorem~\ref{thm:mix277}
gives \(g=(s,D(\delta)t)\) such that the \(\nu\)-character
integral of \(P(hg)\) has modulus at least
\(\sqrt{m_\nu d_\mu/(2d_\rho)}\,\|P\|_2\).

Apply Proposition~\ref{prop:fiber274} at \(hg\) to the
large row \(\sum_z\bar\lambda(z)R(za)\widetilde B\).
Integration against the bounded character on the real compact
group commutes with its box limit. Indeed the large unipotent
quotient and the compact group together bound all values of
the row occurring in the averages at the fixed \(g\).
The small compact group commutes with \(n(x)\) in that formula.
Some real \(x\) therefore retains at least half the modulus
of the character integral. Restrict
\(y\mapsto\widetilde B(y n(x)ga)\) to \(U(S_0)\), project
its right real translates to \(\nu\), and apply the
everywhere finitely unramified determinant twist of
Proposition~\ref{prop:spacing276}. These operations preserve
smoothness and rational left invariance. The terminal
scalars commute with \(g\) because \(D(\delta)\) centralizes
the small row, and the other finite translate is supported at
\(v_*\). Thus the sum and small Fourier integral of this row
give the selected value at one, up to a scalar of modulus one.
The twist turns \(\nu\) into \(\mu\), keeps compact finite
levels, and is trivial on the unipotent. The real evaluation
bound \(d_\mu\|p\|_2\) from Theorem~\ref{thm:types276}
now gives the displayed inequality.

At a terminal place conjugation of the small graph group by
\(g_v\) fixes it pointwise. Conjugation by the inverse local
component of the graph change \(a\) then places
it in the original large lattice group by
Theorem~\ref{thm:graph274}. At \(v\nmid T\) put \(L=L_v\)
and \(b=a_v+2j_v\). If \(k\) lies in the asserted small row
group, the two inclusions of Proposition~\ref{prop:lattices277}
give \((k-1)L\subset\pi_v^b L\). Indeed the displacement
acts on the projection to \(S_0\), loses at most \(e_v\) on
that projection, and loses at most \(e_v\) on inclusion back
into \(L\). For \(k\in K_{\Lambda_{W,v}}(b)\) the same
inclusion costs no \(e_v\), since the column splits off
orthogonally in \(L\). The inverse satisfies the same
inclusion in both cases.

In conjugating \(k\) through \(g_v\) the determinant section
cancels. The remaining \(t_v\) and its inverse take \(L\)
into \(\pi_v^{-j_v}L\), since they act in the core summand
alone. Hence \(t_v^{-1}kt_v\) belongs to \(K_L(a_v)\),
including when \(a_v=0\). It fixes the translated row there.
The real operations do not change that calculation.
\end{proof}

\section{Compact coefficients and projection norms}\label{nav:g-VIII-3}

\subsection{The rectangular compact type at the last hyperplane}\label{nav:V-30}
For \(1\le r<s\) let \(V_b\) be the span of ordered products
\(P_T\) of \(b\) maximal minors of an \(s\)-by-\(r\) matrix,
with the Fischer inner product of Theorem~\ref{thm:frame55}.
Use the coordinate hyperplane of the terminal inclusion
in Proposition~\ref{prop:ranks47}. Let \(V_{b,l}\) be the
span of the products in which exactly \(l\) of the minors
use row \(s\). In particular their degree in that row is \(l\).
Write \(A_0\) for the minor on rows \(1,\ldots,r\), and \(A_1\)
for the minor on rows \(1,\ldots,r-1,s\), both with rows in
increasing order.

\begin{lemma}[The integral vectors in the restricted rectangular type]
\label{lem:branch243}
For \(0\le l\le b\) the spaces \(V_{b,l}\) are pairwise
orthogonal and irreducible under \(U(s-1)\). In the covariant
row convention their highest weights are
\((b^{r-1},b-l,0,\ldots,0)\), where \(b^{r-1}\) denotes
\(r-1\) repeated entries. Dualizing the row action or twisting
by its vacuum character transports the whole decomposition.
Put \(Q_l=A_0^{b-l}A_1^l\). Then
\[
 \|Q_l\|_{\rm F}^2=\frac{C_{r,b}}{\binom bl},\qquad
 \sum_{T:\,\deg_s P_T=l}|L(P_T)|^2=C_{r,b}\|L\|^2
                         \quad (L\in V_{b,l}^\vee),
 \qquad C_{r,b}=\prod_{j=0}^{r-1}\frac{(b+j)!}{j!}.
\]
Here the dual norm is that of a complex linear functional and
the ordered indices retain their multiplicities. On the full group
\(U(s)\) the corresponding space is \(V_b\), with distinguished
vector \(A_0^b\) of squared norm \(C_{r,b}\).
\end{lemma}
\begin{proof}
The unitary translates of \(A_0^b\) span \(V_b\) by
Theorem~\ref{thm:frame55}. Its diagonal-torus weight has
multiplicity one. Indeed an indexed product using any row beyond
\(r\) has positive degree in some such row, and a product using
only the first \(r\) rows equals \(A_0^b\). A unitary invariant
orthogonal projection sends this cyclic vector to a scalar
multiple of itself. As the projection is idempotent the scalar
is zero or one, and cyclicity gives irreducibility. Additions of
an earlier row to a later row in the first minor either leave its
determinant fixed or change a row outside the minor. Thus its
extremal weight, choosing the Borel for this convention, is
\((b^r,0,\ldots,0)\).

The Schur character of this rectangular representation is the
sum of monomials of semistandard tableaux of shape \((b^r)\).
In the alphabet \(1,\ldots,s\) all entries \(s\) lie at the
right end of the bottom row, since columns strictly increase
and rows weakly increase. With \(l\) entries \(s\), deleting
them gives exactly the tableaux of shape \((b^{r-1},b-l)\)
on \(1,\ldots,s-1\). Extension by those \(l\) entries is
inverse to deletion. Taking the coefficient of the \(l\)-th
power of the last diagonal variable in the character therefore
gives one irreducible character of \(U(s-1)\). The corresponding
weight space for the last circle is precisely \(V_{b,l}\).
Distinct last-row degrees are orthogonal by circle invariance.
This also covers \(b=0\).

A real unitary rotation in rows \(r,s\) takes \(A_0\) to
\(\cos\theta\, A_0+\sin\theta\, A_1\). The terms of its \(b\)-th
power with different \(l\) are orthogonal. Its squared norm is
constantly \(C_{r,b}\) by the full-group calculation. Equating
coefficients of \(u^{b-l}v^l\), with \(u=\cos^2\theta\) and
\(v=\sin^2\theta\), gives the first formula. In detail the
homogeneous polynomial for the squared norm agrees with
\(C_{r,b}(u+v)^b\) on \(u+v=1\) in the positive quadrant,
so the two homogeneous polynomials agree identically.
For the second formula the full frame operator
\(\sum_T P_T P_T^*\) is \(C_{r,b}\) times identity on \(V_b\)
by Theorem~\ref{thm:frame55}. Its restriction to the indicated
orthogonal summand uses exactly the terms of degree \(l\).
\end{proof}

For a list \(b_\tau,l_\tau\) use the tensor products \(V,Q\)
of these spaces and vectors over the real places. Their frame
constant and the squared norm of \(Q\) are respectively
\[
 C_{\boldsymbol b}=\prod_\tau C_{r,b_\tau},\qquad
 \|Q\|_{\rm F}^2=C_{\boldsymbol b}/H,\qquad
 H=\prod_\tau\binom{b_\tau}{l_\tau}.
\]
For the full target use \(V=\bigotimes_\tau V_{b_\tau}\),
\(Q=\prod_\tau A_0^{b_\tau}\), and \(H=1\). In either case
write \(d_V=\dim_{\CC}V\). It is the product of the dimensions
of the stated highest weights. Vacuum characters have dimension one.

\subsection{Compact matrix coefficients for an arbitrary partition}\label{nav:V-31}
Let \(n\le s\) and \(\lambda_1\ge\cdots\ge\lambda_n\ge0\)
be integers. Write \(e_j=\lambda_j-\lambda_{j+1}\), with
\(\lambda_{n+1}=0\), and \(d=|\lambda|\). For \(m=n,s\) form
\[
 T_m=\bigotimes_{j=1}^n(\bigwedge\nolimits^j\CC^m)^{\otimes e_j}.
\]
Coordinate wedges are orthonormal. Let \(t_m\) be the tensor of
the first \(j\) coordinate wedges in the indicated slots, and
\(V_m\) its unitary cyclic span in \(T_m\). This is the Cartan
component of highest weight \(\lambda\), padded by zeros, as in
Lemma~\ref{lem:shuffle107}. Its integral lattice here is
its intersection with the coordinate tensor lattice. For complex
inner products in this subsection use linearity in the first entry.

On \(T_n\otimes T_s\) define a linear map \(M\) to degree-\(d\)
matrix polynomials as follows. A coordinate tensor has subsets
\(J_a\) and \(I_a\) of the same size in each matched slot. Its
image is \(\prod_a\det Z_{I_a,J_a}\), with the coordinates in
each subset in increasing order. In particular \(M(t_n\otimes t_s)
=p_\lambda\), the product of leading minors in
Lemma~\ref{lem:hook86}. Use the covariant actions on the tensor
factors and on the polynomial coefficients. In an oscillator
where matrices or their conjugates act contragrediently, pull back
these actions along the corresponding compact group automorphism
on both sides. Include the same vacuum character on the two sides.

\begin{proposition}[Integration against the two specified coefficients]
\label{prop:cartan244}
Let \(\rho_n,\rho_s\) denote the actions on \(V_n,V_s\), including
the respective vacuum characters, and \(\omega\) the action on
polynomials. Write \(h_\lambda=\|p_\lambda\|_{\rm F}^2\) and
\(D_m=\dim V_m\). For \(x,y\in V_n\), \(x',y'\in V_s\) and
a degree-\(d\) polynomial \(P\), the compact average for Haar
probabilities is
\begin{align*}
 &\int\overline{\langle\rho_n(g)x,y\rangle
                 \langle\rho_s(h)x',y'\rangle}\,
                 \omega(g,h)P\,dg\,dh\\
 &\hspace{12mm}=
 \frac{\langle P,M(x\otimes x')\rangle_{\rm F}}
             {D_nD_s h_\lambda}\ M(y\otimes y').
\end{align*}
For \(P=p_\lambda\) the numerator divided by \(h_\lambda\)
is \(\langle t_n,x\rangle\langle t_s,x'\rangle\).
In particular it is one on the two coordinate highest vectors.
The displayed polynomial \(M(y\otimes y')\) has integral
coefficients whenever \(y,y'\) lie in the stated integral lattices.
\end{proposition}
\begin{proof}
Cauchy--Binet on every slot shows that \(M\) intertwines the
source and target actions. More explicitly, a change on a wedge
acts by its minors, which are precisely the coefficients obtained
by expanding that same minor after a matrix substitution. The
irreducible tensor product \(V_n\otimes V_s\) maps nontrivially
since \(p_\lambda\) has value one on the top identity frame.
By the Cauchy decomposition in Lemma~\ref{lem:degree181},
its joint type has multiplicity one in the matrix polynomials.
Equivalently the degree-\(d\) polynomial representation is obtained
by taking the simultaneous permutation invariants in the two
\(d\)-fold tensor powers. Schur--Weyl duality and Schur's lemma
pair each partition with itself once. The invariant positive form
induced by \(M\) on the two factors is a scalar times their tensor
form, and its value on \(t_n\otimes t_s\) gives that scalar
\(h_\lambda\).

On one irreducible unitary representation of dimension \(D\),
the integral of \(\overline{\langle\rho(g)x,y\rangle}\rho(g)\)
maps \(z\) to \(\langle z,x\rangle y/D\). Indeed averaging
\(\rho(g)zx^*\rho(g)^*\) gives \(\langle z,x\rangle I/D\)
by invariance and trace, and evaluation on \(y\) gives the
formula. Coefficients of inequivalent representations give zero
by the same averaging of an intertwiner between them. Apply
the two matrix units on the unique joint summand. Converting
its tensor inner product to the Fischer inner product divides
by \(h_\lambda\), giving the assertion for any \(P\). The
formula for \(p_\lambda\) follows by the same isometry. Each
coordinate minor product has integer coefficients, so linearity
in the coordinate tensor basis gives the integral assertion.
\end{proof}

\subsection{The integral tests at the first seesaw}\label{nav:VI-34}
Fix one of the period fields \(E=E_x\), let \(F=E^+\), and use
the first step of the ascent from the restricted column period.
The two targets there are the initial row space \(Y\) and the
larger space \(X=Y\perp L\), with \(L=Ez\) a line. The bottom
lattice in \(X\) is
\(\Lambda_X=\iota\Lambda_Y\oplus\mathcal O_E z\). The source
\(S\) of the two oscillator spaces is the first core space.
Its map into the first row step is an isometry of lattices:
the extra hyperbolic summand at this step has rank zero, so
the row lattice is exactly the image of \(\Lambda_S\).
These are the bottom lattice and the first row lattice in the
standard frame. The identification uses no choice of a prime.

Here are the finite tests in coordinates. For \(B=X,Y,L\), start
with the additive lattice in \(\operatorname{Hom}_E(S,B)(\mathbb A_{F,f})\)
generated by the maps
\[
   v\longmapsto \alpha a\,t\,h_S(s,v),\qquad
   \alpha\in\widehat{\mathcal O}_F,\quad
   a\in\mathfrak a,\quad s\in\Lambda_S,\quad t\in\Lambda_B .
\]
Here \(h_S\) is linear in its second argument and \(\mathfrak a\)
is the fixed fractional ideal in the oscillator data. Impose in addition
\(e u\in(\mathfrak m\operatorname{Hom}_{\mathcal O_E}
(\Lambda_S,\Lambda_B))^{\wedge}\) on a map \(u\), for every idempotent
\(e\) prescribed by the grid. Such an \(e\) selects the specified finite
components above the split places. The hat denotes finite completion and
\(\mathfrak m\) is the original modulus. Denote the resulting lattice
by \(\Gamma_B\). On the line put \(\Lambda_L=\mathcal O_E z\).
Each \(\Gamma_B\) is compact and open. Indeed the rank-one
generators before the extra conditions span a full finite lattice:
take a finite set of generators of the two integral lattices and of
\(\mathfrak a\), and use nondegeneracy of \(h_S\). The additional
congruences affect only finitely many places and contain a full open
sublattice there.

Use a rational Schr\"odinger polarization respecting the sum \(Y\perp L\),
and the usual symmetric Weyl operators \(R_B\) for the fixed additive
character. The finite depth tests \(\mathcal V_B\) are the Schwartz
vectors fixed by \(R_B(2u)\) for \(u\in\Gamma_B\) and by \(R_B(u)\)
for \(u\in\Gamma_B\) supported on one half of a split place.
On \(-L\) reverse the form and conjugate the oscillator. Let
\(\mathcal V_B^{\rm int}\) be the functions in this space taking values in
\(\overline{\ZZ}\), the ring of algebraic integers in \(\CC\).

\begin{proposition}[Integral factorization at the original depth]
\label{prop:depth245}
Under restriction to the two target summands, \(\Gamma_X\) is
\(\Gamma_Y\oplus\Gamma_L\). The finite Schr\"odinger tensor map
induces an isomorphism
\[
 \mathcal V_Y^{\rm int}\otimes_{\overline{\ZZ}}
      \overline{\mathcal V_{-L}^{\rm int}}
        \ \simeq\ \mathcal V_X^{\rm int}.
\]
It carries product bases to an integral basis, including above \(2\).
In each factor the coefficient of a basis function is obtained by
evaluation at a finite position.
\end{proposition}
\begin{proof}
The projection \(X\to Y\) sends \(\iota y+c z\) to \(y\) and the
projection \(X\to L\) sends it to \(cz\). They preserve the displayed
integral lattices. They split each rank-one generator into the same
two generators with projected target vectors. Conversely the two
inclusions preserve all generators. The group of integral homomorphisms
also splits by composing with the same projections. They commute with
scalar multiplication by every \(e\), so the congruences split with
both inclusions. Precomposition with the isometry from the rank-zero
step preserves the source lattice and its Hermitian form. This proves
the first assertion for the precise first row and core spaces.

In the direct-sum polarization a Weyl displacement is the product
of the two displacements, since their cross pairing is zero.
A doubled displacement splits into independent doubled displacements.
A displacement supported on one half of a split place projects to
two displacements supported on that same half. Each separate
displacement is also present by extension by zero. Thus the fixed
equations on the two variables can be imposed separately. Complex
conjugation reverses the phases on the line and identifies its fixed
equations with those for the opposite form.

For integral coefficients one can solve these equations explicitly.
In a single factor write the finite position and frequency spaces
as \(Z\) and \(Z^\vee\). Choose compact open lattices \(U\subset Z\)
and \(V\subset Z^\vee\) such that \(U\oplus V\subset 2\Gamma_B\),
\(U\) annihilates the frequency projection of \(\Gamma_B\), and \(V\)
annihilates its position projection. Existence follows by shrinking
two compact open lattices inside the indicated open subgroups.
The depth equations for \(U,V\) give functions supported on
\(V^\perp\) and constant modulo \(U\). The other Weyl operators
preserve this finite coordinate space. A position displacement
permutes its coordinates, and a frequency displacement multiplies
them by roots of unity, including the phase of the symmetric Weyl
operator. On each permutation orbit choose a representative. If a
word fixing this representative has a nontrivial phase, its coordinate
must be zero. On every other orbit propagate the value one at the
representative by the fixed equations. This is well defined because
two paths differ by a word fixing the representative. The functions
so obtained have disjoint supports and root-of-unity entries.
Evaluation recovers their coefficients over both \(\CC\) and
\(\overline{\ZZ}\). In the product the two systems of equations
act on separate coordinates, so the product functions give exactly
the same construction on pairs of surviving orbits. This proves
the integral isomorphism. Keeping the doubled and split-half
equations throughout covers the dyadic case.
\end{proof}

\subsection{An integral terminal coefficient on one definite class}\label{nav:VI-37}
On a field of Theorem~\ref{thm:downcore254}, put \(G=G_m\),
\(K_f=K_{m,f}\), \(K_\infty=G(F\otimes_\QQ\RR)\) and let
\(h_m=\#(G(F)\backslash G(\mathbb A_{F,f})/K_f)\).
The real group is compact. Let \(\rho\) be the representation
of \(K_\infty\) of parameter \(P_m\), of dimension \(d_\rho\).
At each selected complex embedding its highest weight is the
decreasing Harish-Chandra parameter minus the ordered half sum
of roots. Denote by \(W\) the resulting algebraic representation
at all these embeddings over a number field \(J\) containing
their images.

\begin{proposition}[A terminal form supported on one class]\leavevmode
\label{prop:terminalcore255}
Fix one representative \(t_c\) of each finite class. There are
finite lists of terminal forms, depending on the field and the
core lattices, which span \(\mathcal A_m\). Each is supported on
a single class \(c\), where its real function is
\(k\mapsto\ell(\rho(k)v)\). Both \(v\) and \(\ell\) may be taken
from generating sets of mutually dual integral coefficient
lattices on that class. With probability measure on the
quotient its squared norm is
\[
      \frac{\|v\|^2\|\ell\|_*^2}{h_m d_\rho}.
\]
For any bottom row vector in Theorem~\ref{thm:downcore254}
one such terminal form gives its nonzero first period after a
single descent using the stated monomials.
\end{proposition}
\begin{proof}
The class set is finite by the finiteness theorem for definite
Hermitian lattices, or by compactness of the anisotropic
adelic quotient and openness of the level in the finite adeles.
Let \(\Lambda_c=E^{n_m}\cap t_c\widehat\Lambda_m\) in rational
coordinates. A rational stabilizer of the class representative
belongs to the automorphism group of this definite lattice.
The integral endomorphisms form a discrete lattice in the product
of the archimedean endomorphism spaces, and the isometries lie in
a compact group there. Their intersection is finite. At the
chosen split auxiliary prime the stabilizer is conjugate into
\(1+p_0\operatorname{Mat}_{n_m}(\ZZ_{p_0})\). This uses complete
splitting and the positive congruence depth. The prime \(p_0\)
is odd by its defining congruence. This principal congruence
group is torsion-free. Indeed, if \(s\ge1\) is the least entry
valuation of \(g-1\ne0\), raising \(g\) to a prime
exponent \(\ell\) leaves that valuation
at \(s\) for \(\ell\ne p_0\) and raises it to \(s+1\) for
\(\ell=p_0\), by the binomial expansion. In the latter case all
terms beyond the linear one have valuation at least \(s+2\).
Iteration precludes finite order. Localization is injective on
rational matrices, so the rational stabilizer is trivial.

The quotient at level \(K_f\) is consequently a disjoint union
of \(h_m\) copies of \(K_\infty\). In the quotient measure
each copy has the same mass, since its finite representative
translates the same compact open set and the stabilizers are
trivial. Haar probability conditional on a copy is compact Haar
probability. On one copy the right \(\rho\)-isotypic functions
are exactly the span of its matrix coefficients by the
Peter--Weyl theorem. Compact Schur orthogonality gives the
stated squared norm on that copy and division by \(h_m\)
gives the global probability norm.

At each embedding take the integral Schur module of the lattice
\(\Lambda_c\) with the indicated highest weight, using the
determinant twist and the image of the integral Young
symmetrizer as in Theorem~\ref{thm:proj}. Tensor these modules
over the embeddings to obtain a projective \(\mathcal O_J\)-lattice
\(W_c\) spanning \(W\). Choose finite module generators of
\(W_c\) and its integral dual. Over \(\CC\) the resulting
coefficients span all matrix coefficients since the two lists
span the representation and its full linear dual. The class
stabilizer imposes no further conditions. In particular
\(\ell(v)\) at the identity is integral in these coefficient
coordinates. In the descent of Theorem~\ref{thm:downcore254}
expand its terminal form in this finite spanning list. With
the tests and the row held fixed the period is linear in the
terminal form, so at least one listed term has nonzero period.
\end{proof}

\subsection{A norm estimate for the tensor projections}\label{nav:VI-51}
Use the tensor projections \(P_i=\iota_i p_i\) of
Proposition~\ref{prop:transfer286}. The invariant form on the
dual defining representation is indefinite. Summing the
positive Hilbert norms of the projections therefore requires
an estimate. The number of row entries below the interlacing
part is \(m=2s+2\).

\begin{lemma}[The sum of the positive Hilbert norms]\leavevmode
\label{lem:signs286}
On the fields and spacings of Section~\ref{sec:absolute163}
let \(U\) be the lowest compact type of the initial row.
After \(a\le 4r^2\) unit lowerings, let \(w\) be a
compact-finite vector in the shifted row whose compact types
occur in \((F^*)^{\otimes a}\otimes U\). For every Euclidean
unit dual basis vector \(\zeta\) one has
\[
       \sum_i\|p_i(\zeta\otimes w)\|^2
                    \le \exp(C s\log b)\|w\|^2,\qquad
                     b=\exp(x^\varepsilon).
\]
The constant is uniform in the prescribed word, place and
shift. The compact types of the projected vectors occur in
\((F^*)^{\otimes(a+1)}\otimes U\).
\end{lemma}
\begin{proof}
We first determine enough of the signs \(\epsilon_i\) to
control the positive sum. They depend only on the fixed
tensor form and the row. For this calculation choose an
auxiliary compact column with decreasing half-integral
parameters \(u_a'\). Place \(u_1'\) above \(\mu_1\), and
place \(u_a'\) between \(\mu_{a-1}\) and \(\mu_a\) for
\(2\le a\le k\), at distance greater than one from both
adjacent entries. Choose \(u_1'\) also at distance greater
than one from \(\mu_1\). These choices are possible because
the row gaps after the allowed shifts tend to infinity.
The signs of the row are unchanged. The packet calculation of
Theorem~\ref{thm:src109} gives a nonzero Bessel model for
this auxiliary column, since the order of the two lists,
the chamber signs and their parities agree with those in
its proof. This auxiliary choice is used only to determine the
sign of the invariant form on a row summand.

For each \(i\), the signed scalar of
Proposition~\ref{prop:cofactor286} for the auxiliary column
is nonzero. The adjoint identity in
Proposition~\ref{prop:transfer286} identifies its sign
with \(\epsilon_i\). For \(i\le k-1\) the numerator product
has \(k-i\) negative factors and the denominator \(r-i\).
Consequently the signs of the first \(k-1\) summands are
the same. The multiplier \(\eta\) is common and is real,
since the signed scalar for any one summand is nonzero and
real while the product quotient is nonzero and real.
Call the common sign \(\epsilon_0\), and let \(J_-\)
be the indices whose sign differs from it. Then
\(|J_-|\le r-k+1=2s+2\).

The split Casimir \(C_0\) has eigenvalues
\(\theta_i=-\mu_i-(r-1)/2\), pairwise at distance at least one.
Use the real polynomial
\[
 Q(X)=2\prod_{i\in J_-}(X-\theta_i)^2-1 .
\]
It equals \(-1\) at the exceptional eigenvalues and is
at least one at all other eigenvalues, also when the exceptional
set is empty. Orthogonality for the invariant form gives,
for \(z=\zeta\otimes w\),
\[
 \sum_i\|p_i z\|^2\le
      \epsilon_0[Q(C_0)z,z]\le\|Q(C_0)z\|_{\rm Euc}\,\|z\|_{\rm Euc}.
\]
Here \(\|\cdot\|_{\rm Euc}\) is the positive tensor norm from
the Euclidean norm on \(F^*\) and the shifted Hilbert norm.
The matrix of the invariant form in this norm is a unitary
involution. The first inequality follows term by term since
\(\epsilon_0\epsilon_i Q(\theta_i)\ge1\).

We estimate the polynomial on \(z\) using
Lemma~\ref{lem:cas96}. On the lowest compact type \(U\) the
norms of all unit compact Lie directions are bounded by the
parameter \(L_*\) of Theorem~\ref{thm:localmass156}. On the
unitary compact tensor power \((F^*)^{\otimes a}\otimes U\)
the bound is \(L_*+a\), by the tensor differentiation rule.
It holds on any compact span having only these types, since
each compact type is unitary and compact isotypic summands
are orthogonal. The square root of the absolute Casimir
eigenvalue of the shifted row is at most
\(C\sqrt r\,(\max_j|\mu_j|+r)\), by the formula in
Proposition~\ref{prop:cofactor286}. Together these furnish
a parameter \(L'\) in Lemma~\ref{lem:cas96} with
\(\log(2+L')=O(\log b)\). Here \(a,r\) are polynomial in \(x\),
and \(L_*\ll\sqrt r(2b+r)\).

In the Euclidean Witt frame \(C_0\) has entries
\(-E_{ij}\) acting on the row. Each complexified matrix
unit is a sum of a real unitary Lie direction and \(i\)
times another one, both with Frobenius norm at most one.
Indeed the real structure \(X\mapsto-JX^*J\) preserves
the Frobenius norm. Expanding \(C_0^d z\) uses at most
\((2r)^d\) vectors in single dual coordinates when \(\zeta\)
is a basis vector. Each is a word of \(d\) such real
derivatives applied to \(w\). Thus for \(d\le 4s+4\),
\[
 \|C_0^d z\|_{\rm Euc}\le
     \{C r^2(L'+4s+5)\}^{d}\|w\|.
\]
By expanding the product defining \(Q\), its norm on \(z\)
is bounded by
\[
 \left[1+2\{C r^2(L'+4s+5)+\max_i|\theta_i|\}^{\,2|J_-|}\right]\|w\|.
\]
The logarithm of the bracket is \(O(s\log b)\), giving the
claimed estimate. Finally \(p_i\) is compact-equivariant.
The compact span of \(\zeta\otimes w\) has exactly the
required containment of compact types, which passes to its
image under \(p_i\).
\end{proof}

\subsection{Prime exponents in compact unitary averages}\label{nav:VII-63}
For \(h\ge1\), \(L\ge0\), put
\[
 R_h(L)=\frac{\operatorname{lcm}(h,h+1,\ldots,h+L)}{h},
 \qquad R(n,L)=\prod_{h=1}^{n-1}R_h(L)^{\,n-h}.
\]
Let \(D(n,L)=\gcd(R(n,L),\mathcal D(n,L))\), where
\(\mathcal D(n,L)\) is the compact moment denominator of
Proposition~\ref{prop:moments242}. Empty products are one.

\begin{theorem}[A logarithmic bound for each prime exponent]\leavevmode
\label{thm:compact313}
Suppose \(P(g,\bar g)\) has degree at most \(L\) in the entries of
\(g\) and degree at most \(L\) in their conjugates, with coefficients
in a fractional ideal \(I\) of a number field. For Haar probability
on \(U(n)\) one has
\[
 D(n,L)\int P(g,\bar g)\,dg\ \in I,\qquad
 v_p(D(n,L))\le \binom n2\log_p(n+L)
\]
at every prime \(p\). For \(n=2\) the smallest positive integer
clearing every such moment with integer coefficients is
\(\operatorname{lcm}(1,\ldots,L+1)\).
\end{theorem}
\begin{proof}
We give an integration formula in which every angular integration
has integer coefficients. For dimension \(s\ge2\) start the first
column at \(e_1\). Successively act on coordinates \(1,k\) for
\(2\le k\le s\) by the block
\[
      \begin{pmatrix} z&-\bar w\\ w&\bar z\end{pmatrix},
        \qquad |z|^2=T,\qquad |w|^2=1-T .
\]
In each block choose the phases of \(z,w\) independently and
uniformly on the circle. Choose \(T\) with density
\(h t^{h-1}\,dt\) on \([0,1]\), where \(h=s-k+1\). Use independent
parameters for different blocks. The squared modulus in coordinate
\(k\) then receives the fraction \(1-T\) of the mass left in
coordinate one.

The resulting first column is uniform on the complex unit sphere.
Indeed the coordinate squared moduli of a vector of independent
standard complex Gaussians are independent exponentials. Dividing
by their sum gives uniform density on the simplex of masses. This
follows by the change from the \(s\) exponentials to their sum \(r\)
and their mass fractions, with Jacobian \(r^{s-1}\). In a uniform
simplex with \(h+1\) coordinates the sum \(T\) of the remaining
\(h\) coordinates has density \(h t^{h-1}\), and their normalized
fractions are independent of \(T\), by the same simplex change
of variables. Iteration gives exactly the chosen masses. Given
the phases of the \(z\)'s, each emitted coordinate has the
independent uniform phase of its \(w\) shifted by a fixed angle.
These phases are independent of the \(z\)'s. The final phase in
coordinate one is the product of the \(z\)-phases and is uniform.
Thus the phases have the Gaussian angular law as well.

Multiply this product of blocks on the right by an independent
\(\operatorname{diag}(1,U)\), with \(U\) Haar in dimension \(s-1\).
For a fixed first column any two unitary completions differ on
the right by its \(U(s-1)\) stabilizer in these coordinates.
The conditional law of the completion is therefore uniform on
the fiber. Since the first column is uniform, the resulting law
is left invariant, hence Haar. Recurse down to one uniform phase
in dimension one. For the original \(n\), the parameter with
density \(h t^{h-1}\) occurs once in each dimension
\(s=h+1,\ldots,n\). The independent beta angles in unitary
integration also appear in
\href{https://arxiv.org/abs/1512.09229}{P.~Diaconis and
P.~J.~Forrester, \emph{A.~Hurwitz and the origins of random
matrix theory in mathematics}}.

Substitute the product into a monomial of \(P\). Each unitary
entry has degree at most one in each individual block, so at
most \(2L\) factors of a given block occur in a term of the
expansion. Averaging its two phases gives zero unless both
the \(z,\bar z\) counts and the \(w,\bar w\) counts match.
Let the numbers of pairs be \(a,b\), respectively. Then
\(a+b\le L\), and the radial average is
\[
 h\int_0^1 t^{h+a-1}(1-t)^b\,dt
      =\sum_{r=0}^b(-1)^r\binom br\frac{h}{h+a+r}.
\]
Its denominator divides \(R_h(L)\) since
\(h R_h(L)\) is divisible by every \(h+a+r\) in the sum.
The final circle phase contributes zero or one. Independence
and expansion with integer coefficients prove the assertion
with \(R(n,L)\). Proposition~\ref{prop:moments242} proves it
with \(\mathcal D(n,L)\). An integral B\'ezout combination gives
the assertion with their gcd, also for a fractional ideal.
Each \(v_p(R_h(L))\) is at most \(\log_p(n+L)\), and there are
\(\binom n2\) blocks.

For \(U(2)\), \(R(2,L)=\operatorname{lcm}(1,\ldots,L+1)\).
For each \(0\le j\le L\), the polynomial \(|g_{11}|^{2j}\)
has mean \(1/(j+1)\) because \(|g_{11}|^2\) is uniform on
\([0,1]\). Every common denominator must therefore contain
the whole least common multiple.
\end{proof}

\section{Whittaker normalizations and terminal restriction spaces}\label{nav:terminal-whittaker}

\subsection{The terminal scalar in the normalized input}\label{sec:terminalscalar357}

Keep the fields, the coded weights and the split generator
\(\pi\in\mathcal O_{K_d}\) of Theorem~\ref{thm:common111}.
Write \(D_0=t_{\rm r}+1\) for the length of each appended
block, \(n_j=k+jD_0\), and \(h\) for the number of high
embeddings. Both \(m\) and \(D_0\) are even. Each embedding
in the CM type has the same restriction to \(K_d\), as in
Section~\ref{nav:I-29}. We write \(\pi,\bar\pi\) also for
the images of its generators in the coefficient field \(J\)
under that common restriction. The terminal prime is
\(T=\pi\bar\pi\). Normalize \(v_w(T)=1\) at \(w\mid T\).

We specify the finite normalization used after the Hecke
operation. At a coefficient prime \(w\) over \(p\) let
\(e_w\) be the ramification index. If the prescribed total gain
minus total error is \(s_w\), the source is multiplied locally
by \(\mathfrak p_w^{-\lfloor e_w s_w\rfloor}\), where
\(\mathfrak p_w\) is the maximal ideal of the local integers.
The sign comes from the inverse fractional ideal in
\srcloc{Formalism/SizeWindowBrackets/SlopesHypothesisFinitePlaces}{27}.
The data giving \(s_w\) are recorded in
\srcloc{Formalism/LocalToGlobal/GlobalTheory/IsResidueClassZero}{441},
and their use on the input in
\srcloc{Formalism/LocalToGlobal/Columns/FiniteSumNonzeroLargeDepth}{1370}.
The following calculation gives the entire part supported over
this word's own \(T\).

\begin{theorem}[The terminal normalization on the coded family]\leavevmode
\label{thm:terminalscalar357}
The multiplier supported over \(T\) is the principal ideal
\((\pi^{A})\) of \(\mathcal O_J\), where
\[
 A=m\left(2\sum_{\tau}e_{m,\tau}+fD_0+
                          \frac{hD_0(m+1)}4\right).
\]
Here \(f\) is the number of embeddings in the CM type and
\(e_{j,\tau}\) are the prescribed count weights of
\eqref{eq:letters111}. The nonnegative integer \(A\) is
the same for every word and translate at this \(x\).
If \(S'_T\) is the lattice of Theorem~\ref{thm:neighborindex356},
the input with all the prescribed changes at this \(T\) and
ordinary input elsewhere is exactly \(\pi^{A}S'_T\).
\end{theorem}
\begin{proof}
For an embedding \(\tau\), let \(W_\tau\) be the maximum
absolute coordinate of its integral weight over all stages
\(0\le i\le m\), omitting the first coordinate at each stage
when \(\tau\) is high. The terminal error per lift in the
prescription is zero at \(w\nmid\pi\), and at \(w\mid\pi\)
it is
\[
 E_T=2\sum_\tau W_\tau+f n_m+\frac{hD_0(m+1)}4 .
\]
This is the numerical rule in
\srcloc{Formalism/LocalToGlobal/GlobalTheory/IsResidueClassZero}{394},
lines 394--408 and 441--449 there. To evaluate the maximum,
write \(u_a-\rho_k(a)\) for an initial coordinate, where
\(\rho_k(a)=(k+1)/2-a\). At stage \(i\) it gains the shift
\(-iD_0/2\). The block introduced at step \(j\le i\) has
constant coordinate
\[
       -e_{j,\tau}+\frac{k}{2}
                      +\frac{(2j-1-i)D_0}{2}.
\]
Indeed at its introduction the unshifted block has value
\(-e_{j,\tau}+k/2+(2j-1)D_0/2\) and the shift applies to
every coordinate. These are the same lists as in
Corollary~\ref{cor:normaldegree347} at the last stage.
The minimum \(B_\tau\) in the coded family grows faster than
every fixed polynomial of \(x\), by \eqref{eq:letters111} and
the choices following it. Thus all block coordinates are negative
throughout the chosen intervals for large \(x\). For fixed \(j\)
their absolute value is largest at \(i=m\). Increasing \(j\)
there increases the absolute value by \(2H_0-D_0>0\). The
largest block value in absolute value is consequently
\[
                     e_{m,\tau}-k/2-(m-1)D_0/2 .
\]
Every retained initial coordinate at a non-high embedding has
absolute value bounded by a fixed polynomial in \(x\), even
after the stage shift. At a high embedding the same is true of
the interior coordinates. The last initial coordinate there
has absolute value at most \(b_*+mD_0/2\), while
\(B_\tau-b_*\ge\lfloor b'/x\rfloor\) grows faster than any
fixed polynomial. Hence the displayed block gives exactly
\(W_\tau\) at every embedding. Substitution with \(n_m=k+mD_0\)
gives \(m E_T=A\).

It remains to identify which local terms contribute to \(s_w\).
The generator gives a prime split completely in \(E_x\),
different from \(d,g_1,g_2\). At such a prime the count gain
at each lift is zero since only the inert and signal classes
carry this gain. Deleting primes from the gain leaves this
value zero. The second arithmetic summand at a lift is also
zero. It is supported on the two auxiliary primes, on the
prime ramified in \(K_d\), and on the specified inert or
signal primes below the cutoff. The formulas are given in
\srcloc{Formalism/LocalToGlobal/GlobalTheory/IsResidueClassZero}{333}
and \srcloc{Formalism/LocalToGlobal/Divisibility/Numbers}{197}.
The terminal error above is subtracted once for each of the
\(m\) lifts. It follows that \(s_w=-A\) above \(\pi\)
and \(s_w=0\) above \(\bar\pi\).

The first two summands in \(A\) are integers. The last is an
integer since \(mD_0\) is divisible by four. The prime ideals
of \(\pi,\bar\pi\) are coprime, and their product is \((T)\).
Consequently \(\operatorname{ord}_w(\pi)=e_w\) at \(w\mid\pi\),
and zero at every other place. The multiplier has exponent
\(e_w A\) above \(\pi\), with zero at all other places. This
proves the principal-ideal formula and the lattice identity.
Finally \(e_{m,\tau}=B_\tau+(2m-1)H_0\). The three sets of
high, distinguished and remaining places have fixed
cardinalities. Their preliminary \(B_\tau\) values are fixed
on each set, and the parity correction adds the same total
zero or one at each translate. Therefore \(\sum_\tau e_{m,\tau}\)
is independent of the word and the translate.
\end{proof}



\subsection{Whittaker coefficients on a prescribed pair of residue flags}
\label{sec:flagwhittaker526}

Fix the field, the translate and the terminal marking of
Section~\ref{sec:missing}. We examine the two auxiliary primes
in the finite factor of the projected Poincar\'e series.
At the places of \(E^+\) over \(g_1g_2\) put
\[
 G=\prod_v\GL_k(\mathbb F_{q_v}),\qquad U=\prod_v U_k(\mathbb F_{q_v}),
 \qquad \theta=\prod_v\theta_{k,v},\qquad d=[G:U].
\]
Here \(U_k\) is upper unitriangular and \(\theta_{k,v}\) is the
character of its first superdiagonal from Section~\ref{sec:whit125}.
Use counting measure. The matrix \(A\) in
\eqref{eq:actual125} has diagonal one. For arrays \(z_i\) on
\(G\) with values in a complex Hilbert space, write
\(Q(z)=\sum_{i,j}A_{ij}\langle z_i,z_j\rangle\), with the inner
product linear in its second argument. The entries of \(A\) are
real and symmetric by Theorem~\ref{thm:char125}.

\begin{lemma}[An integral matrix for the normalized finite pairing]
\leavevmode\label{lem:whittakermatrix526}
Choose representatives \(a\) for \(G/U\) and \(h\) for \(U\backslash G\).
Define the matrix \(H\), with entries in the integers of a field
containing the values of \(\theta\), by
\[
          (Hz)_{a,h}=\sum_{u\in U}\theta(u) z_{auh}.
\]
Then \(H^*H=dA\). Changing representatives multiplies each output
coordinate by a root of unity. These identities hold for any family
of columns, including signed sums evaluated at fixed tests.
\end{lemma}
\begin{proof}
For fixed \(a\) the sets \(aUh\) partition \(G\). Two indices
\(i,j\) occur in the same summand of the squared norm just when
\(a^{-1}ji^{-1}a\in U\). If \(i=auh\) and \(j=au'h\), their
coefficient is \(\overline{\theta(u)}\theta(u')
=\theta(a^{-1}ji^{-1}a)\). Sum over \(aU\). The result is the
induced character \(\operatorname{Ind}_U^G\theta\) evaluated at
\(ji^{-1}\), by the finite induction formula. This element is
conjugate to \(i^{-1}j\). The character is real, and replacing its
argument by the inverse conjugates its value. Dividing by its
degree \(d\) thus gives the entry of \eqref{eq:actual125}.
Replacing \(a,h\) by \(a u_0,u_1 h\) multiplies the coordinate by
\(\theta(u_0u_1)^{-1}\) by a change of variable in \(U\).
Polarization gives the identity for all pairs of columns.
\end{proof}

Let \(B\) be the product of the upper triangular subgroups.
Fix at every place an input and an output complete flag. The set
\(S\) of tuples of matrices sending the former flags to the latter
has the form \(xBy^{-1}\). There is no restriction on the
coefficients at different indices of \(S\).
Set
\[
 \beta=\prod_v\frac{(1-q_v^{-1})^k}{\prod_{j=1}^k(1-q_v^{-j})},
 \qquad
 \alpha=\prod_v\prod_{j=1}^{k-1}(1-q_v^{-j})^{-1}.
\]

\begin{theorem}[Integral recovery and norms on a pair of flags]
\leavevmode\label{thm:flagrecovery526}
For arrays supported on \(S\), a proportion \(\beta\) of the
choices \(aU\) have the following property. The single row family
\(((Hz)_{a,h})_{Uh}\) recovers every \(z_i\) by selecting an entry and
multiplying it by a root of unity. The same selections work for
all columns and all coefficient modules. In particular the map on
this support has an integral left inverse at every coefficient prime.
Its normalized form satisfies
\[
           \beta\sum_{i\in S}\|z_i\|^2
                   \ \le\ Q(z)\ \le\ \alpha\sum_{i\in S}\|z_i\|^2 .
\]
For any matrix of such columns the ideals generated by minors of
each order before and after applying \(H\) are equal over the indicated
coefficient ring. Fractional entries are allowed.
\end{theorem}
\begin{proof}
For the standard flag in one factor let its opposite be the flag
of the reversed coordinate basis. Its stabilizer intersects the
upper unitriangular group trivially. The flags opposite the standard
flag are exactly the translates of this opposite flag by upper
unitriangular matrices. Indeed in each successive subspace of an
opposite flag there is a unique new basis vector with its reversed
pivot normalized to one, its later coordinates zero after elimination,
and with arbitrary entries above its pivot. These vectors give
exactly one upper unitriangular matrix acting on the reversed
basis. Thus their number is \(q^{k(k-1)/2}\).
All flags number \([\GL_k(\mathbb F_q):B_k]\).
Since \(|\GL_k(\mathbb F_q)|=q^{k^2}\prod_{j=1}^k(1-q^{-j})\)
and \(|B_k|=(q-1)^k q^{k(k-1)/2}\), division gives the
local factor in \(\beta\).

Select \(a\) whose flag is opposite the prescribed output flag at
every place. Each flag has the same number of lifts to \(G/U\).
Its unipotent group \(aUa^{-1}\) meets \(xBx^{-1}\) trivially:
in opposite bases the two groups are respectively lower unitriangular
and upper triangular, up to a simultaneous change of frame.
If \(i,j\in S\) occur in the same coordinate for this \(a\), then
\(ji^{-1}\) lies in that intersection. Hence \(i=j\).
For each \(i\in S\) its unique coordinate contributes exactly
\(\theta(u)z_i\), and different indices use distinct coordinates.
This gives the left inverse. Cauchy--Binet in both directions
gives equality of the determinantal ideals.

Each selected \(a\) contributes \(\sum_{i\in S}\|z_i\|^2\) to the
squared norm. Every other contribution is a squared norm as well.
Use Lemma~\ref{lem:whittakermatrix526} and divide by \(d\).
The upper estimate is Proposition~\ref{prop:amplitudes126}.
\end{proof}

\begin{corollary}[Several output flags at the marked tests]
\leavevmode\label{cor:severalflags526}
Fix the input flags, and split any signed array on \(G\) according
to the tuple \(b\) of output flags, writing \(z=\sum_b z_b\).
With \(S_2=\sum_b\|z_b\|^2\) one has
\[
  0\ \le\ \alpha S_2-Q(z)
             \ \le\ (\alpha-\beta)\left(\sum_b\|z_b\|\right)^2 .
\]
In particular for any family of columns supported on the same
\(h\) tuples of flags its normalized Gram matrix is bounded below by
\((\alpha-h(\alpha-\beta))\) times its original Gram matrix.
On the prescribed intervals
\[
 0\le \alpha-\beta \ll f k(g_1^{-1}+g_2^{-1}),\qquad
  1-\beta\ll f k(g_1^{-1}+g_2^{-1}).
\]
These estimates are uniform in the word, the terminal prime and
the marking. For instance, for \(h\le b_x^4\), where
\(b_x=\exp(x^\varepsilon)\), the normalized form on this entire
support differs from the original form by relative
\(O(f k (k b_x)^{-8} b_x^4)\).
\end{corollary}
\begin{proof}
The positive operator \(D=\alpha I-A\) has quadratic form on each
individual flag support bounded above by \((\alpha-\beta)\|z_b\|^2\).
Apply the triangle inequality to \(\sum_b D^{1/2}z_b\) and square.
Disjoint supports give \(\|z\|^2=S_2\). Apply Cauchy--Schwarz to
the sum of \(h\) norms and test linear combinations of the columns.
At a residue size \(q\ge3\), the upper factor has log at most
\(3/q\) by Section~\ref{sec:moment126}, and the lower factor has
negative log at most \(-k\log(1-q^{-1})\le 2k/q\).
Each size occurs \(f\) times and \(g_i\ge(k b_x)^8\).
The exponents tend uniformly to zero since \(f,k\) have polynomial
growth. The stated estimates follow by exponentiation.
\end{proof}

In the orbital expansion the array at index \(i\) is obtained by
summing the weighted amplitudes with that common relative residue
matrix, with a separate component for each ordered pair of lattices.
The definition of these three indices is in
\srcloc{Formalism/LocalToGlobal/SiegelLatticeTheory/ClosedLocalOrbits/TermOp}{322}.
One may first hold the terminal data fixed, sum the amplitudes at
each index and then apply the corollary. Absolute convergence, when
an orbital series is used, permits this finite regrouping.
For a single flag support the theorem recovers the resulting
coefficients including their cancellations within each index. For
the full prescribed period the distribution among output flags,
the sums within the individual indices, and their relation to
the integral theta lattice still enter the comparison in
Section~\ref{sec:missing}.


\subsection{The integral image of the finite Whittaker pairing}
\label{sec:whittakerimage527}

We now allow every relative residue matrix at the two auxiliary
primes of Section~\ref{sec:flagwhittaker526}. Put \(N=|G|\).
Call an irreducible complex representation of a local factor
\(G_q=\GL_k(\mathbb F_q)\) generic if it contains the character
\(\theta_q\) on the upper unitriangular subgroup \(U_q\). On the
regular representation of \(G\) let \(P\) be the orthogonal
projection onto the sum of the isotypic spaces whose factors at
every auxiliary place are generic. It acts on the residue index
alone, including for vector-valued arrays.

\begin{theorem}[The spectrum on all the relative residue matrices]
\leavevmode\label{thm:whittakerspectrum527}
In the notation of Lemma~\ref{lem:whittakermatrix526},
\[
       \ker H=\ker P,\qquad \beta P\ \le\ A\ \le\ \alpha P .
\]
On each surviving isotypic space of \(G\), of irreducible degree
\(d_\pi\), the scalar of \(H^*H\) is \(N/d_\pi\).
In particular the relative norm estimate applies to arbitrary signed
amplitudes after projection to \(P\), with
\(\alpha-\beta\ll fk(g_1^{-1}+g_2^{-1})\) on the period interval.
\end{theorem}
\begin{proof}
First work over one finite field. The intertwining algebra of
\(\operatorname{Ind}_{U_q}^{G_q}\theta_q\) is commutative. To see
this, apply the anti-involution \(g\mapsto w_0 g^{\mathsf t}w_0\),
where \(w_0\) reverses the basis. It preserves \(U_q\) and
\(\theta_q\). By the double coset argument of
Lemma~\ref{lem:intertwine126}, the functions in this algebra
have the same character law on both sides and are supported on
cosets with representatives \(tw\) of the following form.
The permutation maps consecutive intervals increasingly onto
intervals in reverse block order, and the diagonal \(t\) is constant
on each image interval. Each nonzero block of the monomial matrix is
therefore a scalar identity block. Reflection across the antidiagonal
preserves each such block and reverses its entries along the block
diagonal. It fixes \(tw\). The two character laws show that the
anti-involution fixes every function in the algebra. It reverses
convolution by a change of summation variable. Thus convolution
commutes. By complete reducibility the endomorphism algebra has one
matrix factor of size the multiplicity of each constituent, so every
constituent has multiplicity one.

The central convolution matrix \(A_q\) acts on an irreducible
\(\pi\) by \(|U_q|m_\pi/d_\pi\), by character orthogonality as in
Proposition~\ref{prop:amplitudes126}. Here \(m_\pi\) is the
multiplicity just computed. Conjugating \(\theta_q\) leaves the
set of constituents fixed since the inverse character is obtained
by conjugation with a diagonal matrix. This also covers either
regular action convention. The upper bound for this scalar is
the local factor of \(\alpha\), by that proposition.

For the lower bound we use Green's character degree formula
\cite{XVIII-XIII-green527}. Its parameters are partitions \(\lambda_\phi\)
at Frobenius orbits \(\phi\) of degrees \(e_\phi\), with
\(\sum_\phi e_\phi |\lambda_\phi|=k\).
Use the indexing in which a one-row partition for the trivial
degree-one orbit gives the trivial representation.
Writing \(n(\lambda)=\sum_{i\ge1}(i-1)\lambda_i\) and \(h(c)\)
for the hook length of a cell, the degree is
\[
  d_\pi=\prod_{j=1}^k(q^j-1)
       \prod_\phi \frac{q^{e_\phi n(\lambda_\phi)}}
           {\prod_{c\in\lambda_\phi}(q^{e_\phi h(c)}-1)} .
\]
In a partition the sum of the hooks is the number of cells plus
\(n(\lambda)+n(\lambda')\), by counting cells to the right and
below. On extracting powers of \(q\) we obtain
\[
 \frac{d_\pi}{|U_q|}
  =q^{-\sum_\phi e_\phi n(\lambda'_\phi)}
     \frac{\prod_{j=1}^k(1-q^{-j})}
          {\prod_{\phi,c}(1-q^{-e_\phi h(c)})}
  \ \le\ \frac{\prod_{j=1}^k(1-q^{-j})}{(1-q^{-1})^k}.
\]
There are at most \(k\) factors in the denominator. Taking the
reciprocal gives the local factor of \(\beta\) for each generic
representation. Irreducibles, characters and the matrices of the
finite pairings tensor over the places. Multiplication of the local
bounds and of the nonzero scalars gives the assertion, since
\([G:U]|U|=N\). The kernel identity follows by positivity.
\end{proof}

\begin{proposition}[The dimension discarded at the two auxiliary primes]
\leavevmode\label{prop:whittakerrank527}
On the same intervals, for all sufficiently large \(x\),
\[
   1-\frac{\operatorname{rank}P}{N}
        \ \le\ 8 f(g_1^{-3}+g_2^{-3}).
\]
The rank here refers to the residue index with scalar coefficients.
\end{proposition}
\begin{proof}
It suffices to work at \(q\ge4k\), as holds for both primes.
For \(k\ge2\), the upper estimate on the eigenvalue of \(A_q\)
shows that a generic representation has degree at least
\(|U_q|\prod_{j<k}(1-q^{-j})\ge |U_q|(1-1/(q-1))\).
If any \(\lambda_\phi\) has two cells in the same row, the exponent
\(\sum e_\phi n(\lambda'_\phi)\) in the preceding proof is at least
one. Its degree is at most \(2|U_q|/q\), since
\((1-q^{-1})^{-k}\le\exp(k/(q-1))\le2\).
Such a parameter is nongeneric.

To count the remaining parameters let \(e_d\) be the number of
degree-\(d\) Frobenius orbits on the multiplicative character
parameters. Frobenius to the \(a\)-th power fixes \(q^a-1\) such
parameters. Indeed, in the direct limit of the character groups
under pullback by norms these are exactly the characters of order
dividing \(q^a-1\), a cyclic group of that size. It follows that
\(\sum_{d\mid a}d e_d=q^a-1\). On taking logarithms of the
Euler product, this identity gives
\[
                    F(t)=\prod_{d\ge1}(1-t^d)^{-e_d}
                              =\frac{1-t}{1-qt}.
\]
The parameters with all partitions columns are counted by \(F(t)\).
Its coefficient of \(t^k\) is \((q-1)q^{k-1}\), the dimension of
the intertwining algebra by Lemma~\ref{lem:intertwine126}. Multiplicity
one identifies this dimension with the number of distinct generic
representations. Since every generic parameter has columns, all
the column parameters are generic.

Allowing arbitrary partitions gives the generating series
\(F(t)J(t)\), where \(J(t)=\prod_{j\ge2}F(t^j)\), by recording
the multiplicities of parts of each size. The number of nongeneric
parameters at degree \(k\) is at most \(q^k(J(1/q)-1)\): the
coefficients of \(J-1\) are nonnegative and the degree-\(a\)
coefficient of \(F\) is at most \(q^a\). The product for \(J(1/q)\)
telescopes to \((1-q^{-1})^{-1}\). The lost dimension in the
regular representation is the sum of squared degrees of at most
\(q^k/(q-1)\) constituents, each of degree at most \(2|U_q|/q\).
Since \(|G_q|=q^k|U_q|^2\prod_{j\le k}(1-q^{-j})\), their
fraction is at most \(4/(q^2(q-2))\le8q^{-3}\). Rank
fractions multiply in the product group. Subtract their product
from one and use the sum of the local deficits. For \(k=1\) every
character is generic and the deficit is zero.
\end{proof}

\begin{theorem}[Integral comparison for arbitrary signed columns]
\leavevmode\label{thm:whittakerintegral527}
Let \(Z\) be any finite matrix of columns in the residue coordinates,
with entries in a number field containing the values of \(\theta\).
A coefficient space with a fixed basis is allowed, with the identity
on this basis understood in all operators. Write \(\mathfrak d_j\)
for the fractional ideal generated by the minors of size \(j\),
including the zero ideal in deficient rank. Then
\[
     N^{2j}\mathfrak d_j(PZ)\ \subseteq\ \mathfrak d_j(HZ)
                          \ \subseteq\ \mathfrak d_j(PZ).
\]
Away from \(N\) the two ideals agree for every \(j\). For independent
projected columns the total logarithmic discrepancy of their
maximal-minor ideals, normalized per column and by the field degree,
is at most \(2\log N\).
\end{theorem}
\begin{proof}
For each complex irreducible use its central idempotent
\[
             e_\pi=\frac{d_\pi}{N}\sum_{g\in G}
                             \chi_\pi(g^{-1})g .
\]
The sum over the surviving constituents is \(P\). Their set is
Galois invariant. Indeed a Galois conjugate of the inducing
character at a place replaces its nontrivial additive character
by a scalar multiple in the additive argument, obtainable by
diagonal conjugation on the first superdiagonal. The same argument
applies place by place. Character values are algebraic integers
because the eigenvalues of finite-order operators are roots of unity.
Consequently \(NP\) and
\(N\sum_{\pi\ {\rm surviving}}d_\pi e_\pi\) belong to
\(\mathbb Z[G]\), by integrality and Galois invariance of every
coefficient. Put \(C=N^{-1}\sum_{\pi\ {\rm surviving}} d_\pi e_\pi\).
The scalar calculation gives \(CH^*H=P\) and \(HP=H\).
Thus \(HZ=H(PZ)\) and \(N^2PZ=(N^2 C)H^*(HZ)\).
All multiplying matrices on these right sides are integral,
including \(H^*\) since conjugates of the roots in \(H\) are
their inverses. Cauchy--Binet gives the inclusions. This argument
also works in local lattice bases with a common coefficient lattice
on the residue indices. Taking ideal norms and using
\(|\operatorname{N}_{J/\mathbb Q}(N)|=N^{[J:\mathbb Q]}\)
proves the logarithmic bound when the minor ideals are nonzero.
\end{proof}

\begin{corollary}[The same terminal primes as the strict groups]
\leavevmode\label{cor:whittakerterminal527}
At each \(x\) on the common intervals one can impose \(T\nmid N\)
in Corollary~\ref{cor:sameprimes517}. There are still
\(Y_x^{1-o(1)}\) eligible choices, with full strict vanishing
on all the subfields and in all the weights asserted there.
For these choices the Whittaker and projected minor ideals in
Theorem~\ref{thm:whittakerintegral527} agree at every place
above the actual terminal prime \(T\), simultaneously for all
sets of columns. Over all coefficient primes the normalized
logarithmic discrepancy is \(o(b_x/x^A)\) for every fixed \(A>0\).
\end{corollary}
\begin{proof}
Select the two auxiliary primes by Corollary~\ref{cor:sameprimes517}
before selecting \(T\). With \(f\) factors of each residue size,
\[
             \log N\le f k^2\log(g_1g_2)\ll_{\psi,\varepsilon} x^C
\]
for a fixed \(C\), by the polynomial field and rank bounds and the
auxiliary intervals. At most \(\log N/\log2\) more primes are
excluded. The stated terminal window has \(\log Y_x\gg m x^\varepsilon\),
so this exclusion preserves its count in logarithmic order
uniformly for the prescribed data. Apply the two inclusions above
at \(T\), where \(N\) is a unit. At all places use \(2\log N\)
and \(b_x=\exp(x^\varepsilon)\).
\end{proof}

These results apply to the signed arrays in the original relative
residue coordinates at a fixed marking. They give the same rank
as the explicit projection \(P\). The small dimension fraction
lost by \(P\) speaks to the full regular representation.
A prescribed family of columns may lie in its kernel. The rank
and height of the surviving arithmetic amplitudes and the remaining
strict groups on the common interval are still required in
Section~\ref{sec:missing}.


\subsection{Alternating sums of the signed amplitudes on complete flags}
\label{sec:apartmentperiod528}

At each of the two auxiliary primes we extract from the pairing in
Theorem~\ref{thm:whittakerspectrum527} a set of integer linear forms
in the signed amplitudes. They can be evaluated with the given
terminal marking. Their coefficients are integers of absolute value one on the indicated cosets.

Work first with \(G_q=\GL_k(\mathbb F_q)\). Write
\(X_q=G_q/B_q\) for the set of complete flags, with base flag \(C_0\).
A panel is a flag with one subspace of dimension between \(1\) and
\(k-1\) omitted. Call a function on \(X_q\) harmonic if its sum over
the \(q+1\) complete flags on each panel is zero. For \(u\in U_q\)
write
\[
   v_u=\sum_{w\in\mathfrak S_k}(-1)^{\ell(w)}\delta_{uwC_0}.
\]
This is the signed sum on the apartment given by the basis \(u e_i\).
Here \(\ell(w)\) counts inversions. Use counting inner products on all
finite sets. The relative permutation of \(aC_0,bC_0\) is the unique
\(w\) with \(a^{-1}b\in B_q w B_q\).

\begin{theorem}[The finite pairing on the complete flag space]
\leavevmode\label{thm:flagprojector528}
The \(v_u\) form an integral basis of the harmonic functions. In the
space of functions on \(G_q\) constant on each right \(B_q\)-coset,
the local matrix \(A_q\) of Section~\ref{sec:whit125} acts as
orthogonal projection onto the harmonic functions. In flag coordinates
this projector has entries
\[
 P_q(C,D)=\frac{|U_q|}{|X_q|}(-q)^{-\ell(w(C,D))}.
\]
In particular \(|X_q|P_q\) has integer entries.
\end{theorem}
\begin{proof}
Every apartment containing a chamber on a given panel contains exactly
two such chambers, whose permutations differ by an adjacent exchange.
Their signs cancel. With \(w_0\) the reversing permutation, the
opposite chambers are \(uw_0C_0\), one per \(u\in U_q\), as in
Theorem~\ref{thm:flagrecovery526}. The restriction of the displayed
family to these chambers is the identity matrix times the sign of
\(w_0\).

To see why restriction is injective on harmonic functions, take a
chamber of position \(w\ne w_0\) from \(C_0\). Choose an ascent
\(w(j)<w(j+1)\). On removing dimension \(j\), the other \(q\)
ways of filling that panel have position \(w s_j\), of length
\(\ell(w)+1\). Indeed one may move the chamber by an upper
triangular matrix to \(wC_0\); in the two-dimensional quotient the
other lines have a nonzero coefficient on \(e_{w(j+1)}\) and an
arbitrary coefficient on \(e_{w(j)}\). Adding the latter coordinate
to the former pivot is an upper triangular basis change.
The zero panel sum therefore determines the value at the original
chamber from larger lengths. Descending induction proves injection.
Together with the opposite restrictions this proves the basis
assertion over the integers, and the same argument works with
coefficients in any abelian group.

Over \(\mathbb C\) the harmonic space is an irreducible
\(G_q\)-representation. For a \(B_q\)-invariant harmonic vector its
values depend only on \(w\), and the same panel equation reads
\(h(w)=-q h(ws_j)\) at an ascent. Thus its fixed space under \(B_q\)
has dimension at most one. Each irreducible summand of the nonzero
harmonic space occurs in \(\mathbb C[G_q/B_q]\) and has a nonzero
\(B_q\)-fixed vector by finite group reciprocity. Complete
reducibility gives the assertion. Its degree is \(|U_q|\).

The \(\theta_q\)-space in all flag functions has dimension one.
Indeed every \(U_q\)-orbit has a representative \(wC_0\), by
Bruhat decomposition and diagonal conjugation. If \(w\ne w_0\)
there is an \(i\) with \(w^{-1}(i)<w^{-1}(i+1)\). The simple
root group \(1+tE_{i,i+1}\) fixes this chamber, and \(\theta_q\)
is nontrivial on it. Equivariance forces zero on its orbit.
On the opposite orbit the action is free and transitive. A
\(\theta_q\)-function there is determined by one value and can be
extended by zero. This one-dimensional character space already
occurs in the harmonic representation. In fact restriction to
opposite chambers intertwines \(U_q\) and is an isomorphism onto
its regular representation by the integral basis just proved.
It follows that the harmonic summand is the only generic occurrence
in the flag space. The scalar formula of
Theorem~\ref{thm:whittakerspectrum527} gives one on it and zero on
its orthogonal complement. Lifting flag values to all matrices in
each coset multiplies their inner product by the fixed number
\(|B_q|\), so this is also the orthogonal projection for counting
measure on flags.

Each diagonal entry of the projector is its rank divided by
\(|X_q|\), by transitivity. Its column at \(C_0\) is invariant
under \(B_q\) and harmonic. Starting from its diagonal value, the
ascent relation gives the displayed factor on every Bruhat cell.
Translation proves the general formula; the inverse permutation
has the same length. Finally \(|U_q|=q^{\ell(w_0)}\) clears all
powers of \(q\) in this formula after multiplication by \(|X_q|\).
\end{proof}

Return to the product of all auxiliary places, and use
\(X=\prod_v X_{q_v}\), \(B=\prod_v B_{q_v}\) and the groups
\(G,U\) from Section~\ref{sec:flagwhittaker526}.
Let \(J\) copy flag values to their right cosets in \(G\).
For \(u=(u_v)\in U\), put \(v_u=\bigotimes_v v_{u_v}\) and
let \(V\) be the matrix with these columns on \(X\). Put \(K=V^*V\).
For any matrix \(Z\) of signed arrays in the original relative
residue coordinates define \(M=V^*J^*Z\). Thus each coordinate of
a column is an explicit alternating sum
\[
 M_u=\sum_{w\in\prod_v\mathfrak S_k}
       (-1)^{\sum_v\ell(w_v)}\sum_{i\in u w B}Z_i .
\]
The inner sums can use the signed coefficients at the prescribed
marking in Section~\ref{sec:flagwhittaker526}. The formulas act
on the residue coordinates and preserve any additional coefficient
index. Write \(C_u=(u_v w_0 C_{0,v})_v\).

\begin{theorem}[A lower bound at the fixed signed amplitudes]
\leavevmode\label{thm:markedapartments528}
For any such \(Z\) the Gram matrices satisfy
\[
 Z^* A Z\ \ge\ \frac{1}{|B|} M^* K^{-1} M,\qquad
 (K^{-1})_{ut}=\prod_v P_{q_v}(C_{u,v},C_{t,v}).
\]
There is equality for columns constant on the right cosets of \(B\).
In particular for arbitrary columns one has
\[
 Z^* A Z\ \ge\ \frac{M^*M}{|B|\,|U|\,(k!)^{2f}} .
\]
The logarithm of the denominator is polynomially bounded in \(x\)
on the same fields and intervals, uniformly in the translate.
\end{theorem}
\begin{proof}
The orthogonal projection onto the tensor product of harmonic spaces
is \(V K^{-1}V^*=\bigotimes_v P_{q_v}\). On the opposite flags
the rows of \(V\) form the identity times a common sign.
Restricting the projector to those rows and columns gives the
claimed entries of \(K^{-1}\).
The columns of \(JV\) lie in the eigenvalue-one space of \(A\).
Their Gram matrix is \(|B|K\) since \(J^*J=|B|I\).
The projection onto their span is dominated by \(A\), which is
self-adjoint and positive on the orthogonal complement. Its
compression by \(Z\) is exactly the first right side.
For a right \(B\)-invariant column, the tensor product form of the
preceding theorem identifies the action of \(A\) with this projection.

Each column of \(V\) has \((k!)^{2f}\) entries of modulus one and
all its other entries vanish. At each row there are at most \(|U|\)
nonzero entries. Cauchy--Schwarz row by row, followed by summing
over the columns, gives \(K\le |U|(k!)^{2f}I\). Inverting gives
the second estimate. The bounds on \(f,k\) and the auxiliary primes
in Corollary~\ref{cor:whittakerterminal527} give
\(\log(|B|\,|U|\,(k!)^{2f})\le 2\log|G|+2fk\log(2k)\ll x^C\)
for a fixed \(C\).
\end{proof}

\begin{corollary}[Integral tests at the common terminal prime]
\leavevmode\label{cor:apartmentideals528}
Write \(N=|G|\), \(d=[G:U]\), and use the minor ideals from
Theorem~\ref{thm:whittakerintegral527}. For every \(j\) one has
\[
             d^j\mathfrak d_j(M)\ \subseteq\ \mathfrak d_j(HZ).
\]
If the columns are right \(B\)-invariant one also has
\(N^j\mathfrak d_j(HZ)\subseteq\mathfrak d_j(M)\).
At every place over a terminal prime from
Corollary~\ref{cor:whittakerterminal527}, a unit minor of the
alternating sums gives a unit minor of the same size in the
Whittaker matrix of integral columns. For right \(B\)-invariant
columns all the localized minor ideals agree there.
\end{corollary}
\begin{proof}
Since \(H^*H=dA\) and \(A(JV)=JV\), one has
\(dM=(HJV)^*HZ\) with integral coefficients on the left of \(HZ\).
For invariant columns \(H\) kills their component orthogonal to
\(JV\) in the flag space by Theorem~\ref{thm:flagprojector528}
and positivity of \(H^*H\). Hence
\[
            HZ=(HJV)(|B|K)^{-1}M .
\]
The formula for \(K^{-1}\) shows \(|X|K^{-1}\) is integral.
Multiplying by \(|B||X|=N\) clears the matrix on the right.
Cauchy--Binet proves both inclusions. At the selected \(T\), both
\(N\) and \(d\) are units. Over a local ring the ideal generated by
a finite list of integral elements can contain a unit only if some
element of the list is a unit, which gives the unit-minor assertion.
\end{proof}

The sign test sums all contributions at each residue index and uses
the terminal prime at which the strict groups vanish. To apply it
to the full period, one still has to prove that the prescribed
columns have a nonzero alternating minor of the required size, or
establish their independence in other generic representations. The
arithmetic sizes of the surviving period coefficients also enter
Section~\ref{sec:missing}.


\subsection{Rational isometries at the two auxiliary primes}
\label{sec:isometryseparation530}

Let \(E=E_x\), \(F=E^+\), \(f=[F:\mathbb Q]\), and retain the two
split primes \(g_1,g_2\) and the matrix \(H\) of
Section~\ref{sec:flagwhittaker526}. At each \(v\mid g_1g_2\) of \(F\)
choose the place of \(E\) giving the residue coordinate in \(G\).
We examine the part of the finite pairing in which these coordinates
come from maps between two fixed totally definite Hermitian
\(E\)-spaces of dimension \(k\). Write \(b_x=\exp(x^\varepsilon)\).
The case \(k=1\) has \(U=1\) and each residue coordinate is read
directly. We assume \(k\ge2\).

Fix lattices \(\Lambda,\Lambda'\) in the two spaces. A rational
isometry means an \(E\)-linear isometry \(m\) between the spaces.
At the chosen places we consider isometries inducing isomorphisms
between the two local lattices. In fixed bases their residues define
an index \(i(m)\in G\). One can multiply all the residue matrices
by the same invertible matrices on the left and on the right at
each place. This permits changes of the two marked frames.

\begin{theorem}[Exact separation for bounded denominators]
\leavevmode\label{thm:definiteseparation530}
Let \(h\ge1\) be a rational integer prime to \(g_1g_2\). Use any
finite set of distinct rational isometries satisfying
\[
                  h\Lambda'\ \subseteq\ m\Lambda\
                                   \subseteq\ h^{-1}\Lambda'.
\]
If \(4kh^2<g_1g_2\), their residue indices are distinct and the
restriction of the normalized Gram matrix \(A\) to them is the
identity. For \emph{every} fixed \(aU\in G/U\), the coordinates
\((Hz)_{a,t}\), \(t\in U\backslash G\), recover every coefficient
of an array on these indices by a selection of entries and
multiplication by roots of unity. In particular, the ideals of
minors of each size agree before and after the finite pairing,
over every coefficient ring of integers containing the character
values. This holds throughout the period intervals for every
\(h\le b_x^8\) prime to \(g_1g_2\).
\end{theorem}
\begin{proof}
Consider distinct maps \(m,n\) and put \(\gamma=mn^{-1}\) on the
target space. Define
\[
        u=2k-\operatorname{Tr}_E(\gamma)
                  -\operatorname{Tr}_E(\gamma^{-1})
          =\operatorname{Tr}_E\big((1-\gamma)(1-\gamma^{-1})\big),
\]
where the trace is the matrix trace over \(E\). The adjoint of
\(\gamma\) is \(\gamma^{-1}\), so \(u\in F\). At each embedding
of \(F\) choose an orthonormal complex basis. The image of \(u\)
is the squared Hilbert--Schmidt norm of \(1-\gamma\), the sum of
the squared absolute values of its entries. It is positive because
\(\gamma\ne1\) and the embedding is injective. Each eigenvalue of
\(\gamma\) has absolute value one, so this norm squared is at
most \(4k\).

The lattice inclusions imply \(h^2\gamma\Lambda'\subseteq\Lambda'\)
and \(h^2\gamma^{-1}\Lambda'\subseteq\Lambda'\). For instance
\(n^{-1}\Lambda'\subseteq h^{-1}\Lambda\), and
\(m\Lambda\subseteq h^{-1}\Lambda'\). Traces of lattice
endomorphisms are integral, as one sees in a basis at each finite
place. Thus \(h^2u\) is a nonzero integer of \(F\) and
\[
            1\le |N_{F/\mathbb Q}(h^2u)|\le(4kh^2)^f .
\]
Suppose the relative residue matrix \(i(m)i(n)^{-1}\) were
unipotent at every marked place. In the chosen bases it is the
reduction of \(\gamma\) up to conjugacy, including after the
two common frame changes. Its trace and the trace of its inverse
are both \(k\) in the residue field. Hence \(u\) reduces to zero
at every \(v\mid g_1g_2\) of \(F\). Since the primes split
completely, \((g_1g_2)^f\) divides the displayed absolute norm,
contrary to the bound.

There is therefore a nonunipotent local relative matrix for
every distinct pair. For fixed \(a\), two entries of an input
array occur in the same output coordinate precisely when their
relative matrix belongs to \(aUa^{-1}\), by
Lemma~\ref{lem:whittakermatrix526}. This would require unipotence
at all places. Each input entry consequently occurs by itself,
multiplied by a root of unity, and distinct inputs have distinct
indices. In this row family the squared norm is exactly the sum
of the input squared norms. Averaging over \(aU\) and polarizing
gives the identity Gram matrix. The entry selections give an
integral left inverse. Applying Cauchy--Binet to both maps gives
the assertion on minors, also for several columns at once.
Finally \(g_i\ge(kb_x)^8\) on these intervals, and
\(4k b_x^{16}<k^{16}b_x^{16}\) for \(k\ge2\).
\end{proof}

In particular the full set of isometries of one definite integral
lattice onto another is covered with \(h=1\), independently of the
orders of their stabilizers. Such a set is finite. Indeed the images
of a fixed \(\mathbb Z\)-basis lie in a discrete lattice and in a
bounded set for the sum of the positive Hermitian norms at infinity.
The same argument applies with target \(h^{-1}\Lambda'\).
An individual residue index in the period may already be a signed
sum. The theorem recovers that sum.

\begin{lemma}[The largest probability of a flag collision]
\leavevmode\label{lem:flagcollision530}
Put \([j]_q=(q^j-1)/(q-1)\). For a nonidentity unipotent matrix
in \(\operatorname{GL}_k(\mathbb F_q)\), the proportion of
complete flags it fixes is at most
\[
        \rho_k(q)=
        \frac{\sum_{s=0}^{k-2}(s+1)q^s}{[k]_q[k-1]_q}
        \ \le\ \frac{k-1}{[k]_q}.
\]
Equality in the first bound holds for a transvection.
\end{lemma}
\begin{proof}
Write the matrix as \(1+N\) with \(N\) nilpotent. A fixed flag
starts with a line in \(\ker N\), and its completion is a fixed
flag for the induced nilpotent on the quotient by that line.
Let \(B_j\) be the number of flags fixed by a transvection in
dimension \(j\), and \(D_j\) the total flag count. In dimension
two a nonzero nilpotent has rank one and fixes just its kernel
line. Assume in dimension \(k-1\ge2\) that \(B_{k-1}\) is the
maximum for nonzero nilpotents. For rank at least two in
dimension \(k\) there are at most \([k-2]_q\) choices of line
and the quotient nilpotent is always nonzero. A rank-one
nilpotent allows \([k-1]_q\) lines, including its image line.
The quotient is zero at that line and rank one at every other.
Thus \(B_k=D_{k-1}+([k-1]_q-1)B_{k-1}\), which dominates
\([k-2]_q B_{k-1}\). This proves maximality by induction.

For rank one put \(L=\operatorname{im}N\) and \(K=\ker N\).
The line lies in the hyperplane. For a fixed flag \(F_i\), let
\(i\) be the first step containing \(L\), and \(j\) the first
step outside \(K\). Stability means \(i<j\). Indeed
\(N F_j=L\), and \(F_j\cap K=F_{j-1}\). Conversely \(L\subseteq
F_{j-1}\) makes every step stable. Relative to the fixed flag
there are \(q^{i-1}\) lines at the indicated step and
\(q^{k-j}\) hyperplanes first missed at step \(j\), by
normalizing the first nonzero coefficient of their defining
functional. For \(i<j\) incidence is automatic. Summing gives
\(\sum_{s=0}^{k-2}(s+1)q^s\) such pairs \((L,K)\).
There are \([k]_q[k-1]_q\) incident pairs altogether.
Transitivity on flags and on incident pairs permits counting in
either order. Bounding each \(s+1\) by \(k-1\) gives the last
inequality.
\end{proof}

\begin{proposition}[The number of places distinguishing rational maps]
\leavevmode\label{prop:heightcollision530}
Retain the two spaces, and require the rational isometries to be
integral isomorphisms at the marked places. For a distinct pair,
use \(u\) from the proof of Theorem~\ref{thm:definiteseparation530}.
Let \(\mathfrak d(u)\) be its denominator ideal in \(F\), and put
\[
 W=f\log(g_1g_2),\qquad
 B_{mn}= f\log(4k)+\log N_{F/\mathbb Q}\mathfrak d(u).
\]
If the relative residues are unipotent everywhere, then
\[
  B_{mn}\ge W,\qquad
  \sum_{\substack{v\mid g_1g_2\\ i(m)_v\ne i(n)_v}}\log q_v
                                      \ \ge\ 2W-B_{mn}.
\]
For a uniformly chosen \(aU\), the probability that the two
indices occur in the same coordinate of the row family
\((Hz)_{a,t}\) is zero if some relative residue is nonunipotent.
Otherwise it is at most
\(\prod_{i(m)_v\ne i(n)_v}\rho_k(q_v)\).
\end{proposition}
\begin{proof}
The positive valuations of \(u\), with weights \(\log Nv\),
sum to \(\log|N_{F/\mathbb Q}u|+\log N_{F/\mathbb Q}\mathfrak d(u)\)
by factorization of its fractional ideal. This is at most \(B_{mn}\).
At every marked unipotent residue the trace calculation gives
\(v(u)\ge1\), with valuation one on a uniformizer. Where the two
indices agree, both \(1-\gamma\) and \(1-\gamma^{-1}\) vanish
modulo the corresponding place of \(E\) in the target lattice
basis. Their product has entries divisible by its square. The
place splits above \(F\), so \(v(u)\ge2\) in \(F\).
Summing gives the two inequalities.

At a local unipotent residue, membership in \(a_vU_qa_v^{-1}\)
is equivalent to preserving the flag of \(a_v\), since every
diagonal eigenvalue in a triangular form is one. The map from
\(G_q/U_q\) to the flags has fibers of the same size.
Lemma~\ref{lem:flagcollision530} bounds the probability when
the residue differs from the identity. The choices at distinct
places are independent in the product \(G/U\).
\end{proof}

\begin{corollary}[Integral recovery at larger denominators]
\leavevmode\label{cor:markeddenominators530}
Suppose that \(s\) distinct maps obey the lattice inclusions of
Theorem~\ref{thm:definiteseparation530} and that, for some
\(0<\xi\le1\),
\[
      4kh^2\le(g_1g_2)^{2-\xi},\qquad
      e=(s-1)(g_1g_2)^{-\xi f(k-1)/2}<1 .
\]
Their indices are distinct. On these indices \(H\) has an integral
left inverse, and the eigenvalues of the normalized Gram matrix
lie in \([1-e,1+e]\). The choices are uniform in the common marked
frames and work for signed columns at the prescribed terminal prime.
\end{corollary}
\begin{proof}
Here \(h^2u\) is integral, so \(B_{mn}\le f\log(4kh^2)\).
Whenever all relative residues are unipotent, the places of
disagreement have total log size at least \(\xi W>0\).
At every marked prime \(q\) one has
\(\rho_k(q)\le(k-1)q^{-(k-1)}\le q^{-(k-1)/2}\), by the
prime-size condition above. Proposition~\ref{prop:heightcollision530}
bounds the probability of collision in a row family by
\((g_1g_2)^{-\xi f(k-1)/2}\) for any distinct pair, also when
some relative residue is nonunipotent.

Fix an input index. The fraction of row families where any other
input shares its coordinate is at most \(e\), by summing over the
other \(s-1\) choices. Some row therefore isolates this input.
Selecting an isolating row separately for each input and undoing
the unit character factors gives the left inverse on the entire
support. The absolute value of an off-diagonal entry of \(A\)
is at most the corresponding collision probability, because it
is an average of roots of unity on the colliding families and
zero on the others. The row sums of the absolute off-diagonal
entries are at most \(e\). The matrix is Hermitian with unit
diagonal, so the eigenvalue bound follows, or equivalently one
can apply \(2|z_i z_j|\le |z_i|^2+|z_j|^2\) to its quadratic form.
\end{proof}

\begin{theorem}[Definite theta coefficients with prescribed Gram matrix]
\leavevmode\label{thm:rectangulartheta530}
Let \(A_0,M\) be integral totally definite Hermitian lattices over
\(\mathcal O_E\), of ranks \(k\) and \(D\ge k\). Let
\(A_0^\#=\{v:\langle v,A_0\rangle\subseteq\mathcal O_E\}\), and
choose positive integers \(d_0,h\), prime to \(g_1g_2\), such that
\(d_0 A_0^\#\subseteq A_0\). Consider all isometric embeddings
\(X:A_0\to h^{-1}M\). At each marked residue place retain the image
subspace of the reduced frame. For each tuple of image subspaces,
choose bases of those subspaces and form \(H\) on the resulting
square frame coordinates. If \(4k d_0 h^2<g_1g_2\), this operation
recovers the individual coefficients of all the embeddings over the
integers of the character field. Its normalized squared norm is
the sum of their squared absolute values.

For free \(A_0\) with basis and Gram matrix \(S\), the assertion
applies to the full \(S\)-coefficient of the definite lattice theta
series with these residue data retained. In particular, on each of
the period fields take \(A_0=\mathcal O_E^k\) with its identity
form, \(M=A_0\perp M_0\) for any integral totally definite \(M_0\),
and \(h=1\). For any assignment of a sign to every embedding, the
refined coefficient after \(H\) has an entry which is a unit at
every coefficient prime.
\end{theorem}
\begin{proof}
Write \(X^*\) for the Hermitian adjoint on the rational spaces.
One has \(X^*X=1\), and \(h X^*M\subseteq A_0^\#\), by taking
pairings with \(A_0\). At each marked place the chosen integers
are units. The reductions of \(X\) and \(X^*\) are therefore
defined on their lattices and their composition in the order
\(X^*X\) is the identity. Thus the frames have rank \(k\).

For distinct embeddings \(X,Y\) set
\[
      u_{XY}=\operatorname{Tr}_E\big((X-Y)^*(X-Y)\big)
             =2k-\operatorname{Tr}_E(X^*Y+Y^*X)\in F.
\]
At an infinite place the two maps are isometric embeddings
between ordinary Hermitian spaces. Each has squared
Hilbert--Schmidt norm \(k\). The triangle inequality gives
\(0<u_{XY}\le4k\) there. At finite places
\(d_0h^2 X^*Y A_0\subseteq A_0\), since \(hY A_0\subseteq M\)
and \(h X^*M\subseteq A_0^\#\). The same inclusion holds for
\(Y^*X\). Thus \(d_0h^2 u_{XY}\) is integral in \(F\).

Different tuples of image subspaces already lie in different
rows. Within a tuple, a collision in a row family of \(H\)
would give, at each chosen place, \(\overline X=\overline Y t\)
for some unipotent \(k\)-by-\(k\) matrix \(t\). Here the bars
mean reduction, and unipotence is invariant under the changes
from the chosen subspace bases. Reducing the adjoint identities
shows
\[
     \overline{Y^*X}=t,\qquad
     \overline{X^*Y}\,t=\overline{X^*X}=1.
\]
These identities use the integral adjoint maps at that place.
The traces of \(t\) and \(t^{-1}\) are both \(k\), so every
place of \(F\) above \(g_1g_2\) divides \(u_{XY}\). The positive
integer \(|N_{F/\mathbb Q}(d_0h^2 u_{XY})|\) would be at least
\((g_1g_2)^f\) and at most \((4k d_0h^2)^f\). This contradicts
the stated inequality. The distinct embeddings therefore occupy
separate coordinates in every row family, and entry selections
and the norm identity follow as in
Theorem~\ref{thm:definiteseparation530}.

In a basis of \(A_0\) an embedding is exactly an ordered tuple in
\((h^{-1}M)^k\) with Hermitian Gram matrix \(S\): the equalities
of the pairings extend by sesquilinearity, and definiteness
implies injectivity. Grouping the defining theta sum over tuples
by this Gram matrix gives its stated Fourier coefficient.
At a fixed \(S\), all images of a \(\mathbb Z\)-basis are
bounded at all infinite places, so only finitely many embeddings
occur. Any weights on these tuples can thus be used in the sum.
For the last assertion \(d_0=1\) is admissible, and the inclusion
of the first summand gives an embedding. Its signed coefficient
is isolated up to a root of unity in every row family.
\end{proof}

\begin{proposition}[Cancellation on singular frame orbits]
\leavevmode\label{prop:isotropickernel530}
Over the same \(E\), take a full frame \(X\) of \(k\) vectors.
Its first \(l\ge1\) vectors are the first Witt basis vectors of
mutually orthogonal integral hyperbolic planes and the remaining \(k-l\) form an
integral basis of a nondegenerate orthogonal summand. At any
given positive precision above \(T\) there is a finite family
of integral frames congruent to \(X\) there, with its column
lattice, Gram matrix and residue image subspaces. On this family
the coefficient array with all weights one is killed by \(H\) if
\(l\ge2\). For \(l=1\) the family constructed below has an
entry equal to a root of unity after applying \(H\).
\end{proposition}
\begin{proof}
Write \(\pi\) for the coordinate projection onto the last \(k-l\)
entries. The determinant-one changes of basis fixing \(\pi\)
are precisely the matrices with block rows \((A,B)\), \((0,I)\),
where \(A\in\operatorname{SL}_l\) and \(B\) is arbitrary. Since
the Gram form factors through \(\pi\), integral changes \(s\)
of this shape preserve both the Gram matrix and the column
lattice of \(Xs\). At each chosen place use the finite group
\[
             S_v=\{s\in\operatorname{SL}_k(\mathbb F_{q_v}):
                                  \pi s=\pi\},\qquad S=\prod_v S_v .
\]
Choose one integral lift of this shape per element of \(S\),
equal to the identity at the specified terminal precision.
For the block \(B\) use the Chinese remainder theorem, taking it
zero at the terminal places. For the block \(A\), over a field
determinant-one matrices are products of elementary row additions
by Gaussian elimination when \(l\ge2\), and \(A=1\) when \(l=1\).
The signed interchange
and the diagonal scaling needed in elimination follow from
\[
  \begin{pmatrix}1&t\\0&1\end{pmatrix}
  \begin{pmatrix}1&0\\-t^{-1}&1\end{pmatrix}
  \begin{pmatrix}1&t\\0&1\end{pmatrix}
       =\begin{pmatrix}0&t\\-t^{-1}&0\end{pmatrix}
       \qquad(t\ne0)
\]
and the product of two such matrices with parameters \(t,-1\).
Lift each elementary parameter by the same Chinese remainder
argument, making it zero also at the other auxiliary places.

In the bases of the reduced frame \(X\) the resulting array
is \(\mathbf1_S\), with a single tuple of image subspaces.
Suppose \(l\ge2\). Fix a row \(a,t\) of \(H\) and one auxiliary
place, suppressing its subscript. The first \(k-1\) columns of
\(\pi a\) lie in dimension \(k-l<k-1\). Some linear dependence
has last nonzero coefficient \(b_i\), with \(i\le k-1\).
Complete its coefficient vector \(b\) by zeros after \(i\).
For \(c\in\mathbb F_q\), the matrix \(u_c=I+c b e_{i+1}^{\mathsf t}\)
is upper unitriangular and \(\pi a u_c=\pi a\).
Hence \(a u_c a^{-1}\in S_v\). Its character before conjugation
is \(\theta_v(u_c)=\psi_v(c b_i)\), which differs from one for
some \(c\). Extend \(u_c\) by the identity at the other places.
The set of \(u\in U\) with \(aut\in S\) is stable under left
multiplication by \(u_c\); its character sum equals that sum
times \(\theta(u_c)\). Every coordinate of \(H\mathbf1_S\) is
therefore zero. This proof includes totally isotropic frames
and frames with a single nonzero Gram entry when \(k\ge3\).

If \(l=1\), each \(S_v\) consists of shears into the first
coordinate line and fixes that line pointwise. Choose \(a_v\)
so that this line is last in its ordered basis. Then
\(a_v U_q a_v^{-1}\cap S_v=1\), since the shears are lower
unitriangular in that basis. Choose the row with \(t=a^{-1}\).
Its support meets \(S\) just in the identity, giving the unit
entry. Changing representatives only multiplies entries by roots
of unity. Common left and right changes of frame relabel \(a,t\)
in the formula for \(H\), so both conclusions hold in those
coordinates as well.
\end{proof}

These statements concern one fixed pair of lattice rows at a
time, as in the index of
\srcloc{Formalism/LocalToGlobal/SiegelLatticeTheory/ClosedLocalOrbits/TermOp}{322}.
They can be applied directly to a restricted period sum whose
residue labels are the reductions of the indicated rational
isometries in common frames. The rectangular theta coefficient in
Theorem~\ref{thm:rectangulartheta530} retains its residue image
subspaces. The first-plane orbital calculation in
Lemma~\ref{lem:realflag265} also uses singular Gram matrices.
Proposition~\ref{prop:isotropickernel530} exhibits cancellation
on its orbit types with radical of dimension at least two, even
with fixed terminal congruences.
The full orbital formula uses adelic Levi matrices and prescribed
path weights. A lower bound for the theta lattice in
Section~\ref{sec:missing} requires control of those actual
weights on the singular orbits and of the signed sums when
residue rows are combined.


\subsection{Integral lifts of a rational partner flag}
\label{sec:partnerlifts531}

Fix two words and an ordered pair of their lattice rows
\(\Lambda',\Lambda\) on the field \(E=E_x\) of
Section~\ref{sec:missing}. Write \(F=E^+\), \(f=[F:\QQ]\), and
\(Y'=(e'_1,\ldots,e'_s,z')\), \(Y=(e_1,\ldots,e_s,z)\)
for the source and target frames. Let \(I'\) be the span of the
\(e'_j\). These vectors give an ordered isotropic flag. The vector
\(z'\) lies in \((I')^\perp\) and has totally negative norm.
Put
\[
 J'=(I')^\perp/I',\qquad C'=(E\bar z')^\perp\subset J',\qquad
 L'=C'\cap\operatorname{im}\big(\Lambda'\cap(I')^\perp\longrightarrow J'\big).
\]
Here and below a bar on \(z'\) means its image in \(J'\).
Removing the \(s\) hyperbolic planes leaves signature \((k,1)\)
at every real place. Hence \(C'\) is totally positive of rank \(k\).
The image in this formula spans \(J'\) and is finitely generated
over \(\mathcal O_E\). Its intersection with the rational subspace
\(C'\) is a full lattice there, as one sees by clearing the
denominators of a basis of \(C'\).

Let \(N'\) be the integral source isometries acting trivially on
the successive lines of the isotropic flag and on \(J'\), and use
\(N\) for the corresponding target group. Let \(S'\) be the integral
source isometries fixing every vector of \(Y'\) pointwise, and write
\(D\) for their image in \(\operatorname{Aut}(L')\). A lifted term
is a rational isometry \(\gamma:\Lambda'\longrightarrow\Lambda\)
with admissible transported frame, taken modulo \(N\) on the left
and \(N'\) on the right. Its geometric partner is the class of the
transported frame \(\gamma Y'\). Two transported frames define the
same partner if one is obtained from the other by first applying
the integral unipotent group of its own isotropic flag and then
an element of \(N\). This definition includes the terminal vector
in both actions. It is the equivalence used in
\srcloc{Formalism/LocalToGlobal/SiegelLatticeTheory/ClosedLocalOrbits/TermDef}{18}.

At each \(v\mid g_1g_2\) take the chosen split factor, of residue
size \(q_v\). Use the original marked bases on the two rows and
write \(i(\gamma)_v\) for the induced matrix on the \(k\)-dimensional
residual complement. There are \(f\) places of each residue size.
Write \(G=\prod_v\operatorname{GL}_k(\mathbb F_{q_v})\),
\(U=\prod_v U_k(\mathbb F_{q_v})\), and retain the character
\(\theta\), the matrix \(H\) and \(A=H^*H/[G:U]\) of
Lemma~\ref{lem:whittakermatrix526}. Work at \(x\) large enough
that \(g_1,g_2>2fk+1\), as holds on the common intervals.

\begin{lemma}[The stabilizer in the marked middle quotient]
\leavevmode\label{lem:middlelattice531}
The homomorphism \(S'\to D\) has kernel \(S'\cap N'\).
The group \(D\) is finite. In each single marked factor its
reduction \(d\mapsto d_v\) is injective and only its identity
reduces to a unipotent matrix. Writing
\(\delta(d)=(d_v)_v\), one has
\[
       i(n\gamma n')=i(\gamma),\qquad
       i(\gamma h)=i(\gamma)\delta(d)
       \quad(n\in N,\ n'\in N',\ h\in S',\ h|_{C'}=d).
\]
\end{lemma}
\begin{proof}
Each member of \(S'\) preserves \(J'\), the vector \(\bar z'\)
and the lattice \(L'\). Since \(J'=E\bar z'\perp C'\), trivial
action on \(C'\) is equivalent to trivial action on \(J'\).
The pointwise flag condition gives the assertion about the kernel.
The sum of the Hermitian norms over all real places is positive
definite on the underlying real space of \(L'\). The image of
each vector of a fixed integral basis under a lattice isometry
lies on a compact sphere in a discrete lattice. This gives
finitely many possibilities. In particular \(D\) has a faithful
rational representation of dimension \(2fk\).

If \(g_i\) divided \(|D|\), Cauchy's theorem would give an
element of order \(g_i\). In characteristic zero its minimal
polynomial would contain the cyclotomic polynomial of degree
\(g_i-1>2fk\), since at least one eigenvalue would have that
prime order. Thus \(g_i\nmid |D|\). If \(d_v\) is unipotent,
its order is a power of \(g_i\) and divides \(|D|\), so it is
the identity. Lift to the localization of \(L'\) at the chosen
place of \(E\). If \(d^m=1\) and \(d\equiv1\) there, the sum
\(1+d+\cdots+d^{m-1}\) reduces to \(mI\) and is invertible
for \(m=|D|\). Multiplying it by \(d-1\) gives zero. Thus
\(d=1\) already on the rational lattice.

For the multiplication formulas, reduce the saturated isotropic
lattice and its orthogonal complement at \(v\). The latter reduces
to the orthogonal of the former. Indeed the pairing map onto the
dual of the saturated summand is surjective by self-duality, and
its exact sequence splits over the local integers. At a split
place this uses the perfect pairing between the two factors.
The norm of the marked last vector is a unit at \(v\), so taking
its perpendicular in the middle quotient commutes with reduction.
Admissibility identifies the two reduced isotropic flags up to an
upper unitriangular change and identifies their last vectors
modulo the isotropic summand. The chosen basis and its dual on
the remaining \(k\)-space therefore read the matrix of the
induced map by ordinary pairings. This is exactly the entry
definition in
\srcloc{Formalism/LocalToGlobal/SiegelLatticeTheory/ClosedLocalOrbits/TermOp}{300}.
On this quotient \(n,n'\) induce the identity, and \(h\) induces
\(d_v\). Composing the induced maps proves both formulas.
\end{proof}

\begin{theorem}[All lifts of a fixed geometric partner]
\leavevmode\label{thm:partnerlifts531}
Let \(\gamma_0\) be one lift of a geometric partner in the fixed
pair of rows. Its lifts are indexed bijectively by \(d\in D\)
and occupy precisely \(i(\gamma_0)\delta(D)\). Each single marked
place distinguishes them. For an array supported on this coset,
every row family \(aU\) of \(H\) recovers all its coefficients by
entry selection and multiplication by roots of unity. Its
normalized squared norm is exactly the sum of the squared norms
of those coefficients.
\end{theorem}
\begin{proof}
Each \(h\in S'\) fixes the transported frame of \(\gamma_0 h\)
exactly. Hence the transported frame satisfies every original
admissibility condition. If \(\gamma Y'\) has the same partner
as \(\gamma_0Y'\), write
\(\gamma Y'=n m\gamma_0Y'\) with \(n\in N\) and \(m\) in the
integral unipotent group of the transported flag. The latter
group is \(\gamma_0N'\gamma_0^{-1}\), by transport of the
lattice, the flag and its middle quotient. Thus
\(\gamma=n\gamma_0 n'h\) for some \(n'\in N'\) and \(h\in S'\).
The group \(S'\) normalizes \(N'\), since it preserves the
lattice and the defining filtration. Move \(n'\) to the right
of \(h\) and pass to the double quotient. Two choices of \(h\)
with the same \(d\) give the same lift by the kernel assertion.
Different \(d\)'s give different indices at each marked place
by Lemma~\ref{lem:middlelattice531}.

In a row family \(aU\) two indices coincide in one entry of
\(H\) precisely if their relative matrix lies in \(aUa^{-1}\).
For two distinct lifts the relative matrix at each place is a
conjugate of the reduction of a nonidentity element of \(D\).
It cannot be unipotent. Hence each coefficient occurs in its own
entry multiplied by its character value. This proves the integral
recovery and the norm identity, including for operator coefficients
with their Hilbert--Schmidt norm.
\end{proof}

\begin{proposition}[The signed orbital coefficients on a coset]
\leavevmode\label{prop:partnerweights531}
Let \(Z(\gamma)\) include the finite outer unipotent character,
the real orbital integral and the common volume factor in the
unfolded period. Let \(\rho'\) be the unitary source column type
on the compact group of \(C'\) at infinity. With the choices
of the preceding theorem the full operator satisfies
\[
              Z(\gamma_0 h)=Z(\gamma_0)\rho'(d).
\]
For finitely many geometric partners, sum their contributions at
each actual residue index and denote the resulting operator by
\(Z_i\). For every coset \(C=i_C\delta(D)\) in this sum there
is a single operator \(Z_{i_C}\) such that
\[
       Z_{i_C\delta(d)}=Z_{i_C}\rho'(d),\qquad
       \sum_{i\in C}\|Z_i\|_{\rm HS}^2=|D|\,\|Z_{i_C}\|_{\rm HS}^2 .
\]
The same conclusions hold for an absolutely summable expansion
grouped by partner. The assertions use the full source operator
before evaluation on any particular test vector.
\end{proposition}
\begin{proof}
Use representatives \(\gamma_0h\) with \(h\) fixing \(Y'\)
pointwise. In the source marked finite basis the matrix of \(h\)
fixes each reduced \(e_j\) and the reduced \(z_0\). They are the
images of its fixed rational vectors after the prescribed marking.
The finite outer character pairs the dual flag vector \(f_j\)
with the image of the next vector in the list
\((e_1,\ldots,e_s,z_0)\), sums over \(j\), divides by \(g_1g_2\)
and applies the finite additive character. Right multiplication by
the matrix of \(h\) preserves these pairings modulo \(g_1g_2\)
at every marked place. At the other finite places the difference
is integral and \(g_1g_2\) is a unit. The finite character is
trivial on the resulting integral differences, by its conductor
normalization. This gives the same finite factor for both terms.

The group determining the real domain is
\(N\cap\gamma_0N'\gamma_0^{-1}\). It is unchanged on replacing
\(\gamma_0\) by \(\gamma_0 h\), so take the same fundamental
domain in the two unfolded integrals. The real outer character
of \(h\) is one, since it fixes every next vector and each
corresponding simple off-diagonal coordinate is zero. The
equivariance in \eqref{eq:real127} consequently replaces
\(W'(\gamma_0^{-1}y)\) by
\(\rho'(d)^{-1}W'(\gamma_0^{-1}y)\). In the integrand
\(W(y)W'(\gamma_0^{-1}y)^*\) this multiplies the operator on
the right by \(\rho'(d)\), by unitarity. The quotient-integral
calculation is independent of the compatible representatives.
Indeed in the unfolded adelic integral left multiplication by an
integral target unipotent changes the variable by that
translation, and right multiplication by an integral source
unipotent gives the source equivariance factor. In either
case its finite and infinite character values on the rational
unipotent multiply to one, by triviality of the global additive
character on \(E\). In particular translating a fundamental
domain by an element of the intersection leaves its integrand
unchanged after combining the characters. The volume factor
depends only on the rows and words. This proves the operator
formula.

For each partner whose coset is \(C\), there is a unique lift
at \(i_C\). Add its operator there, with the sign in the original
sum. The same right action \(\rho'(d)\) transports each summand
to \(i_C\delta(d)\). Unitarity and summation give the second
formula. Absolute summability allows passage from finite sums
by continuity in the finite-dimensional operator space.
\end{proof}

For two distinct cosets \(C=i\delta(D)\), \(C'=j\delta(D)\)
write \(J(C,C')=\{v:i_vD_v\ne j_vD_v\}\), where
\(D_v\) is the image of \(D\) at \(v\), and let \(\mathcal V\)
denote the set of the \(2f\) places. For numbers \(0<t_v\le1\)
define
\[
 w_t(C,C')=
 \begin{cases}
 0,&J(C,C')=\varnothing,\\
 \prod_{v\in J(C,C')} t_v,&\varnothing\ne J(C,C')\ne\mathcal V,\\
 \min_{v_0\in\mathcal V}\prod_{v\ne v_0}t_v,&J(C,C')=\mathcal V.
 \end{cases}
\]
Set \(p_v=\rho_k(q_v)\) as in Lemma~\ref{lem:flagcollision530}
and \(c_v=(q_v^k-1)^{-1}\). Assume \(k\ge2\). For \(k=1\)
the groups \(U_v\) are trivial and the entries are read directly.

\begin{theorem}[Comparison for several partners in the same two rows]
\leavevmode\label{thm:cosetcomparison531}
For an array on any set \(\mathcal C\) of distinct cosets
of \(\delta(D)\), let \(Z_C\) denote its whole restriction
to \(C\), with the sum-of-squares norm. Then
\[
 \left|\langle Z,AZ\rangle-\sum_C\|Z_C\|^2\right|
 \le \sum_{C\ne C'} w_c(C,C')\|Z_C\|\,\|Z_{C'}\|.
\]
The sum counts ordered pairs. Put
\(\epsilon_p=\max_C\sum_{C'\ne C}w_p(C,C')\), with value zero
for an empty family. If \(\epsilon_p<1\), all raw coefficients
on the union can be recovered from \(HZ\) by integral entry
selections. Each index is isolated in at least a proportion
\(1-\epsilon_p\) of the row families. For any matrix of columns
on this support all minor ideals before and after \(H\) agree
over the coefficient ring containing the character values.
\end{theorem}
\begin{proof}
Fix a particular index \(i'\) in \(C\). In one factor and for
a fixed row family \(a_vU_v\), at most one \(d\in D\) can place
\(j_vd_v\) in the same row as \(i'_v\). Two such elements
would have unipotent relative matrix, contrary to
Lemma~\ref{lem:middlelattice531}. If the two local cosets
agree, this single choice is forced by \(j_vd_v=i'_v\) and
works for every row family there. If all local cosets agree
and these forced choices were consistent, the two global
cosets would meet and be equal. This proves the zero case.

When the local cosets agree at some places and differ at
others, any agreeing place thus specifies at most one possible
global partner. In a differing factor its relative matrix
is nonidentity. The fraction of the local row families giving a
collision is zero for a nonunipotent and at most \(p_v\) for a
nonidentity unipotent. Indeed \(a_vU_v\) induces a uniformly
distributed complete flag, and a unipotent preserving that flag
lies in its upper unitriangular group. The choices of local row
families are independent. If all local cosets differ, expose any
one row family first. Conditional on it there is at most one
candidate, to which the product bound at the other places applies.
Minimizing over the exposed place proves the bound \(w_p\) in
all cases. A union bound over \(C'\) now gives one isolating row
for every \(i'\) as soon as \(\epsilon_p<1\). Each chosen output
entry, multiplied by the inverse of its character unit, retrieves
the input. These selections define a left inverse independent of
the columns. Cauchy--Binet applied to \(H\) and to this inverse
gives both inclusions of minor ideals, also for fractional entries.

For the norm estimate the local normalized Gram entry has
absolute value at most the collision probability by the row
formula. On a nonidentity relative matrix it is also at most
\(c_v\). Indeed it vanishes off the unipotents, and for a
unipotent with fixed dimension \(l<k\) Theorem~\ref{thm:char125}
gives absolute value \(\prod_{j=l+1}^k(q_v^j-1)^{-1}\).
The global entry is the product of the local entries. When some
local cosets agree, the forced choices leave at most one nonzero
entry in the sum over \(j\delta(d)\), and the differing places
give the product of the \(c_v\). When all differ, keep the
collision-probability bound at one chosen place. Its sum over
\(d\in D\) is at most one since the collision events there are
disjoint. At the other places use the bounds \(c_v\). This proves
that each absolute row sum of the block \(A_{C,C'}\) is at most
\(w_c(C,C')\). The same holds for column sums by symmetry.
To bound the operator norm by this common bound, apply weighted
Cauchy--Schwarz to each row of a matrix \(B\). This gives
\(\|\sum_jB_{ij}v_j\|^2\le
 (\sum_j|B_{ij}|)\sum_j|B_{ij}|\|v_j\|^2\), and one sums over
\(i\). The diagonal blocks are identities by
Theorem~\ref{thm:partnerlifts531}. Summing the off-diagonal
blocks proves the estimate.
\end{proof}

\begin{corollary}[Uniform recovery on the period interval]
\leavevmode\label{cor:partnerinterval531}
On the common fields and intervals, suppose the arrays in each
ordered pair of rows occupy at most \(b_x^{4(k-1)}\) such cosets,
where \(b_x=\exp(x^\varepsilon)\). The integral comparison of
Theorem~\ref{thm:cosetcomparison531} then holds, in particular
at each terminal prime of Corollary~\ref{cor:whittakerterminal527}.
The normalized squared norm differs from the raw squared norm
by at most \(o(b_x^{-2})\) times the latter, uniformly in the
arrays and independently of \(|D|\).
\end{corollary}
\begin{proof}
Write \(q_*=\min(g_1,g_2)\ge(kb_x)^8\). There are at least
two marked factors. Every nonzero product in \(w_p\) contains
at least one factor at most \((k-1)q_*^{-(k-1)}\) by
Lemma~\ref{lem:flagcollision530}. The factors \(c_v\) satisfy
the same bound. Hence each of the two row sums of weights over
the other cosets is at most
\[
       b_x^{4(k-1)}(k-1)(kb_x)^{-8(k-1)}
            \le b_x^{-4(k-1)} .
\]
This is \(o(b_x^{-2})\) uniformly for \(k\ge2\). Apply the
theorem, using \(2ab\le a^2+b^2\) and symmetry for its norm
bound. Distinct ordered pairs of rows are separate coordinates
of \(H\), so their squared norms add and the integral selections
can be made in each pair independently.
\end{proof}

For a single geometric partner the comparison has no support
hypothesis. When several partners have the same coset, the operator
at its chosen base index is their transported signed sum. A
lower bound for the prescribed projected theta periods in
Section~\ref{sec:missing} still requires control of these sums,
their evaluation on the prescribed tests, and interference across
the occupied cosets in the full expansion.


\subsection{The two frame stabilizers in the signed period}
\label{sec:twostabilizers532}

Use the same two words, rows and markings as in
Section~\ref{sec:partnerlifts531}. Attach the finite groups
\(D_t,D_s\) of Lemma~\ref{lem:middlelattice531} to the target
and source respectively. Denote their embeddings in \(G\) by
\(\delta_t,\delta_s\). The group \(P=D_t\times D_s\) acts on
relative residue indices by
\[
            (a,b)i=\delta_t(a)i\delta_s(b)^{-1}.
\]
On a source-to-target operator \(B\) let
\(\tau(a,b)B=\rho_t(a)B\rho_s(b)^{-1}\), where \(\rho_t,\rho_s\)
are the two real column types. This is a unitary action on the
Hilbert--Schmidt space. We first bound the orders of both groups
in terms of the dimension over \(E\).

\begin{lemma}[A uniform bound for the column stabilizers]
\leavevmode\label{lem:conductorstabilizer532}
Suppose \(E\subseteq\QQ(\zeta_N)\), \(4\mid N\), and
\(D_0\) is any finite subgroup of \(\operatorname{GL}_k(E)\).
Then, with an absolute implied constant,
\[
           \log|D_0|\le k\log N+O(k\log^2(2k)).
\]
On the polynomial-conductor intervals of Theorem~\ref{thm:conductor232}
this gives \(\log|D_0|=O_{\psi,\varepsilon}(k\log^2(2x))\).
The same estimate for \(A_*\) applies in Theorems~\ref{thm:norm211}
and~\ref{thm:norm212}.
\end{lemma}
\begin{proof}
Use reduction at rational primes that split completely in
\(\QQ(\zeta_N)\), following Minkowski's reduction method
\cite{XVIII-XIII-IV-serre211}. Outside a finite set all matrices in \(D_0\)
and their inverses are integral, and their reductions remain
distinct. For the last assertion choose a nonzero entry of each
difference and exclude all rational primes below the prime
divisors of these entries. At any remaining split prime \(l\),
\[
          |D_0|\ \bigm|\ |\operatorname{GL}_k(\mathbb F_l)|
            =l^{k(k-1)/2}\prod_{j=1}^k(l^j-1).
\]
Fix an odd prime \(p\nmid N\). Choose a generator modulo \(p\)
and lift it modulo \(p^2\) so that its \((p-1)\)-st power is
different from one there. Among its \(p\) lifts precisely one
has power one, by binomial expansion and the unit derivative.
Impose this residue on \(l\), together with \(l\equiv1\bmod N\).
The Chinese remainder theorem and Dirichlet's theorem on primes
in arithmetic progressions supply such an \(l\) outside the
finite excluded set. Then \(p\mid l^j-1\) precisely for
\(j=(p-1)h\), and \(v_p(l^j-1)=1+v_p(h)\). Indeed raising a
unit \(1+p^m u\) to a power prime to \(p\) preserves \(m\),
and raising to the \(p\)-th power increases it by one, by
binomial expansion. Thus
\[
       v_p(|D_0|)\le\sum_{a\ge0}
               \left\lfloor\frac{k}{(p-1)p^a}\right\rfloor.
\]
For \(p\mid N\) put \(s=v_p(N)\). Choose
\(l\equiv1+p^s\bmod p^{s+1}\) and \(l\equiv1\bmod N/p^s\)
by the same argument. This splits \(l\) completely.
Now \(v_p(l^j-1)=s+v_p(j)\), also for \(p=2\) since \(s\ge2\).
The bound is \(ks+v_p(k!)\). The contributions at all \(p\mid N\)
are therefore at most \(k\log N+\log(k!)\). At the other primes
only \(p\le k+1\) contribute, by the preceding floor sum.
Their total is at most
\(\sum_{p\le k+1} kp\log p/(p-1)^2=O(k\log^2(2k))\),
by extending the sum of \(\log p/p\) to all integers in that
range. This proves the first estimate. On the stated intervals
\(N\) and \(k\) are bounded by fixed powers of \(x\).
Each coarse column stabilizer acts faithfully on its definite
\(k\)-dimensional \(E\)-space, so the same argument applies to
the maximum of the orders of these stabilizers.
\end{proof}

For a single lifted term choose a representative \(\gamma\) and
write \(i=i(\gamma)\). On the double quotient of lifted terms
use \((a,b)[\gamma]=[\tilde a\gamma\tilde b^{-1}]\), with tildes
denoting pointwise integral frame stabilizers. Define two
subgroups of \(P\) by
\[
  K_\gamma=\operatorname{Stab}_P([\gamma]),\qquad
  M_i=\operatorname{Stab}_P(i).
\]
For a subgroup \(M\) of \(P\), let
\(\Pi_M B=|M|^{-1}\sum_{m\in M}\tau(m)B\).

\begin{theorem}[The signed sum above one residue in a two-sided orbit]
\leavevmode\label{thm:orbitaverage532}
The action on the lifted terms is well defined and the residue map
is equivariant. Write \(B_\gamma\) for the single orbital operator
of Proposition~\ref{prop:partnerweights531}, and sum all distinct
terms of \(P\cdot[\gamma]\) at their actual residue indices.
Then the entries and the squared norm of the resulting array satisfy
\begin{align*}
    Z_{p\cdot i}&=[M_i:K_\gamma]\,\tau(p)\Pi_{M_i}B_\gamma
                                  &&(p\in P),\\
    \sum_{j\in P\cdot i}\|Z_j\|_{\rm HS}^2
       &=[P:M_i][M_i:K_\gamma]^2\|\Pi_{M_i}B_\gamma\|_{\rm HS}^2 .
\end{align*}
The range of \(\Pi_{M_i}\)
is the space of operators intertwining the two restrictions to
\(M_i\). These formulas also give the contributions of any
absolutely summable collection of such orbits to a common residue
by addition at that residue.
\end{theorem}
\begin{proof}
The pointwise integral stabilizer of either frame normalizes its
unipotent group. Two lifts of \(a\) or \(b\) differ by this
unipotent group, by Lemma~\ref{lem:middlelattice531}. Thus the
action descends to the lifted terms. Its right half preserves
the transported frame. For the left half, \(\tilde a\) preserves
the target lattice and fixes the target frame exactly. Applying
it to an admissibility congruence preserves the congruence to
the same upper unitriangular combination of target vectors.
It carries the transported frame's saturated isotropic lattice,
the specified fractional-ideal multiples of its vectors and its
integral primitive terminal class to the corresponding data for
the new frame. This preserves the remaining row conditions.
The induced maps on the reduced middle quotient compose, proving
equivariance and \(K_\gamma\subseteq M_i\).

For the left action, the reduced images under \(\gamma\) of the
successive source flag vectors, including the last vector, are
linear combinations of the fixed target frame vectors. Hence
\(\tilde a\) fixes these images at every marked place. In the
formula for the finite outer character the simple-coordinate
pairings are therefore unchanged modulo \(g_1g_2\). At the
other finite places the differences after division are integral.
The real domain for \(\tilde a\gamma\) is obtained from the one
for \(\gamma\) by left translation by \(\tilde a\), since the
intersection subgroup is conjugated by \(\tilde a\). Make this
change of variable. The target real vector then acquires the
factor \(\rho_t(a)\) by \eqref{eq:real127}, its outer character
being one on the pointwise stabilizer. Together with
Proposition~\ref{prop:partnerweights531} on the source this
gives \(B_{p[\gamma]}=\tau(p)B_\gamma\).
In particular \(B_\gamma\) is fixed by \(K_\gamma\).
The source actions are also displayed in
\srcloc{Formalism/LocalToGlobal/SiegelLatticeTheory/ClosedLocalOrbits/GaugeLatticeShifts/SatGramActR}{58}.

The orbit is \(P/K_\gamma\), its residue set is \(P/M_i\),
and the fibre above \(i\) is \(M_i/K_\gamma\). When summing
over all elements of \(M_i\) each term in that fibre is counted
\(|K_\gamma|\) times. This gives the first identity at \(p=1\)
and equivariance gives it at every index. The group average is
self-adjoint since inversion permutes the summands. It fixes
every invariant operator and has invariant image by reindexing.
It is consequently the orthogonal projection to the operators
satisfying \(\rho_t(a)B=B\rho_s(b)\) for every \((a,b)\in M_i\).
All entries have equal Hilbert--Schmidt norm, giving the second
identity. Absolute convergence permits the specified regrouping.
\end{proof}

\begin{lemma}[Equality of the compact characters on the residue stabilizer]
\leavevmode\label{lem:matchedcompact532}
At every archimedean embedding the defining \(k\)-dimensional
representations of \(M_i\) on the two positive complements are
unitarily equivalent. In particular if the two column types have
the same highest weight at each real place, there is a unitary
intertwiner \(U\) of those types on \(M_i\). For any such \(U\)
and every operator \(B\) between the two column types,
\[
       \|\Pi_{M_i}B\|_{\rm HS}^2
           \ \ge\ \frac{|\operatorname{Tr}(U^*B)|^2}{\dim V_s}.
\]
\end{lemma}
\begin{proof}
For \((a,b)\in M_i\) the two finite matrices satisfy
\(a_v=i_v b_v i_v^{-1}\) at every marked place. Their rational
traces in the defining spaces belong to \(\mathcal O_E\) since
they preserve lattices. If their difference \(z\) were nonzero,
each of the \(f\) chosen primes of \(E\) over \(g_1\) would
divide \(z\). All have norm \(g_1\), so that
\[
                g_1^f\le |N_{E/\QQ}(z)|\le(2k)^{2f}.
\]
The upper bound follows because all eigenvalues of these finite
order transformations have modulus one at every embedding.
On our intervals \(g_1>(2k)^2\), a contradiction. Hence the
traces agree for every element of \(M_i\) under each embedding.
Character orthogonality gives equal multiplicities for all
irreducible complex representations of this finite group.
Taking orthonormal bases in their unitary decompositions gives
a unitary intertwiner on the defining spaces. Conjugation by
such a unitary identifies the two ambient compact unitary groups.
Irreducible representations of the same highest weight are
equivalent under this identification, which is independent up
to inner automorphism of the choices of orthonormal basis.
Take the tensor product of the local equivalences for the column
types. The resulting \(U\) lies in the range of the orthogonal
projection \(\Pi_{M_i}\) and \(\|U\|_{\rm HS}^2=\dim V_s\).
Taking its inner product with \(B\) and applying Cauchy--Schwarz
proves the inequality.
\end{proof}

\begin{theorem}[Integral comparison on full two-sided residue orbits]
\leavevmode\label{thm:doubleorbit532}
Use the period fields of Theorem~\ref{thm:conductor232} and
the two marked primes of Corollary~\ref{cor:whittakerterminal527}.
Fix \(0<\lambda<8\) and \(k\ge2\). On the union of at most
\(b_x^{(8-\lambda)(k-1)}\) two-sided residue orbits in one
ordered pair of lattice rows, \(H\) has an integral left inverse.
For every array \(Z\) on that union one has, uniformly,
\[
  \left|\frac{\|HZ\|^2}{[G:U]}-\|Z\|^2\right|
       \le b_x^{-(\lambda-o(1))(k-1)}\|Z\|^2 .
\]
For a single full orbit the exponent may be taken to be
\(-(8-o(1))(k-1)\). In particular this compares the full
signed array of Theorem~\ref{thm:orbitaverage532} with its
explicit projected operator. At the terminal prime \(T\) the
averaging and the factor \([M_i:K_\gamma]\) in that formula
have unit denominators and the latter factor is a unit.
\end{theorem}
\begin{proof}
The orbit of \(i\) is \(\delta_t(D_t)i\delta_s(D_s)\), a union
of at most \(|D_t|\) left cosets of \(\delta_s(D_s)\).
Lemma~\ref{lem:conductorstabilizer532} gives
\(|D_t|\le b_x^{o(k-1)}\) uniformly, since
\(k/(k-1)\le2\) and \(\log^2(2x)/x^\varepsilon\to0\).
Every off-diagonal weight of Theorem~\ref{thm:cosetcomparison531}
is at most \((k-1)(kb_x)^{-8(k-1)}\), for both its unsigned and
its signed estimate. Multiply by the number of cosets just found.
The maximum sum of off-diagonal weights is at most the relative
error in the statement, which is less than one for large \(x\).
Entry selection gives the integral left inverse. The quadratic
estimate follows by the inequality \(2uv\le u^2+v^2\) on the
off-diagonal blocks. For one orbit use only \(|D_t|\) in the
count. Both arithmetic groups embed in \(G\) and \(T\nmid|G|\),
so all subgroup orders in the average and the orbit indices are
units at \(T\).
\end{proof}

More explicitly, take several columns of orbital operators on a
single \(P\)-orbit of terms with the same equivariance law.
At a place over \(T\) where these operators are defined over a
coefficient field, use a \(P\)-stable operator lattice. The
minor ideals of the stacked columns in \(HZ\) then equal those
of the base columns \(\Pi_{M_i}B_\gamma\), in local lattice
coordinates. Indeed the left inverse in the theorem first
recovers all residue blocks. Taking the block at \(i\) and
dividing by the unit \([M_i:K_\gamma]\) recovers the base
columns. Conversely each residue block is an integral translate
of the base columns multiplied by this unit. Cauchy--Binet in
both directions proves the equality. A stable lattice can be
formed by summing the finitely many translates of a lattice in
any field of definition of the action.

The formulas determine exactly which component of each orbital
operator contributes at the fixed marking. The trace in
Lemma~\ref{lem:matchedcompact532} depends on the real functions
and on the rational partner. Its size, and the signed sum of
these components over distinct rational orbits at a prescribed
residue, still enter the global theta-period estimate.


\subsection{The residue of the common definite quotient}
\label{sec:fullrankresidue533}

Retain a fixed pair of lattice rows and the marking of
Section~\ref{sec:partnerlifts531}. Assume \(k\ge2\) and put
\(Q_0=g_1g_2\). For a
lift \(\gamma\), write \(Y_0,Y_1\) for the spans of the target
frame and the transported source frame in the target space, and
\(X_0,X_1\) for their isotropic parts. In this section a
\emph{full-rank cell} is one for which \(Y_0\cap Y_1\) contains
a totally negative vector. The term refers to the dimension of
the common positive quotient. In the notation
\[
 I=X_0\cap X_1,\qquad A=(Y_0+Y_1)^\perp,\qquad
 P_j=Y_j^\perp/X_j ,
\]
the maps \(\pi_j:A/I\to P_j\) are then rational isometries onto
spaces of dimension \(k\), by Theorem~\ref{thm:cross132}.
Let \(P_s,P_t=P_0\) be the positive spaces of the two fixed rows.
The rational isometry attached to the cell is
\[
          j_\gamma=\pi_0\pi_1^{-1}\bar\gamma:P_s\longrightarrow P_t ,
\]
where \(\bar\gamma:P_s\to P_1\) is induced by \(\gamma\).
Passage to the lifted double quotient preserves this map,
since the unipotent groups act trivially on the two endpoint
quotients and carry \(A,I\) to the corresponding subspaces.

\begin{lemma}[The comparison map at the marked places]
\leavevmode\label{lem:markedtransport533}
Put \(U_j=(\Lambda\cap X_j)/(\Lambda\cap I)\). Let
\(\mathfrak a_\gamma\) be the annihilator of the cokernel of
the integral pairing map
\[
       U_1\longrightarrow U_0^\sharp,\qquad
            x\longmapsto\bigl(y\longmapsto h(y,x)/Q_0\bigr),
\]
where the dual consists of conjugate-linear integral functionals.
For the zero-rank map set \(\mathfrak a_\gamma=\mathcal O_E\).
If \(L_s,L_t\) are the positive column lattices, one has
\[
  \mathfrak a_\gamma j_\gamma L_s\subseteq L_t,\qquad
  \overline{\mathfrak a_\gamma}j_\gamma^{-1}L_t\subseteq L_s .
\]
At every \(v\mid Q_0\) prime to
\(\mathfrak b_\gamma=N_{E/F}(\mathfrak a_\gamma)\),
\(j_\gamma\) induces an isomorphism of local lattices and its
matrix in the marked residue bases equals \(i(\gamma)_v\).
\end{lemma}
\begin{proof}
The common negative vector \(z_*\) gives \(Y_j=X_j+Ez_*\).
The pairing between \(X_0,X_1\) has kernel \(I\) on each side.
Indeed an isotropic vector in \(X_1\) orthogonal to \(X_0\)
is also orthogonal to \(z_*\). Its image in the positive
quotient \(P_0\) must vanish. Interchanging the indices gives
the other kernel. Thus the displayed map is nonsingular over
\(E\). The two isotropic lattices have the same reduction in
the target row at the marked places. A vector of \(\Lambda\cap
Y_1^\perp\) therefore pairs with \(\Lambda\cap X_0\) in
\(Q_0\mathcal O_E\), by reducing in both split factors.
Away from \(Q_0\) this follows from integrality. In particular
the displayed pairing map is integral.

Take \(w\in\Lambda\cap Y_1^\perp\) and
\(\beta\in\mathfrak a_\gamma\). The functional
\(y\mapsto h(y,\beta w)/Q_0\) on \(U_0\) belongs to the image.
Hence some \(u\in\Lambda\cap X_1\) makes \(\beta w-u\)
orthogonal to \(X_0\), and it is already perpendicular to
\(X_1\) and \(z_*\). Thus \(\beta w-u\in\Lambda\cap A\).
Its image in \(P_1\) is \(\beta[w]\) and the same vector now
represents its image in \(P_0\). This proves the first
denominator claim after transport by \(\gamma\). The reversed
pairing map is the conjugate transpose in dual local bases.
Transposition preserves the elementary divisors, and conjugation
exchanges the conjugate primes. The annihilator of its cokernel
is consequently \(\overline{\mathfrak a_\gamma}\), proving the
inverse assertion. This is the local argument of
Proposition~\ref{prop:denom235} in both directions.

At a marked \(v\) prime to \(\mathfrak b_\gamma\), take
\(\beta=1\) locally in each split factor. The vectors \(w\)
and \(w-u\) above define the same class in the common residual
middle quotient, since \(u\) reduces into the common isotropic
space. The projections of the same integral vector \(w-u\)
therefore identify the positives at both ends by the identity
in this common quotient. The negative line splits off modulo
\(v\) since its norm is a unit. Composing with \(\bar\gamma\)
now gives exactly the matrix read from the marked basis and
its dual in Lemma~\ref{lem:middlelattice531}.
\end{proof}

Define the positive real number
\[
       h_\gamma=(N_{E/\QQ}\mathfrak a_\gamma)^{1/(2f)} .
\]
It can be read from the largest local elementary divisors of
the pairing in the lemma, including when its ideal is
nonprincipal.

\begin{theorem}[Separation of distinct rational transports]
\leavevmode\label{thm:fullrankseparation533}
Consider two full-rank cells \(\gamma_1,\gamma_2\) in the
fixed rows. If their indices share any output entry of \(H\),
then either \(j_{\gamma_1}=j_{\gamma_2}\), or
\[
                  Q_0\le 4k(h_{\gamma_1}h_{\gamma_2})^4 .
\]
If their indices agree, the alternative for unequal transports
is \(Q_0^2\le4k(h_{\gamma_1}h_{\gamma_2})^6\).
In particular, in any collection with \(h_\gamma\le b_x^2\)
the projected arrays belonging to distinct rational transports
occupy disjoint output entries in every row family. There is
no restriction on the number of cells in this assertion.
\end{theorem}
\begin{proof}
Write \(j_l=j_{\gamma_l}\) and let \(u=j_2^{-1}j_1\) act on
the positive source space. Define
\(\mathfrak b=\mathfrak b_{\gamma_1}\mathfrak b_{\gamma_2}\)
in \(F\). Its extension to \(E\) is the product of both
\(\mathfrak a_{\gamma_l}\overline{\mathfrak a_{\gamma_l}}\),
by the ideal norm formula for a quadratic extension.
Lemma~\ref{lem:markedtransport533} then shows that this extended
ideal clears the denominators of both \(u\) and \(u^{-1}\) on
the source lattice. Consider the matrix trace
\[
       t=2k-\operatorname{tr}_E u-\operatorname{tr}_E u^{-1}
        =\operatorname{tr}_E((1-u)(1-u^{-1}))\in F .
\]
The two traces are conjugate because \(u\) is unitary. If
\(u\ne1\), at every embedding of \(F\) the number \(t\) is the
squared Hilbert--Schmidt norm of \(1-u\), positive and at most
\(4k\). Moreover \(\mathfrak b(t)\) is an integral ideal of \(F\).
Indeed the traces of integral lattice endomorphisms are locally integral,
and integrality of an element of \(F\) can be checked in \(E\).
Thus, putting \(B=N_{F/\QQ}\mathfrak b\), one has
\(N_{F/\QQ}(\mathfrak b(t))\le B(4k)^f\).

At a marked place prime to \(\mathfrak b\), a collision in a
row of \(H\) makes the reduction of \(u\) unipotent by
Lemma~\ref{lem:markedtransport533} and the row criterion of
Lemma~\ref{lem:whittakermatrix526}. Both traces reduce to \(k\),
giving \(v(t)\ge1\). If the indices agree, both \(1-u\)
and \(1-u^{-1}\) vanish modulo the chosen split prime. Their
product gives \(v(t)\ge2\). Write \(e=1\) or \(2\) respectively.
The product of norms of the marked primes dividing \(\mathfrak b\)
is at most \(B\). The product over all marked primes is \(Q_0^f\),
as both rational primes split completely. It follows that
\[
         (Q_0^f/B)^e\le N_{F/\QQ}(\mathfrak b(t))
                   \le B(4k)^f .
\]
Since \(B=(h_{\gamma_1}h_{\gamma_2})^{2f}\), this proves the
two bounds. For \(h_{\gamma_l}\le b_x^2\) their first right
side is at most \(4k b_x^{16}<Q_0\) for \(k\ge2\), using
\(g_i\ge(kb_x)^8\). Comparing each pair of cells proves the
disjointness, including after summing any coefficients at
repeated indices with the same transport.
\end{proof}

\begin{corollary}[The signed array separated by its rational isometry]
\leavevmode\label{cor:transportarray533}
For the cutoff \(h_\gamma\le b_x^2\), let \(Z^{(j)}\) be the
raw contribution of all cells with \(j_\gamma=j\). For arbitrary
coefficients in finite sums one has
\[
       \|HZ\|^2=\sum_j\|HZ^{(j)}\|^2 .
\]
Each \(HZ^{(j)}\) is read integrally by entry selection from
\(HZ\). If all the denominator ideals in the collection are
prime to \(Q_0\), every row family of \(H\) recovers the single
raw operator for each \(j\). The normalized form equals the
raw sum of squares and all corresponding minor ideals agree.
These conclusions extend to absolutely summable collections.
\end{corollary}
\begin{proof}
Apply disjointness in Theorem~\ref{thm:fullrankseparation533}
and set entries outside the designated output support equal
to zero. For denominator ideals prime to \(Q_0\), every cell
with the same \(j\) has the same input index at every marked
place by Lemma~\ref{lem:markedtransport533}. Such an index
appears in its own row with coefficient a character unit.
Selecting it and undoing the unit gives a left inverse.
Averaging the exact single-family norms and applying
Cauchy--Binet to the two integral maps gives the remaining
assertions. The sets of input and output coordinates are
finite, allowing passage to absolutely convergent sums.
\end{proof}

\begin{theorem}[Survival of the individual full-rank orbital operator]
\leavevmode\label{thm:orbitnoncancellation533}
For an absolutely convergent full-rank cell of identical words
suppose that
\[
                       Q_0^2>4k h_\gamma^6 .
\]
Then every \((a,b)\in M_{i(\gamma)}\) satisfies \(a j_\gamma
=j_\gamma b\) over \(E\). In Theorem~\ref{thm:orbitaverage532}
one has \(\Pi_{M_i}B_\gamma=B_\gamma\).
Thus a nonzero single-cell operator survives the signed sum over
its entire two-sided arithmetic orbit, followed by \(H\).
This applies on the common intervals for every \(h_\gamma\le b_x^5\).
At the terminal prime the individual operator is recovered from
the orbit array over any stable coefficient lattice on which the
operators are defined.
\end{theorem}
\begin{proof}
Write \(j=j_\gamma\) and take \(u=j^{-1}ajb^{-1}\) on \(P_s\).
The integral actions of \(a,b\) preserve their column lattices.
Thus in the trace argument of Theorem~\ref{thm:fullrankseparation533}
the single ideal \(\mathfrak b_\gamma\) clears the denominators
of \(u,u^{-1}\), since each composition uses \(j,j^{-1}\)
once. At every marked prime outside this ideal Lemma~\ref{lem:markedtransport533}
and the definition of \(M_i\) give \(u\equiv1\). The same
argument with \(e=2\) now has \(B=h_\gamma^{2f}\). For
\(u\ne1\) it would give \(Q_0^2\le4k h_\gamma^6\).

The archimedean cell intertwines the unitary actions on the
common \(A/I\), by Theorem~\ref{thm:common133}. Under its two
projections these are the full compact groups of \(P_s,P_t\)
identified by \(j\). In particular
\(\rho_t(a)B_\gamma=B_\gamma\rho_s(b)\) whenever \(aj=jb\).
The common finite phase and volume scalar preserve this equality.
Therefore the entire average fixes \(B_\gamma\). The raw
operator at \(i\) is exactly \([M_i:K_\gamma]B_\gamma\),
whose scalar factor is positive and a unit at the common terminal
prime. Theorem~\ref{thm:doubleorbit532} supplies the integral
left inverse for the full orbit. Finally
\(Q_0^2\ge(kb_x)^{32}>4k b_x^{30}\) for \(k\ge2\).
\end{proof}

This calculation permits signed weights and every lift in the
chosen arithmetic orbit. Distinct geometric orbits with the
same rational transport can contribute different scalars. Their
sum, the cells with common positive dimension \(k-1\), and the
large-denominator contribution are still present in the projected
theta periods of Section~\ref{sec:missing}.


\subsection{The marked projection on a codimension-two subspace}
\label{sec:codimtwo624}

The small-rank terms of Theorem~\ref{thm:markedranks623} include
collections of matrices with a common restriction to a subspace
of codimension two. On each such collection we determine the kernel
of the finite projection, the precision of integral recovery, and
the loss in its complex norm. All statements concern the finite
matrix actually used at the two marked primes.

First work at one place. Let \(V=\mathbb F_q^k\), \(k\ge2\), let
\(W\subset V\) have dimension \(k-2\), and let \(P\) be the
pointwise stabilizer of \(W\) in \(G=\mathrm{GL}(V)\). Put
\(P^1=P\cap\mathrm{SL}(V)\). Take \(U\) upper unitriangular in a
reference basis and \(\theta\) a character of \(U\) nontrivial on
every simple root subgroup, as in Section~\ref{sec:whit125}.
For a table \(z\) supported on \(P\) use zero extension to \(G\).
Its finite projection is the unnormalized matrix \(H\) of
Lemma~\ref{lem:whittakermatrix526}. Choose a base point \(b_e\in P\)
of each determinant \(e\in\mathbb F_q^\times\). Write
\[
   (Dz)_y=z_y-z_{b_{\det y}}\quad (y\in P,\ y\ne b_{\det y}).
\]

\begin{theorem}[The integral differences detected by one marked projection]
\leavevmode\label{thm:codimtwo624}
Over the integers of a number field containing the values of
\(\theta\), there are integral matrices \(A,B\) with
\[
               H|_P=A D,\qquad qD=B(H|_P).
\]
In characteristic zero the kernel on \(P\) consists exactly of
the functions constant on each determinant fiber. In particular
the indicator of \(P^1\) is a nonzero integral table killed by
the projection. Every two indices in its support differ by rank
at most two. The assertions also hold for a left and right
translate of \(P\), using its transported coordinates.
\end{theorem}
\begin{proof}
Fix a row family \(aU\) of \(H\), let \(F_\bullet\) be its flag,
and put \(U_F=aUa^{-1}\). A nonempty row entry on \(P\) is a
sum over \((U_F\cap P)y\) for \(y\in P\). Up to a constant root
of unity its coefficients are the restriction of the character
\(\theta_F=\theta(a^{-1}(\,\cdot\,)a)\). Choose a flag-adapted
basis in which a subset of the basis spans \(W\). To do this,
at each step where \(\dim(W\cap F_j)\) increases choose the
next basis vector there, and at the other steps extend the
previous basis in \(F_j\). Rescaling the vectors gives the
prescribed decorations. Two positions of the basis are outside
the chosen subset. At least one such position \(j\) is at least
two. The simple shear \(1+tE_{j-1,j}\), \(t\in\mathbb F_q\),
lies in \(U_F\cap P^1\), and \(\theta_F\) is nontrivial on it.
Nontriviality is unchanged by the chosen basis lifts, since
changing lifts conjugates by a unitriangular matrix. The sum
of this nontrivial additive character is zero. Partitioning
the entry into cosets of the shear subgroup kills every function
of the determinant. Subtracting \(z_{b_{\det y}}\) at \(y\)
and inserting zero at the base points therefore yields the
same output. This gives the integral matrix \(A\).

Next choose a complement \(C\) of \(W\), a line \(l\subset C\),
and flags with \(F_1=l,\ F_2=C\). Put \(S(C,l)\) equal to the
unitriangular group of \(C\) fixing \(l\) pointwise and acting
as one on \(C/l\), extended by the identity on \(W\). We have
\(U_F\cap P=S(C,l)\). Indeed such an element preserves \(C\)
and fixes \(W\), and every later flag space is the sum of \(C\)
and its intersection with \(W\). The group \(S(C,l)\) is
additively isomorphic to \(\mathbb F_q\). Varying the decorations
on the first two quotients of the flag realizes all its nontrivial
additive characters, since nonzero scalar multiples of one
nontrivial additive character give all the others. Thus for each
left orbit \(S(C,l)y\) the rows of \(H\) contain all its nontrivial
Fourier sums, up to known root-of-unity factors. For a parametrization
\(s_t\) of this group, Fourier orthogonality gives
\[
 q(z_y-z_{s_t y})=
   \sum_{\chi\ne1}(1-\overline{\chi(t)})
                     \sum_{u\in\mathbb F_q}\chi(u)z_{s_u y}.
\]
All coefficients on the right are integral.

We record a bound on the number of shears needed to join two
indices. Every element of \(P^1\) is a product of at most
\(B_0=13\) elements of groups \(S(C,l)\), allowing the
complement and the line to change. Fix one splitting \(W\oplus C\).
On \(C\) every matrix in \(\mathrm{SL}_2(\mathbb F_q)\) is a
product of at most nine upper or lower elementary shears.
One row addition makes the upper left entry nonzero if necessary,
and two more clear the lower left and upper right entries.
For unit \(u\) put \(w(u)=U(u)L(-u^{-1})U(u)\), using the upper
and lower elementary matrices. Then \(w(u)w(-1)\) equals
\(\operatorname{diag}(u,u^{-1})\), which accounts for the other
six factors.

The remaining part of \(P^1\) consists of the translations
\(t_X=\left(\begin{smallmatrix}1&X\\0&1\end{smallmatrix}\right)\),
\(X:C\to W\). A translation \(X=w\otimes\phi\) of rank one
is a product of two allowed shears. Choose \(0\ne v\in\ker\phi\)
and \(Y:C\to W\) with \(Yv=w\). With \(s=1+v\otimes\phi\)
on \(C\), the product \(t_Ys t_Y^{-1}s^{-1}\) is
\(t_{Y(1-s^{-1})}=t_X\). Its first shear is taken in the
complement \(t_Y C\). Every \(X\) is a sum of at most two
rank-one maps, by expanding on a basis of the two-dimensional
domain. This proves the bound of thirteen, including \(W=0\).

For two indices of equal determinant their left quotient belongs
to \(P^1\). Telescoping the Fourier difference identity along
a product of the shears just constructed expresses \(q\) times
their difference by integral combinations of rows of \(H\).
This proves the formula for \(B\) and the kernel assertion.
All indices in \(P\) agree on \(W\), giving the rank bound.
For translated support \(xPy\) the substitutions \(a\mapsto
x^{-1}a,\ h\mapsto h y^{-1}\) relabel the two systems of
representatives in the formula for \(H\). Changes back to
the chosen representatives multiply entries by units.
\end{proof}

\begin{proposition}[A uniform lower bound on the detected norm]
\leavevmode\label{prop:codimnorm624}
For Hilbert-space-valued coefficients on \(P\), let \(\Pi z\)
be the table of their means on the determinant fibers, using
counting measure. There is an absolute \(c=3/(8B_0^2)>0\) such
that
\[
    c\,\|z-\Pi z\|_2^2
       \le \frac{\|Hz\|_2^2}{[G:U]}
       \le \left(\prod_{j=1}^{k-1}(1-q^{-j})^{-1}\right)
                  \|z-\Pi z\|_2^2 .
\]
\end{proposition}
\begin{proof}
The kernel in Theorem~\ref{thm:codimtwo624} allows subtraction
of \(\Pi z\). The upper bound then follows from the upper
operator bound in Theorem~\ref{thm:whittakerspectrum527}.
For the lower bound take a uniform decorated flag. The probability
that \(F_2\) is complementary to \(W\) is
\[
   p=\frac{q^{2(k-2)}}{{k\brack2}_q}
     =\frac{(1-q^{-1})(1-q^{-2})}
                 {(1-q^{-k})(1-q^{1-k})}.
\]
Indeed each complement is the graph of one linear map from a
fixed complement to \(W\). Counting ordered independent pairs
in \(V\) and dividing by their number in a plane gives the
Gaussian binomial in the denominator. Conditional on this
event, the complement \(C\) and the line \(l\) in it are
uniform and, conditional on the ordinary flag, the restricted
character is uniform among the \(q-1\) nontrivial ones.

Write \(E_S\) for the average over the left translations of
\(S=S(C,l)\) on \(\ell^2(P)\). On each orbit \(Sy\) the sum
of the squares of all nontrivial Fourier sums is \(q\) times
the sum of squared distances to the orbit mean, by the ordinary
finite Fourier identity. For a fixed row family the output
coordinates meeting \(P\) index these orbits once each. Averaging
the squared row norms over all families and retaining only the
complements bounds the normalized square below by
\[
              p\,\frac{q}{q-1}\,
                      \mathbb E_{C,l}\|(1-E_S)z\|_2^2.
\]
For an element \(u\in S\) and its left translation \(L_u\),
\(\|L_uz-z\|\le2\|(1-E_S)z\|\), since \(L_uE_S=E_S\).
Every \(h\in P^1\) is a product of at most \(B_0\) such elements
by the preceding proof. Conjugate this product by a uniform
element of \(P\). For each factor its conjugated pair \((C,l)\)
is uniform because \(P\) is transitive on complements and their
lines. The triangle inequality for the product of translations
and Cauchy--Schwarz give
\[
  \mathbb E_{a\in P}\|L_{aha^{-1}}z-z\|_2^2
                 \le 4B_0^2\,\mathbb E_{C,l}\|(1-E_S)z\|_2^2 .
\]
Averaging \(h\) over \(P^1\) gives \(2\|z-\Pi z\|_2^2\) on
the left by expanding the square. Finally \(p\,q/(q-1)\ge
1-q^{-2}\ge3/4\). These inequalities prove the lower bound.
\end{proof}

\begin{corollary}[Differences at both auxiliary primes]
\leavevmode\label{cor:codimmarked624}
Fix one codimension-two source subspace and its prescribed
injective restriction at each place over \(Q=g_1g_2\). Consider
any matrix \(Z\) of signed coefficient tables supported there,
with a common choice of restrictions for its columns. These may
be the sums of the actual orbital coefficients at every
supported index. After translating each local support into
\(P_v\), write \(D_*=\bigotimes_vD_v\),
\(P_*=\bigotimes_v(1-\Pi_v)\), and \(d_*=[G:U]\) for the
product marked group. For each complex column \(z\),
\[
    c^{2f}\|P_*z\|_2^2
    \ \le\ d_*^{-1}\|Hz\|_2^2
    \ \le\ \prod_{v\mid Q}\prod_{j=1}^{k-1}(1-q_v^{-j})^{-1}
                                  \|P_*z\|_2^2 .
\]
For algebraic coefficients let \(I_h\) denote the fractional
ideal generated by the \(h\)-rowed minors. At every finite
place of the coefficient field,
\[
     Q^{fh} I_h(D_*Z)\ \subseteq\ I_h(HZ)\ \subseteq\ I_h(D_*Z).
\]
Thus the minor ideals agree away from the two marked rational
primes, including at the common terminal prime \(T\). After
division by \(h\), the logarithmic loss from the displayed
rational multiplier, averaged over the coefficient field, is at
most \(f\log Q=o(b/x^M)\) for every fixed \(M>0\) on the
original intervals.
\end{corollary}
\begin{proof}
An injective restriction extends to an isomorphism by completing
bases. All matrices with this restriction form a translate of
the pointwise stabilizer, and missing entries can be filled by
zero. The marked matrix is the tensor product of the local
matrices by its character-sum formula. The positive local
Gram inequalities in Proposition~\ref{prop:codimnorm624} tensor
in the same order, since the difference of two such products
telescopes as a sum of tensor products of positive factors.
Thus the bounds hold for arbitrary product-space vectors, also
with Hilbert-valued entries. The integral factorizations in
Theorem~\ref{thm:codimtwo624} tensor with multiplier
\(\prod_{v\mid Q}q_v=Q^f\). Taking exterior powers, or
expanding minors by Cauchy--Binet, proves both inclusions.
The logarithmic assertion follows from the norm of a rational
principal ideal, \(f\ll_\psi x^{2+8\varepsilon}\), and
\(\log Q=O(\log b)\).
\end{proof}

This result compares the finite projected table with its actual
mixed differences whenever the restriction on a codimension-two
space is prescribed. A determinant-constant contribution at
just one place can vanish even in comparison rank two throughout
that place. For the full period, contributions from different
restrictions and from different orbits at one index must still
be combined with their coefficients in Section~\ref{sec:missing}.


\subsection{The spectrum after fixing a codimension-two subspace}
\label{sec:markedfourier625}

We compute the normalized Gram operator on the support of
Section~\ref{sec:codimtwo624}. This improves its lower norm
bound to a factor tending to one over the entire marked field.
In particular the finite projection on the detected subspace
incurs a vanishing relative error even when the number of
archimedean places grows.

At one marked place write \(V=W\oplus C\),
\(\dim_{\mathbb F_q}C=2\), \(h=\dim W=k-2\).
The elements of the pointwise stabilizer \(P\) of \(W\) are
\[
 g(X,A)=\begin{pmatrix}1&X\\0&A\end{pmatrix},
 \qquad X\in N=\operatorname{Hom}(C,W),\quad A\in\mathrm{GL}(C).
\]
Let \(\mathcal G\) be the compression of \(H^*H/[G:U]\) to
functions supported on \(P\). At the rank-two group
\(\mathrm{GL}(C)\) write \(\mathcal G_2\) for its full normalized
finite Gram matrix. For a line \(l\subset C\) let \(U_l\)
be the subgroup fixing \(l\) and \(C/l\) pointwise.
It has order \(q\). Fix a nontrivial additive character
\(\psi\) of \(\mathbb F_q\), and put
\[
   \eta_k=\{(1-q^{-k})(1-q^{1-k})\}^{-1},
             \qquad p_0=(1-q^{-1})(1-q^{-2}).
\]

\begin{theorem}[Fourier decomposition of the restricted Gram operator]
\leavevmode\label{thm:markedfourier625}
For each \(D\in\operatorname{Hom}(W,C)\), the functions
\(z(g(X,A))=\psi(\operatorname{tr}(DX))\,v(A)\) form an
invariant space for \(\mathcal G\). These spaces are mutually
orthogonal and exhaust the arrays on \(P\). The operator induced
on \(v\) is as follows.
\begin{enumerate}
\item[\textup{(i)}] For \(D=0\) it is \(\eta_k p_0\mathcal G_2\).
\item[\textup{(ii)}] For \(\operatorname{rank}D=1\) put \(l=\operatorname{im}D\).
On each set \(A U_l\) it is \(\eta_k(I-q^{-2}\mathbf J_q)\),
where \(\mathbf J_q\) has every entry equal to one.
\item[\textup{(iii)}] For \(\operatorname{rank}D=2\) it is \(\eta_k I\).
\end{enumerate}
Only the ranks realizable in \(\operatorname{Hom}(W,C)\) occur.
\end{theorem}
\begin{proof}
The trace pairing between \(N\) and \(\operatorname{Hom}(W,C)\)
is perfect, as can be checked on rectangular matrix units.
Multiplication by a nonzero field element is surjective on
\(\mathbb F_q\), so the displayed characters are distinct.
Finite Fourier orthogonality proves the asserted decomposition
of the underlying vector space.

For two indices, their relative matrix in the target has blocks
\[
       g(X',A')g(X,A)^{-1}=g(Y,M),\qquad
                  Y=(X'-X)A^{-1},\quad M=A'A^{-1}.
\]
By Corollary~\ref{cor:content125} the Gram entry vanishes unless
\(M\) is unipotent. For unipotent \(M\) its value depends on
the rank of the stacked map \(v\mapsto(Yv,(M-1)v)\) from \(C\).
At ranks zero, one and two the values are respectively
\[
       1,\quad a_1=-\frac1{q^k-1},\quad
              a_2=\frac1{(q^k-1)(q^{k-1}-1)} .
\]
If rank two is absent, the formulas below only evaluate entries
at the ranks that occur. Substitution \(X'=X+YA\) in a Fourier
mode removes its common factor \(\psi(\operatorname{tr}(DX))\).
The remaining character on \(Y\) is that indexed by \(B=AD\).
This proves invariance of each Fourier mode.

Set \(N_0=q^h\). For a fixed line \(l'\subset C\), the maps
\(Y\) that kill \(l'\) form a group of order \(N_0\).
Its character sum at \(B\) is \(N_0\) if
\(\operatorname{im}B\subseteq l'\), and zero otherwise, by
the trace pairing on \(C/l'\). Every map of rank one is
counted once by its kernel line, and the zero map is counted
\(q+1\) times. Hence the Fourier sum of the indicator of
\(\operatorname{rank}Y\le1\) is \(N_0 n_B-q\), where \(n_B\)
is the number of lines containing \(\operatorname{im} B\).
For ranks zero, one and two of \(B\) it equals \(q+1,1,0\),
respectively.

At \(M=1\) the entry as a function of \(Y\) is
\(a_2+(a_1-a_2)\mathbf1_{\operatorname{rank}Y\le1}
+(1-a_1)\mathbf1_{Y=0}\). The diagonal coefficient in the
mode \(D\) is consequently
\[
 a_2 N_0^2\mathbf1_{D=0}
       +(a_1-a_2)(N_0 n_B-q)+(1-a_1).
\]
For \(M\ne1\) unipotent in dimension two, put
\(l'=\ker(M-1)\). The stacked map has rank one exactly when
\(Y\) kills \(l'\). Its Fourier sum is therefore
\[
    a_2 N_0^2\mathbf1_{D=0}
          +(a_1-a_2)N_0\,\mathbf1_{\operatorname{im}B\subseteq l'} .
\]
The three scalar simplifications are
\[
  1-a_1-q(a_1-a_2)=\eta_k,\qquad
  (a_1-a_2)N_0=-\eta_k/q^2,\qquad
  a_2N_0^2=\eta_k/q^3 .
\]
For rank two this gives the scalar operator in (iii).
For rank one, an off-diagonal coefficient occurs exactly
when \(M\) belongs to \(U_{Al}\setminus\{1\}\), or equivalently
when \(A'\in A U_l\setminus\{A\}\). Together with the diagonal
coefficient this gives (ii). In rank zero the diagonal is
\(\eta_k p_0\), and the value at each nonidentity unipotent is
\(-\eta_k p_0/(q^2-1)\). All remaining values are zero.
These are the entries in (i) by the same character formula
in dimension two.
\end{proof}

\begin{corollary}[Near-isometry on the determinant differences]
\leavevmode\label{cor:markednearisometry625}
Let \(\Pi\) denote the determinant-fiber average on \(P\).
In the order on Hermitian forms,
\[
     \eta_k(1-q^{-1})^2(I-\Pi)
       \ \le\ \mathcal G\ \le\ \eta_k(I-\Pi).
\]
On any product of the prescribed restriction cosets at the
two auxiliary primes use \(P_*\) of
Corollary~\ref{cor:codimmarked624}. The corresponding bounds
are \(a_*P_*\) and \(b_*P_*\), where
\[
      a_*=\prod_{v\mid Q}\eta_k(q_v)(1-q_v^{-1})^2,\qquad
      b_*=\prod_{v\mid Q}\eta_k(q_v).
\]
Here \(\eta_k(q_v)\) denotes the local value defined above.
Uniformly on the period intervals,
\[
       0\le1-a_*\ll f/q_*,
       \qquad 0\le b_*-1\ll f q_*^{1-k},
        \qquad q_*=\min(g_1,g_2).
\]
\end{corollary}
\begin{proof}
In the zero Fourier mode the determinant depends only on
\(A\in\mathrm{GL}_2\). By Theorem~\ref{thm:codimtwo624}
with \(k=2\) the kernel of \(\mathcal G_2\) is exactly the
space of functions of this determinant. On its orthogonal
complement the two bounds from
Theorem~\ref{thm:whittakerspectrum527} specialize to
\((1-q^{-1})/(1-q^{-2})\) and \((1-q^{-1})^{-1}\).
Multiplication by \(\eta_k p_0\) gives the lower coefficient
in the statement and the upper bound \(\eta_k(1-q^{-2})\)
on that complement. Each nonzero mode has mean zero on every
determinant fiber by summation over \(X\). In rank one its
two eigenvalues are \(\eta_k(1-q^{-1})\) and \(\eta_k\), by
the constant and zero-sum decomposition on \(A U_l\).
Rank two has the single eigenvalue \(\eta_k\). This proves
the local inequalities.

Tensor the positive Gram inequalities over the \(2f\) places,
as in Corollary~\ref{cor:codimmarked624}. Each lower factor is
at most one since \(k\ge2\), and is at least \((1-q_v^{-1})^2\).
Each upper factor has logarithm \(O(q_v^{1-k})\), by
\(-\log(1-t)\le 2t\) for \(0\le t\le1/2\).
The prime bound \(q_*\ge(kb)^8\) and polynomial growth of \(f\)
give the claimed estimates upon summing and exponentiating.
\end{proof}

\begin{proposition}[Effect on the Gram blocks of signed inputs]
\leavevmode\label{prop:markedgram625}
On each target row prescribe one of the product supports in
Corollary~\ref{cor:codimmarked624}. Allow a distinct support
on each target row. Write \(Z_i\) for any incoming column
operator to the direct sum of these tables, including their
original signed coefficients. Let \(Y_i=P_*Z_i\), using the
appropriate projection on each row, and
\(V_i=[G:U]^{-1/2}HZ_i\). Put
\(\delta=\max(1-a_*,b_*-1)\). Then
\[
  \|V_i^*V_j-Y_i^*Y_j\|_{\rm HS}
   \le\delta\min\{\|Y_i\|_{\rm op}\|Y_j\|_{\rm HS},
                   \|Y_i\|_{\rm HS}\|Y_j\|_{\rm op}\}.
\]
For the full matrix \(Y\) of these incoming columns and the
corresponding \(V\), the two ranks agree. On every subspace
of dimension \(u\) on which \(Y\) is injective the ratio of
the Gram determinants for \(V\) and \(Y\) lies between
\(a_*^u\) and \(b_*^u\).
\end{proposition}
\begin{proof}
On each target support let \(B\) be its compressed normalized
finite Gram operator. The preceding corollary implies
\(B=P_*BP_*\) and \(\|B-P_*\|_{\rm op}\le\delta\).
The same holds on the direct sum. Thus the mixed difference
is \(Y_i^*(B-P_*)Y_j\). Applying the operator norm on the
left and the Hilbert--Schmidt norm on the right proves the
first product bound. Taking adjoints proves the other.
The forms for the full maps satisfy
\(a_*Y^*Y\le V^*V\le b_*Y^*Y\), giving equality of kernels.
On the indicated subspace conjugate its restricted Gram form
for \(V\) by the inverse positive square root of the restricted
form for \(Y\). All eigenvalues lie in \([a_*,b_*]\).
Their product is the claimed determinant ratio.
\end{proof}

These estimates and the integral inclusions of
Corollary~\ref{cor:codimmarked624} apply on exactly the same
supports. The computation controls the contribution of the
finite projection to the norm there. The sizes of the signed
differences \(Y_i\) supplied by the arithmetic periods and
the cross terms between different restriction supports on one
target row still enter the comparison in Section~\ref{sec:missing}.


\subsection{Mixed periods on different restriction subspaces}
\label{sec:restrictionrecursion626}

Two signed inputs to one target row can have distinct restrictions at a
marked place. The dimension on which those restrictions agree controls
the size of the finite groups needed to compute their mixed inner product.
We first prove the reduction at one place. All inner products use
counting measure.

Write \(G_j=\mathrm{GL}_j(\mathbb F_q)\), with \(G_0\) the trivial
group. For \(0\le d\le k\) let \(P_{k,d}\subset G_k\) fix pointwise a
specified subspace \(W\) of dimension \(w=k-d\). Let \(\mathcal G_{k,d}\)
be the compression to \(P_{k,d}\) of the normalized Gram matrix of
Lemma~\ref{lem:whittakermatrix526}. Thus \(\mathcal G_{j,j}\) acts on
the full group \(G_j\), and \(\mathcal G_{0,0}=1\). Write
\[
 \eta_{k,d}=\prod_{i=k-d+1}^k(1-q^{-i})^{-1},
 \qquad p_e=\prod_{i=1}^e(1-q^{-i}).
\]
Empty products equal one. Translating the support by any fixed element
of \(G_k\) gives the same operator in the translated coordinates, since
its entries are values of a class function on relative matrices.

\begin{theorem}[Reduction to the dimension of the moving subspace]
\leavevmode\label{thm:restrictionrecursion626}
There is an orthogonal decomposition of \(\mathcal G_{k,d}\) into
blocks, each unitarily equivalent to
\[
                  \eta_{k,d}\,p_e\,\mathcal G_{e,e}
                    \qquad (0\le e\le d).
\]
The following rule specifies the blocks and the transform on the inputs.
Choose a complement \(C\) of \(W\). Fourier transform the coordinate
\(X\in\operatorname{Hom}(C,W)\) in the matrices
\(g(X,A)=\left(\begin{smallmatrix}1&X\\0&A\end{smallmatrix}\right)\).
For frequency \(D:W\to C\) put \(t=\operatorname{rank}D\).
Group the remaining indices \(A\in G_d\) by left cosets of the
pointwise stabilizer of \(\operatorname{im}D\).
For \(t=0\) stop with \(e=d\). For \(t>0\) continue inside each coset
with ambient dimension \(d\) and moving dimension \(d-t\).
A group of dimension zero contributes the identity. On the orthogonal
complement of the kernel the eigenvalues of the original operator lie in
\[
            [\,\eta_{k,d}(1-q^{-1})^d,\ \eta_{k,d}\,].
\]
\end{theorem}
\begin{proof}
Fix a nontrivial additive character \(\psi\) of the residue field.
The characters \(\psi(\operatorname{tr}(DX))\) give every Fourier mode
by the perfect pairing of rectangular matrices. For two indices put
\(M=A'A^{-1}\) and \(Y=(X'-X)A^{-1}\). The relative matrix is
\(g(Y,M)\). By Corollary~\ref{cor:content125}, its entry is zero
unless \(M\) is unipotent. Otherwise put
\[
    K=\ker(M-1)\cap\ker Y,\qquad
    a_k(r)=\frac{(-1)^r}{\prod_{i=k-r+1}^k(q^i-1)}.
\]
The entry is \(a_k(d-\dim K)\), since the relative matrix minus the
identity has rank \(d-\dim K\). Acting on one Fourier mode replaces
the sum over \(X'\) by the sum over \(Y\) with character indexed by
\(B=AD\), times the original character at \(X\).

For \(0\le j\le d\) set
\[
 c_j=a_k(d)(-1)^j q^{wj+j(j+1)/2}.
\]
For a subspace \(K\) of dimension \(s\) one has
\(a_k(d-s)=\sum_{L\subseteq K}c_{\dim L}\).
Indeed, writing \({s\brack j}_q\) for the number of \(j\)-dimensional
subspaces of an \(s\)-space, the sum divided by \(a_k(d)\) equals
\(\prod_{i=0}^{s-1}(1-q^{w+1+i})\).
Here we use the polynomial identity
\[
 \sum_{j=0}^{s}{s\brack j}_q(-1)^j q^{j(j-1)/2}z^j
     =\prod_{i=0}^{s-1}(1-zq^i).
\]
It follows by induction on \(s\), counting subspaces inside a fixed
hyperplane and those meeting it in codimension one. The latter number
with prescribed intersection of dimension \(j-1\) is \(q^{s-j}\).
Cancelling the factors of \(a_k(d)\) proves the asserted subspace sum.

For \(L\subseteq\ker(M-1)\) of dimension \(j\), the sum of the
character at \(B\) over maps \(Y\) killing \(L\) is \(q^{w(d-j)}\)
if \(\operatorname{im}B\subseteq L\), and zero otherwise. This follows
by applying the trace pairing to \(C/L\). Thus the transformed entry
can be nonzero only if \(M\) fixes \(\operatorname{im}B\) pointwise.
This is precisely \(A'\in A\operatorname{Stab}_{\rm pt}
(\operatorname{im}D)\). In that case let \(s=\dim\ker(M-1)\). Use
\[
 c_j q^{w(d-j)}
   =\eta_{k,d}(-1)^{d-j}q^{-d(d+1)/2+j(j+1)/2}.
\]
There are \({s-t\brack j-t}_q\) eligible \(L\) for each \(j\ge t\).
The same polynomial identity, now at \(z=q^{t+1}\), gives the sum
\[
 \eta_{k,d}(-1)^{d-s}q^{-d(d+1)/2+t(t+1)/2}
       \prod_{i=t+1}^{s}(q^i-1)
      =\frac{\eta_{k,d}}{\eta_{d,d-t}}\,a_d(d-s).
\]
For nonunipotent \(M\) both matrix entries are zero. Conjugating by
the representative of the coset identifies the result with
\(\eta_{k,d}\mathcal G_{d,d-t}/\eta_{d,d-t}\). Unitarity follows by
normalizing each Fourier transform by the inverse square root of the
size of its additive group. The successive factors cancel. At the
last step \(1/\eta_{e,e}=p_e\), proving the stated formula.

For \(e\ge1\), Theorem~\ref{thm:whittakerspectrum527} gives the
lower bound \((1-q^{-1})^e/p_e\) on the positive spectrum of
\(\mathcal G_{e,e}\). Proposition~\ref{prop:amplitudes126} gives the
upper bound \(\prod_{i=1}^{e-1}(1-q^{-i})^{-1}\).
Multiplying both bounds by \(\eta_{k,d}p_e\) and using \(e\le d\)
proves the inequalities. The scalar at \(e=0\) satisfies them as well.
\end{proof}

\begin{corollary}[A common transform for mixed signed periods]
\leavevmode\label{cor:mixedrestrictions626}
At each marked place \(v\) suppose the two tables are supported,
respectively, on \(g_{1,v}P_{W_{1,v}}\) and \(g_{2,v}P_{W_{2,v}}\).
Here \(g_{i,v}\) are invertible and \(P_W\) fixes \(W\) pointwise.
Put
\[
 W_v=(W_{1,v}\cap W_{2,v})\cap\ker(g_{1,v}-g_{2,v}),\qquad
 d_v=k-\dim W_v .
\]
Apply the theorem on the coset \(g_{1,v}P_{W_v}\), tensor over the
marked places, and let \(\mathcal F\) denote the resulting unitary
transform. On each final block apply the orthogonal projection onto
the generic part of each \(G_{e_v}\), as defined in
Theorem~\ref{thm:whittakerspectrum527}, with the identity for \(e_v=0\).
Transport the product of these projections back by \(\mathcal F\)
and take their direct sum, denoted by \(\Pi\) in the original coordinates.
For the two tables \(z,z'\), extended by zero to the common support,
\[
 \left|\frac{\langle Hz,Hz'\rangle}{[G:U]}
                -\langle\Pi z,\Pi z'\rangle\right|
 \ll \frac{\sum_v d_v}{q_*}\,\|\Pi z\|\,\|\Pi z'\|,
 \qquad q_*=\min(g_1,g_2).
\]
Here \(g_1,g_2\) in \(q_*\) are the two auxiliary primes. The exact
mixed product is obtained by substituting the blocks of the theorem
in place of the projections. If each \(W_{i,v}\) has codimension two,
then
\(\sum_v d_v\le 8f+\sum_v\operatorname{rank}(g_{1,v}-g_{2,v})\).
All statements hold for Hilbert-space-valued signed coefficients on
the original period intervals.
\end{corollary}
\begin{proof}
Both supports prescribe the map \(g_{1,v}|_{W_v}\) on \(W_v\).
Translate by \(g_{1,v}^{-1}\) and apply the theorem at each place.
By Lemma~\ref{lem:whittakermatrix526} the mixed product is exactly
the inner product for the tensor of these compressed Gram operators.
Its kernel is the orthogonal complement of the indicated surviving
blocks. Its positive eigenvalues are bounded below by
\(\prod_v(1-q_v^{-1})^{d_v}\) and above by
\(\prod_v\eta_{k,d_v}(q_v)\). The logarithms of these bounds are
\(O(\sum_v d_v/q_*)\) in absolute value. This tends to zero uniformly
since \(\sum_v d_v\le 2fk\), \(q_*\ge(kb)^8\), \(f,k\) are
polynomially bounded in \(x\), and \(b\) grows faster than every
power of \(x\). Subtract the identity on the surviving blocks and
use Cauchy--Schwarz. Tensoring with the identity on the coefficient
space gives the vector-valued assertion. Finally codimension is
subadditive under intersection, and imposing the equality of the two
maps costs at most the rank of their difference. Summing over the
\(2f\) places gives the dimension estimate.
\end{proof}

The same construction applies to any collection on a single target
row by taking the common subspace of agreement of all its prescribed
maps. It retains every sign before the last projections. The size of
the resulting generic components of the arithmetic tables still
enters the period comparison in Section~\ref{sec:missing}.


\subsection{Integral matrices for mixed periods at the terminal prime}
\label{sec:integralrestrictions627}

The unitary transform in Section~\ref{sec:restrictionrecursion626}
computes mixed inner products. Integral periods require the matrix
\(H\) itself, before division by \([G:U]\). We give a transform
for its rows using only character values and Fourier sums. It
compares the ideals of minors even at a prime dividing \([G:U]\),
provided that it differs from the residue characteristic of the mark.

At one place use \(W,C,P_{k,d}\) as in
Theorem~\ref{thm:restrictionrecursion626}, with \(w=k-d\).
Let \(R\) be the ring of integers in a number field containing the
additive character values. Let \(\Gamma_j\) denote the induced
representation \(R[G_j]\otimes_{R[U_j]} R_\theta\), where
\(\theta\) is the nondegenerate upper unitriangular character.
Use the same additive character in all dimensions, and set
\(\Gamma_0=R\). In its coset basis the matrix
\(\sum_g z_g\Gamma_j(g)\), with its entries listed as rows of a
linear map in \(z\), is \(H_jz\). Indeed its entry at the basis
vectors \(h^{-1}U_j,aU_j\), in source and target respectively,
is the sum of \(\theta(u)z_{auh}\) over \(u\in U_j\).

\begin{theorem}[Integral reduction at a single mark]
\leavevmode\label{thm:integralrestrictions627}
For arrays on \(P_{k,d}\) define the following collection of
smaller arrays. Take the orbits of \(G_d\) acting by
postcomposition on the maps \(B:W\to C\). In each orbit choose
\(B\), list its members \(B_i=a_iB\) with \(a_i\in G_d\), and
put \(S_B=\{s\in G_d:sB=B\}\). For every ordered pair \(i,j\) form
\[
   z_{ij}(s)=\sum_{X:C\to W}
      \psi\bigl(\operatorname{tr}(B_jX)\bigr)
      z\bigl(g(X,a_i s a_j^{-1})\bigr),\qquad s\in S_B .
\]
For \(B=0\) retain \(H_d z_{11}\). Otherwise repeat the operation
on \(S_B\), the pointwise stabilizer of a \(t\)-space in dimension
\(d\), where \(t=\operatorname{rank} B\). A terminal trivial
group gives its coefficient itself. Let \(L_{k,d}z\) be the stack
of all the retained rows. Over \(R[1/q]\) the row modules of
\(H_k|_{P_{k,d}}\) and \(L_{k,d}\) coincide inside the dual of
the same coefficient lattice. In particular there are matrices
over this ring expressing either set of rows in terms of the other.
The assertion holds after left translation of the support.
\end{theorem}
\begin{proof}
Work first over the completion \(O\) of the coefficient integers
at a prime of residue characteristic different from that of
\(\mathbb F_q\). We may extend it to the integers of a finite
splitting extension for the finite groups under consideration,
in both the ordinary and the residual characteristic. Descent
will be justified at the end. Write \(n(X)=g(X,1)\),
\(l(A)=g(0,A)\), so \(g(X,A)=l(A)n(X)\).
Since \(q\) is a unit, the additive group of the \(n(X)\) has
orthogonal character idempotents
\[
 e_B=q^{-wd}\sum_X \psi(-\operatorname{tr}(BX))\,n(X),
                      \qquad \sum_B e_B=1 .
\]
Conjugation by \(l(A)\) sends \(e_B\) to \(e_{AB}\). Identify all
the summands \(e_{B_i}\Gamma_k\) in one orbit with \(e_B\Gamma_k\)
using \(l(a_i)\). These summands are free over the discrete valuation
ring as direct summands of a free lattice. The block with source
\(j\) and target \(i\) in the matrix of the group algebra element
\(\sum_g z_g g\) is now precisely
\[
                  \sum_{s\in S_B} z_{ij}(s)\,l(s)|_{e_B\Gamma_k}.
\]
Indeed \(n(X)\) acts on the source by its character at \(B_j\),
and \(l(A)\) carries this source to \(i\) exactly when
\(A=a_i s a_j^{-1}\). Other orbits give separate blocks.

We identify \(e_B\Gamma_k\) as an \(O[S_B]\)-module.
In characteristic zero its character at \(s\in S_B\) is the trace
of \(e_Bl(s)\) on \(\Gamma_k\). In the expression for \(e_B\),
\(n(X)l(s)=g(Xs,s)\) and \(sB=B\). Thus replacing \(Xs\) by
the summation variable leaves its character factor unchanged.
The Fourier sum computed in Theorem~\ref{thm:restrictionrecursion626},
at frequency \(-B\), gives
\[
 \operatorname{tr}(e_Bl(s)\mid\Gamma_k)
  =\frac{\gamma_k(1)}{q^{wd}\gamma_d(1)}
           \frac{\eta_{k,d}}{\eta_{d,d-t}}\,
                    \operatorname{tr}(s\mid\Gamma_d),
 \qquad
 \gamma_j(1)=[G_j:U_j]=\prod_{a=1}^j(q^a-1).
\]
The rank and image used in that sum are unchanged by the sign.
The expression is valid at nonunipotent elements as both
character values vanish there. The identity for the degree counts
ordered bases and divides by \(|U_j|=q^{j(j-1)/2}\).
Cancelling the factors in the ratio gives the positive integer
\[
              m=q^{t(t+1)/2}\prod_{a=t+1}^{w}(q^a-1).
\]
Here \(t\le w\) because \(B\) has domain \(W\).

Both representations being compared are projective over
\(O[S_B]\). In fact \(\Gamma_j\) is induced from a rank-one
summand of the regular module of \(U_j\), using the character
average over the group of unit order in \(O\). Induction
therefore makes it projective over \(O[G_j]\), and restriction
to a subgroup preserves projectivity because the group algebra
is free over the subgroup algebra. Apply this to \(S_B\subseteq G_d\)
and to \(l(S_B)\subseteq G_k\). The idempotent \(e_B\) commutes
with \(l(S_B)\), so its image remains projective.

We use here the usual character theorem for projective modules
over a splitting modular system. The reductions of indecomposable
projective lattices give the indecomposable projective modules
over the residue field, each with a unique projective lift.
The ordinary characters of these lifts are linearly independent.
Consequently equality of ordinary characters determines all the
multiplicities in a direct sum of projectives. Applying this to
the character equality just proved yields the integral isomorphism
\[
           e_B\Gamma_k \simeq
                (\Gamma_d|_{S_B})^{\oplus m}.
\]
These statements about projectives are Proposition~9.4.5 and
Corollary~10.2.4 of Webb.~\cite{XVIII-XIII-projectivecharacters627}

Choose \(O\)-bases using this isomorphism in each orbit.
Vectorizing an integral change of basis on source and target
gives an invertible coordinate change on matrix entries.
The transformed output consists of zeros and repeated copies
of the entries of \(\sum_s z_{ij}(s)\Gamma_d(s)\), with at least
one copy for every \(i,j\). Row span is unchanged on removing
these zeros and repetitions. This proves the recursion over the
splitting completion. When \(t=0\) the smaller group is all of
\(G_d\) and the retained matrix is \(H_d\). When \(t>0\) its
moving dimension \(d-t\) is smaller. When the moving dimension
is zero the support is a singleton, and the identity matrix has
a unit entry. Induction proves equality with the stated final
rows, including the case of trivial groups.

All the specified arrays and both row modules are defined over
\(R[1/q]\). A finite extension of complete discrete valuation
rings is faithfully flat, as is completion of a discrete
valuation ring. Equality of the two row modules descends along
these extensions, by applying faithful flatness to the quotient
of their sum by each one. Equality at all maximal ideals of
\(R[1/q]\) gives equality over the ring itself. Finally left
translation multiplies the output matrix by the invertible
matrix of the translating group element, which preserves its
row module.
\end{proof}

\begin{corollary}[The terminal ideals for inputs with different restrictions]
\leavevmode\label{cor:integralmixed627}
Use the fields, two auxiliary primes \(g_1,g_2\), and common
supports of Corollary~\ref{cor:mixedrestrictions626}. Put
\(Q=g_1g_2\). On each target row take the tensor of the
corresponding \(L_{k,d_v}\) and denote the resulting map, summed
directly over targets, by \(L\). For any matrix \(Z\) of integral
signed input columns on these supports and each minor size \(u\),
\[
           I_u(HZ)\,R[1/Q]= I_u(LZ)\,R[1/Q].
\]
In particular the two matrices have exactly the same determinantal
valuations and nonzero Smith elementary divisors up to units at
every prime over the terminal prime \(T\). The factors on the
right are computed on groups of dimensions at most \(d_v\), for
the very same columns and for every \(x\) on the period interval.
\end{corollary}
\begin{proof}
Use the two row transformations from the theorem at each place.
They still express either tensor product through the other over
\(R[1/Q]\). Take their direct sums and compose with \(Z\) on the
right. Cauchy--Binet gives the two inclusions of determinantal
ideals. At a prime above \(T\) the auxiliary primes are units.
Over the resulting discrete valuation ring the valuations of
the nonzero determinantal ideals are the successive partial
sums of the Smith exponents. Their differences determine the
individual exponents. This also preserves the rank.
\end{proof}

Thus the denominators in a normalized complex Gram matrix cause
no loss at \(T\) in this reduction of the finite projection. The
remaining estimates concern the sizes and divisibility of the
signed arithmetic coefficients after the indicated Fourier sums,
including the other places in the integral theta period.


\subsection{Global heights of the reduced signed period matrices}
\label{sec:reducedheights628}

We compare the matrices of Section~\ref{sec:integralrestrictions627}
at the two auxiliary primes themselves and at every complex embedding.
The comparison includes the mixed inner products used in
Proposition~\ref{prop:rank11}. Write \(q=p^a\) at one mark,
\(w=k-d\), \(P=P_{k,d}\), and
\[
 \gamma_j=\prod_{i=1}^j(q^i-1),\qquad
 b_P=wd+\tbinom d2=kd-\tbinom{d+1}2,\qquad
 U_{k,d}=\gamma_k/\gamma_w .
\]
The symbol \(b_P\) here gives the exponent in the \(p\)-part
\(q^{b_P}\) of \(|P|\). We use the original unnormalized
\(H_k\) and the integral transform \(L=L_{k,d}\) of
Theorem~\ref{thm:integralrestrictions627}.

\begin{lemma}[Multiplicities in the reduced period matrix]
\leavevmode\label{lem:reducedmultiplicities628}
The rows of \(L\) are the matrix entries, with zeros omitted
if necessary, of a representation \(\tau\) of \(P\) by monomial
matrices with roots of unity as nonzero entries. In characteristic
zero \(\tau\) is multiplicity-free. Its constituents are exactly
the constituents of \(\Gamma_k|_P\).
Group its output rows \(L_\nu\) by the successive choices of
character orbits in the reduction. At each stage keep all ordered
source and target pairs from that orbit in the same group.
For every \(z,z'\) on \(P\) and at every complex embedding one has
\[
 \frac{\langle H_kz,H_kz'\rangle}{\gamma_k}=
       \sum_\nu \frac{m_\nu}{\gamma_k}
                       \langle L_\nu z,L_\nu z'\rangle,\qquad
 \gamma_w\le m_\nu\le\gamma_k,\qquad
 v_p(m_\nu)\le a\tbinom{d+1}2 .
\]
The inner products on the right are the usual ones on matrix entries.
Every constituent in group \(\nu\) occurs \(m_\nu\) times in
\(\Gamma_k|_P\).
\end{lemma}
\begin{proof}
If \(d=0\), \(P\) is trivial. Use its trivial representation and
multiplicity \(\gamma_k\), a \(p\)-adic unit. Otherwise write
\(N=\operatorname{Hom}(C,W)\) for the normal translation subgroup.
For a dual character \(B\) with stabilizer \(S_B\subseteq G_d\),
its character \(\chi_B\) extends to \(N\rtimes S_B\) by giving it
value one on \(S_B\). With the notation of
Theorem~\ref{thm:integralrestrictions627}, the summand of
\(\Gamma_k|_P\) belonging to the orbit of \(B\) is induced from
\(\chi_B\otimes e_B\Gamma_k\) on \(N\rtimes S_B\), where the
second factor denotes the action of \(S_B\) on that space.
Indeed its \(N\)-eigenspaces are indexed by that orbit and are
carried one to another by \(l(a_i)\). In characteristic zero the
calculation of the trace on \(e_B\Gamma_k\) in that theorem gives
\(m(k,d,t)\) copies of \(\Gamma_d|_{S_B}\), where
\[
             m(k,d,t)=q^{t(t+1)/2}\frac{\gamma_w}{\gamma_t},
                    \qquad t=\operatorname{rank}B .
\]
Define \(\tau\) recursively by retaining one copy at each such
step. End with \(\Gamma_d\) when \(t=0\); in moving dimension
zero use the trivial representation. Induction of a representation
by monomial matrices and tensor product with \(\chi_B\) again give
monomial matrices in the coset basis. Their block from \(j\) to
\(i\) is precisely the action of the array \(z_{ij}\) in the
theorem there. Induction through the steps identifies their entries
with the rows asserted in the lemma.

For irreducible representations of \(S_B\), the corresponding
induced representations are irreducible and are distinguished by
their orbit and their original \(S_B\)-representation. Indeed
every invariant subspace splits into its \(N\)-eigenspaces.
Its piece at \(B\) is stable under \(S_B\), and determines all the
other pieces by translation. The same argument applied to an
intertwiner shows that inequivalent choices at \(B\), or different
orbits, cannot give equivalent representations. At the last stage
\(\Gamma_e\) is multiplicity-free by
Theorem~\ref{thm:whittakerspectrum527}, also in dimension one.
This proves the assertion on multiplicities.

Along an orbit path the factors \(m(k,d,t)\) multiply, with
the factor \(\gamma_j\) if the path ends at a singleton in
ambient dimension \(j\). The first factor is at least \(\gamma_w\)
because \(\gamma_t\le q^{t(t+1)/2}\); all later factors are
positive integers. A nonzero representation of multiplicity
\(m_\nu\) consumes at least \(m_\nu\) dimensions of \(\Gamma_k\).
Finally the moving dimension decreases by \(t\) at each positive
step, so the successive \(t\)'s sum to at most \(d\). Since all
\(\gamma_j\) are \(p\)-adic units and
\(\sum_i t_i(t_i+1)/2\le d(d+1)/2\), this proves the last bound.

The action at any embedding is unitary in the original monomial
bases. Character equality gives a unitary isomorphism with
the indicated orthogonal sum: an intertwining isomorphism \(A\)
can be made unitary using \(A(A^*A)^{-1/2}\).
Unitary conjugation preserves the inner product of matrices.
Each copy contributes the same mixed product of entries.
This gives the identity, including for signed inputs.
\end{proof}

\begin{theorem}[The loss at the residue characteristic]
\leavevmode\label{thm:residueconductor628}
Write \(\mathcal R_H,\mathcal R_L\) for the modules generated by
the respective integral rows in the dual of the coefficient lattice
on \(P\). Over the full coefficient ring of integers one has
\[
       q^{2b_P+\binom{d+1}2}\mathcal R_L\ \subseteq\ \mathcal R_H,
       \qquad q^{2b_P}\mathcal R_H\ \subseteq\ \mathcal R_L .
\]
Away from \(p\) the two modules are equal. The bounds apply to
a translated support and to arrays with a common coefficient
lattice at the mark.
\end{theorem}
\begin{proof}
Work locally at a prime over \(p\), first after a finite splitting
extension. Let \(\rho\) be either \(\Gamma_k|_P\) or \(\tau\) and
let \(T_\rho\) take the entries of \(\sum z_g\rho(g)\).
The coefficientwise complex conjugates in \(T_\rho^*\) are again
roots of unity or zero. On functions on \(P\), the matrix
\(T_\rho^*T_\rho\) is central convolution with the character of
\(\rho\), with inversion according to the convolution convention.
In fact its \(g,h\) entry is \(\operatorname{tr}\rho(g^{-1}h)\).
Character orthogonality gives its nonzero scalars
\(|P|m_\pi/d_\pi\), where \(d_\pi\) is the irreducible degree and
\(m_\pi\) the multiplicity in \(\rho\). Index the central
projectors in this convention, replacing \(\pi\) by its dual if
necessary. Each projector \(e_\pi\) has coefficients of the form
\(d_\pi\chi_\pi(g)/|P|\), possibly with \(g\) inverted.
The character values are algebraic integers since each eigenvalue
of a finite-order matrix is a root of unity.

On the common support of the two representations take
\[
               C_\rho=\sum_{\pi:\,m_\pi>0}
                        \frac{d_\pi}{|P|m_\pi}e_\pi .
\]
Thus \(C_\rho T_\rho^*T_\rho\) is the identity on those
constituents and zero on their complement. If \(\rho'\) is the
other representation, its matrix of entries satisfies
\[
                T_{\rho'}=T_{\rho'}C_\rho T_\rho^*T_\rho .
\]
The identity first holds over the splitting field in characteristic
zero by the scalar calculation. All its entries are algebraic, so
it holds also in the chosen completion. Each denominator in
\(C_\rho\) divides \(|P|^2m_\pi\) for its summand. Since
\(|P|=q^{wd}|G_d|\), its \(p\)-part is \(q^{b_P}\).
Lemma~\ref{lem:reducedmultiplicities628} now clears all
denominators at \(p\) by \(q^{2b_P+\binom{d+1}2}\) for the original
representation and \(q^{2b_P}\) for the reduced one. The displayed
factorization gives the two inclusions. Faithful flatness descends
them to the original local integers. At primes away from \(p\)
Theorem~\ref{thm:integralrestrictions627} gives equality. Localizing
at all maximal ideals proves the global inclusions. Left translation
changes the original output by an invertible monomial matrix, and
tensoring row transformations with the identity preserves them.
\end{proof}

\begin{corollary}[All-place comparison on the original period fields]
\leavevmode\label{cor:globalreducedheight628}
Use the same marked places, translates and fields as in
Corollary~\ref{cor:integralmixed627}. On each target row \(\alpha\)
choose one common support for all its incoming signed columns, with
moving dimensions \(d_{\alpha v}\). Form the product of the local
transforms on that row and then the direct sum over rows. Denote
the full original and reduced maps by \(H,L\), and normalize \(H\)
at infinity by dividing on each row by
\((\prod_v\gamma_k(q_v))^{1/2}\); write \(H^\circ\) for that metric
normalization. In the \(L\)-coordinates put the squared-norm
weight \(\prod_v m_{\nu_v}/\gamma_k(q_v)\) on each final block
of that row. With this metric, the complex Gram matrices of
\(H^\circ Z\) and \(LZ\) coincide for every family \(Z\) of signed
input columns and at every embedding.

For the unweighted usual metric on \(L\), let \(h\) be the
logarithmic height of the exterior product of \(u\) independent
output columns. Use the original integral matrices at finite
places. Uniformly for every such choice over a coefficient
field \(J\),
\[
  \frac{|h(H^\circ Z)-h(LZ)|}{u}
     \ \ll\ k\sum_{v\mid g_1g_2}
                   \bigl(\max_\alpha d_{\alpha v}\bigr)\log q_v
     \ \ll\ f k^2\log(g_1g_2).
\]
This bound is \(o(b_x/x^A)\) for each fixed \(A>0\) at all \(x\)
on the period intervals, where \(b_x=\exp(x^\varepsilon)\).
The ranks of the two outputs agree. At the terminal prime their
determinantal ideals agree exactly.
\end{corollary}
\begin{proof}
The mixed-product identity of the lemma tensors over the marks
and adds over target rows, proving the exact Gram assertion.
With the unweighted metric on \(L\), put
\(U=\max_\alpha\prod_v U_{k,d_{\alpha v}}(q_v)\).
The squared norm for \(H^\circ\) on any input then lies between
\(U^{-1}\) times the squared norm for \(L\) and the squared norm
for \(L\). This proves equality of kernels, also after taking
any subfamily. Restrict the forms to the span of \(u\) independent
outputs using their common input quotient. Taking determinants
there bounds the ratio of the exterior norms between
\(U^{-u/2}\) and one.

Let \(D\) be the product over \(v\) of \(q_v\) to the largest
exponent \(2b_P+\binom{d_{\alpha v}+1}2\) on a target row at \(v\).
Theorem~\ref{thm:residueconductor628} gives integral row
transformations in both directions after multiplication by \(D\),
by tensoring on each row and taking the direct sum.
Cauchy--Binet therefore gives \(D^u I_u(LZ)\subseteq I_u(HZ)\)
and \(D^u I_u(HZ)\subseteq I_u(LZ)\) for the fractional ideals
of the \(u\)-rowed minors. The bound for the normalized difference
of their logarithmic ideal norms is \(u\log D\). These ideals agree
at \(T\) by Corollary~\ref{cor:integralmixed627}.

Explicitly, the height used here is the sum of the logarithms
of the exterior norms over all embeddings of \(J\), less the
logarithm of the absolute norm of the minor ideal, all divided
by \([J:\QQ]\). The product formula makes it independent of a
common change of basis of the input columns. Combining the finite
and infinite estimates costs at most \(\log D+\tfrac12\log U\)
per rank. Now \(2b_P+\binom{d+1}2=2kd-\binom{d+1}2\le 2kd\)
and \(\log U_{k,d}\le kd\log q\).
There are \(f\) places over each of \(g_1,g_2\), with the residue
sizes specified in Corollary~\ref{cor:whittakerterminal527}.
The last bound there is polynomial in \(x\) for fixed field
data and \(\varepsilon\), giving the asserted little-oh estimate.
\end{proof}

In particular one can form the diagonal and off-diagonal second
moments in Proposition~\ref{prop:rank11} with the weighted reduced
matrices and their original signed coefficients. Their numerical
values at infinity are unchanged. To establish that proposition's
hypotheses for the theta maps one still needs estimates for the
arithmetic coefficients in these matrices and the comparison with
the prescribed global period lattice.

\clearpage
\section{References}\label{nav:bibliography-XVIII}
\begingroup\def\section*#1{}\begin{thebibliography}{99}\small
\bibitem{XVIII-XIII-IV-etingof86} P. Etingof et al., \emph{Introduction to Representation Theory}, Student Mathematical Library 59, American Mathematical Society, 2011. \href{https://math.mit.edu/~etingof/representationtheorybook.pdf}{Theorem 5.22.1}.
\bibitem{XVIII-XIII-IV-kudla29} S. S. Kudla, \emph{Notes on the local theta correspondence}, European School of Group Theory, 1996. \href{https://www.math.toronto.edu/skudla/castle.pdf}{Chapter III, Section 6, pp.~53--54}.
\bibitem{XVIII-XIII-IV-ghy209} W. T. Gan, J. P. Hanke and J.-K. Yu, \emph{On an exact mass formula of Shimura}, Duke Math. J. \textbf{107} (2001), 101--133. \href{https://web.math.princeton.edu/~jonhanke/Web-02/Mass-Formula-of-Shimura/shimura-mass.pdf}{Section 4, formulas (4.2)--(4.3) and Propositions 4.4--4.5}.
\bibitem{XVIII-XIII-I-kott91} R. E. Kottwitz, \emph{Tamagawa numbers}, Ann. of Math. (2) \textbf{127} (1988), 629--646. \href{https://annals.math.princeton.edu/1988/127-3/p05}{Invariance under inner twisting}.
\bibitem{XVIII-XIII-I-olbrich91} M. Olbrich, \emph{\(L^2\)-invariants of locally symmetric spaces}, Doc. Math. \textbf{7} (2002), 219--237. \href{https://arxiv.org/pdf/math/0009039}{Section 5, Corollary 5.2}.
\bibitem{XVIII-XIII-green527} J. A. Green, \emph{The characters of the finite general linear groups}, Trans. Amer. Math. Soc. \textbf{80} (1955), 402--447, Theorem~14. \href{https://doi.org/10.1090/S0002-9947-1955-0072878-2}{Original paper}.
\bibitem{XVIII-XIII-IV-serre211} J.-P. Serre, \emph{Bounds for the orders of the finite subgroups of \(G(k)\)}, in \emph{Group Representation Theory}, EPFL Press, 2007. \href{https://arxiv.org/pdf/1011.0346}{Lecture I, Section 1.1, Theorem 1}.
\bibitem{XVIII-XIII-projectivecharacters627} P. Webb, \emph{A Course in Finite Group Representation Theory}, Cambridge University Press, 2016, Proposition~9.4.5 and Corollary~10.2.4. \href{https://www-users.cse.umn.edu/~webb/RepBook/RepBookLatex.pdf}{Author's text}.
\end{thebibliography}\endgroup
\end{document}
